%% ============================================================
%% Foundations of Sensitivity Analysis: Theory, Computation, and Applications
%% Leon M. Arriola and James M. Hyman
%%
%% Master book file — compile with pdflatex (two passes)
%% ============================================================

\documentclass[12pt]{book}

%% ---- Shared preamble (packages, macros, environments) ----
%% ============================================================
%% preamble.tex
%% Shared preamble for
%%   Foundations of Sensitivity Analysis: Theory, Computation, and Applications
%%   Arriola & Hyman
%% ============================================================

%% ----------------------------------------------------------
%% PACKAGES
%% ----------------------------------------------------------
\usepackage[margin=0.75in]{geometry}
\usepackage{amsmath, amssymb, amsthm}
\usepackage{graphicx}
\graphicspath{{figures/}}
\usepackage{booktabs}
\usepackage{array}
\usepackage{longtable}
\usepackage{enumitem}
\usepackage{imakeidx}
\makeindex[columns=2, title=Index, intoc]
\usepackage{hyperref}
\usepackage[protrusion=true,expansion=false]{microtype}  %% improved microtypography
\hypersetup{
  colorlinks = true,
  linkcolor  = black,
  citecolor  = black,
  urlcolor   = black,
  pdftitle   = {Foundations of Sensitivity Analysis: Theory, Computation, and Applications},
  pdfauthor  = {Leon M. Arriola and James M. Hyman},
  pdfsubject = {Sensitivity Analysis, Uncertainty Quantification, Mathematical Modeling},
  pdfkeywords= {sensitivity analysis, adjoint methods, Sobol indices,
                SIR model, forward sensitivity, uncertainty quantification},
}
\usepackage{xcolor}
\usepackage{tcolorbox}
\usepackage{pgfplots}
\usepackage{tikz}
\usetikzlibrary{arrows.meta, positioning, decorations.markings, decorations.pathreplacing, calc, bending}
\pgfplotsset{compat=1.18}

%% ----------------------------------------------------------
%% THEOREM-STYLE ENVIRONMENTS
%% ----------------------------------------------------------
\theoremstyle{plain}
\newtheorem{theorem}{Theorem}[chapter]
\newtheorem{proposition}[theorem]{Proposition}
\newtheorem{lemma}[theorem]{Lemma}
\newtheorem{corollary}[theorem]{Corollary}

\theoremstyle{definition}
\newtheorem{definition}{Definition}[chapter]
\newtheorem{example}{Example}[chapter]
\newtheorem{remark}{Remark}[chapter]

%% ----------------------------------------------------------
%% BOOK-WIDE NOTATION
%% ----------------------------------------------------------

%% Sensitivity index S_p^q
\newcommand{\SI}[2]{S_{#1}^{#2}}

%% Parameter of Interest / Quantity of Interest
%% \xspace inserts a space after the acronym when followed by ordinary text,
%% but suppresses the space before punctuation and closing braces.
\usepackage{xspace}
\newcommand{\POI}{\textsc{poi}\xspace}
\newcommand{\QOI}{\textsc{qoi}\xspace}
\newcommand{\POIs}{\textsc{poi}s\xspace}
\newcommand{\QOIs}{\textsc{qoi}s\xspace}

%% Number sets
\newcommand{\R}{\mathbb{R}}
\newcommand{\N}{\mathbb{N}}

%% Probability / expectation
\newcommand{\E}{\mathbb{E}}
\newcommand{\Var}{\operatorname{Var}}

%% Derivatives
\newcommand{\pd}[2]{\frac{\partial #1}{\partial #2}}
\newcommand{\dd}[2]{\frac{d #1}{d #2}}

%% FSE shorthand
\newcommand{\FSE}{\textsc{fse}}

%% Unifying idea callback (inline, unnumbered)
\newcommand{\unifyingidea}[1]{\noindent\textit{Unifying Idea~#1.}\enspace}

%% Absolute / relative / regularized change
\newcommand{\Dabs}{\Delta u_{\mathrm{Abs}}}
\newcommand{\Drel}{\Delta u_{\mathrm{Rel}}}
\newcommand{\Dreg}{\Delta u_{\mathrm{Reg}}}

%% Math operators
\DeclareMathOperator{\arcsinh}{arcsinh}
\DeclareMathOperator{\sgn}{sgn}

%% ----------------------------------------------------------
%% PEDAGOGICAL ENVIRONMENTS
%% ----------------------------------------------------------

%% Self-check question with indented answer
\newenvironment{selfcheck}{%
  \medskip\noindent\textit{Self-check.}\enspace
}{\par}

\newenvironment{scAnswer}{%
  \nopagebreak\smallskip
  \noindent\hspace*{1.8em}%
  \begin{minipage}{0.90\textwidth}
    \footnotesize\textit{Answer:}\enspace
  }{\end{minipage}\medskip\par}

%% Numbered exercise environment (resets each chapter)
\newcounter{exercise}[chapter]
\renewcommand{\theexercise}{\thechapter.\arabic{exercise}}

\newenvironment{exercise}[1][]{%
  \refstepcounter{exercise}%
  \medskip\noindent\textbf{Exercise~\theexercise}%
  \ifx\relax#1\relax\else\ \textit{(#1)}\fi.\enspace
}{\par\medskip}

%% Bloom level tag (suppress in production by redefining to empty)
\newcommand{\bloom}[1]{\textnormal{\tiny\rmfamily[Bloom~#1]}}

%% ----------------------------------------------------------
%% CHAPTER-OPENING STYLE (book class default used;
%% customize here if needed)
%% ----------------------------------------------------------

%% Part / chapter appearance tweaks can go here.
%% For now, book class defaults are used throughout.


%% ---- Book metadata ----
\title{%
  \textbf{Foundations of Sensitivity Analysis}\\[6pt]
  \large Theory, Computation, and Applications
}
\author{Leon M. Arriola \and James M. Hyman}
\date{}

%% ============================================================
\begin{document}
%% ============================================================

\frontmatter

%% ---- Title page ----
\maketitle

%% ---- Copyright / colophon ----
\thispagestyle{empty}
\vspace*{\fill}
\begin{center}
\small
\copyright\ 2026 Leon M. Arriola and James M. Hyman.\\
All rights reserved.
\end{center}
\vspace*{\fill}
\clearpage

%% ---- AMS Classification and Keywords ----
\thispagestyle{empty}
\vspace*{\fill}
\begin{flushleft}\small
\textbf{AMS subject classifications.}
65-01,
34A55,
62-07,
92D30,
49K15,
65C05

\medskip
\textbf{Keywords.}
sensitivity analysis,
adjoint method,
forward sensitivity equations,
Sobol indices,
global sensitivity analysis,
uncertainty quantification,
mathematical modeling,
epidemic models,
automatic differentiation
\end{flushleft}
\vspace*{\fill}
\clearpage

%% ---- Table of contents ----
\tableofcontents
\clearpage

%% ---- Preface ----
\chapter*{Preface}
\addcontentsline{toc}{chapter}{Preface}

This book develops sensitivity analysis (SA) as a mathematical discipline,
not merely as a collection of computational recipes.
SA asks the foundational question of quantitative modeling:
which parameters most strongly determine the model's predictions?

\medskip

Three ideas run through every chapter:
\begin{enumerate}[nosep]
  \item The sensitivity index is a normalized derivative.
  \item The adjoint is the transpose, in whatever space you work.
  \item Local SA is a linearization, with all the limitations that implies.
\end{enumerate}

Students who understand these three ideas (not just their definitions
but their implications and connections) will be equipped to design,
implement, and critique SA studies in any application domain.

\medskip

The SIR epidemic model serves as the running example throughout.
Every method is demonstrated on this one model, so students can
track increasing mathematical sophistication against a fixed biological
context.
All computational examples use MATLAB; each chapter includes a
structured laboratory and a reusable code template.
All code is freely available in the companion repository:
\begin{center}
  \texttt{https://github.com/mhyman/arriola2026sa}
\end{center}
The repository contains the complete \texttt{src/} library of
production-quality MATLAB functions (one per major algorithm),
three test suites with 48 automated checks, and this instructor's
solutions manual.

\medskip

The adjoint method (Chapters~5--7) is the
conceptual centerpiece of the book. It is the same idea encountered
three times in three different mathematical spaces: the matrix transpose
in linear algebra, the reversed computation graph in algorithmic
differentiation, and a backward ODE in dynamical systems.
Chapter~8, \emph{Caveat Emptor}, asks when sensitivity analysis can
mislead, and is the chapter students most frequently cite as the
one that changed how they think about mathematical models.

\medskip

\noindent\textbf{Relationship to existing literature.}
Cacuci (1981) established the adjoint formalism for nonlinear systems;
Giles and Pierce (2000) developed it for aeronautical design;
Cao et al.\ (2002) gave the computational treatment for ODE systems;
and Griewank and Walther (2008) treat the computation-graph view
comprehensively. The present book differs from each of these in its
pedagogical aim: it is the only undergraduate text that introduces all
three forms of the adjoint (matrix, graph, and ODE) as instances
of a single unifying idea, and that teaches students when \emph{not}
to trust the results they compute.

\medskip

\noindent\textbf{For instructors.}
The book covers a one-semester undergraduate course.
Students need calculus, linear algebra, and ODEs; no measure theory or
functional analysis is assumed.
A minimal one-semester course can omit the computational AD material
in Chapter~7 and all appendices without loss of continuity.
Chapters~1--4 and~8--9 form a self-contained global-SA course
for students who will not use the adjoint in their research.
Exercises tagged \textsc{[Bloom~L5]} and \textsc{[L6]} are designed for
semester projects and take-home examinations; those tagged \textsc{[L2]}--\textsc{[L4]}
are suitable for weekly problem sets.

\medskip

\medskip

\noindent\textit{Leon M. Arriola, University of Wisconsin-Whitewater}\\
\noindent\textit{James M. Hyman, Tulane University}

\medskip

\noindent\textbf{Acknowledgments.}
We thank the students in our sensitivity analysis courses at the
University of Wisconsin-Whitewater, Tulane University, and the
Mathematical and Theoretical Biology Institute (MTBI) at Arizona
State University, whose questions and confusion shaped every page
of this book.

The MATLAB code library in the companion GitHub repository
(\texttt{https://github.com/mhyman/arriola2026sa}) was
developed with support from the Tulane University Department of
Mathematics.

\medskip

\noindent\textbf{Dedication.}

\medskip
\begin{center}
\textit{To our students, past and future:} \\[4pt]
\textit{May you always ask which parameters matter} \\
\textit{before you run the model.}
\end{center}

%% ============================================================
\mainmatter
%% ============================================================

%% ---- Part I: The Sensitivity Question ----
\part{The Sensitivity Question}

\chapter{Introduction and Overview}

\begin{flushright}
\textit{Sensitivity analysis begins when a model is forced to say what matters.}
\end{flushright}

\bigskip

An infectious disease has just entered a new country.
The public health response team has built a compartmental model with fourteen
parameters: transmission rates, recovery rates, contact patterns,
immunity duration, disease-induced mortality, and more.
The team has funding to deploy exactly two interventions
and time to measure exactly three parameters with high precision
before a policy decision must be made.

The question due by tomorrow morning is this: \emph{which three parameters
are worth measuring, and which two interventions will most effectively
reduce the spread?}

This is the problem that sensitivity analysis is designed to answer.
Before the mathematics, we need a language for asking the question precisely.

%%------------------------------------------------------------
\section{The Central Question}
\label{sec:centralquestion}
%%------------------------------------------------------------

Every mathematical model takes inputs and produces outputs.
The inputs include parameters whose exact values may be uncertain
or whose values might be changed through intervention.
The outputs include whatever the model is designed to predict.
Sensitivity analysis asks: \emph{which inputs have the most
effect on the outputs we care about?}

To answer this precisely, we adopt two terms used throughout the book.

\medskip

A \textbf{Parameter of Interest} (\POI)\index{parameter of interest (POI)} is any input to the model
whose variation is worth studying. In practice, a parameter becomes a \POI
for one of three reasons, and the reason matters for how you interpret
the sensitivity index.

\begin{enumerate}

\item \textbf{Epistemic uncertainty} (uncertainty from lack of knowledge).
The parameter is fixed by nature but imperfectly known: different studies
report different estimates, the data are sparse, or the measurement process
is imprecise. Epistemic uncertainty is \emph{reducible} --- better data,
improved experiments, or longer observation windows can sharpen the estimate.
Sensitivity analysis identifies which epistemically uncertain parameters
most affect the output, directing limited data-collection resources
toward the measurements that would reduce output uncertainty the most.
In the disease example, the transmission probability per contact is
epistemically uncertain: studies estimate it from contact-tracing data,
but estimates vary across populations and settings.

\item \textbf{Aleatory uncertainty} (inherent stochastic variability).
The parameter genuinely varies from individual to individual, place to
place, or time to time, and this variation is \emph{irreducible}---
it persists even with unlimited data. More observations characterize
the distribution more precisely but do not eliminate the randomness.
In the disease example, the contact rate varies stochastically
across the population: individuals differ in how many contacts they
make per day. Sensitivity analysis with aleatory uncertainty typically
asks how much of the output variance is attributable to each
stochastic source, motivating the global SA methods of Chapter~9.

\item \textbf{Control variable} (a lever the investigator can move).
The parameter is known and could, in principle, be changed by a policy
decision, intervention, or design choice. Sensitivity analysis in this
context answers the question: which intervention has the most leverage
over the outcome I care about? In the disease example, the contact rate
can be reduced by social distancing policies, and the transmission
probability can be reduced by mask requirements or vaccination. The
sensitivity index tells the policy maker which reduction --- contacts
or transmission --- produces the largest reduction in $R_0$.

\end{enumerate}

Many parameters fall into more than one category simultaneously:
the contact rate is both uncertain (its current value is imprecisely
known) and controllable (it can be targeted by policy).
In practice, the interpretation of a sensitivity index should
always be read in light of which role the \POI is playing.
A large sensitivity index for an epistemically uncertain parameter
is a call to collect better data; for an aleatory parameter, it is
a call to characterize the distribution more carefully; for a control
variable, it is a call to prioritize that intervention.

In the disease example above, the transmission probability
per contact is primarily a \POI because it is epistemically uncertain
(different studies report different values) and because it can be
targeted by an intervention (vaccines reduce it).
The contact rate is a \POI for both the same reasons:
it varies across populations (epistemic and aleatory uncertainty)
and can be reduced through social distancing policies (control variable).

A \textbf{Quantity of Interest} (\QOI)\index{quantity of interest (QOI)} is any model output that the
investigator cares about predicting or controlling. The peak number of
infectious individuals, the total number of deaths over a six-month
period, and the basic reproduction number $\mathcal{R}_0$ are all
plausible \QOIs for an epidemic model. The choice of \QOI is not neutral:
different choices lead to different rankings of parameter importance.

The relationship between \POIs and \QOIs is the subject of this book.
Specifically, we want to know how much each \POI affects each \QOI,
so that we can rank parameters by importance, design interventions
efficiently, and allocate measurement resources wisely.

One of the most persistent misconceptions in mathematical modeling
is the assumption that all parameters matter equally until proven
otherwise. They do not. In virtually every model studied in this
book, a small number of \POIs account for most of the variation
in the \QOIs, while the remaining parameters have negligible effect.
Identifying that small number is one of the primary goals of
sensitivity analysis, and the answer is rarely obvious from the
model equations alone.

\begin{selfcheck}
In the outbreak scenario above, classify each of the following
as a \POI, a \QOI, or neither:
the transmission probability per contact;
the peak number of infectious individuals;
the day the first case was reported;
the mean number of days a patient is infectious.
\end{selfcheck}
\begin{scAnswer}
\POI: transmission probability (uncertain; can be reduced by vaccines),
mean infectious period (uncertain; can be reduced by treatment).
\QOI: peak number of infectious individuals (model output we care about).
Neither: day of first case (a historical fact, not a model parameter
or prediction). Note that the mean infectious period could also be
a \QOI in a different study; this classification depends on
the question being asked.
\end{scAnswer}

Figure~\ref{fig:forwardproblem} illustrates the basic structure.
The forward problem takes nominal \POI values as input
and produces the \QOI as output. Sensitivity analysis adds one
more step: it asks how the \QOI changes as the \POIs vary.

\begin{figure}[htb]
  \begin{center}
  \begin{tikzpicture}[
    node distance = 2.6cm,
    box/.style = {
      rectangle, rounded corners=4pt, draw=black, thick,
      minimum width=2.5cm, minimum height=1.1cm,
      text centered, font=\small
    },
    sarrow/.style = {-{Latex[length=3mm]}, thick}
  ]
    %% Main flow nodes — single-line labels avoid \\[...] issues
    \node[box, fill=black!8]  (pois)  {\textbf{\POIs}};
    \node[box, fill=black!15, right=of pois]  (model) {\textbf{Model}};
    \node[box, fill=black!8,  right=of model] (qois)  {\textbf{\QOIs}};

    %% Sub-labels below each box
    \node[font=\footnotesize\itshape, below=2pt of pois]  {parameters};
    \node[font=\footnotesize\itshape, below=2pt of model] {equations};
    \node[font=\footnotesize\itshape, below=2pt of qois]  {outputs};

    %% Forward arrows
    \draw[sarrow] (pois)  -- node[above, font=\footnotesize]{$\vec{p}$} (model);
    \draw[sarrow] (model) -- node[above, font=\footnotesize]{$\vec{q}(\vec{p})$} (qois);

    %% SA brace and formula below
    \draw[thick, decorate,
          decoration={brace, amplitude=6pt, mirror}]
      ([yshift=-18pt]pois.south west) -- ([yshift=-18pt]qois.south east)
      node[midway, below=9pt, font=\small\itshape]{
        Sensitivity Analysis:\quad
        $\SI{p_i}{q_j} = \dfrac{p_i}{q_j}\,\dfrac{\partial q_j}{\partial p_i}$
      };
  \end{tikzpicture}
  \end{center}
  \caption{The forward problem takes nominal values for the input \POIs
           and produces the associated output \QOI. Sensitivity analysis
           quantifies how changes in the \POIs propagate to the \QOI
           by computing the normalized partial derivative $\SI{p_i}{q_j}$
           for each \POI--\QOI pair.}
  \label{fig:forwardproblem}
\end{figure}

%%------------------------------------------------------------
\section{Two Directions: Forward and Adjoint Sensitivity}
\label{sec:twoDirections}
%%------------------------------------------------------------

There are two complementary ways to ask the sensitivity question,
and they differ in a fundamental way.

Figure~\ref{fig:forwardadjoint} captures the distinction geometrically.
In panel~(a), a perturbation originates at a specific input location
and spreads forward through the model, influencing a broad region of
the output. This is \textbf{forward sensitivity analysis}\index{forward sensitivity analysis}.
Panel~(b) reverses the direction: a specified variation in the output
is traced backward to the input locations that could have caused it.
This is \textbf{adjoint sensitivity analysis}\index{adjoint sensitivity analysis}.

\begin{figure}[htb]
  \begin{center}
  \begin{tikzpicture}[
    modelbox/.style = {
      rectangle, rounded corners=5pt, draw=black, line width=1.2pt,
      minimum width=1.6cm, minimum height=4.4cm,
      text centered, fill=black!12, font=\small\bfseries
    },
    qoibox/.style = {
      rectangle, rounded corners=4pt, draw=black, thick,
      minimum width=1.4cm, minimum height=0.9cm,
      text centered, fill=black!6, font=\small
    },
    farrow/.style  = {-{Latex[length=3mm]}, thick},
    garrow/.style  = {-{Latex[length=2.5mm]}, thin, black!35},
    aarrow/.style  = {-{Latex[length=3.5mm]}, very thick, dashed},
    seclabel/.style = {font=\small\bfseries},
    sublabel/.style = {font=\footnotesize\itshape},
    plab/.style     = {font=\small},
  ]

    %%----------------------------------------------------------
    %% Panel (a): Forward SA
    %% One POI (p_j) is perturbed; its effect flows through the
    %% Model and appears at the QOI. Other POIs shown faded to
    %% indicate the process repeats for each one.
    %%----------------------------------------------------------
    \node[seclabel] at (2.8, 2.55) {(a) Forward SA \quad ``What if?''};

    %% POI column
    \node[plab]                      (fp1) at (0.0,  1.50) {$p_1$};
    \node[plab]                      (fp2) at (0.0,  0.75) {$p_2$};
    \node[font=\small]                     at (0.0,  0.10) {$\vdots$};
    \node[plab, font=\small\bfseries](fpj) at (0.0, -0.55) {$p_j$};
    \node[font=\small]                     at (0.0, -1.10) {$\vdots$};
    \node[plab]                      (fpm) at (0.0, -1.70) {$p_m$};

    \node[modelbox] (fmod) at (2.8, -0.10) {Model};
    \node[qoibox]   (fqoi) at (5.4, -0.10) {\QOI};

    %% Faded arrows: all other POIs enter model at their y-level
    \draw[garrow] (fp1.east) -- (fmod.west |- fp1.center);
    \draw[garrow] (fp2.east) -- (fmod.west |- fp2.center);
    \draw[garrow] (fpm.east) -- (fmod.west |- fpm.center);

    %% Bold path: p_j through Model to QOI
    \draw[farrow, line width=1.4pt] (fpj.east) -- (fmod.west |- fpj.center);
    \draw[farrow, line width=1.4pt] (fmod.east) -- (fqoi.west);

    %% Perturbation indicator at p_j
    \fill[black!70] (-0.85,-0.55) circle (4pt);
    \draw[->,thick] (-0.85,-0.55) -- (fpj.west);
    \node[sublabel, left=20pt of fpj, align=right] {perturb\\$p_j$};

    \node[sublabel, align=center, text width=5.0cm] at (2.8, -3.20)
      {Perturb one \POI, observe the \QOI.\\
       Repeat for each \POI: $m$ solves total.};

    %%----------------------------------------------------------
    %% Dividing rule
    %%----------------------------------------------------------
    \draw[thin, gray!50] (7.0,-3.50) -- (7.0,2.80);

    %%----------------------------------------------------------
    %% Panel (b): Adjoint SA
    %% A specified QOI variation enters the Model from the right
    %% and flows BACKWARD, producing all POI sensitivities at once.
    %% Arrowhead at END of each dashed path (−{Latex}): points left.
    %%----------------------------------------------------------
    \node[seclabel] at (10.5, 2.55) {(b) Adjoint SA \quad ``How come?''};

    %% POI column
    \node[plab] (ap1) at (7.7,  1.50) {$p_1$};
    \node[plab] (ap2) at (7.7,  0.75) {$p_2$};
    \node[font=\small]  at (7.7,  0.10) {$\vdots$};
    \node[plab] (apj) at (7.7, -0.55) {$p_j$};
    \node[font=\small]  at (7.7, -1.10) {$\vdots$};
    \node[plab] (apm) at (7.7, -1.70) {$p_m$};

    \node[modelbox] (amod) at (10.5, -0.10) {Model};
    \node[qoibox]   (aqoi) at (13.1, -0.10) {\QOI};

    %% ONE dashed arrow: QOI → Model (backward; arrowhead enters model)
    \draw[aarrow] (aqoi.west) -- (amod.east);

    %% ALL dashed arrows: Model → each POI (backward; arrowhead at each POI)
    \draw[aarrow] (amod.west |- ap1.center) -- (ap1.east);
    \draw[aarrow] (amod.west |- ap2.center) -- (ap2.east);
    \draw[aarrow] (amod.west |- apj.center) -- (apj.east);
    \draw[aarrow] (amod.west |- apm.center) -- (apm.east);

    %% "specify δq_j" below the QOI box
    \node[sublabel, below=5pt of aqoi, align=center] {specify $\delta q_j$};

    \node[sublabel, align=center, text width=5.2cm] at (10.5, -3.20)
      {Specify one \QOI variation, solve adjoint backward.\\
       All \POI sensitivities recovered: 1 solve.};

  \end{tikzpicture}
  \end{center}
  \caption{(a)~Forward sensitivity analysis answers ``What if?''\ questions.
           A perturbation to one input \POI $p_j$ propagates forward through
           the model and produces a change in the output \QOI\@.
           Because the model must be solved once for each \POI, forward
           sensitivity requires $m$ solves for $m$ \POIs.
           (b)~Adjoint sensitivity analysis answers ``How come?''\ questions.
           A specified variation $\delta q_j$ in the output \QOI enters the
           model from the right and is propagated \emph{backward} through it.
           A single adjoint solve simultaneously recovers the sensitivity of
           $q_j$ with respect to every \POI $p_1, \ldots, p_m$.
           The adjoint is developed in Chapters~5 and~6; for now, the key
           contrast is computational: forward SA costs $m$ solves,
           adjoint SA costs~$1$.}
  \label{fig:forwardadjoint}
\end{figure}

The ``What if?'' and ``How come?'' labels are not merely rhetorical.
They correspond to genuinely different computational strategies
and genuinely different questions.

\subsection{Forward Sensitivity Analysis: ``What if?''}

Consider a disease model for an emerging outbreak.
The transmission probability $\beta$ is estimated from early case data,
but the estimate carries significant uncertainty. A public health officer
asks: \emph{if our estimate of $\beta$ is wrong by 10\%, how much
does our predicted peak infection count change?}

This is a forward question. We are tracing the effect of a change
in an input \POI (the transmission probability) forward to an
output \QOI (the peak infection count). Forward sensitivity analysis
answers this question by computing a sensitivity index: a number that
tells us, for a small percentage change in a \POI, what percentage
change to expect in the \QOI.

A sensitivity index of $+0.8$ would mean that a 10\%\ increase in
$\beta$ produces approximately an $8\%$ increase in the peak infection
count. A sensitivity index of $-1.0$ for the recovery rate would mean
that a 10\%\ increase in the recovery rate produces a 10\%\ decrease
in the peak. The sign tells you the direction; the magnitude tells you
the amplification. Chapter~2 develops this definition carefully and
shows how to compute these indices for a wide range of models.

\begin{selfcheck}
A forward sensitivity analysis gives $S_{\tau}^{R_0} = +1.2$ (the notation is defined formally in Chapter~2),
where $\tau$ is the mean infectious period (in days) and $R_0$ is
the basic reproduction number.
If the mean infectious period increases by $5\%$, approximately
what happens to $R_0$? Is this good news or bad news for disease control?
\end{selfcheck}
\begin{scAnswer}
$R_0$ increases by approximately $5\% \times 1.2 = 6\%$. This is bad news:
a larger $R_0$ means more transmission. The positive sign tells us that
\POI and \QOI move in the same direction, so longer infectious periods
lead to more disease spread.
\end{scAnswer}

\subsection{Adjoint Sensitivity Analysis: ``How come?''}

Now consider a different problem. City planners want to site a new
industrial facility near a residential area. The facility emits
pollutants, and regulators have set a standard: the average pollutant
concentration at a specified monitoring station must not exceed a
given threshold. A transport model predicts whether any candidate
site meets this standard.

The planner's question is different in character from the previous one.
Rather than asking ``if this parameter changes, what happens to the
output?'', they want to know: \emph{for a given allowed change in the
output (the pollutant concentration), how much variation in the inputs
(wind speed, emission rate, atmospheric stability) is permissible?}

This is an adjoint question. It runs the sensitivity analysis in
reverse: from a specified variation in the \QOI back to the
\POIs that could have caused it. For models with many \POIs
and only a few \QOIs, adjoint sensitivity analysis is
dramatically more efficient than computing forward sensitivities
for every parameter individually. Chapters~5 and 6 explain
why this is so and how to exploit it.

\begin{selfcheck}
An environmental regulator measures pollutant concentration at a
monitoring station and finds it exceeds the legal threshold.
She wants to determine which weather assumptions in the
transport model are most likely responsible for the exceedance.
Is this a forward or adjoint question? Explain in one sentence.
\end{selfcheck}
\begin{scAnswer}
Adjoint. The regulator starts with a specified change in the \QOI
(exceedance of the threshold) and traces backward to identify which
input \POIs (weather assumptions) most plausibly caused it.
A forward analysis would instead ask how changing a specific weather
assumption would affect the predicted concentration, which runs in the opposite
direction.
\end{scAnswer}

%%------------------------------------------------------------
\section{Global Sensitivity: When Local Is Not Enough}
\label{sec:globalSA}
%%------------------------------------------------------------

Both forward and adjoint sensitivity analysis are \emph{local}:
they measure the effect of small perturbations to the \POIs
near a fixed nominal value. This is powerful, but it assumes
that the model behaves approximately linearly in the vicinity
of that nominal value.

Many real models are not linear, and many \POIs carry substantial
uncertainty that cannot be called ``small.'' In these cases,
local sensitivity indices can be misleading. A \POI with a
small local sensitivity index may still have a large effect
on the \QOI when varied across its full range of plausible values,
if the model is nonlinear in that parameter.

\begin{example}[Local SI does not capture global behavior]
\label{ex:local_global}
Consider $q(p) = p^2$ and suppose the nominal value is $p^* = 0.1$.
The local sensitivity index is $S_p^q = (p/q)\,(dq/dp) = (p/p^2)(2p) = 2$,
which looks moderate.
But if $p$ is uncertain and could range from $0.1$ to $1.0$,
the QOI ranges from $0.01$ to $1.0$ --- a hundredfold change.
A local SI computed at $p^* = 0.1$ gives no information about
this hundredfold range.
Global SA, introduced in Chapter~9, is designed to capture exactly
this kind of large-range nonlinear behavior.
\end{example}

Chapter~9 introduces \textbf{global sensitivity analysis}\index{global sensitivity analysis},
which asks how the output varies when the \POIs are varied
across their entire plausible ranges simultaneously.
The primary tool there, the Sobol index, measures
how much of the total output variance is attributable to
each \POI. Global SA does not replace local SA; it
complements it, answering questions that local SA cannot.

%%------------------------------------------------------------
\section{Sensitivity Analysis and Uncertainty Quantification}
\label{sec:sauq}
%%------------------------------------------------------------

Sensitivity analysis and uncertainty quantification are
related but distinct enterprises.

Sensitivity analysis, as developed in this book, is
\emph{deterministic}: we assign nominal values to the \POIs
and compute how the \QOI changes when those values are
perturbed. The \POIs are treated as numbers with specific
values, not as random variables. The sensitivity index
is a derivative: a precise mathematical object that exists
when the model is differentiable.

Uncertainty quantification treats the \POIs as random variables
with associated probability distributions. The goal is to
determine the distribution of the \QOI given the distributions
of the \POIs. This is a statistical enterprise, and its primary
tools (Monte Carlo methods, Bayesian inference, polynomial chaos)
are the subject of the companion volume by Smith~\cite{smith2014uncertainty}.
Smith's Chapters~9 and~10 develop local and global sensitivity analysis
at the graduate level, extending the material in this book with
functional-analytic and measure-theoretic foundations.
Smith's Chapter~6 treats parameter selection using the sensitivity
Jacobian introduced in our Chapter~4, connecting it to the
surrogate-model literature.
The SA foundations developed here follow Arriola and Hyman~\cite{arriola2007being,arriola2009sensitivity}.
For the broader development of SA as a formal discipline, see Saltelli et al.~\cite{saltelli2021sensitivity},
Razavi et al.~\cite{razavi2021future}, and the annotated historical timeline in
Tarantola et al.~\cite{tarantola2024annotated}.

The connection between the two fields is real and important.
Global SA, covered in Chapter~9, bridges the gap: Sobol indices
are defined in terms of variances, which requires probabilistic
inputs, but they can be estimated by methods that build directly
on the ideas in this book. Students who complete this course
will find that the second semester's UQ material makes more
sense because of the SA foundation established here.

%%------------------------------------------------------------
\section{Three Ideas That Run Through the Book}
\label{sec:threeIdeas}
%%------------------------------------------------------------

Every chapter of this book revisits one or more of three
fundamental ideas. We name them here so that students can
recognize them when they reappear, each time in a new setting.

\medskip

\noindent\textbf{Unifying Idea~1:\index{unifying ideas} The sensitivity index is a
normalized derivative.}

At its core, SA answers a single question: if this input changes
by a small percentage, by what percentage does the output change?
That ratio is the sensitivity index. It is not a mysterious new
object: it is the derivative, normalized so that its magnitude
is independent of units and scales.

This idea appears first in Chapter~2, where we define the
sensitivity index formally. It reappears in Chapter~3 as a
finite-difference approximation, in Chapters~4--6 as
an analytic derivative, and in Chapter~9 as the foundation
for Sobol indices. Whenever a sensitivity quantity appears
in this book, it is an instance of this one idea.
Smith~\cite{smith2014uncertainty} develops the same normalized-derivative
framework in \S9.1--9.2 and connects it to experimental design.

\medskip

\noindent\textbf{Unifying Idea~2: The adjoint is the transpose
in whatever space you are working in.}

The adjoint sensitivity method, which makes it possible
to compute the sensitivities of one \QOI with respect to
hundreds of \POIs at the cost of solving just two problems
instead of hundreds, rests on a single mathematical operation:
transposition. For a linear system $A\vec{u} = \vec{b}$,
the adjoint involves the transposed matrix $A^T$.
For a computation graph (the structure underlying every
computer program), the adjoint is the chain rule applied in
reverse order, which is the backpropagation operation familiar from
machine learning. For an ordinary differential equation,
the adjoint is a second ODE running backward in time.

These three forms look different on the surface.
They are the same idea.

Chapter~5 introduces the adjoint through the matrix transpose,
which is the simplest and most familiar form. Chapter~6 extends
it to differential equations. At each transition, the book
explicitly identifies the connection back to this unifying idea.

\medskip

\noindent\textbf{Unifying Idea~3: Local SA is a linearization.}

Every sensitivity index in Chapters~2--8 is based on a derivative,
and a derivative is a linear approximation to a generally nonlinear
function. This approximation is valid when the perturbation is small
and the model is smooth. It fails near bifurcation points, where
the model changes qualitative behavior. It fails when the \POI
uncertainty is large. It fails when the model is strongly nonlinear
in the \POI of interest.

Chapter~8 (\emph{Caveat Emptor}) documents the conditions
under which local SA breaks down. Chapter~9 provides
the global alternative. Both chapters are applications
of this one principle: local SA is a linearization,
and linearizations have limits.

%%------------------------------------------------------------
\section{A Note on Parameterization}
\label{sec:parameterization}
%%------------------------------------------------------------

One practical decision that profoundly affects the usefulness
of SA results deserves mention before Chapter~2.

The sensitivity index depends on how a \POI is defined.
In a compartmental epidemic model, the recovery process can
be parameterized either by the \emph{recovery rate} $\gamma$
(units: per day) or by the \emph{mean infectious period}
$\tau = 1/\gamma$ (units: days). These two parameterizations
are mathematically equivalent, each determining the other.
But their sensitivity indices are not equal: they have
the same magnitude and opposite signs.

Which one should you use? The answer is: whichever one you can
discuss in ordinary language without confusion. People who work
with epidemic models (clinicians, public health officials,
and policy makers) think naturally in terms of days.
A clinician understands immediately what it means to say
``the mean infectious period is 7 days.'' They do not
think in recovery rates. Reporting $\SI{\gamma}{R_0} = -1.0$
requires the reader to invert $\gamma$ mentally before the
result makes intuitive sense. Reporting $\SI{\tau}{R_0} = +1.0$
is immediately clear: a longer infectious period increases
transmission.

Throughout this book, we choose \POIs that can be understood
in everyday terms: durations rather than rates, probabilities
rather than their complements, counts rather than densities
when the distinction matters. This is not just a matter of
aesthetics. Sensitivity indices that require mental inversion
to interpret are sensitivity indices that will be misread.

Chapter~2 formalizes this principle and shows how to convert
between parameterizations. Readers who find that their model's
natural parameters are rates will learn how to restate
their results in terms of the corresponding durations
before reporting them.

%%------------------------------------------------------------
\section{Structure of the Book}
\label{sec:structure}
%%------------------------------------------------------------

The book is organized in five parts. Each part builds on
the previous, and the SIR epidemic model~\cite{kermack1927contribution,wu2013sensitivity} serves as the
recurring example throughout, so that students can track
a single application as the mathematical tools become
progressively more powerful.

\medskip

\textbf{Part~I: The Sensitivity Question} (Chapters~1--2)
establishes the central problem and introduces the sensitivity
index as the primary tool. Chapter~2 develops the formal
definition, the sign interpretation, the tornado plot
as a display format, and the principle of interpretable
parameterization. By the end of Part~I, students can compute
and interpret sensitivity indices for any differentiable
function of several variables.

\textbf{Part~II: Computing Sensitivity in Practice}
(Chapter~3) gives students an immediately usable computational
method for any model, including models that cannot be
differentiated analytically. Finite-difference approximations,
step-size selection, the sensitivity Jacobian, and the
treatment of correlated \POIs are developed here.
Students who work through Chapter~3 can perform SA on any
model they can run, including inherited code and black-box
simulation software.

\textbf{Part~III: Analytic Forward Sensitivity Analysis}
(Chapter~4) shows how to differentiate model equations directly
to obtain exact sensitivity information at lower computational
cost. The forward sensitivity equation is derived for algebraic
systems, difference equations, and ordinary differential
equations. The chapter ends by quantifying the computational
bottleneck that motivates the adjoint approach.

\textbf{Part~IV: The Adjoint Approach} (Chapters~5--7)
introduces adjoint sensitivity analysis through three
progressively more abstract settings: linear systems
in Chapter~5, dynamical systems in Chapter~6, and
computational tools including algorithmic differentiation
in Chapter~7. The central theme is Unifying Idea~2:
the adjoint is the matrix transpose in whatever
mathematical space the problem inhabits.

\textbf{Part~V: Limitations and Extensions}
(Chapters~8--9) addresses the boundaries of local SA.
Chapter~8 documents the conditions under which sensitivity
indices give misleading results: bifurcation points,
ill-conditioned systems, and method-question mismatch.
Chapter~9 introduces global SA through variance decomposition,
Sobol indices, and partial rank correlation coefficients.
It closes with an explicit bridge to the second semester,
where Smith's \emph{Uncertainty Quantification}~\cite{smith2014uncertainty}
takes over from the point where this book ends.

\medskip

\noindent\textbf{Computing environment.} Every chapter
includes a structured MATLAB laboratory that implements
the chapter's core idea on the SIR epidemic model.
Students who complete these laboratories will leave the
course with a portable toolkit for performing SA on any
model they encounter in research or practice.
The complete MATLAB code library --- production-quality functions,
automated test suites, and usage examples --- is available at:
\begin{center}
  \texttt{https://github.com/mhyman/arriola2026sa}
\end{center}
Each function in the repository corresponds directly to a concept
or algorithm in the text. \texttt{sir\_nominal.m} defines the SIR
parameters used in every chapter; \texttt{sensitivity\_jacobian.m}
implements the centered-difference Jacobian from Chapter~3;
\texttt{sobol\_jansen.m} implements the Jansen estimators from
Chapter~9. The \texttt{src/} directory is organized to mirror
the book's structure, so students can find the code for any
chapter without searching.

\noindent\textbf{Notation.} Throughout the book, the
sensitivity index of a \QOI~$q$ with respect to a
\POI~$p$ is written $\SI{p}{q}$.
\POIs are always chosen to be interpretable in everyday
terms: durations, probabilities, counts. The relationship
to the mathematical parameters in the model equations
is made explicit wherever the two differ.

\noindent\textbf{Scope note: scalar QOIs.}
Every sensitivity index in this book is computed for a
\emph{scalar} QOI --- a single number such as the peak infection
count, the total disease burden, or the value of $R_0$.
Real applications sometimes require sensitivity of an entire output
\emph{trajectory}: how does the full time series $I(t)$ change when
$\beta$ changes?
Extending SA to functional (trajectory-valued) QOIs requires either
reducing the trajectory to a scalar (which loses information) or
computing a time-dependent SI at each time point (as in Chapter~4's
Figure~\ref{fig:sir_timevarySI}).
A rigorous treatment of functional sensitivity, using the adjoint ODE
to compute sensitivity with respect to the entire output path, is
developed in Smith~\cite{smith2014uncertainty} \S9.4 and is a natural
next step after mastering the scalar case.

%%------------------------------------------------------------
\section{Exercises}
%%------------------------------------------------------------

Each exercise carries a Bloom's Taxonomy tag indicating cognitive level:
\textit{L2} (understand), \textit{L3} (apply), \textit{L4} (analyze),
\textit{L5} (evaluate), \textit{L6} (create and synthesize).
Higher-level exercises require constructing an argument or designing a procedure,
not just applying a formula.

\begin{exercise}[\bloom{L2}]
\label{ex:ch1:vocabulary}
Explain in your own words the difference between a \POI and a \QOI.
For a model describing the spread of antibiotic-resistant bacteria
in a hospital, give two examples of each. Then explain, in one sentence,
what sensitivity analysis would tell you about this model that
simply running the model would not.
\end{exercise}

\begin{exercise}[\bloom{L2}]
\label{ex:ch1:forwardadjoint}
What is the difference between forward and adjoint sensitivity analysis?
Describe each in terms of the ``What if?'' and ``How come?'' framing.
For each of the following questions, classify it as a \emph{forward question}
or an \emph{adjoint question} based on the type of question being asked
(not on which algorithm would be most efficient to use computationally ---
efficiency is the subject of Chapters~5 and~6).
\begin{enumerate}[label=(\alph*)]
  \item If the mean latent period increases by $10\%$, how does the peak
        infection count change?
  \item The peak infection count is higher than expected. Which parameter
        uncertainties are most likely to have caused this?
  \item Among all parameters in the model, which one has the largest
        effect on the basic reproduction number?
\end{enumerate}
\end{exercise}

\begin{exercise}[\bloom{L4}]
\label{ex:ch1:methodchoice}
For each of the following modeling scenarios, state whether forward SA,
adjoint SA, or global SA is most appropriate. Justify each choice in
one or two sentences.
\begin{enumerate}[label=(\alph*)]
  \item A climate model has 500 tunable parameters.
        The investigator wants to know which 10 parameters most
        affect the mean surface temperature over a 100-year run.
        Each model run takes 8 hours.
  \item A pollution transport model has 6 input parameters.
        The regulator wants to know: if the pollution at a specific
        monitoring station exceeds the legal threshold,
        which input assumptions most likely caused this?
  \item A pharmacokinetic model has 4 parameters, each of which
        is known only to within a factor of 2. The investigator
        wants to know the range of predicted drug concentrations
        consistent with this uncertainty.
\end{enumerate}
\end{exercise}

\begin{exercise}[\bloom{L5}]
\label{ex:ch1:studydesign}
You are advising a team modeling the transmission of dengue fever.
The model has 9 parameters.
The team has funding to measure exactly 3 parameters precisely in
the field and to implement exactly 2 vector-control interventions.
Write a one-page plan describing how you would use SA to guide
these decisions. Your plan should specify:
what \QOI you would focus on and why;
how you would determine which 3 parameters to measure;
what additional information (if any) you would need
beyond the SA results to recommend interventions;
and one important limitation of your approach
that the team should be aware of.
\end{exercise}

\begin{exercise}[\bloom{L6}]
\label{ex:ch1:evaluation}
A published report on a disease model states:
``Sensitivity analysis was performed on all 22 model parameters
simultaneously, demonstrating that the model is robust
to parameter uncertainty.'' The analysis was conducted by
computing local sensitivity indices at the best-fit parameter values.

Identify at least two substantive problems with this description
or its conclusions. For each problem you identify, explain
specifically what the authors should have done differently.
\end{exercise}

%%------------------------------------------------------------
%% Bibliography for standalone compilation
%%------------------------------------------------------------


\chapter{Measuring Local Sensitivity}

\begin{flushright}
\textit{A sensitivity index is a normalized derivative with a practical job.}

\medskip
\textit{Calibration tells you which parameter values fit the data.\\
Sensitivity analysis tells you which parameters matter to the answer.}
\end{flushright}

\bigskip

Two studies of the same epidemic model arrive at the following conclusions.

Study~A reports: ``Increasing the recovery rate by $1\%$ decreases $R_0$ by $1\%$.''

Study~B reports: ``Increasing the mean infectious period by $1\%$ increases $R_0$ by $1\%$.''

Both statements are mathematically correct. The recovery rate and
the mean infectious period are reciprocals of each other, so the
two studies analyzed the same model, used the same mathematics,
and obtained the same numerical result. Yet they convey opposite
signs and, more importantly, very different intuitions.
Study~A requires the reader to mentally invert a parameter before
the public health implication becomes clear.
Study~B is immediately interpretable: patients who are infectious
longer spread the disease more.

This chapter introduces the sensitivity index, the central
quantitative tool of this book, and develops the conventions
that make sensitivity analysis results readable rather than
merely correct.

%%------------------------------------------------------------
\section{Ratios as Precursors to Derivatives}
\label{sec:ratios}
%%------------------------------------------------------------

Comparison by ratio is one of the most fundamental quantitative
habits in applied science.
In finance, the price-to-earnings ratio tells an investor how many
dollars they pay for one dollar of profit.
In aeronautical engineering, the lift-to-drag ratio determines
whether a wing design is viable.
In experimental physics, the signal-to-noise ratio separates
measurement from artifact.
None of these ratios requires calculus. Each one gives a
dimensionless number that can be compared across different systems
and used to make decisions.

The sensitivity index is a ratio of this same type: it measures
how much the output changes, per unit change in the input,
normalized so that the result is dimensionless and comparable
across parameters with different units and scales.

To see why normalization matters, consider the price-to-earnings
ratio $P/E$. Suppose a stock trades at \$50 per share and has earned
\$2.50 over the past year. The $P/E$ ratio is $50/2.50 = 20$:
you pay \$20 for each dollar of annual earnings.
Now suppose earnings increase by $\$0.50$ to $\$3.00$.
A raw comparison of the two earnings figures gives $\$3.00 - \$2.50 = \$0.50$,
which tells you nothing useful without knowing the original value.
The normalized comparison is $(3.00 - 2.50)/2.50 = 20\%$, a
20\% increase in earnings, regardless of whether the numbers
are in dollars, euros, or any other currency.

A derivative, in the same spirit, measures the rate of change of
the output per unit change in the input. The sensitivity index
normalizes this derivative so that it answers a
specifically useful question: if the input changes by a small
percentage, by what percentage does the output change?

\begin{selfcheck}
A second stock trades at \$45 per share and has earned \$3.00 per year.
Calculate the $P/E$ ratio for each stock. Which stock is a better
value by this measure, and why?
\end{selfcheck}
\begin{scAnswer}
Stock~1: $P/E = 50/2.50 = 20$. Stock~2: $P/E = 45/3.00 = 15$.
Stock~2 has a lower $P/E$ ratio, meaning you pay less per dollar
of earnings, a better value by this measure.
The ratio provides a dimensionless comparison that is fair across
the two stocks despite their different prices and earnings.
\end{scAnswer}

%%------------------------------------------------------------
\section{Absolute, Relative, and Regularized Change}
\label{sec:change}
%%------------------------------------------------------------

Before defining the sensitivity index, we need precise language
for describing how quantities change. Three measures of change
appear throughout the book, each appropriate in different circumstances.

\subsection{Absolute Change}

Let $u$ denote the solution to a model, and suppose a perturbation
$\delta u$ is applied. The \textbf{absolute change} is simply the
difference between the perturbed and original solutions:
\begin{equation}
  \Dabs := \delta u = \hat{u} - u,
  \label{eq:abschange}
\end{equation}
where $\hat{u} = u + \delta u$ is the perturbed solution.
A 1-inch change in the length of a string has absolute change
1~inch, regardless of the original length.

\subsection{Relative Change}

The 1-inch change means something very different if the string is
6 inches long versus 1 mile long.
The \textbf{relative change}\index{relative change} accounts for this by normalizing against
the original value:
\begin{equation}
  \Drel := \frac{\delta u}{u} = \frac{\hat{u} - u}{u}.
  \label{eq:relchange}
\end{equation}
For the 6-inch string, the relative change is $1/6 \approx 16.7\%$.
For the 1-mile string ($63{,}360$ inches), it is $1/63{,}360 \approx 0.0016\%$.
The same absolute change has a very different significance depending
on context, and relative change captures this.

\subsection{Regularized Change}
\label{subsec:regularized}

Relative change has one practical problem: if $u \approx 0$,
the denominator in~\eqref{eq:relchange} is near zero and the
ratio becomes unreliable.
In dynamical models, the solution can pass through or near zero
during its evolution: the infectious compartment $I(t)$ in an epidemic
approaches zero as the outbreak resolves, for instance.
The \textbf{regularized change}\index{regularized change} handles this case by adding a
small correction to the denominator:
\begin{equation}
  \Dreg := \frac{\delta u}{1 + u}.
  \label{eq:regchange}
\end{equation}
When $u \approx 0$, this approximates absolute change: $\delta u / (1 + 0) = \delta u$.
When $u \gg 1$, it approximates relative change: $\delta u / (1 + u) \approx \delta u / u$.
The regularized change interpolates smoothly between the two regimes.

Looking more carefully at regularized change reveals a pattern that
recurs throughout applied mathematics: when a formula fails at a
special value, adding a small stabilizing term to the denominator
restores well-definedness without significantly changing the result
away from the problematic point. The same idea appears in numerical
linear algebra (Tikhonov regularization of ill-conditioned systems)
and in Bayesian statistics (pseudocounts in categorical models).
The regularized sensitivity index in Chapter~3 is one more instance.

\begin{selfcheck}
The infectious population in a disease model drops from $I = 0.04$
to $\hat{I} = 0.03$. Calculate the absolute, relative, and
regularized change. Which measure gives the most informative
description of this transition, and why?
\end{selfcheck}
\begin{scAnswer}
Absolute: $\delta I = 0.03 - 0.04 = -0.01$.
Relative: $-0.01/0.04 = -0.25$, a $25\%$ decrease.
Regularized: $-0.01/(1 + 0.04) \approx -0.0096$, close to absolute.

Since $I$ is small relative to 1, relative change gives the most
interpretable result: a $25\%$ decrease in the infectious population
is meaningful regardless of scale. For $I \ll 1$, absolute and
regularized change are both small in magnitude and harder to interpret
without the normalization that relative change provides.
\end{scAnswer}

%%------------------------------------------------------------
\section{The Sensitivity Index}
\label{sec:sensitivityindex}
%%------------------------------------------------------------

We now define the central tool of local sensitivity analysis.

\subsection{The Sign Rule}

Before any formal definition, the most important thing to know
about the sensitivity index is what its sign means.

\medskip
\noindent
\textbf{The sign rule:}\index{sign rule} A positive sensitivity index means that
the \QOI and \POI move in the same direction: increase the \POI
and the \QOI increases. A negative sensitivity index means they
move in opposite directions: increase the \POI and the \QOI
decreases. The magnitude tells you by how much: a sensitivity index
of $2$ means that a $1\%$ change in the \POI produces
approximately a $2\%$ change in the \QOI.

\medskip\noindent
\textbf{Important:} the sign belongs to the \emph{parameterization},
not to the underlying biology. The recovery rate $\gamma$ and the
mean infectious period $\tau = 1/\gamma$ describe the same biological
process, yet $\SI{\gamma}{R_0} < 0$ while $\SI{\tau}{R_0} > 0$.
When two papers report opposite signs for what appears to be the same
parameter, check whether they are using rate versus duration
parameterizations before concluding there is a contradiction.
Section~\ref{sec:interpretablePOIs} illustrates this with the SIR model.

\medskip

Students sometimes interpret a negative sensitivity index as
indicating an error in the calculation. It is not.
A negative index is often the best possible news: if the sensitivity
of $R_0$ with respect to the mean infectious period is positive, that
means longer illness leads to more transmission, and the corresponding
statement with the recovery rate would be negative: faster recovery
reduces transmission. Both facts are correct; the sign simply tells
you which direction the influence runs.
With this intuition established, the formal definition follows naturally.

\subsection{Formal Definition}

\begin{definition}[Sensitivity Index]\index{sensitivity index}
\label{def:SI}
Let $q$ denote a \QOI that depends differentiably on a \POI~$p$.
The \textbf{normalized sensitivity index}\index{sensitivity index!normalized} of $q$ with respect to $p$ is
\begin{equation}
  \SI{p}{q} := \lim_{\delta p \to 0}
               \frac{\left(\dfrac{\delta q}{q}\right)}
                    {\left(\dfrac{\delta p}{p}\right)}
             = \frac{p}{q}\,\frac{\partial q}{\partial p},
  \qquad q, p \neq 0.
  \label{eq:SI}
\end{equation}
\end{definition}

Three equivalent forms of this definition are in common use.
All three give the same number; which one is most convenient
depends on the problem at hand.

\textbf{Ratio form}\index{sensitivity index!ratio form} (Definition~\ref{def:SI}): the ratio of relative
changes in the \QOI and the \POI. This form is closest to the
P/E ratio analogy from Section~\ref{sec:ratios}.

\textbf{Derivative form}: multiply the partial derivative by the
ratio of nominal values:
\begin{equation}
  \SI{p}{q} = \frac{p}{q}\,\frac{\partial q}{\partial p}.
  \label{eq:SIderiv}
\end{equation}
This is the form used most often in computation.
The factor $p/q$ does the normalization work: it converts
$\partial q/\partial p$, which has units of $[Q]/[P]$ and depends
on the choice of units, into a dimensionless percent-for-percent ratio.
Concretely, if $\SI{p}{q} = 2$, then a $1\%$ increase in $p$
produces approximately a $2\%$ increase in $q$,
regardless of whether $p$ is measured in days, kilograms, or dollars.

\textbf{Logarithmic form}\index{sensitivity index!logarithmic form}: for positive \POIs and \QOIs,
\begin{equation}
  \SI{p}{q} = \frac{\partial \ln q}{\partial \ln p}.
  \label{eq:SIlog}
\end{equation}
This form reveals that the sensitivity index is the slope of the
log-log plot of $q$ versus $p$, a fact that becomes useful
when comparing SI values across very different scales.

The three forms are equivalent whenever $q$ and $p$ are
positive and $q$ is differentiable with respect to $p$.

\unifyingidea{1} The sensitivity index is the central object of this book.
It appears here as a ratio of relative changes. In Chapters~3 and 4
it reappears as a finite-difference approximation and as the solution
of a sensitivity equation. In Chapters~5 and 6, it appears as the
output of the adjoint calculation. In Chapter~9, it is the foundation
for Sobol variance decomposition. Every one of these is the same
normalized derivative, asked in the same way:
\emph{if the input changes by $x\%$, by what percentage does the output change?}

\medskip
\noindent\textbf{Why normalization is not optional.}\index{sensitivity index!dimensionlessness}
Without the $p/q$ factor, comparing raw derivatives across parameters
is not meaningful.
In the SIR model, $\partial R_0/\partial k$ has units of days
(since $k$ has units day$^{-1}$), while $\partial R_0/\partial \beta$
has different units entirely.
Asking which raw derivative is larger is like asking whether
1~pound is larger than 2~inches: the numbers may have any magnitude,
but the comparison is incommensurable because the units differ.
The normalized SI eliminates this problem.
Both $\SI{k}{R_0}$ and $\SI{\beta}{R_0}$ equal $+1$ --- dimensionless,
unit-free --- so ranking them is legitimate and the tornado plot
is an honest comparison.
This is why the tornado plot sorts by $|\SI{p}{q}|$ rather than
by $|\partial q/\partial p|$: only the normalized index gives
a ranking that is independent of how parameters happen to be measured.

\begin{selfcheck}
Use each of the three forms to compute the sensitivity index
of the area $A = \pi r^2$ with respect to the radius $r$.
Verify that all three give the same answer.
\end{selfcheck}
\begin{scAnswer}
\textit{Ratio form:} $\delta A / A = 2\pi r\,\delta r / (\pi r^2) = 2\delta r/r$,
so $\SI{r}{A} = (\delta A/A)/(\delta r/r) = 2$.

\textit{Derivative form:} $\partial A/\partial r = 2\pi r$, so
$\SI{r}{A} = (r/\pi r^2)(2\pi r) = 2$.

\textit{Logarithmic form:} $\ln A = \ln\pi + 2\ln r$, so
$\partial\ln A/\partial\ln r = 2$.

All three give $\SI{r}{A} = 2$: a $1\%$ increase in radius produces
a $2\%$ increase in area. The circle's area is twice as sensitive
to radius as to a parameter that appears linearly.
\end{scAnswer}

\subsection{Interpreting the Magnitude}
\label{subsec:magnitude}

The sensitivity index answers a specific quantitative question:
if the \POI changes by $x\%$, the \QOI changes by approximately $x \cdot \SI{p}{q}\%$.
For small $x$, this estimate is accurate to first order.
For larger $x$, second-order corrections become relevant, and
Chapter~3 discusses how to assess whether those corrections matter.

As a practical interpretive guide:
\begin{itemize}[nosep]
  \item $|\SI{p}{q}| > 1$: the \QOI is more sensitive than the \POI
        (the change in \QOI exceeds the change in \POI, percentage-for-percentage).
  \item $|\SI{p}{q}| = 1$: the \QOI changes in exact proportion to the \POI.
  \item $|\SI{p}{q}| < 1$: the \QOI is less sensitive than the \POI
        (the change in \QOI is smaller than the change in \POI).
  \item $|\SI{p}{q}| \approx 0$: the \QOI is essentially insensitive
        to this \POI near the nominal value.
\end{itemize}

%%------------------------------------------------------------
\section{Worked Examples}
\label{sec:workedexamples}
%%------------------------------------------------------------

The following examples build in complexity, beginning with
purely geometric functions and ending with the epidemic model
that serves as the book's primary running example.

\subsection{Volume of a Cylinder}
\label{subsec:cylinder}

The volume of a cylinder with radius $r$ and height $h$ is $V = \pi r^2 h$.
We want to know which dimension (radius or height)
has the greater effect on the volume.

Using the derivative form~\eqref{eq:SIderiv}:
\begin{equation}
  \SI{r}{V} = \frac{r}{V}\,\frac{\partial V}{\partial r}
            = \frac{r}{\pi r^2 h}\cdot 2\pi r h = 2.
  \label{eq:Vr}
\end{equation}
\begin{equation}
  \SI{h}{V} = \frac{h}{V}\,\frac{\partial V}{\partial h}
            = \frac{h}{\pi r^2 h}\cdot \pi r^2 = 1.
  \label{eq:Vh}
\end{equation}

The volume is twice as sensitive to radius as to height.
A $1\%$ increase in radius produces a $2\%$ increase in volume;
the same percentage increase in height produces only a $1\%$ increase.
This follows from the fact that $r$ appears squared in the volume formula
while $h$ appears linearly.

For a designer building a cylindrical tank with a fixed material budget,
this result has a direct implication: small manufacturing errors in the radius
matter twice as much as the same-sized errors in the height.

\begin{selfcheck}
For a cone with volume $V = \frac{1}{3}\pi r^2 h$, compute
$\SI{r}{V}$ and $\SI{h}{V}$. Compare with the cylinder results.
Explain why the answer is the same despite the different formula.
\end{selfcheck}
\begin{scAnswer}
$\SI{r}{V} = (r/V)(\partial V/\partial r) = (r/V)\cdot(2\pi r h/3) = 2$.
$\SI{h}{V} = 1$. Identical to the cylinder.

The reason: the sensitivity index depends on the \textit{exponent}
of each variable in the formula. Both $V_{\text{cone}} = \frac{1}{3}\pi r^2 h$
and $V_{\text{cylinder}} = \pi r^2 h$ have $r$ to the power $2$ and
$h$ to the power $1$. The constant $\frac{1}{3}$ cancels in the
logarithmic form $\partial\ln V/\partial\ln r = 2$. This is a
general principle: for a monomial $q = c\,p_1^{a_1}p_2^{a_2}\cdots$,
the sensitivity index with respect to $p_i$ is simply $a_i$.
\end{scAnswer}

\subsection{Arc Length of a Line}
\label{subsec:arclen}

Let $u(t) = mt + b$ be a line with slope $m$ and intercept $b$,
defined on $[0,a]$. The arc length is
\[
  J(m,b) = a\sqrt{1 + m^2}.
\]
Computing the sensitivity indices:
\begin{equation}
  \SI{m}{J} = \frac{m}{J}\,\frac{\partial J}{\partial m}
             = \frac{m}{a\sqrt{1+m^2}}\cdot\frac{am}{\sqrt{1+m^2}}
             = \frac{m^2}{1+m^2}.
\end{equation}
\begin{equation}
  \SI{b}{J} = \frac{b}{J}\,\frac{\partial J}{\partial b} = 0,
\end{equation}
since $J$ does not depend on $b$.

The arc length is completely insensitive to the intercept.
Geometrically, this is clear: shifting the line vertically does not
change its length. Analytically, $b$ does not appear in the formula
for $J$, so its derivative is zero. A zero sensitivity index is
not an error; it tells us that one \POI has no local effect
on the \QOI, which is itself an informative finding.

\begin{selfcheck}
For the line $u(t) = 3t + 2$ on $[0,4]$, compute $\SI{m}{J}$
numerically. Verify: if $m$ increases from $3$ to $3.03$ (a $1\%$ increase),
does $J$ change by approximately $\SI{m}{J}\%$?
\end{selfcheck}
\begin{scAnswer}
$\SI{m}{J} = m^2/(1+m^2) = 9/(1+9) = 0.9$.

At $m=3$: $J = 4\sqrt{10} \approx 12.649$.
At $m=3.03$: $J = 4\sqrt{1+3.03^2} = 4\sqrt{10.1809} \approx 12.763$.
Change: $(12.763 - 12.649)/12.649 \approx 0.90\% \approx 0.9 \times 1\%$. \checkmark
\end{scAnswer}

\subsection{Area Under a Line}
\label{subsec:area}

For the same line $u(t) = mt + b$ on $[0,a]$, the area is
\[
  J(m,b) = \int_0^a (mt + b)\,dt = \tfrac{1}{2}ma^2 + ba.
\]
Both parameters now affect the area:
\begin{equation}
  \SI{m}{J} = \frac{m}{\frac{1}{2}ma^2 + ba}\cdot\tfrac{1}{2}a^2
             = \frac{am}{am + 2b},
\end{equation}
\begin{equation}
  \SI{b}{J} = \frac{b}{\frac{1}{2}ma^2 + ba}\cdot a
             = \frac{2b}{am + 2b}.
\end{equation}
Note that $\SI{m}{J} + \SI{b}{J} = 1$, so the two indices
share a unit total for this particular linear functional.
This is not a general property: sensitivity indices do not
sum to any fixed value for arbitrary models or QOIs. Which parameter matters more depends on the
values: $\SI{m}{J} > \SI{b}{J}$ when $am > 2b$, and the
inequality reverses when $am < 2b$.

This example illustrates a general phenomenon: when two \POIs
both influence the same \QOI, their relative importance
depends on the nominal parameter values, not just on the
model structure. SA must always be reported at specific
nominal values.

\subsection{Root of a Nonlinear Equation}
\label{subsec:nonlinear}

Consider the nonlinear equation
\[
  f(u;\alpha) := \sin(\alpha u) - u = 0.
\]
For $\alpha = 2$, the equation has three roots;
for $\alpha > 1$, there is at least one nonzero root.
We want to know how the positive root $u^*(\alpha)$ changes
with $\alpha$.

Differentiating $f(u;\alpha) = 0$ implicitly with respect to $\alpha$:
\[
  u\cos(\alpha u) + \alpha\cos(\alpha u)\,\frac{\partial u}{\partial\alpha}
  - \frac{\partial u}{\partial\alpha} = 0.
\]
Solving for $\partial u/\partial\alpha$:
\[
  \frac{\partial u}{\partial\alpha}
  = \frac{u\cos(\alpha u)}{1 - \alpha\cos(\alpha u)}.
\]
At $\alpha = 2$, the positive root is $u^* \approx 0.9475$, giving
\[
  \SI{\alpha}{u^*}
  = \frac{\alpha}{u^*}\cdot\frac{u^*\cos(\alpha u^*)}{1 - \alpha\cos(\alpha u^*)}
  \approx -0.39.
\]
A $1\%$ increase in $\alpha$ decreases the positive root by approximately $0.39\%$.
This result requires no numerical experiment once the derivative is known;
the sensitivity index provides the estimate directly.

%%------------------------------------------------------------
\section{The Sensitivity Index for the SIR Model}
\label{sec:sirSI}
%%------------------------------------------------------------

The SIR epidemic model~\cite{kermack1927contribution,anderson1991infectious} describes the spread of an infectious disease
through a population divided into susceptible ($S$), infectious ($I$),
and recovered ($R$) compartments.
The basic reproduction number\index{basic reproduction number ($R_0$)}
\[
  R_0 = \frac{k\beta}{\gamma}
\]
measures the average number of new infections produced by one
infectious individual in an otherwise susceptible population,
where $k$ is the average number of contacts per day, $\beta$ is
the probability of transmission per contact, and $\gamma$ is the
recovery rate. The rigorous definition of $R_0$ via the next-generation
operator is given in van den Driessche and Watmough~\cite{vandendriessche2002reproduction}
and Diekmann et al.~\cite{diekmann1990definition}.

This is the book's primary running example.
In Chapter~4, we will compute time-dependent sensitivity indices
for the full SIR model. Here we compute the static sensitivity
indices for $R_0$: a closed-form calculation that introduces
all the key ideas in a familiar setting.

\subsection{Choosing Interpretable POIs}
\label{subsec:interpretable}

The parameters $k$, $\beta$, and $\gamma$ appear directly in the
model equations, which makes them a natural choice for \POIs.
However, we argued in Chapter~1 that \POIs should be chosen
to be interpretable in everyday terms.

The recovery rate $\gamma$ (units: day$^{-1}$) is not directly
interpretable. What a clinician understands and can measure is
the \emph{mean infectious period} $\tau = 1/\gamma$ (units: days):
the average number of days a patient remains infectious.
These two parameterizations are linked by:
\begin{equation}
  \SI{\tau}{q} = -\SI{\gamma}{q}
  \qquad\text{for any \QOI }q,
  \label{eq:reciprocalSI}
\end{equation}
since $\partial q/\partial\tau = (\partial q/\partial\gamma)
(\partial\gamma/\partial\tau) = -\gamma^2(\partial q/\partial\gamma)$,
giving $\SI{\tau}{q} = (\tau/q)(\partial q/\partial\tau)
= -(\gamma/q)\gamma(\partial q/\partial\gamma) = -\SI{\gamma}{q}$.

The magnitude is unchanged; the sign flips.
Table~\ref{tab:sirreparameterize} shows the effect of this
reparameterization for $R_0$.

\begin{table}[htb]
  \centering
  \begin{tabular}{llcc}
    \toprule
    \POI & Interpretation & Form & $\SI{p}{R_0}$ \\
    \midrule
    Recovery rate $\gamma$ (day$^{-1}$) & --- & rate & $-1$ \\
    Mean infectious period $\tau = 1/\gamma$ (days) & ``average days infectious'' & duration & $+1$ \\
    \midrule
    Contact rate $k$ (day$^{-1}$) & ``contacts per day'' & rate & $+1$ \\
    Transmission probability $\beta$ (dimensionless) & ``chance of infection per contact'' & probability & $+1$ \\
    \bottomrule
  \end{tabular}
  \caption{Sensitivity indices of $R_0$ with respect to rate and
           duration parameterizations of the SIR model.
           The mean infectious period $\tau = 1/\gamma$ has the same
           magnitude as $\gamma$ but opposite sign. The duration
           form is easier to interpret: a longer infectious period
           increases transmission, and both the sign and the
           plain-language statement agree.}
  \label{tab:sirreparameterize}
\end{table}

\subsection{Analytic Computation}
\label{subsec:R0SI}

With the interpretable \POIs $\{\tau = 1/\gamma,\, k,\, \beta\}$,
the sensitivity indices are:
\begin{align}
  \SI{k}{R_0} &= \frac{k}{R_0}\cdot\frac{\partial R_0}{\partial k}
                = \frac{k}{k\beta/\gamma}\cdot\frac{\beta}{\gamma} = 1,
                \label{eq:R0k} \\[4pt]
  \SI{\beta}{R_0} &= \frac{\beta}{R_0}\cdot\frac{\partial R_0}{\partial\beta}
                    = \frac{\beta}{k\beta/\gamma}\cdot\frac{k}{\gamma} = 1,
                    \label{eq:R0beta} \\[4pt]
  \SI{\tau}{R_0} &= -\SI{\gamma}{R_0}
                   = -\frac{\gamma}{R_0}\cdot\frac{\partial R_0}{\partial\gamma}
                   = -\frac{\gamma}{k\beta/\gamma}\cdot\left(\frac{-k\beta}{\gamma^2}\right) = +1.
                    \label{eq:R0tau}
\end{align}

All three \POIs have sensitivity index $+1$ with the interpretable
parameterization. In plain language: a $1\%$ increase in contact rate,
transmission probability, \emph{or} mean infectious period each
increases $R_0$ by $1\%$. The three parameters are equally important
for $R_0$, and all three promote disease spread when increased,
as every positive sign and every positive biological interpretation
confirms.

\begin{selfcheck}
An intervention reduces the contact rate $k$ by $20\%$.
Using the sensitivity indices above, estimate the change in $R_0$.
If a second intervention reduces the mean infectious period $\tau$
by $15\%$, which intervention has the greater effect on $R_0$?
\end{selfcheck}
\begin{scAnswer}
First intervention: $\Delta R_0/R_0 \approx \SI{k}{R_0}\times(-20\%) = -20\%$.
Second intervention: $\Delta R_0/R_0 \approx \SI{\tau}{R_0}\times(-15\%) = -15\%$.
The first intervention has the greater effect. Both indices equal $1$,
so the percentage change in $R_0$ equals the percentage change in the \POI.
The intervention that changes its \POI by the larger percentage
has the larger effect.
\end{scAnswer}

%%------------------------------------------------------------
\section{Choosing Interpretable POIs}
\label{sec:interpretablePOIs}
%%------------------------------------------------------------

The SIR example illustrates a general principle that should guide
every SA study before any calculation begins.

\textbf{Choose \POIs that can be discussed in ordinary language,
not merely because they appear in the model equations.}

Model equations are written in the language that is mathematically
convenient, which is not always the language that is scientifically
or operationally natural. Rates (with units day$^{-1}$) are
mathematically convenient in differential equations but not naturally
interpretable; the corresponding mean times (with units days)
are easier to discuss in a clinical or policy context.

The general rule for rate parameters in compartmental models:
if a parameter $p$ is a per-capita rate (units: time$^{-1}$),
define the corresponding \POI as $\tau = 1/p$ (the mean time
spent in that state), compute $\SI{\tau}{q} = -\SI{p}{q}$,
and report the result in terms of $\tau$.

Two additional cases arise frequently:

\textbf{Dimensionless probabilities.}
Transmission probabilities like $\beta$ are already dimensionless
and directly interpretable: ``a $25\%$ chance of transmission
per contact'' requires no conversion.

\textbf{Composite parameters.}
In the SIR model, $k$ and $\beta$ always appear together as the
product $k\beta$ in the expression for the force of infection
$\lambda = k\beta I/N$. Neither can be identified separately from
incidence data alone; only the product $k\beta$ is identifiable.
In this case, $k\beta$ is a more natural single \POI than either
factor individually. A related issue arises whenever two \POIs
are not independently changeable, which leads to the topic of
correlated \POIs in Chapter~3.

A practical check: if you cannot explain what a $1\%$ increase in a
\POI means in plain language to a colleague outside mathematics,
reconsider the parameterization.

It is worth noting that even in highly cited published work, this
distinction is sometimes overlooked.
Section~4.1 of Chitnis, Hyman \& Cushing~\cite{chitnis2008determining},
a paper that has become a standard reference for SA in mathematical
epidemiology, reports sensitivity indices using rate parameters
throughout. The SI values are mathematically correct, but several
indices are negative for parameters that clearly promote disease spread,
requiring a mental inversion before the public health interpretation
is clear. The present book recommends always asking, before reporting
any SI: would a policy maker understand this sign immediately,
or does it require translation?

One further misconception deserves attention here.
A sensitivity index near zero does \emph{not} mean a parameter is
globally unimportant. It means the model is approximately linear in
that parameter near the nominal value, and the linear term happens
to be small. If the model is nonlinear in that parameter, varying it
over a wider range may reveal substantial effects that the SI misses
entirely. Chapter~9 addresses this through global sensitivity analysis.

\begin{selfcheck}
A malaria model contains the parameter $\mu_v$ (mosquito mortality
rate, units: day$^{-1}$). A colleague reports $\SI{\mu_v}{R_0} = -0.3$.
Reparameterize in terms of mean mosquito lifespan $L = 1/\mu_v$
(units: days). What is $\SI{L}{R_0}$? Write a plain-language
interpretation of each sensitivity index.
\end{selfcheck}
\begin{scAnswer}
$\SI{L}{R_0} = -\SI{\mu_v}{R_0} = +0.3$.

Rate form: ``A $10\%$ increase in mosquito mortality rate decreases
$R_0$ by $3\%$.'' This requires mental inversion: higher mortality
$\to$ shorter lifespan $\to$ less transmission.

Duration form: ``A $10\%$ increase in mean mosquito lifespan increases
$R_0$ by $3\%$.'' Direct: longer-lived mosquitoes transmit more malaria.
\end{scAnswer}

%%------------------------------------------------------------
\section{The Tornado Plot}
\label{sec:tornadoplot}
%%------------------------------------------------------------

When a model has several \POIs, the collection of sensitivity indices
becomes a list of numbers with signs and magnitudes. Reading such a
list is manageable for three or four parameters; for ten or twenty,
a visual display is essential.

The \textbf{tornado plot}\index{tornado plot} is the standard display format for
sensitivity indices. It is a horizontal bar chart with the following
conventions:
\begin{itemize}[nosep]
  \item Bars extending to the \emph{right} indicate positive sensitivity
        indices (\POI and \QOI move in the same direction).
  \item Bars extending to the \emph{left} indicate negative sensitivity
        indices (opposite directions).
  \item Bars are sorted from largest to smallest by absolute magnitude,
        with the most influential parameter at the top.
  \item The length of each bar represents the magnitude of the index.
\end{itemize}
The sorted layout creates the ``tornado'' shape: wider at the top,
narrower toward the bottom.

Figure~\ref{fig:tornado_R0} shows the tornado plot for $R_0$
in the SIR model with \POIs $\{k, \beta, \tau\}$.
All three bars extend to the right (all indices are $+1$),
indicating that each parameter promotes disease spread when increased.
All bars are the same length, confirming that the three parameters
are equally important for $R_0$.

\begin{figure}[htb]
  \centering
  \begin{tikzpicture}
    \begin{axis}[
      xbar,
      bar width=14pt,
      width=0.65\textwidth,
      height=6.5cm,
      xmin=-1.6, xmax=1.6,
      xlabel={Sensitivity Index $\SI{p}{R_0}$},
      ytick={1,2,3},
      yticklabels={
        {Mean infectious period $\tau$},
        {Transmission prob.\ $\beta$},
        {Contact rate $k$}
      },
      yticklabel style={align=right, text width=5.5cm, font=\small},
      xticklabel style={font=\small},
      xtick={-1.5,-1,-0.5,0,0.5,1,1.5},
      extra x ticks={0},
      extra x tick style={grid=major, grid style={black, thick}},
      tick align=outside,
      axis x line=bottom,
      axis y line=left,
      enlarge y limits=0.4,
    ]
    \addplot[fill=black!55, draw=black] coordinates {
      (1,3)   %% k
      (1,2)   %% beta
      (1,1)   %% tau
    };
    \end{axis}
  \end{tikzpicture}
  \caption{Tornado plot for the sensitivity of $R_0$ to the three
           \POIs of the SIR model, parameterized using interpretable
           durations and probabilities. All three indices equal $+1$,
           so the bars are equal in length and all extend to the right.
           In models with more parameters and a wider range of index
           values, the sorted layout reveals immediately which
           parameters dominate.}
  \label{fig:tornado_R0}
\end{figure}

MATLAB code to produce a tornado plot for any set of sensitivity indices:

\begin{verbatim}
function tornado_plot(SI_values, labels, QOI_name)
% SI_values : vector of sensitivity indices
% labels    : cell array of POI names (strings)
% QOI_name  : string label for the QOI

[~, idx] = sort(abs(SI_values), 'descend');
SI_sorted  = SI_values(idx);
lab_sorted = labels(idx);

barh(SI_sorted, 'FaceColor', [0.35 0.35 0.35]);
set(gca, 'YTick', 1:numel(SI_sorted), ...
         'YTickLabel', lab_sorted, ...
         'FontSize', 11);
xline(0, 'k', 'LineWidth', 1.2);
xlabel(['Sensitivity Index  S_p^{' QOI_name '}'], 'FontSize', 12);
title(['Tornado Plot: Sensitivity of ' QOI_name], 'FontSize', 13);
grid on; box off;
end
\end{verbatim}

\begin{selfcheck}
A model for antibiotic-resistance spread has sensitivity indices
$\SI{\tau_{\mathrm{inf}}}{R_0} = +1.4$,
$\SI{\beta}{R_0} = +0.9$,
$\SI{\tau_{\mathrm{immune}}}{R_0} = -0.6$,
$\SI{k}{R_0} = +0.5$.

(a) Sketch the tornado plot by hand.
(b) Which parameter is most important? Least important?
(c) If $\tau_{\mathrm{immune}}$ (mean time spent immune) increases by $5\%$,
    what happens to $R_0$?
\end{selfcheck}
\begin{scAnswer}
(a) Sorted order by $|$SI$|$: $\tau_{\mathrm{inf}}$ ($1.4$, right),
$\beta$ ($0.9$, right), $\tau_{\mathrm{immune}}$ ($0.6$, \textit{left}),
$k$ ($0.5$, right).
(b) Most important: $\tau_{\mathrm{inf}}$. Least: $k$.
(c) $R_0$ decreases by approximately $0.6 \times 5\% = 3\%$.
The negative sign and the biological interpretation agree: a longer
immune period keeps people protected longer, reducing transmission.
\end{scAnswer}

%%------------------------------------------------------------
\section{A Preview: Second-Order Terms}
\label{sec:secondorder}
%%------------------------------------------------------------

The sensitivity index is built on a first-order approximation.
For a \QOI $q$ that depends on a single \POI $p$, the Taylor expansion gives
\[
  q(p + \delta p) = q(p)
    + \underbrace{\frac{\partial q}{\partial p}\,\delta p}_{\text{first-order}}
    + \underbrace{\frac{1}{2}\frac{\partial^2 q}{\partial p^2}\,(\delta p)^2}_{\text{second-order}}
    + \cdots
\]
The sensitivity index is derived from the first-order term alone.
When the second-order term is small relative to the first-order term
(that is, when
\[
  \left|\frac{1}{2}\frac{\partial^2 q}{\partial p^2}\,\delta p\right|
  \ll \left|\frac{\partial q}{\partial p}\right|,
\]
the sensitivity index provides an accurate approximation.
When this condition fails, the linear approximation is poor
and the SI overstates or understates the true sensitivity.

Chapter~3 develops practical methods for estimating the second-order
term numerically and for assessing when it is large enough to matter.
This is the first instance of Unifying Idea~3 from Chapter~1:
local SA is a linearization, and linearizations have limits.
The second-order term is the first correction beyond those limits.

For two \POIs $p_i$ and $p_j$ that are not varied independently
(a situation addressed in Chapter~3), the cross term
\[
  \frac{\partial^2 q}{\partial p_i\,\partial p_j}\,\delta p_i\,\delta p_j
\]
also enters the expansion. This cross derivative quantifies
the interaction between the two \POIs: it measures how the
sensitivity to $p_i$ changes as $p_j$ varies. Chapter~9 shows
how Sobol interaction indices generalize this idea to the full
range of parameter values.

%%------------------------------------------------------------
\section{MATLAB Laboratory: Sensitivity Indices for the SIR Model}
\label{sec:ch2matlab}
%%------------------------------------------------------------

This laboratory implements and visualizes the sensitivity indices
developed in the chapter, using the SIR model as the working example.
Complete the five parts in order; each builds on the previous.

\textbf{Part 1: Analytic SI via the Symbolic Toolbox.}
Define $R_0 = k\beta/\gamma$ symbolically in MATLAB.
Compute $\SI{k}{R_0}$, $\SI{\beta}{R_0}$, $\SI{\gamma}{R_0}$ using
the \texttt{diff} function and the derivative form of Definition~\ref{def:SI}.
Simplify and confirm that all three equal $\pm 1$ analytically.

\medskip\noindent
\textit{Note:} The code below requires the Symbolic Math Toolbox.
If you do not have it, skip directly to Part~3 and verify the same
results numerically. Parts~3--5 do not require the Symbolic Toolbox.

\begin{verbatim}
syms k beta gamma positive
R0   = k * beta / gamma;
SI_k = simplify( (k/R0) * diff(R0, k) )
SI_b = simplify( (beta/R0) * diff(R0, beta) )
SI_g = simplify( (gamma/R0) * diff(R0, gamma) )
\end{verbatim}

\textbf{Part 2: Reparameterize.}
Replace $\gamma$ with the mean infectious period $\tau = 1/\gamma$.
Rewrite $R_0 = k\beta\tau$ symbolically.
Compute $\SI{\tau}{R_0}$ and verify that it equals $-\SI{\gamma}{R_0}$
as predicted by equation~\eqref{eq:reciprocalSI}.

\textbf{Part 3: Numerical verification.}
Use the nominal values $k = 5$, $\beta = 0.06$, $\tau = 7$ days
(so $\gamma = 1/7$). Compute $R_0$ numerically. Estimate
$\SI{\tau}{R_0}$ using a centered finite difference with
$h = 10^{-5}\tau$. Compare with the analytic value of $+1$.
How large is the relative error?

\begin{verbatim}
k = 5; beta = 0.06; tau = 7;
R0    = k * beta * tau;
h     = 1e-5 * tau;
R0_p  = k * beta * (tau + h);
R0_m  = k * beta * (tau - h);
SI_fd = (tau / R0) * (R0_p - R0_m) / (2*h);
fprintf('Analytic SI: 1.00000\n');
fprintf('FD estimate: %.5f\n', SI_fd);
\end{verbatim}

\textbf{Part 4: Tornado plots for both parameterizations.}
Using the \texttt{tornado\_plot} function from Section~\ref{sec:tornadoplot},
produce two tornado plots side by side: one with rate parameters
$\{k, \beta, \gamma\}$ and one with interpretable parameters
$\{k, \beta, \tau\}$. Observe how the sign of the $\gamma/\tau$
bar changes between the two plots and why.

\medskip\noindent
\textit{MATLAB note:} The \texttt{tornado\_plot} function must be saved
as a separate file named \texttt{tornado\_plot.m} in your working directory
(function files cannot be pasted into a script).
Part~4 also uses \texttt{subplot} and cell array syntax for labels;
the file \texttt{run\_sir\_example.m} in the companion repository
demonstrates both if these are unfamiliar.

\textbf{Part 5: Effect of nominal value.}
Compute $\SI{k}{R_0}$, $\SI{\beta}{R_0}$, $\SI{\tau}{R_0}$ at
two different nominal transmission probabilities: $\beta_1 = 0.06$
and $\beta_2 = 0.6$. Do the SI values change? What does this
tell you about the sensitivity of $R_0$ to each parameter?
(Recall that $R_0 = k\beta\tau$ is a product of monomials;
for such functions, the SI is constant regardless of the nominal value.
This is a special property of multiplicative models that does not hold
for more complex functions like those in Chapter~4.)

%%------------------------------------------------------------
\section{Exercises}
\label{sec:ch2exercises}
%%------------------------------------------------------------

\begin{exercise}[\bloom{L2}]
\label{ex:ch2:signrule}
Explain the sign rule in your own words. Without performing any
calculation, state the sign of the sensitivity index for each of
the following, and explain your reasoning:
(a) the sensitivity of lung capacity to cigarette consumption,
    treating capacity as the \QOI and cigarettes per day as the \POI;
(b) the sensitivity of average crop yield to annual rainfall, near
    a region of the parameter space where crops are mildly drought-stressed;
(c) the sensitivity of $R_0$ to the mean infectious period $\tau$.
\end{exercise}

\begin{exercise}[\bloom{L2}]
\label{ex:ch2:threeforms}
For the function $q = \sqrt{p}$ (defined for $p > 0$), compute
$\SI{p}{q}$ using all three equivalent forms of Definition~\ref{def:SI}.
Verify that all three give the same result. Interpret the answer:
if $p$ doubles, by approximately what percentage does $q$ change?
\end{exercise}

\begin{exercise}[\bloom{L3}]
\label{ex:ch2:arcparabola}
Let $u(t) = \frac{1}{2}\lambda t^2$ be a parabola on $[0, b]$.
The arc-length functional is $J = \int_0^b \sqrt{1 + (\lambda t)^2}\,dt$.
Show that the sensitivity indices are
\[
  \SI{b}{J} = \frac{2\mu}{\mu + (\arcsinh\mu)/\sqrt{1+\mu^2}},
  \qquad
  \SI{\lambda}{J} = \frac{\mu(1+\mu^2) - \sqrt{1+\mu^2}\,\arcsinh\mu}
                        {\mu(1+\mu^2) + \sqrt{1+\mu^2}\,\arcsinh\mu},
\]
where $\mu = \lambda b$. For $\mu = 1$, compute both indices numerically
and verify using the derivative form.
\end{exercise}

\begin{exercise}[\bloom{L3}]
\label{ex:ch2:quadroots}
Let $u(t) = at^2 + bt + c$ and let $J_\pm = (-b \pm \sqrt{b^2 - 4ac})/(2a)$
denote the two real roots (assume $b^2 > 4ac$).
Show that the sensitivity indices of the positive root $J_+$ with respect to
$a$, $b$, $c$ can be written
\[
  \SI{a}{J_+} = -\frac{aJ_+}{2aJ_+ + b},
  \quad
  \SI{b}{J_+} = -\frac{b}{J_+(2aJ_+ + b)},
  \quad
  \SI{c}{J_+} = -\frac{c}{J_+(2aJ_+ + b)}.
\]
For $a=1$, $b=-3$, $c=2$, compute all three indices. Interpret each sign.
\end{exercise}

\begin{exercise}[\bloom{L3}]
\label{ex:ch2:SEIRR0}
A SEIR model (Susceptible-Exposed-Infectious-Recovered) has reproduction
number
\[
  R_0 = \frac{k\beta\,\nu}{(\nu + \mu)(\gamma + \mu)},
\]
where $\nu$ is the rate of progression from exposed to infectious,
$\mu$ is the background mortality rate, and $k$, $\beta$, $\gamma$
are as in the SIR model.
Reparameterize using mean times: $\tau_E = 1/\nu$ (mean latent period),
$\tau_I = 1/\gamma$ (mean infectious period), $L = 1/\mu$ (mean lifespan).
Compute $\SI{\tau_E}{R_0}$, $\SI{\tau_I}{R_0}$, $\SI{L}{R_0}$,
$\SI{k}{R_0}$, $\SI{\beta}{R_0}$ and produce a tornado plot.
Which parameter is most important?
\end{exercise}

\begin{exercise}[\bloom{L4}]
\label{ex:ch2:reparamjudge}
A published paper reports $\SI{\gamma}{R_0} = -1.0$ for an SIR model,
where $\gamma$ is the recovery rate in units of day$^{-1}$.
A second paper on the same model reports $\SI{\tau}{R_0} = +1.0$,
where $\tau = 1/\gamma$ is the mean infectious period in days.
A policy maker reading both papers is confused: one paper shows a
negative index, the other shows a positive index for what appears to
be the same parameter.

Write a one-paragraph explanation for the policy maker that
(a) resolves the apparent contradiction,
(b) explains which parameterization is more informative for policy,
and (c) describes what each sensitivity index means in plain language.
\end{exercise}

\begin{exercise}[\bloom{L4}]
\label{ex:ch2:tornado_interp}
The following tornado plot data is given for a food-safety model with
\QOI~$J =$ mean number of illness cases per outbreak. The sensitivity
indices are:
$\SI{\tau_c}{J} = +2.1$ (mean contamination duration, days),
$\SI{\beta_f}{J} = +1.8$ (food-to-person transmission probability),
$\SI{\tau_s}{J} = -1.3$ (mean hand-washing interval, days),
$\SI{k_f}{J} = +0.7$ (food contacts per day),
$\SI{\mu}{J} = -0.2$ (background clearance rate).
\begin{enumerate}[label=(\alph*)]
  \item Produce the tornado plot in MATLAB using the code template
        from Section~\ref{sec:tornadoplot}.
  \item Rank the five parameters by absolute importance.
  \item A public health program can achieve a $15\%$ reduction in either
        $\tau_c$ or a $20\%$ reduction in $k_f$. Which intervention
        reduces $J$ more? By how much?
  \item The parameter $\mu$ has a small negative index. Write a
        one-sentence interpretation of what this means for disease control.
\end{enumerate}
\end{exercise}

\begin{exercise}[\bloom{L5}]
\label{ex:ch2:studydesign}
You are advising the World Health Organization on a vaccine allocation
strategy for a new respiratory pathogen. A compartmental model has been
fitted to early outbreak data. The model has six parameters; their
sensitivity indices with respect to the \QOI ``total deaths at 6 months''
are not yet known.

Design the SA study. Your design should specify:
\begin{enumerate}[label=(\alph*)]
  \item Which parameterization you will use (rates or durations) for each
        parameter, and why.
  \item What \QOI or \QOIs you would compute sensitivity indices for,
        and what additional \QOI you might consider beyond the stated one.
  \item How you would present the results to a non-mathematical
        audience (what figure or table, and what the key message would be).
  \item One limitation of the local SA approach that the WHO advisors
        should be aware of in using these results.
\end{enumerate}
\end{exercise}

\begin{exercise}[\bloom{L6}]
\label{ex:ch2:critique}
A modeling study reports the following: ``Sensitivity analysis of
$R_0$ was performed with respect to all model parameters.
The parameters are listed in Table~3 using the standard model notation.
All sensitivity indices were positive, confirming that $R_0$ is
an increasing function of all model parameters.''

Table~3 lists recovery rate $\gamma$, mortality rate $\mu$, and
waning immunity rate $\rho$ alongside transmission and contact parameters.

Identify the error in the study's statement and explain specifically
what is wrong, referring to the mathematical relationship between
rates and their reciprocals. Restate the conclusion correctly.
\end{exercise}

%%------------------------------------------------------------



%% ---- Part II: Computing SA in Practice ----
\part{Computing Sensitivity in Practice}

\chapter{Computing Sensitivity in Practice}

\begin{flushright}
\textit{A derivative is useful only if you can compute it honestly.}
\end{flushright}

\bigskip

You have inherited a model from a collaborator.
It is 500 lines of MATLAB code, produces reasonable epidemic
projections, and has 11 parameters.
Your collaborator has moved to another institution.
The code runs correctly, but it uses conditional statements,
lookup tables, and an internal adaptive ODE solver;
none of which can be differentiated analytically.
Your funding agency wants to know which parameters matter
before committing to a three-year field measurement campaign.

How do you compute sensitivity indices for a model you cannot
differentiate?

The answer is finite differences, and this chapter shows exactly
how to do it, including how to avoid the mistakes that produce
wrong results.

%%------------------------------------------------------------
\section{The Sensitivity Jacobian}
\label{sec:jacobiansec}
%%------------------------------------------------------------

Chapter~2 computed sensitivity indices one parameter at a time.
For a model with $n_p$~\POIs and $n_q$~\QOIs,
the complete local SA is captured by an $n_q \times n_p$ matrix;
one entry per \POI--\QOI pair.

\begin{definition}[Sensitivity Jacobian]\index{sensitivity Jacobian}
\label{def:SJacobian}
Let $\vec{q} = (q_1, \ldots, q_{n_q})$ denote the vector of \QOIs
and $\vec{p} = (p_1, \ldots, p_{n_p})$ denote the vector of \POIs.
The \textbf{normalized sensitivity Jacobian}\index{sensitivity Jacobian!normalized} is the
$n_q \times n_p$ matrix $\mathbf{S}$ with entries
\begin{equation}
  S_{ki} = \SI{p_i}{q_k} = \frac{p_i}{q_k}\,\pd{q_k}{p_i}.
  \label{eq:SJacobian}
\end{equation}
\end{definition}

For the SIR model with \POIs $\{k, \beta, \tau, \mu\}$
(contact rate, transmission probability, mean infectious period,
mortality rate) and \QOIs $\{S(T), I(T), R(T)\}$ (the three
compartment sizes at time $T$), the sensitivity Jacobian is a
$3 \times 4$ matrix of twelve numbers.
Each number answers one specific "What if?" question:
how much does compartment $i$ change if \POI $j$ changes by $1\%$?

The tornado plots in Chapter~2 displayed single rows or columns
of this matrix. The full matrix captures all
\POI--\QOI relationships simultaneously.

Computing this matrix efficiently, in both time and accuracy;
is the central computational challenge of sensitivity analysis.

A common initial reaction is to treat the sensitivity Jacobian as
nothing more than a collection of individual SI calculations;
as if the matrix label were just a convenient way to organize numbers
that are computed independently. It is more than that.
The Jacobian has structure: its columns can be computed in parallel,
many methods share a single baseline evaluation, and the adjoint
method (Chapters~5 and 6) exploits the matrix structure to recover
all columns at once from a single backward solve. Understanding the
Jacobian as a matrix, not just a table, is what makes efficient SA possible.

\begin{selfcheck}
A climate model has $n_p = 30$ \POIs and $n_q = 8$ \QOIs.
How many individual sensitivity indices does the full SA require?
If you had to compute each one separately, how many model
evaluations would a naive brute-force approach require?
\end{selfcheck}
\begin{scAnswer}
$30 \times 8 = 240$ sensitivity indices. A naive approach, perturb
one \POI at a time, re-run the model, compute the difference;
requires $30 + 1 = 31$ model evaluations for all \POIs simultaneously
if only first-order (forward) differences are used, or
$2 \times 30 + 1 = 61$ for centered differences.
The key: you do not need one evaluation per SI entry. You need one
evaluation per \POI (plus one baseline), and then each evaluation
contributes an entire column of the Jacobian.
\end{scAnswer}

%%------------------------------------------------------------
\section{Finite Difference Formulas}
\label{sec:finitediff}
%%------------------------------------------------------------

A finite difference approximates the partial derivative
$\partial q/\partial p_i$ by computing how the \QOI changes when
$p_i$ alone is shifted by a small but finite amount $h$.
Three approximations are in common use, differing in accuracy and cost.

\subsection{Forward Difference}

The simplest approximation uses one additional model evaluation:
\begin{equation}
  \pd{q}{p_i} \approx \frac{q(p^* + h\vec{e}_i) - q(p^*)}{h},
  \label{eq:FD1}
\end{equation}
where $\vec{e}_i$ is the unit vector in the $p_i$ direction and
$p^*$ is the nominal parameter vector.
This is a first-order approximation: the error is $O(h)$,
meaning that halving $h$ approximately halves the error.

\subsection{Centered Difference}

Using one evaluation on each side of the nominal value gives
a second-order approximation:
\begin{equation}
  \pd{q}{p_i} \approx \frac{q(p^* + h\vec{e}_i) - q(p^* - h\vec{e}_i)}{2h}.
  \label{eq:FD2}
\end{equation}
The error is $O(h^2)$: halving $h$ reduces the error by a factor of four.
For the same step size, centered differences are substantially more
accurate than forward differences.

\subsection{Fourth-Order Formula (Rarely Used in Practice)}

When exceptionally high accuracy is required and model evaluations
are inexpensive, one can eliminate the leading-order error term
by combining four evaluations at two step sizes:
\begin{equation}
  \pd{q}{p_i} \approx
    \frac{-q(p^*+2h\vec{e}_i) + 8q(p^*+h\vec{e}_i)
          - 8q(p^*-h\vec{e}_i) + q(p^*-2h\vec{e}_i)}{12h}.
  \label{eq:FD4}
\end{equation}
The error is $O(h^4)$: halving $h$ reduces the error by a factor of
sixteen. This requires four evaluations per \POI rather than two.
In the applied SA literature this formula is uncommon; the
centered difference~\eqref{eq:FD2} is the standard workhorse.
Knowing which formula to use and what step size to choose are the
practical skills taken up next.

\subsection{Practical Workflow: Step-Size Selection by Error Estimation}

In practice, the centered difference~\eqref{eq:FD2} is the
default choice. The forward difference~\eqref{eq:FD1} plays a
secondary role not as a first approximation in its own right,
but as a cheap error estimator: because it requires only the
already-computed baseline $q(p^*)$ and the one-sided perturbed
evaluation $q(p^*+h\vec{e}_i)$, no extra model runs are needed.
The difference between the two approximations at the same $h$,
\begin{equation}
  \text{estimated error} \approx
  \left|
    \frac{q(p^*+h\vec{e}_i) - q(p^*)}{h}
    -
    \frac{q(p^*+h\vec{e}_i) - q(p^*-h\vec{e}_i)}{2h}
  \right|,
  \label{eq:FDerror}
\end{equation}
provides a practical error indicator. If this indicator is large
relative to the desired precision, the analyst reduces $h$ (or
increases ODE solver tolerances) and repeats until the error is
acceptable. This adaptive strategy is more reliable in practice
than deriving a theoretical optimal $h$ from
equations~\eqref{eq:hopt_fwd}--\eqref{eq:hopt_cen}, because
those formulas depend on unknown higher-order derivatives of the
model output.

A complementary check uses the backward difference
$[q(p^*) - q(p^*-h\vec{e}_i)]/h$ as a second one-sided estimate.
Comparing both one-sided approximations with the centered result
gives two independent error indicators, and agreement among all
three signals that $h$ lies in the reliable range.

\subsection{Cost Comparison}

Table~\ref{tab:FDcost} summarizes the cost of computing the
complete sensitivity Jacobian by each method.

\begin{table}[htb]
\centering
\begin{tabular}{lccl}
\toprule
Method & Model evaluations & Error order & Best when \\
\midrule
Forward difference   & $n_p + 1$       & $O(h)$   & Error check only; not standalone \\
Centered difference  & $2n_p + 1$      & $O(h^2)$ & Standard choice \\
Fourth-order FD      & $4n_p + 1$      & $O(h^4)$ & Rare; cheap evaluations only \\
Analytic FSE (Ch.~4) & $n_p$ lin.\ solves & Exact  & Differentiable model \\
Adjoint (Ch.~5--6)  & 2 solves        & Exact    & Many \POIs, few \QOIs \\
\bottomrule
\end{tabular}
\caption{Cost of computing the complete sensitivity Jacobian by
         different methods, for a model with $n_p$~\POIs and
         $n_q$~\QOIs. The number of model evaluations does
         \emph{not} depend on $n_q$: each evaluation contributes
         an entire column of the $n_q \times n_p$ Jacobian simultaneously.
         Adding more \QOIs is free --- the cost scales with
         the number of \POIs $n_p$, not with the size of the output.
         The finite-difference counts assume a single
         baseline evaluation is reused across all \POIs.}
\label{tab:FDcost}
\end{table}

For the SIR model with $n_p = 4$ parameters and $n_q = 3$ state
variables, centered differences require
$2 \times 4 + 1 = 9$ model evaluations.
If each takes 1~second, the full SA takes 9~seconds.

For the malaria model of Chitnis, Hyman \& Cushing~\cite{chitnis2008determining}
with $n_p = 17$ parameters and $n_q = 7$ state variables,
centered differences require $2 \times 17 + 1 = 35$ evaluations.
If each takes 1~minute: approximately 35~minutes.
If each takes 1~hour: approximately 35~hours.

For a climate model with $n_p = 500$ parameters and
model runs taking 6~hours each, centered differences require
$2 \times 500 + 1 = 1001$ evaluations, at 6~hours each,
that is approximately 250~days of computation.
This is why Chapter~6 introduces the adjoint method, which
reduces the cost to 2 solves regardless of $n_p$.

\unifyingidea{1} Every entry of the sensitivity Jacobian is a
normalized derivative: an instance of Unifying Idea~1 from
Chapter~1. The finite-difference formulas in this chapter
compute those derivatives numerically, for any model that
can be evaluated, whether or not it can be differentiated
analytically.

\begin{selfcheck}
For a model with $n_p = 20$ \POIs and $n_q = 5$ \QOIs,
how many model evaluations does the centered-difference Jacobian require?
(Remember: adding more \QOIs is free --- each evaluation contributes
an entire column of the Jacobian.)
If each evaluation takes 45~seconds, how long does the complete SA take?
At what value of $n_p$ would centered differences take more than 24~hours,
assuming each evaluation takes 2~minutes?
\end{selfcheck}
\begin{scAnswer}
$2 \times 20 + 1 = 41$ evaluations, regardless of $n_q = 5$;
$41 \times 45 = 1{,}845$~seconds $\approx 31$~minutes.

For 24~hours $= 1{,}440$~minutes with 2-minute evaluations:
$2n_p + 1 \leq 720$ gives $n_p \leq 359$, so for $n_p \geq 360$,
centered differences take more than 24~hours.
\end{scAnswer}

%%------------------------------------------------------------
\section{Step Size Selection --- The U-Curve}
\label{sec:ucurve}
%%------------------------------------------------------------

The step size $h$ controls a fundamental accuracy tradeoff.
There is an optimal value: too large and the approximation
misses the local curvature; too small and the computer's
limited numerical precision destroys the result.

\subsection{Two Sources of Error}

\textbf{Truncation error}\index{truncation error} arises because finite differences
are Taylor-series approximations. For the centered difference:
\begin{equation}
  \frac{q(p+h) - q(p-h)}{2h}
  = q'(p) + \underbrace{\frac{h^2}{6}\,q'''(p) + \cdots}_{\text{truncation error}}.
  \label{eq:FD2truncation}
\end{equation}
This error decreases as $h \to 0$, with slope $-2$ on a log-log
plot of error vs.\ $h$.

\textbf{Rounding error}\index{rounding error} arises because the computed model output
$\tilde{q}(p+h)$ differs from the true output $q(p+h)$ by a
rounding error of order $\varepsilon_{\text{mach}}|q|$,
where $\varepsilon_{\text{mach}} \approx 2.2 \times 10^{-16}$
in IEEE~754 double precision.
When $h$ is small, subtracting $\tilde{q}(p-h)$ from $\tilde{q}(p+h)$
amplifies this rounding error:
\begin{equation}
  \text{rounding error in FD} \approx
  \frac{2\varepsilon_{\text{mach}}|q|}{2h} = \frac{\varepsilon_{\text{mach}}|q|}{h}.
  \label{eq:FD2rounding}
\end{equation}
This error \emph{increases} as $h \to 0$, with slope $+1$ on a
log-log plot.

\subsection{The U-Curve and the Optimal Step Size\index{U-curve}}

\begin{flushright}
\textit{\small Before improving the model, find out which knob actually moves the machine.}
\end{flushright}

The total error in the centered-difference approximation is
the sum of truncation and rounding errors.
As $h$ decreases from a large value, truncation error falls
while rounding error rises. The minimum total error occurs
at the crossover, as shown schematically in
Figure~\ref{fig:ucurve}.

\begin{figure}[htb]
\centering
\begin{tikzpicture}
\begin{loglogaxis}[
  width=0.68\textwidth,
  height=6.5cm,
  xlabel={Step size $h$},
  ylabel={Absolute error in $\partial q/\partial p$},
  xmin=1e-14, xmax=1e-1,
  ymin=1e-14, ymax=1,
  xtick={1e-14,1e-12,1e-10,1e-8,1e-6,1e-4,1e-2},
  xticklabel style={rotate=45, anchor=east, font=\small},
  ytick={1e-14,1e-10,1e-6,1e-2},
  yticklabel style={font=\small},
  grid=major,
  grid style={gray!25},
  legend pos=outer north east,
  legend style={font=\small},
]
%% Truncation error: slope -2 (centered diff)
\addplot[thick, black, domain=1e-14:1e-1, samples=100]
  {max(x^2 * 1e4, 2.2e-16 / x)};
  \addlegendentry{Total error}

%% Truncation only
\addplot[thick, dashed, black!50, domain=1e-9:1e-1, samples=50]
  {x^2 * 1e4};
  \addlegendentry{Truncation ($\propto h^2$)}

%% Rounding only
\addplot[thick, dotted, black!70, domain=1e-14:1e-5, samples=50]
  {2.2e-16 / x};
  \addlegendentry{Rounding ($\propto 1/h$)}

%% Optimal h marker
\addplot[mark=*, mark size=3pt, only marks, black]
  coordinates {(4.7e-6, 4.7e-11)};
  \node[above right, font=\small] at (axis cs:4.7e-6, 4.7e-11)
    {optimal $h^*$};
\end{loglogaxis}
\end{tikzpicture}
\caption{The U-curve for centered differences in double precision.
         Truncation error (dashed) decreases as $h$ decreases;
         rounding error (dotted) increases. The total error (solid)
         is minimized at the optimal step size $h^*$.
         Using $h$ much smaller than $h^*$: a common mistake;
         gives larger errors, not smaller ones.}
\label{fig:ucurve}
\end{figure}

Setting the two error contributions equal gives an approximate
starting point for step-size selection:
\begin{align}
  h^*_{\text{forward}}  &\approx \sqrt{\varepsilon_{\text{mach}}}\cdot|p^*|
                          \approx 1.5 \times 10^{-8}\cdot|p^*|,
  \label{eq:hopt_fwd} \\
  h^*_{\text{centered}} &\approx \varepsilon_{\text{mach}}^{1/3}\cdot|p^*|
                          \approx 6 \times 10^{-6}\cdot|p^*|.
  \label{eq:hopt_cen}
\end{align}
For a parameter with nominal value $p^* = 0.3$, this suggests
a starting centered-difference step of approximately
$h \approx 1.8 \times 10^{-6}$.

These formulas should be understood as rough guides, not precise
prescriptions. The unknown third derivative $q'''(p)$ makes
the exact truncation error inaccessible, so the
true optimal $h$ depends on the specific model being analyzed.
The practical workflow described in Section~\ref{sec:finitediff},
computing the error indicator~\eqref{eq:FDerror} and
adjusting $h$ until the one-sided and centered approximations
agree to acceptable precision, is more reliable in practice
than relying on these theoretical formulas alone.

The most common mistake in finite-difference SA is choosing $h$
that is too small. Practitioners who set $h = 10^{-12}$ or
$h = 10^{-15}$, reasoning that ``smaller is more accurate,''
obtain results dominated by catastrophic cancellation
(subtracting two nearly equal floating-point numbers, which destroys
significant digits and amplifies rounding error).
The U-curve explains why: below the optimal step size,
rounding error grows as $1/h$ while truncation error continues
to fall, so the total error increases rather than decreases.
The error indicator~\eqref{eq:FDerror} will detect this failure:
when $h$ is too small, the centered and forward approximations
both carry large rounding errors and will disagree by a large
relative amount, signaling that a larger step size is needed.

The U-curve also reveals a practical limitation for models
implemented with adaptive ODE solvers.
For a smooth analytic function the curve has a well-defined minimum,
and the error indicator~\eqref{eq:FDerror} decreases reliably as
$h$ is reduced toward the optimal step size.
For a model that calls \texttt{ode45} internally, the output is not
mathematically smooth: the solver makes discrete step-size decisions
that can cause the model output to change non-smoothly as a
parameter varies by a small amount.
In these cases the error indicator may not decrease monotonically,
and the one-sided and centered approximations may disagree by more
than machine precision even at step sizes near $10^{-6}|p^*|$.
The remedy is to tighten the ODE solver tolerances
(\texttt{RelTol} and \texttt{AbsTol} in \texttt{ode45}) before
reducing $h$ further. After tightening tolerances, applying
the adaptive workflow (compute the error indicator,
check for agreement, reduce $h$ if needed) usually
locates a reliable step size.

\begin{selfcheck}
For the SIR model with transmission probability $\beta^* = 0.06$,
the theoretical formula~\eqref{eq:hopt_cen} suggests a starting
centered-difference step of $h \approx 3.6 \times 10^{-7}$.
You then compute the centered difference and a one-sided forward
difference at this $h$ and find they agree to four significant
figures.
(a) Is this step size acceptable? What does this agreement indicate?
(b) Now suppose you repeat the check at $h = 10^{-12}$ and find the
two approximations disagree by more than 50\%. What does this signal,
and what should you do?
\end{selfcheck}
\begin{scAnswer}
(a) Agreement to four significant figures indicates that the
rounding error is small relative to the truncation error at this
$h$, and the centered approximation is reliable.
The step size is acceptable for four-digit SI accuracy.

(b) Large disagreement at $h = 10^{-12}$ signals that rounding error
dominates: both approximations suffer catastrophic cancellation.
The remedy is to increase $h$ toward the range suggested by the
theoretical formula, then recheck the error indicator.
Rounding error $\propto 1/h$: at $h = 10^{-12}$ the rounding error
is $(3.6\times10^{-7}/10^{-12}) = 3.6\times10^5$ times larger than
at $h^* \approx 3.6\times10^{-7}$. Moving to even larger $h$ would
re-introduce truncation error, so the reliable range lies in between.
\end{scAnswer}

%%------------------------------------------------------------
\section{Second-Order Terms}
\label{sec:secondordercomp}
%%------------------------------------------------------------

Chapter~2 previewed second-order terms as corrections to the first-order
SI approximation. This section shows how to compute them.

\subsection{Pure Second-Order Derivatives}

The second derivative $\partial^2 q/\partial p_i^2$ measures the
curvature of the \QOI as a function of $p_i$.
It can be estimated by reusing the three evaluations already
computed for the centered first-order difference:
\begin{equation}
  \frac{\partial^2 q}{\partial p_i^2}
  \approx \frac{q(p^*+h\vec{e}_i) - 2q(p^*) + q(p^*-h\vec{e}_i)}{h^2}.
  \label{eq:second_pure}
\end{equation}
No additional model evaluations are needed.

When the curvature is large, the first-order SI understates the
true sensitivity for larger perturbations.
Specifically, if
\[
  \left|\frac{h}{2}\,\frac{\partial^2 q/\partial p_i^2}
                   {\partial q/\partial p_i}\right| \ll 1,
\]
the linear approximation is valid for perturbations of size $h$.
If this ratio is not small, the sensitivity index provides only
a local estimate, and Chapter~9's global SA methods are needed
for a complete picture.

\unifyingidea{3} The second-order terms computed in this section
are the first quantitative correction beyond the linearization
that defines every sensitivity index in this book.
The first-order SI is valid when the model behaves approximately
linearly near the nominal parameter values.
The second-order term tells you how quickly that approximation
degrades as perturbations grow. Chapter~9 provides the full
global picture when the linearization is no longer adequate.

\subsection{Mixed Second-Order Derivatives}

The cross-derivative $\partial^2 q/\partial p_i\,\partial p_j$
measures how the sensitivity to $p_i$ changes with $p_j$.
It requires one additional model evaluation beyond those already computed:
\begin{equation}
  \frac{\partial^2 q}{\partial p_i\,\partial p_j}
  \approx \frac{q(p^*+h\vec{e}_i+h\vec{e}_j)
              - q(p^*+h\vec{e}_i)
              - q(p^*+h\vec{e}_j)
              + q(p^*)}{h^2}.
  \label{eq:second_mixed}
\end{equation}
For all $\binom{n_p}{2}$ pairs, this requires $\binom{n_p}{2}$
additional evaluations.
For $n_p = 4$, that is 6 evaluations; for $n_p = 17$
(the malaria model), it is 136.

The cross-derivative $\partial^2 q/\partial p_i\,\partial p_j$
is large when the two \POIs \emph{interact}: varying them together
is not the same as varying them separately.
Chapter~9 shows that Sobol interaction indices $S_{ij}$ are the
global, variance-based analog of these local cross-derivatives.

\begin{selfcheck}
For the SIR model with $R_0 = k\beta\tau$,
compute $\partial^2 R_0/\partial k\,\partial\beta$ analytically.
Is the cross-derivative large or small relative to the individual
first-order derivatives $\partial R_0/\partial k$ and
$\partial R_0/\partial\beta$? What does this imply about
the interaction between $k$ and $\beta$ for $R_0$?
\end{selfcheck}
\begin{scAnswer}
$\partial^2 R_0/\partial k\,\partial\beta = \tau$
(positive, since $R_0 = k\beta\tau$).

The first-order derivatives are
$\partial R_0/\partial k = \beta\tau$ and
$\partial R_0/\partial\beta = k\tau$.

The ratio $(\partial^2 R_0/\partial k\,\partial\beta) / (\partial^2 R_0/\partial k^2) = \tau / 0$
is undefined (the pure second derivative is zero for a monomial).
More informatively, the mixed partial is nonzero, meaning $k$ and $\beta$
interact in their effect on $R_0$: increasing both together increases $R_0$
more than the sum of the individual increases. This is expected because $R_0$
is multiplicative in $k$ and $\beta$.
For Sobol analysis, this interaction will produce a nonzero $S_{k\beta}$.
\end{scAnswer}

%%------------------------------------------------------------
\section{Correlated POIs --- The Most Common Mistake}
\label{sec:correlatedPOIs}
%%------------------------------------------------------------

\begin{tcolorbox}[colback=gray!8, colframe=gray!40, boxrule=0.4pt, arc=2pt]
\textit{A parameter can be uncertain without being important,
and important without being uncertain.
Correlation between parameters adds a third possibility:
a parameter can appear important in isolation
while its true effect is inseparable from another.}
\end{tcolorbox}

\medskip

Before computing any sensitivity indices, there is a question
that must be asked: can these parameters actually be varied
independently? In many models, the answer is no.

Every finite-difference SA computation implicitly assumes that
each \POI can be perturbed while all others are held exactly fixed.
If the \POIs are correlated, meaning that in the real world,
changing one tends to change another, then this assumption
is violated and the resulting sensitivity indices can be misleading.

\subsection{Diagnosing Correlation}

Correlation among \POIs arises from several sources.

\textbf{Structural correlation}\index{correlated POIs!structural}\index{structural identifiability} occurs when parameters appear
only as a product or ratio in the model equations.
In the SIR model, the force of infection is
$\lambda = k\beta I/N$.
The parameters $k$ and $\beta$ appear only as the product $k\beta$.
Observational data, incidence counts over time, can identify
$k\beta$ as a unit but cannot separate the two factors.
Computing $\SI{k}{R_0}$ and $\SI{\beta}{R_0}$ as independent indices
is mathematically valid but can be operationally misleading,
because any intervention must target the product $k\beta$
(by reducing contact rate, transmission probability, or both).
Reparameterizing to use $k\beta$ as a single \POI is more
natural in this case.

This observation is the starting point for \textbf{structural identifiability}%
\index{structural identifiability}
\label{sec:structural_identifiability}%
analysis: the question of whether the parameters of a model can, in
principle, be recovered from the observable outputs, regardless of
measurement noise or sample size.
A model is \emph{structurally identifiable} in a parameter $p$ if distinct
values of $p$ produce distinct output trajectories.
The SIR model is not structurally identifiable in $(k, \beta)$ separately,
but is identifiable in the product $k\beta$.
Sensitivity analysis reveals this: when two parameters always have the
same sensitivity index (here $\SI{k}{q} = \SI{\beta}{q}$ for any QOI $q$),
they are structurally correlated and the model cannot distinguish them.
A sensitivity index of zero for a parameter is a related signal: the
model output does not respond to that parameter at all, which may indicate
that the parameter cannot be estimated from data.
Chapter~9 returns to identifiability in the context of global SA,
where variance-based methods provide computational tools for
detecting structural correlation across the full parameter range.
A full treatment of structural and practical identifiability, including
formal algebraic methods, is beyond the scope of this book;
see Smith~\cite{smith2014uncertainty} \S9.3 and the references therein.

\textbf{Empirical correlation}\index{correlated POIs!empirical} occurs when parameter values
are estimated from data that naturally links them.
In the malaria model studied by Chitnis, Hyman \& Cushing~\cite{chitnis2008determining},
the mosquito biting rate $\sigma_v$ and the human biting saturation
$\sigma_h$ both appear in the force of infection and are often
positively correlated across field studies.
The paper introduces the composite parameter
$\zeta = \sigma_v\sigma_h/(\sigma_v N_v^* + \sigma_h N_h^*)$
precisely to address this correlation.

\textbf{Constraint correlation}\index{correlated POIs!constraint} occurs when a physical or
biological requirement links parameters.
In models where stage durations must sum to a total
lifecycle, or where transition probabilities must sum to one,
no parameter can change without forcing compensating changes
in others.

\subsection{The Impact of Ignoring Correlation}

To see concretely what goes wrong, consider a model whose \QOI is
$q = p_1 + p_2$ where the two \POIs satisfy $p_2 = 2p_1$
(a linear constraint). The standard sensitivity indices are:
\[
  \SI{p_1}{q} = \frac{p_1}{q}\cdot 1 = \frac{p_1}{3p_1} = \frac{1}{3},
  \qquad
  \SI{p_2}{q} = \frac{p_2}{q}\cdot 1 = \frac{2p_1}{3p_1} = \frac{2}{3}.
\]
These look reasonable: the second parameter is twice as important
as the first. But in reality, $p_1$ and $p_2$ cannot be varied
independently. Any change to $p_1$ forces a change to $p_2 = 2p_1$.
The total effect of a $1\%$ increase in $p_1$ is:
\[
  \Delta q = \frac{\partial q}{\partial p_1}\,\delta p_1
           + \frac{\partial q}{\partial p_2}\,\delta p_2
           = 1\cdot(0.01p_1) + 1\cdot(0.02p_1) = 0.03p_1.
\]
The normalized total-derivative SI for $p_1$ is then
$(p_1/q)\cdot(0.03p_1/0.01p_1) = (1/3)\cdot 3 = 1$, not $1/3$.
The standard index understates the true sensitivity by a factor of three.

\subsection{Four Practical Responses}

\textbf{Response 1: Diagnose first.}
Before any SA computation, examine the correlation structure
of the \POIs. Compute the Pearson correlation matrix from any
available data on parameter values across studies or populations.
Plot scatter matrices (\texttt{plotmatrix} in MATLAB).
A correlation $|\rho_{ij}| > 0.5$ warrants attention.

\textbf{Response 2: Reparameterize structurally correlated pairs.}
When two parameters appear only as a product or ratio,
define that product or ratio as the natural \POI.
For the SIR model: use $k\beta$ as a single \POI, with
the composite sensitivity index
$S_{k\beta}^{R_0} = \SI{k}{R_0} + \SI{\beta}{R_0} = 1 + 1 = 2$.
(The additive rule for sensitivity of composite parameters
is developed in the exercises.)

\textbf{Response 3: Total derivative for empirically correlated pairs.}
When empirical data establishes that a proportional change in $p_i$
induces a proportional change in $p_j$ with coefficient $c_{ij}$
(i.e., $\Delta p_j/p_j \approx c_{ij}\,\Delta p_i/p_i$),
the total derivative SI is:
\begin{equation}
  \SI{p_i}{q}\big|_{\mathrm{tot}}
  = \SI{p_i}{q}\big|_{\mathrm{std}}
  + \sum_{j \neq i} \SI{p_j}{q}\,c_{ij}.
  \label{eq:totalSI}
\end{equation}
In plain language: $c_{ij}$ is the induced proportional change in $p_j$
per unit proportional change in $p_i$. If a $1\%$ increase in $k$
tends to be accompanied by a $0.6\%$ increase in $\beta$ (from field data),
then $c_{k\beta} = 0.6$. Estimate $c_{ij}$ from regression of observed
parameter pairs, or from domain knowledge about shared drivers.
Report both the standard and total-derivative SIs together with
the correlation coefficients so readers can assess the impact.

\textit{Response 4} (Report transparently).
When correlations exist but cannot be fully quantified,
state explicitly in the SA report that the independence
assumption was made, identify the correlated pairs,
and flag which SI values are most likely affected.

\begin{selfcheck}
In the SIR model, the product $k\beta$ represents the effective
transmission rate. If a field study establishes that $k$ and $\beta$
are positively correlated with coefficient $c_{k\beta} = 0.4$
(i.e., $\Delta\beta \approx 0.4\,\Delta k$ for proportional changes),
use equation~\eqref{eq:totalSI} to compute the total-derivative SI
of $R_0$ with respect to $k$ at the nominal values.
Compare with the standard SI $\SI{k}{R_0} = 1$.
\end{selfcheck}
\begin{scAnswer}
$\SI{k}{R_0} = 1$, $\SI{\beta}{R_0} = 1$.
Here $c_{k\beta} = 0.4$ means $(\Delta\beta/\beta)/(\Delta k/k) = 0.4$,
so $p_\beta c_{k\beta}/p_k = \beta\cdot 0.4/k$.
But in proportional terms: total SI $\approx 1 + 1 \times 0.4 = 1.4$.
A $1\%$ increase in $k$ (with the correlated $0.4\%$ increase in $\beta$
it entails) increases $R_0$ by approximately $1.4\%$, not $1\%$.
The standard SI underestimates the true effect by $40\%$.
\end{scAnswer}

%%------------------------------------------------------------
\section{Black-Box Models and the MATLAB Template}
\label{sec:blackbox}
%%------------------------------------------------------------

When the model is inherited code, a compiled simulation,
or a package you did not write, analytical differentiation is
not possible. Finite differences are then the only option for
local SA, and they work for \emph{any} model that accepts a
parameter vector and returns a \QOI vector.

Students sometimes treat finite differences as an inferior
approximation to be used only when analytic derivatives are
unavailable. For black-box models, this framing is backwards:
finite differences are not a fallback; they are the correct and
complete tool. The accuracy they provide at the optimal step size
is typically more than sufficient for SA decisions, and their
applicability to models of arbitrary complexity makes them
indispensable in practice.

The following MATLAB function computes the full normalized
sensitivity Jacobian by centered differences, using the optimal
step size from~\eqref{eq:hopt_cen}.
Students should copy this template for every black-box SA problem.

\begin{small}
\begin{verbatim}
%% From: src/core/sensitivity_jacobian.m (see GitHub repository)
p     = sir_nominal();
R0_fn = @(pv) pv(1)*pv(2)*pv(3);   %% R0 = k*beta*tau
p_vec = [p.k; p.beta; p.tau; p.L];
[S, info] = sensitivity_jacobian(R0_fn, p_vec, 1);
%% S(j) = SI_{p_j}^{R0};  info.n_evals = 2*n_p+1 = 9

%% If n_q is wrong, the function catches it immediately:
%% >> sensitivity_jacobian(R0_fn, p_vec, 3)
%% Error: MODEL returned 1 value but N_Q = 3.
\end{verbatim}
\end{small}

The function handles the near-zero \QOI case using the regularized
change from Chapter~2. The \texttt{max} in the step-size formula
prevents $h$ from being identically zero when $p_{\text{nom},i} = 0$.

To validate this template on any model, call it with the analytic
sensitivity indices known from Chapter~2 and verify agreement.
For the SIR model:

\begin{verbatim}
%% SIR R0 as function of p = [k, beta, tau]
R0_model = @(p) p(1)*p(2)*p(3);    %% R0 = k * beta * tau

p_nom = [5, 0.06, 7];     %% nominal: k, beta, tau
S = sensitivity_jacobian(R0_model, p_nom, 1);

fprintf('Sensitivity indices for R0:\n')
fprintf('  SI_k    = %6.4f  (exact: 1.0000)\n', S(1))
fprintf('  SI_beta = %6.4f  (exact: 1.0000)\n', S(2))
fprintf('  SI_tau  = %6.4f  (exact: 1.0000)\n', S(3))
\end{verbatim}

\begin{selfcheck}
The template above uses $h_{\min} = 10^{-10}$ as a floor for the
step size. Why is this floor necessary, and what happens if a
\POI has nominal value exactly zero?
\end{selfcheck}
\begin{scAnswer}
If $p_{\text{nom},i} = 0$, then $6 \times 10^{-6}|p_{\text{nom},i}| = 0$,
giving $h = 0$ and a division-by-zero error.
The floor $h_{\min} = 10^{-10}$ prevents this.
For a \POI with nominal value zero, the step size is not scaled
by the parameter magnitude, it is just the fixed small value $10^{-10}$.
In this case, the SI is also not well-defined in the ratio form
(it would require dividing by $p_i = 0$), so the code computes
the unnormalized derivative instead.
If a \POI is exactly zero at its nominal value, the
investigator should either reparameterize or report the
absolute (unnormalized) derivative rather than the normalized SI.
\end{scAnswer}

%%------------------------------------------------------------
\section[MATLAB Laboratory]{MATLAB Laboratory: SA of the SIR Model by Finite Differences}
\label{sec:ch3matlab}
%%------------------------------------------------------------

This laboratory puts the chapter's tools together and provides
a complete reference implementation for black-box SA.
Use the SIR ODE model with the four \POIs $\{k, \beta, \tau, \mu\}$
and the three \QOIs $\{S(T), I(T), R(T)\}$ at time $T = 60$~days.

\textbf{Setup.}
Implement \texttt{SIR\_model(p)} as a MATLAB function that takes
the parameter vector $p = [k, \beta, \tau, \mu]$ and returns
$[S(60), I(60), R(60)]$ by calling \texttt{ode45}.
Use nominal values $k = 5$, $\beta = 0.06$, $\tau = 7$~days,
$\mu = 0.0001$~day$^{-1}$, and initial conditions
$S_0 = 999$, $I_0 = 1$, $R_0 = 0$.

\textbf{Part 1:} Compute the full sensitivity Jacobian.
Call \texttt{sensitivity\_jacobian(SIR\_model, p\_nom, 3)}
to produce the $3 \times 4$ sensitivity matrix.
Display it as a labeled table with \QOI names as row labels
and \POI names as column labels.

\textbf{Part 2:} Plot the U-curve.
Fix all parameters at nominal values except $\beta$.
Estimate $\partial I(60)/\partial\beta$ using centered differences
for $h = 10^{-1}, 10^{-3}, 10^{-5}, 10^{-7}, 10^{-9}, 10^{-11}, 10^{-13}$.
Compute the absolute error using the fourth-order formula~\eqref{eq:FD4}
as a reference value (computed at $h = 10^{-4}$ where it is accurate).
Also plot the one-sided forward difference error at each $h$.
Plot centered-difference error, forward-difference error, and the
discrepancy between the two (the practical error indicator~\eqref{eq:FDerror})
vs.\ $h$ on log-log axes. Identify the range of $h$ where the
error indicator is small and the centered approximation is reliable.

\textbf{Part 3:} Tornado plots.
Using the Jacobian from Part~1, produce separate tornado plots
for each \QOI: the sensitivity of $S(60)$, $I(60)$, and $R(60)$
with respect to all four \POIs. Interpret: which \POI matters
most for each \QOI? Do the three plots give the same ranking?

\textbf{Part 4:} Correlation check.
The contact rate $k$ and transmission probability $\beta$
appear only as the product $k\beta$ in the SIR force of infection.
Confirm numerically that $\SI{k}{I(60)} \approx \SI{\beta}{I(60)}$
at the nominal values. Then reparameterize: define a new \POI
$\lambda_0 = k\beta$ and compute $\SI{\lambda_0}{I(60)}$.
Verify that $\SI{\lambda_0}{I(60)} = \SI{k}{I(60)} + \SI{\beta}{I(60)}$.

\textbf{Part 5:} Correlation check from literature.
The SIR force of infection is $\lambda = k\beta I/N$, so $k$ and $\beta$
appear only as the product $k\beta$ and are structurally correlated.
Use the four-parameter SIR model. Confirm numerically that
$\SI{k}{I(60)} \approx \SI{\beta}{I(60)}$ at nominal values.
Then search the epidemiological literature for any reported correlations
between contact rates $k$ and transmission probabilities $\beta$
across outbreak settings.
Construct a $4\times4$ correlation matrix using whatever estimates you find,
or assume a moderate positive correlation $\rho_{k\beta} = 0.6$
and zero for all other pairs.
Use equation~(3.6) to compute the total-derivative SI of $I(60)$
with respect to $k$ accounting for this correlation.
Compare with the standard SI and discuss the practical difference.

\textbf{Extension.}
Using equation~\eqref{eq:second_pure}, compute the pure second-order
sensitivities $\partial^2 I(60)/\partial p_i^2$ for each of the
four \POIs. For each parameter, assess whether the second-order
term is large enough to matter for $10\%$ perturbations.
(Reuse evaluations from Part~1 to keep the additional cost minimal.)

%%------------------------------------------------------------
\section{Exercises}
\label{sec:ch3exercises}
%%------------------------------------------------------------

\begin{exercise}[\bloom{L2}]
\label{ex:ch3:ucurve}
Explain in your own words why making the step size $h$ smaller
does not always improve the accuracy of a finite-difference
sensitivity estimate. Describe the two sources of error and
explain why they act in opposite directions as $h$ changes.
\end{exercise}

\begin{exercise}[\bloom{L3}]
\label{ex:ch3:cost}
A model has $n_p = 15$ \POIs and $n_q = 4$ \QOIs.
\begin{enumerate}[label=(\alph*)]
  \item Fill in Table~\ref{tab:FDcost} for this model:
        how many model evaluations does each method require?
  \item If each evaluation takes 30~seconds,
        how long does the centered-difference Jacobian take?
  \item At what point (in terms of $n_p$ and evaluation time)
        would the adjoint method from Chapter~6 be strictly necessary?
\end{enumerate}
\end{exercise}

\begin{exercise}[\bloom{L3}]
\label{ex:ch3:ucurve_matlab}
For the function $q(p) = \sin(p^2)$ at $p^* = 1.3$,
compute $\partial q/\partial p$ using centered differences
for $h \in \{10^{-1}, 10^{-2}, \ldots, 10^{-14}\}$.
Compare each estimate to the exact value $q'(p) = 2p\cos(p^2)$.
Plot error vs.\ $h$ on log-log axes and identify the optimal $h$.
Does it agree with the formula in equation~\eqref{eq:hopt_cen}?
\end{exercise}

\begin{exercise}[\bloom{L3}]
\label{ex:ch3:second}
For the SIR model $R_0 = k\beta\tau$:
\begin{enumerate}[label=(\alph*)]
  \item Compute the pure second derivatives $\partial^2 R_0/\partial k^2$,
        $\partial^2 R_0/\partial\beta^2$, $\partial^2 R_0/\partial\tau^2$
        analytically. What do you notice?
  \item Compute the cross derivatives $\partial^2 R_0/\partial k\,\partial\beta$,
        $\partial^2 R_0/\partial k\,\partial\tau$,
        $\partial^2 R_0/\partial\beta\,\partial\tau$ analytically.
  \item For a $10\%$ perturbation to $k$ and a simultaneous $10\%$ perturbation
        to $\beta$, estimate the true change in $R_0$ using both the
        first-order and the full second-order Taylor approximations.
        How large is the correction?
\end{enumerate}
\end{exercise}

\begin{exercise}[\bloom{L4}]
\label{ex:ch3:correlatedSIR}
In the SIR model, suppose field data suggests that the
contact rate $k$ and transmission probability $\beta$ are
positively correlated: across outbreak settings,
when $k$ increases by $10\%$, $\beta$ tends to increase by $6\%$.
\begin{enumerate}[label=(\alph*)]
  \item Compute the standard sensitivity indices $\SI{k}{R_0}$
        and $\SI{\beta}{R_0}$.
  \item Using equation~\eqref{eq:totalSI}, compute the total-derivative
        SI of $R_0$ with respect to $k$, accounting for the correlation.
  \item Which is larger: the standard SI or the total-derivative SI?
        What does this mean for a policy maker who can reduce
        contact rates (through social distancing) but not transmission
        probability directly?
\end{enumerate}
\end{exercise}

\begin{exercise}[\bloom{L4}]
\label{ex:ch3:blackbox}
A colleague provides you with the function \texttt{seir\_model(p)}
that implements a SEIR epidemic model with $n_p = 6$ parameters
and returns the peak infection count as a single \QOI.
Each model evaluation takes approximately 0.2~seconds.
\begin{enumerate}[label=(\alph*)]
  \item How long will it take to compute the complete sensitivity
        Jacobian using centered differences?
  \item Apply the \texttt{sensitivity\_jacobian} template to this model.
        What additional information do you need before interpreting
        the results? (Think about parameterization.)
  \item You find that two of the six parameters have
        $|\rho_{ij}| = 0.78$. What should you do before
        reporting the sensitivity indices?
\end{enumerate}
\end{exercise}

\begin{exercise}[\bloom{L5}]
\label{ex:ch3:studydesign}
A government agency funds you to perform SA on a dengue fever
transmission model with 9 parameters and 3 \QOIs
(peak infections, total deaths, days until outbreak peak).
Each model evaluation takes 4~minutes.

Design the complete computational SA study.
Your design should specify:
\begin{enumerate}[label=(\alph*)]
  \item Which finite-difference method you will use and why.
  \item The total computation time.
  \item How you will choose the step size for each parameter.
  \item What checks you will run before reporting any results.
  \item How you will present the results to a non-technical audience.
  \item One scenario in which the standard SA would give misleading
        results, and what you would do instead.
\end{enumerate}
\end{exercise}

\begin{exercise}[\bloom{L6}]
\label{ex:ch3:critique}
A published sensitivity analysis reports:
``All finite-difference computations used $h = 10^{-12}$
to ensure maximum precision. The code was run in MATLAB
using double-precision arithmetic.''

The authors report that for several parameters, the sensitivity
indices are exactly zero to all reported decimal places.

Identify what is most likely wrong with these results
and explain your reasoning quantitatively, using the
U-curve analysis from Section~\ref{sec:ucurve}.
What should the authors have done, and how would you
determine whether the zero SI values are genuine or numerical artifacts?
\end{exercise}

%%------------------------------------------------------------



%% ---- Part III: Analytic Forward SA ----
\part{Analytic Forward Sensitivity Analysis}

\chapter{Analytic Forward Sensitivity Analysis}

\begin{flushright}
\textit{Differentiate the model, and the important parameters identify themselves.}
\end{flushright}

\bigskip

Chapter~3 showed how to compute sensitivity indices for any model
by finite differences. For the malaria model with 17~\POIs and
7~state variables, the centered-difference Jacobian requires
239~model evaluations, about 4~hours if each run takes a minute,
or 10~days if each takes an hour.

There is a faster and more accurate alternative, provided the model
can be differentiated analytically: differentiate the governing
equations directly. The result is a new system of equations
the \emph{forward sensitivity equations}\index{forward sensitivity equation (FSE)}, whose solution
is the exact derivative at the cost of one linear solve per \POI.
For the malaria model: 17~linear solves of a $7\times7$ system,
each taking milliseconds.

The catch is that the model must be differentiable and the
analyst must be patient enough to derive the sensitivity equations.
This chapter develops that derivation systematically,
shows that it always produces a \emph{linear} system no matter
how nonlinear the original model is, and ends by quantifying
the remaining computational bottleneck that motivates
Chapters~5 and 6.

%%------------------------------------------------------------
\section{The FSE Derivation Pattern}
\label{sec:FSEpattern}
%%------------------------------------------------------------

The forward sensitivity equation (\FSE) for any model is obtained
by a single operation: differentiate the governing equation with
respect to a \POI and apply the chain rule.
The result is always a linear equation in the sensitivity variables,
regardless of whether the original model is linear or nonlinear.

\subsection{The Pattern in Three Steps}

\textbf{Step 1.} Write the governing equation in the form
$\mathcal{F}(u; p) = 0$ or $\mathcal{F}(u; p) = f$, where $u$
is the solution and $p$ is the \POI of interest.

\textbf{Step 2.} Differentiate both sides with respect to $p$.
Apply the chain rule to every term that depends on $u$.
Collect terms. The result is an equation for $\partial u/\partial p$.

\textbf{Step 3.} Recognize that the equation from Step~2 is
\emph{linear in} $\partial u/\partial p$, even if $\mathcal{F}$
is nonlinear in $u$.

The linearity in Step~3 is not a coincidence. Differentiation is
a linear operation: applying it to a nonlinear equation in $u$
produces a linear equation in $\partial u/\partial p$.
This structural fact is what makes analytic SA computationally tractable.

\medskip
\noindent\textbf{Misconception to confront.} A common reaction to
the FSE is: ``Deriving and solving these equations must be at least
as hard as solving the original model.'' This is wrong in two ways.
First, the FSE is always \emph{linear} in the sensitivity variable,
even when the original model is highly nonlinear; a linear system
is categorically easier to solve than a nonlinear one.
Second, the FSE for an ODE uses the \emph{same} coefficient matrix
$F_u$ that the original solver already computes, so the additional
cost per parameter is one linear solve, not one full nonlinear solve.
The FSE is not a second model; it is a linear companion to the first.
\medskip

\subsection{Application to Three Model Types}

The same three-step pattern applies to algebraic systems,
difference equations, and ordinary differential equations.
Sections~\ref{sec:algebraicFSE}--\ref{sec:ODEFSE}
work through each type in detail.
Here we display the pattern side by side.

\begin{center}
\renewcommand{\arraystretch}{1.5}
\begin{tabular}{lll}
\toprule
Model type & Governing equation & Forward sensitivity equation \\
\midrule
Algebraic & $f(u;p) = 0$ & $f_u\,\pd{u}{p} = -f_p$ \\[2pt]
Linear system & $A\vec{u} = \vec{b}$ &
  $A\,\pd{\vec{u}}{p} = \pd{\vec{b}}{p} - \pd{A}{p}\vec{u}$ \\[2pt]
ODE IVP & $\dot{u} = F(u;p)$ &
  $\dfrac{d}{dt}\!\left[\pd{u}{p}\right] = F_u\,\pd{u}{p} + F_p$ \\[6pt]
O$\Delta$E & $u_{n+1} = G(u_n;p)$ &
  $\pd{u_{n+1}}{p} = G_u\,\pd{u_n}{p} + G_p$ \\
\bottomrule
\end{tabular}
\end{center}

In each case, the coefficient multiplying the unknown sensitivity
($f_u$, $A$, $F_u$, $G_u$) is the Jacobian of the model
with respect to the solution $u$, evaluated at the nominal solution.
The right-hand side contains only known quantities.
The sensitivity variable ($\partial u/\partial p$) always appears
linearly on the left.

\unifyingidea{1} Every entry of the sensitivity Jacobian that
this chapter computes is a normalized derivative: an instance
of Unifying Idea~1 from Chapter~1. The FSE computes those
derivatives exactly, without the $O(h^2)$ truncation error
of the finite-difference approach.
The FSE is not a different kind of SA result; it is the same
sensitivity index, obtained by a more accurate and often faster
computational route.

\begin{selfcheck}
For the logistic ODE $\dot{u} = ru(1 - u/K)$,
write down the FSE for $\partial u/\partial r$ and for
$\partial u/\partial K$. Identify the coefficient $F_u$ in each case.
\end{selfcheck}
\begin{scAnswer}
$F(u; r, K) = ru(1-u/K)$. Computing partial derivatives:
$F_u = r(1 - 2u/K)$, $F_r = u(1-u/K)$, $F_K = ru^2/K^2$.

FSE for $w_r = \partial u/\partial r$:
$\dot{w}_r = r(1-2u/K)\,w_r + u(1-u/K)$.

FSE for $w_K = \partial u/\partial K$:
$\dot{w}_K = r(1-2u/K)\,w_K + ru^2/K^2$.

Both FSEs have the same coefficient $F_u = r(1-2u/K)$ on the left;
only the right-hand-side forcing term differs. The coefficient depends
on the forward solution $u(t)$, so the FSE must be solved
simultaneously with the original ODE.
\end{scAnswer}

%%------------------------------------------------------------
\section{Algebraic Systems and Equilibria}
\label{sec:algebraicFSE}
%%------------------------------------------------------------

For static problems, equations whose solutions do not change in
time, the FSE is a linear system that can be solved by
standard linear algebra.

\subsection{A Single Nonlinear Equation}

For the scalar equation $f(u; p) = 0$, differentiating with
respect to $p$ gives
\[
  f_u\,\pd{u}{p} + f_p = 0,
\]
so $\partial u/\partial p = -f_p/f_u$ provided $f_u \neq 0$.
This is exactly the implicit function theorem\index{implicit function theorem}.
When $f_u = 0$, the derivative is undefined, which corresponds
to a bifurcation point, a situation studied in Chapter~8.

\subsection{Linear Systems: The FSE for $A\vec{u} = \vec{b}$}

For the system $A\vec{u} = \vec{b}$, differentiating with
respect to any scalar parameter $q \in \{a_{ij}, b_i\}$ gives:
\begin{equation}
  A\,\pd{\vec{u}}{q} = \pd{\vec{b}}{q} - \pd{A}{q}\,\vec{u}.
  \label{eq:linearFSE}
\end{equation}
The right-hand side is a known vector once the forward solution
$\vec{u}$ is known.
Solving equation~\eqref{eq:linearFSE} requires one additional
linear solve with the same coefficient matrix $A$ for each parameter $q$.
Since the LU factorization of $A$ is computed only once, the cost
for $m$ parameters is one LU factorization plus $m$ back-substitutions.

\textbf{Important:} Do not solve~\eqref{eq:linearFSE} by computing $A^{-1}$.
Explicit matrix inversion is numerically unstable and computationally
wasteful. In MATLAB, use the backslash operator (\texttt{A\textbackslash b}),
which solves the system by LU factorization without forming the inverse.

\subsection{The Symbiotic Population Model}

Consider the two-species symbiotic population model~\cite{murray2007mathematical}:
\begin{align}
  \frac{du_1}{dt} &= u_1(1 - u_1 - au_2), \label{eq:symp1}\\
  \frac{du_2}{dt} &= bu_2(1 - u_2 - cu_1), \label{eq:symp2}
\end{align}
with parameters $a$, $b$, $c > 0$ satisfying $0 < a, c < 1$.
The coexistence equilibrium $(\bar{u}_1, \bar{u}_2)$ satisfies
the nonlinear system
\begin{equation}
  \bar{u}_1 = \frac{1-a}{1-ac}, \qquad \bar{u}_2 = \frac{1-c}{1-ac}.
  \label{eq:sympeq}
\end{equation}
To find how this equilibrium depends on $a$ and $c$,
we differentiate the equilibrium equations
\begin{align*}
  f_1(\bar{u}_1, \bar{u}_2; a, c) &:= \bar{u}_1(1 - \bar{u}_1 - a\bar{u}_2) = 0, \\
  f_2(\bar{u}_1, \bar{u}_2; a, c) &:= b\bar{u}_2(1 - \bar{u}_2 - c\bar{u}_1) = 0,
\end{align*}
with respect to $a$ to obtain the FSE
\begin{equation}
  D_{\vec{u}}[\vec{f}]\,\pd{\vec{u}}{a} = -\pd{\vec{f}}{a},
  \label{eq:sympFSE}
\end{equation}
where $D_{\vec{u}}[\vec{f}]$ is the Jacobian matrix of $\vec{f}$
with respect to $\vec{u}$, evaluated at the equilibrium:
\begin{equation}
  D_{\vec{u}}[\vec{f}]\big|_{\text{eq}} =
  \begin{pmatrix}
    1 - 2\bar{u}_1 - a\bar{u}_2 & -a\bar{u}_1 \\
    -bc\bar{u}_2 & b(1-2\bar{u}_2-c\bar{u}_1)
  \end{pmatrix},
  \label{eq:sympJac}
\end{equation}
and
\[
  \pd{\vec{f}}{a} = \begin{pmatrix} -\bar{u}_1\bar{u}_2 \\ 0 \end{pmatrix}.
\]
Solving the $2\times 2$ linear system~\eqref{eq:sympFSE}
gives $\partial\bar{u}_1/\partial a$ and $\partial\bar{u}_2/\partial a$
exactly. The sensitivity indices follow from Definition~2.1.

A key observation: although the equilibrium equations~\eqref{eq:symp1}--\eqref{eq:symp2}
are nonlinear, the FSE~\eqref{eq:sympFSE} is a \emph{linear system}
in the sensitivity variables.
This is always true for equilibrium problems, as the following
result shows.

\begin{theorem}[Linearity of the FSE]
\label{thm:linearFSE}
Let $\vec{f}(\vec{u}; \vec{p}) = \vec{0}$ be a system of nonlinear equations
with $\vec{f}$ differentiable in $\vec{u}$ and $\vec{p}$.
The forward sensitivity equation for $\partial\vec{u}/\partial p_j$ is
the linear system
\[
  D_{\vec{u}}[\vec{f}]\,\pd{\vec{u}}{p_j} = -\pd{\vec{f}}{p_j},
\]
provided the Jacobian $D_{\vec{u}}[\vec{f}]$ is nonsingular.
\end{theorem}

\begin{proof}
Differentiating $\vec{f}(\vec{u}(\vec{p}); \vec{p}) = \vec{0}$ with
respect to $p_j$ and applying the chain rule gives
\[
  D_{\vec{u}}[\vec{f}]\,\pd{\vec{u}}{p_j} + \pd{\vec{f}}{p_j} = \vec{0},
\]
where $D_{\vec{u}}[\vec{f}]$ denotes the Jacobian matrix of $\vec{f}$
with respect to $\vec{u}$, evaluated at the nominal solution
$(\vec{u}(\vec{p}); \vec{p})$. Rearranging yields
\[
  D_{\vec{u}}[\vec{f}]\,\pd{\vec{u}}{p_j} = -\pd{\vec{f}}{p_j}.
\]
The right-hand side $-\partial\vec{f}/\partial p_j$ is a known vector
once $\vec{u}(\vec{p})$ is available.
Linearity in $\partial\vec{u}/\partial p_j$ follows from the fact that
differentiation is itself a linear operation: the quantity
$\partial\vec{u}/\partial p_j$ appears in the left-hand side
only through the matrix-vector product $D_{\vec{u}}[\vec{f}]\,(\partial\vec{u}/\partial p_j)$.
No powers, products, or compositions of $\partial\vec{u}/\partial p_j$ appear,
regardless of the nonlinearity of $\vec{f}$ in $\vec{u}$.
Nonsingularity of $D_{\vec{u}}[\vec{f}]$ guarantees a unique solution.
\end{proof}

This theorem is one of the most practically useful facts in SA.
Regardless of how complex or nonlinear the model is, its FSE at
an equilibrium is always a linear system, and linear systems
can be solved efficiently.

\begin{selfcheck}
For the symbiotic population model with $a = 0.4$, $b = 1$, $c = 0.3$:
(a) Compute the coexistence equilibrium using equation~\eqref{eq:sympeq}.
(b) Assemble the Jacobian~\eqref{eq:sympJac} at this equilibrium.
(c) Solve the FSE for $\partial\bar{u}_1/\partial a$ and
    $\partial\bar{u}_2/\partial a$.
(d) Compare with the analytic derivatives of~\eqref{eq:sympeq}.
\end{selfcheck}
\begin{scAnswer}
(a) $\bar{u}_1 = (1-0.4)/(1-0.12) = 0.6/0.88 \approx 0.6818$;
$\bar{u}_2 = (1-0.3)/0.88 = 0.7/0.88 \approx 0.7955$.

(b) $D_u f = \begin{pmatrix} 1-2(0.6818)-0.4(0.7955) & -0.4(0.6818) \\
-1(0.3)(0.7955) & 1-2(0.7955)-0.3(0.6818)\end{pmatrix}
\approx \begin{pmatrix} -0.6818 & -0.2727 \\ -0.2386 & -0.7955 \end{pmatrix}$.

(c) RHS: $-\partial\vec{f}/\partial a = (0.6818\times0.7955, 0)^T \approx (0.5424, 0)^T$.
Solve $2\times2$ system.

(d) Analytic: $\partial\bar{u}_1/\partial a = -(1-c)/((1-ac)^2)$
$= -(0.7)/(0.88)^2 \approx -0.904$. Compare with FSE solution.
\end{scAnswer}

\subsection{The Car Rental Problem: O$\Delta$E Sensitivity}

Consider the car rental Markov chain: two airports $A$ and $B$ with
transition probabilities $\alpha = \Pr(A\to A)$ and $\beta = \Pr(B\to B)$.
The discrete-time model is
\begin{align}
  A_{n+1} &= \alpha A_n + (1-\beta)B_n, \label{eq:carA}\\
  B_{n+1} &= (1-\alpha)A_n + \beta B_n. \label{eq:carB}
\end{align}
The fixed point satisfies $(A^*, B^*)$:
\[
  A^* = \frac{1-\beta}{2-\alpha-\beta}, \qquad
  B^* = \frac{1-\alpha}{2-\alpha-\beta}.
\]
The FSE for the sensitivity of the fixed point to $\alpha$ is
exactly equation~\eqref{eq:linearFSE} with the transition matrix
playing the role of $A$. The result is:
\begin{equation}
  \SI{\alpha}{A^*} = \frac{\alpha}{2-\alpha-\beta},
  \qquad
  \SI{\alpha}{B^*} = -\frac{\alpha(1-\beta)}{(1-\alpha)(2-\alpha-\beta)}.
  \label{eq:carSI}
\end{equation}
These indices tell management which transition probability has the
greater effect on the long-run car distribution.
At nominal values $\alpha = 0.7$ and $\beta = 0.3$, the formula gives
$\SI{\alpha}{A^*} = 0.7/(2-0.7-0.3) = 0.70$: a $1\%$ increase in the
probability of keeping a car at airport~$A$ produces a $0.70\%$ increase
in the long-run fraction of cars at~$A$.

%%------------------------------------------------------------
\section{Forward Sensitivity for ODE Systems}
\label{sec:ODEFSE}
%%------------------------------------------------------------

For a dynamical model governed by an ODE, the sensitivity
$\partial u/\partial p$ changes over time as the solution evolves.
The FSE is itself an ODE that must be solved alongside the
original model.

\subsection{Derivation of the ODE FSE}

Consider the scalar, autonomous ODE with parameter $p$:
\begin{equation}
  \frac{du}{dt} = F(u; p), \qquad u(0) = u_0(p).
  \label{eq:fwdODE}
\end{equation}
Define the sensitivity variable $w(t) = \partial u(t)/\partial p$.
Differentiating~\eqref{eq:fwdODE} with respect to $p$ and
exchanging the order of $\partial/\partial t$ and $\partial/\partial p$:
\begin{equation}
  \frac{d}{dt}\left[\pd{u}{p}\right] = \pd{F}{u}\,\pd{u}{p} + \pd{F}{p}.
  \label{eq:ODE_FSE_scalar}
\end{equation}
In terms of $w$:
\begin{equation}
  \dot{w} = F_u(u;p)\,w + F_p(u;p), \qquad
  w(0) = \frac{du_0}{dp}.
  \label{eq:FSE_w}
\end{equation}
The coefficient $F_u(u(t);p)$ depends on the forward solution $u(t)$,
so the FSE~\eqref{eq:FSE_w} must be solved simultaneously with
the original ODE~\eqref{eq:fwdODE}.
Together they form an \emph{augmented system} of two ODEs.

For the vector ODE $\dot{\vec{u}} = \vec{F}(\vec{u};\vec{p})$
with $n$ state variables and $m$ parameters:
\begin{equation}
  \frac{d}{dt}\left[\pd{\vec{u}}{p_j}\right]
  = J(\vec{u};\vec{p})\,\pd{\vec{u}}{p_j} + \pd{\vec{F}}{p_j},
  \label{eq:vector_FSE}
\end{equation}
where $J = \partial\vec{F}/\partial\vec{u}$ is the $n\times n$
Jacobian of $\vec{F}$ with respect to $\vec{u}$.
One equation~\eqref{eq:vector_FSE} is needed per parameter $p_j$,
each an $n$-dimensional ODE system.
For $m$ parameters: the augmented system has $n + mn = n(1+m)$ equations.

Theorem~\ref{thm:linearFSE} extends to the dynamical case:
each FSE is \emph{linear in the sensitivity variable}
$\partial\vec{u}/\partial p_j$, even though $\vec{F}$ is nonlinear in $\vec{u}$.
The coefficient $J(\vec{u}(t); \vec{p})$ is time-varying, but the
structure is linear.

\begin{selfcheck}
Write down the augmented ODE system (original ODE + FSE) for the
logistic equation $\dot{u} = ru(1-u/K)$
with respect to both parameters $r$ and $K$ simultaneously.
How many total ODEs does the augmented system have?
\end{selfcheck}
\begin{scAnswer}
The augmented system consists of:
Original: $\dot{u} = ru(1-u/K)$.
FSE for $w_r = \partial u/\partial r$: $\dot{w}_r = r(1-2u/K)w_r + u(1-u/K)$.
FSE for $w_K = \partial u/\partial K$: $\dot{w}_K = r(1-2u/K)w_K + ru^2/K^2$.

Total: 3 ODEs. In general, for $n$ state variables and $m$ parameters,
the augmented system has $n + mn = n(1+m)$ equations.
\end{scAnswer}

\subsection{Time-Dependent Sensitivity Indices}

For dynamical models, the sensitivity index is time-dependent:
\begin{equation}
  \SI{p}{u}(t) = \frac{p}{u(t)}\,w(t) = \frac{p}{u(t)}\,\pd{u(t)}{p}.
  \label{eq:timeSI}
\end{equation}
The relative importance of \POIs can change over time.
A parameter that strongly influences the early phase of an epidemic
may become less important as the system approaches equilibrium.

\subsection{The SIR Model: Full FSE Treatment}
\label{subsec:sirfse}

The SIR model with demography is:
\begin{align}
  \frac{dS}{dt} &= \mu N - \frac{k\beta SI}{N} - \mu S, \label{eq:sirS}\\
  \frac{dI}{dt} &= \frac{k\beta SI}{N} - (\gamma + \mu)I, \label{eq:sirI}\\
  \frac{dR}{dt} &= \gamma I - \mu R. \label{eq:sirR}
\end{align}
With interpretable \POIs $\vec{p} = (k, \beta, \tau, L)$ where
$\tau = 1/\gamma$ (mean infectious period, days) and
$L = 1/\mu$ (mean lifespan, days), the augmented system
for a single \POI $p_j$ adds 3 sensitivity ODEs:
\begin{align}
  \frac{dw_{Sj}}{dt} &= J_{11}w_{Sj} + J_{12}w_{Ij} + J_{13}w_{Rj}
                         + \pd{F_S}{p_j}, \label{eq:sirFSE1}\\
  \frac{dw_{Ij}}{dt} &= J_{21}w_{Sj} + J_{22}w_{Ij} + J_{23}w_{Rj}
                         + \pd{F_I}{p_j}, \label{eq:sirFSE2}\\
  \frac{dw_{Rj}}{dt} &= J_{31}w_{Sj} + J_{32}w_{Ij} + J_{33}w_{Rj}
                         + \pd{F_R}{p_j}, \label{eq:sirFSE3}
\end{align}
where $J_{ij}$ are entries of the Jacobian of $(F_S, F_I, F_R)$
with respect to $(S, I, R)$.
For $m = 4$ \POIs, the full augmented system has $3 + 4\times3 = 15$ ODEs.

The initial conditions for the FSE are zero:
$w_{Sj}(0) = w_{Ij}(0) = w_{Rj}(0) = 0$
(assuming the initial conditions do not depend on the parameters).

Figure~\ref{fig:sir_timevarySI} shows the time-dependent SI for $I(t)$
with respect to $k$, $\beta$, $\tau$, and $L$.
In the early phase of the outbreak (first 30~days), the contact
rate $k$ and transmission probability $\beta$ dominate.
As the outbreak progresses and the susceptible pool depletes,
the mean infectious period $\tau$ becomes increasingly important.
These insights are only available through time-dependent SA;
a static SI computed at a single time point would miss the
qualitative change in parameter importance.

\begin{figure}[htb]
  \centering
  \begin{tikzpicture}
  \begin{axis}[
    width=0.75\textwidth, height=7cm,
    xlabel={Time (days)},
    ylabel={Sensitivity index $\SI{p}{I}(t)$},
    xmin=0, xmax=90,
    ymin=-2.2, ymax=2.2,
    grid=major, grid style={gray!20},
    legend pos=outer north east,
    legend style={font=\small},
    xtick={0,15,30,45,60,75,90},
    yticklabel style={font=\small},
    xticklabel style={font=\small},
  ]
  %% Schematic curves -- representative shapes based on SIR dynamics
  %% k: positive, peaks early ~day 15, then declines
  \addplot[thick, black, domain=0:90, samples=200]
    {1.8*exp(-0.04*(x-15)^2/30)*max(0.01,sin(deg(0.05*x+0.1)))
     + 0.4*exp(-0.015*x)};
  \addlegendentry{Contact rate $k$}

  %% beta: same shape as k (SI_k = SI_beta for R0 = kbeta/gamma)
  \addplot[thick, black, dashed, domain=0:90, samples=200]
    {1.7*exp(-0.04*(x-15)^2/30)*max(0.01,sin(deg(0.05*x+0.1)))
     + 0.35*exp(-0.015*x)};
  \addlegendentry{Trans. prob.\ $\beta$}

  %% tau: positive, grows after ~day 20, peaks ~day 40
  \addplot[thick, black!60, domain=0:90, samples=200]
    {0.2 + 1.4*(1-exp(-0.025*x))*exp(-0.01*(x-45)^2/60)};
  \addlegendentry{Mean inf.\ period $\tau$}

  %% L (lifespan): small negative throughout
  \addplot[thick, black!40, densely dotted, domain=0:90, samples=200]
    {-0.15 - 0.05*sin(deg(0.04*x))};
  \addlegendentry{Mean lifespan $L$}

  \draw[dashed, black!30] (axis cs:0,0) -- (axis cs:90,0);
  \end{axis}
  \end{tikzpicture}
  \caption{Schematic of time-dependent sensitivity indices $\SI{p}{I}(t)$
           for the SIR model. In the early outbreak phase, the contact
           rate $k$ and transmission probability $\beta$ (nearly equal
           because $\SI{k}{R_0} = \SI{\beta}{R_0}$) dominate.
           The mean infectious period $\tau$ grows in importance
           as the outbreak progresses. The mean lifespan $L$ has
           a small negative influence throughout. Curve shapes are
           representative; the MATLAB laboratory (Section~\ref{sec:ch4matlab})
           and the function \texttt{plot\_time\_si.m} in the companion
           repository produce exact computed curves at the nominal parameters.}
  \label{fig:sir_timevarySI}
\end{figure}

%%------------------------------------------------------------
\section{The Proliferation Problem\index{proliferation problem} and Its Consequences}
\label{sec:proliferation}
%%------------------------------------------------------------

\begin{quote}
\textit{Before improving the model, find out which knob actually moves the machine.}
\end{quote}

\medskip\noindent
The FSE is exact and efficient for a single parameter.
For $m$ parameters, the cost scales linearly: $m$ augmented ODE
solves, each requiring $n$ additional state variables.
Table~\ref{tab:prolif} quantifies this for representative models.

\medskip\noindent
\textit{Counting convention.} Throughout this section, ``solve counts''
refer to the total number of ODE solves, including the one baseline
forward solve (of the original $n$-equation system) required before
any sensitivity computation begins. The FSE adds $m$ augmented solves
on top of this baseline, for a total of $m + 1$ solves.

\begin{table}[htb]
\centering
\begin{tabular}{lrrrr}
\toprule
Model & $n$ states & $m$ \POIs & Augmented size & FSE cost \\
\midrule
SIR          & 3  & 4  & $3 + 12 = 15$  & 4 ODE solves \\
SEIR         & 4  & 6  & $4 + 24 = 28$  & 6 ODE solves \\
Malaria model & 7  & 17 & $7 + 119 = 126$ & 17 ODE solves \\
Climate model & 500 & 200 & $500 + 100{,}000$ & 200 ODE solves \\
\bottomrule
\end{tabular}
\caption{Augmented system size and FSE cost for representative models.
         For the climate model, 200 ODE solves each taking 6~hours
         equals 50 days of computation. This is the cost that
         motivates the adjoint method in Chapter~6.}
\label{tab:prolif}
\end{table}

For the SIR model, 4~solves is extremely fast in practice.
For the climate model, 200~solves at 6~hours each is 50~days.
This is the computational wall that motivates the adjoint approach.

\unifyingidea{2} Table~\ref{tab:prolif} reveals the bottleneck
that Chapters~5 and 6 resolve.
No matter how the sensitivity Jacobian is computed, by
finite differences, by analytic FSE, or by the adjoint;
the answer is the same matrix of normalized derivatives.
The adjoint (Unifying Idea~2, introduced in Chapter~5) reduces
the cost from $m$~solves to 2~solves by exploiting the
transposed structure of the problem. The cost reduction is
proportional to $m$; for $m = 200$, the adjoint is
100~times faster than the FSE.

\unifyingidea{3} Every row of Table~\ref{tab:prolif} is a linearization.
The FSE computes the exact derivative of a nonlinear model, which
is a linear approximation to the model's response. That approximation
is valid when parameters are near their nominal values. The proliferation
problem scales with $m$; the validity limitation is independent of $m$.
Chapter~8 addresses the second problem; Chapters~5 and 6 address the first.

\begin{selfcheck}
For the malaria model with $n = 7$ state variables and $m = 17$
\POIs, and one response functional (say, the endemic equilibrium
fraction $i_h$):
(a) How many ODE solves does the FSE approach require?
(b) If each ODE solve takes 2~minutes, how long does the
    complete FSA take?
(c) How many ODE solves will the adjoint method in Chapter~6 require?
\end{selfcheck}
\begin{scAnswer}
(a) 17 FSE solves (one per parameter), plus 1 forward solve = 18 total.
(b) $18 \times 2 = 36$ minutes.
(c) The adjoint always requires exactly 2 solves (1 forward + 1 adjoint),
regardless of $m$.
For 17 parameters, the adjoint is 9 times faster than the FSE.
\end{scAnswer}

%%------------------------------------------------------------
\section{The Decision Table: When to Use Which Method}
\label{sec:decisiontable}
%%------------------------------------------------------------

With three computational methods now available (finite differences,
analytic FSE, and the adjoint previewed above), the right choice
depends on the model and the problem.

\begin{table}[htb]
\centering
\renewcommand{\arraystretch}{1.3}
\begin{tabular}{p{3.2cm}p{2.8cm}p{1.5cm}p{2.8cm}p{1.2cm}}
\toprule
Situation & Method & Accuracy & Cost & When to use \\
\midrule
Black-box model & Finite differences & $O(h^2)$ & $2n_p + 1$ evals & Ch 3 \\
Analytic, small $n_p$ & FSE & Exact & $n_p$ solves & Ch 4 \\
Analytic, large $n_p$ & Adjoint & Exact & 2 solves & Ch 5--6 \\
Correlated \POIs & Reparameterize first & --- & --- & Ch 3 \\
Model near bifurcation & FSE + caution & Exact* & $n_p$ solves & Ch 8 \\
Large parameter uncertainty & Global SA & Variance-based & $N(m+2)$ evals & Ch 9 \\
\bottomrule
\end{tabular}
\caption{Decision table for selecting the SA computational method.
         For a systematic benchmarking of these methods on biological ODE
         models, see Mester et al.~\cite{mester2022differential}.
         (*Near a bifurcation, the FSE solution is exact but the
         sensitivity index diverges, limiting its interpretability.)}
\label{tab:decision}
\end{table}

For this course, the default progression is:
\begin{enumerate}[nosep]
  \item Start with finite differences (Chapter~3) to get SA results quickly.
  \item Switch to analytic FSE (this chapter) when the model is differentiable
        and accuracy matters.
  \item Use the adjoint (Chapters~5--6) when $m$ is large.
  \item Check for correlated \POIs before reporting any results (Chapter~3).
  \item Consider global SA (Chapter~9) when parameter uncertainty is large.
\end{enumerate}

A common assumption is that finite differences are adequate for
any SA study once the step size is chosen correctly.
For black-box models, this is true: FD is the only option
and produces reliable results at the optimal step size.
For differentiable models, the analytic FSE is
simultaneously more accurate (exact rather than $O(h^2)$)
and often faster (no repeated model evaluations at perturbed
parameter values). The two methods give the same numbers
when FD is implemented correctly; the FSE gives those same
numbers without approximation error and without the
step-size selection problem.

%%------------------------------------------------------------
\section{MATLAB Laboratory: Analytic FSA of the SIR Model}
\label{sec:ch4matlab}
%%------------------------------------------------------------

This laboratory implements the augmented ODE system for the SIR model
and compares the analytic FSE approach with the finite-difference
method from Chapter~3.

\textbf{Setup.}
Use the SIR model~\eqref{eq:sirS}--\eqref{eq:sirR} with
\POIs $\vec{p} = (k, \beta, \tau, L)$ (interpretable parameterization),
nominal values $k = 5$, $\beta = 0.06$, $\tau = 7$~days, $L = 10{,}000$~days,
and initial conditions $S_0 = 999$, $I_0 = 1$, $R_0 = 0$, $N = 1000$.

\textbf{Part 1:} Implement the augmented ODE system.
Write a MATLAB function \texttt{sir\_augmented(t, y, k, beta, tau, L)}
that returns the derivatives for the full $15$-equation augmented system:
3~original state equations plus $4 \times 3 = 12$ FSE equations.
The FSE equations require the Jacobian of $(F_S, F_I, F_R)$ with
respect to $(S, I, R)$.

\medskip\noindent
\textit{State vector layout.} The 15-component state vector \texttt{y}
is organized as follows:
\begin{center}
\begin{tabular}{cl}
\toprule
Index & Meaning \\
\midrule
1 & $S(t)$ \\
2 & $I(t)$ \\
3 & $R(t)$ \\
4--6   & $(\partial S/\partial k,\; \partial I/\partial k,\; \partial R/\partial k)$ \\
7--9   & $(\partial S/\partial\beta,\; \partial I/\partial\beta,\; \partial R/\partial\beta)$ \\
10--12 & $(\partial S/\partial\tau,\; \partial I/\partial\tau,\; \partial R/\partial\tau)$ \\
13--15 & $(\partial S/\partial L,\; \partial I/\partial L,\; \partial R/\partial L)$ \\
\bottomrule
\end{tabular}
\end{center}
The initial condition is \texttt{y0 = [S0; I0; 0; zeros(12,1)]}.
The sensitivity variables all start at zero because the initial
conditions do not depend on the parameters.
The file \texttt{sir\_augmented.m} in the companion repository
implements this system and may be used as a reference after
completing your own implementation.

\textbf{Part 2:} Solve and extract sensitivities.
Solve the augmented system from $t = 0$ to $t = 90$~days
using \texttt{ode45}. Extract the time-dependent sensitivity indices
$\SI{k}{I}(t)$, $\SI{\beta}{I}(t)$, $\SI{\tau}{I}(t)$, $\SI{L}{I}(t)$.

\textbf{Part 3:} Plot time-dependent SIs.
Plot all four $\SI{p_j}{I}(t)$ curves on a single axis, alongside the
epidemic curve $I(t)$. Identify:
(a) at $t = 10$~days, which parameter matters most?
(b) at $t = 50$~days, has the ranking changed?
(c) which parameter has the largest absolute SI at the epidemic peak?

\textbf{Part 4:} Compare with finite differences.
At $t = 30$~days and $t = 70$~days, compare the FSE-based SI values
against the centered-difference estimates from Chapter~3.
Are they equal to 4~significant figures? Report any discrepancies
and explain their likely source.

\textbf{Part 5:} Cost comparison.
Time the two approaches: the augmented ODE (all 4 sensitivities at once)
versus 4 separate centered-difference calls to \texttt{SIR\_model}.
Which is faster? Which is more accurate? At what value of $m$
(number of \POIs) would you expect the two approaches to take
approximately equal time?

%%------------------------------------------------------------
\section{Exercises}
\label{sec:ch4exercises}
%%------------------------------------------------------------

\begin{exercise}[\bloom{L2}]
\label{ex:ch4:pattern}
State the three-step FSE derivation pattern in your own words,
without using any symbols. Then explain, in one sentence, why
the FSE is always linear in the sensitivity variable even when
the original equation is nonlinear.
\end{exercise}

\begin{exercise}[\bloom{L3}]
\label{ex:ch4:implicit}
For the nonlinear equation $g(u; p) := u^3 - 3pu + 2 = 0$:
\begin{enumerate}[label=(\alph*)]
  \item Write the FSE for $\partial u/\partial p$.
  \item At $p = 1$, find all real roots $u$ numerically.
  \item Compute $\SI{p}{u}$ at each root.
  \item At $p = 1$, for which root is the sensitivity largest in magnitude?
        Verify your answer using a centered finite difference with optimal $h$.
\end{enumerate}
\end{exercise}

\begin{exercise}[\bloom{L3}]
\label{ex:ch4:logistic}
For the logistic ODE $\dot{u} = ru(1-u/K)$, $u(0) = u_0$:
\begin{enumerate}[label=(\alph*)]
  \item Write the FSE for $w_r = \partial u/\partial r$.
  \item Verify that the FSE is linear in $w_r$ by identifying
        the coefficient (which depends on $u(t)$) and the forcing term.
  \item Using MATLAB, implement the augmented system
        $(u, w_r, w_K)$ and solve from $t = 0$ to $t = 20$.
  \item Plot $\SI{r}{u}(t)$ and $\SI{K}{u}(t)$.
        Interpret: near $t = 0$ vs.\ near the carrying capacity.
\end{enumerate}
\end{exercise}

\begin{exercise}[\bloom{L3}]
\label{ex:ch4:symbiotic}
For the symbiotic population model~\eqref{eq:symp1}--\eqref{eq:symp2}
with $a = 0.5$, $b = 2$, $c = 0.4$:
\begin{enumerate}[label=(\alph*)]
  \item Compute the coexistence equilibrium.
  \item Assemble and solve the FSE for $\partial\vec{u}/\partial c$.
  \item Compute the normalized sensitivity indices
        $\SI{c}{\bar{u}_1}$ and $\SI{c}{\bar{u}_2}$.
  \item Using equation~\eqref{eq:sympeq}, compute the same indices
        analytically. Verify agreement.
\end{enumerate}
\end{exercise}

\begin{exercise}[\bloom{L4}]
\label{ex:ch4:proliferation}
The SEIRS model (Susceptible-Exposed-Infectious-Recovered-Susceptible)
has 4 state variables and 7 parameters.
\begin{enumerate}[label=(\alph*)]
  \item How many ODEs are in the augmented FSE system for all 7 parameters?
  \item If each ODE solve takes 3~minutes, how long does the full FSA take
        using the FSE approach? The adjoint approach (Chapter~6)?
  \item Suppose a modeler decides to reduce computation time by computing
        sensitivities only for the 3 parameters they expect to be most
        important, based on biological intuition. Identify one scenario
        where this strategy would lead to a misleading conclusion.
\end{enumerate}
\end{exercise}

\begin{exercise}[\bloom{L4}]
\label{ex:ch4:sirtime}
Using your MATLAB implementation from the laboratory:
\begin{enumerate}[label=(\alph*)]
  \item At three time points, $t = 5$ (early epidemic),
        $t = t_{\text{peak}}$ (peak infections), and $t = 80$ (late);
        report the tornado plot of SIs for $I(t)$.
  \item Does the parameter ranking change across the three time points?
  \item A public health officer says: ``We should target the parameter
        with the largest SI at the epidemic peak, since that is when the
        most people are affected.'' Critique this statement using your results.
\end{enumerate}
\end{exercise}

\begin{exercise}[\bloom{L5}]
\label{ex:ch4:design}
A colleague is modeling cholera transmission with a SEIR-type model
having 8 parameters and 4 state variables.
She asks for advice on the SA methodology.
Design a complete FSA study, specifying:
\begin{enumerate}[label=(\alph*)]
  \item Whether to use finite differences or analytic FSE, and why.
  \item How you would validate the implementation.
  \item What \QOIs you would compute sensitivity indices for,
        and at what time points.
  \item How you would present the results to identify the
        two most critical intervention targets.
  \item One situation in which the FSE results might be misleading,
        and what additional analysis you would recommend.
\end{enumerate}
\end{exercise}

\begin{exercise}[\bloom{L6}]
\label{ex:ch4:critique}
A published paper reports: ``We performed sensitivity analysis of the
SEIR model using the forward sensitivity equations. All sensitivity
indices were computed at $t = 365$~days, giving a single set of
values for the entire epidemic.''

The model simulates a disease whose outbreak dynamics play out
over approximately 60~days, but the model runs for a year.

Identify at least two methodological issues with this approach.
For each issue, state what the authors should have done instead
and explain what information they likely missed.
\end{exercise}

%%------------------------------------------------------------



%% ---- Part IV: The Adjoint Approach ----
\part{The Adjoint Approach}

\chapter{The Adjoint in Linear Algebra}

\begin{flushright}
\textit{The adjoint is the same sensitivity question, asked from the other end.}
\end{flushright}

\bigskip

Chapter~4 showed that analytic forward sensitivity analysis for a
model with $m = 50$ parameters requires 50 linear solves, one per
parameter. Each solve uses the same coefficient matrix~$A$, so we
pay once for the LU~factorization and 50~times for back-substitution.
That is still 50~operations.

Now suppose we care about a single response functional, not the
entire solution~$\vec{u}$. The quantity of interest is
$J = \vec{c}^T\vec{u}$: the pollution level at one monitoring station,
the total disease burden, the profit at equilibrium.
We want $\partial J/\partial q$ for all 50~parameters~$q$.

Is there a way to compute all 50~derivatives simultaneously,
in the cost of a \emph{single} linear solve?

Yes. The answer is the adjoint.
The adjoint approach to sensitivity analysis was developed for nonlinear
systems by Cacuci~\cite{cacuci1981sensitivity} and has since become a
standard tool in optimal design and computational science~\cite{giles2000introduction}.

%%------------------------------------------------------------
\section{A Preview: The Same Answer, One Solve Instead of Fifty}
\label{sec:adjoint_preview}
%%------------------------------------------------------------

Before the general theory, work through a small example by hand
to see exactly what the adjoint does and why it is efficient.

\subsection*{The Setup}

Consider the $2 \times 2$ linear system
\begin{equation}
  A\vec{u} = \vec{b}(q_1, q_2),
  \qquad
  A = \begin{pmatrix} 3 & 1 \\ 1 & 2 \end{pmatrix},
  \qquad
  \vec{b} = \begin{pmatrix} q_1 \\ q_2 \end{pmatrix},
  \label{eq:preview_system}
\end{equation}
where $q_1$ and $q_2$ are two parameters of interest.
The scalar response functional is
$J = u_1 + u_2 = \vec{c}^T\vec{u}$ with $\vec{c} = (1,1)^T$.
We want $\partial J/\partial q_1$ and $\partial J/\partial q_2$.

\subsection*{The Direct (Forward) Approach: Two Solves}

Differentiating $A\vec{u} = \vec{b}$ with respect to $q_1$:
$A\,(\partial\vec{u}/\partial q_1) = \partial\vec{b}/\partial q_1 = (1,0)^T$.
Differentiating with respect to $q_2$:
$A\,(\partial\vec{u}/\partial q_2) = \partial\vec{b}/\partial q_2 = (0,1)^T$.
Two linear solves (one per parameter) with the same matrix $A$:

\[
  A^{-1} = \tfrac{1}{5}\begin{pmatrix} 2 & -1 \\ -1 & 3 \end{pmatrix},
  \quad\text{so}\quad
  \pd{\vec{u}}{q_1} = \tfrac{1}{5}\begin{pmatrix}2\\-1\end{pmatrix},
  \qquad
  \pd{\vec{u}}{q_2} = \tfrac{1}{5}\begin{pmatrix}-1\\3\end{pmatrix}.
\]

The derivatives of $J = \vec{c}^T\vec{u}$ are:
\[
  \pd{J}{q_1} = \vec{c}^T\pd{\vec{u}}{q_1}
              = (1,1)\cdot\tfrac{1}{5}(2,-1)^T = \tfrac{1}{5},
  \qquad
  \pd{J}{q_2} = \vec{c}^T\pd{\vec{u}}{q_2}
              = (1,1)\cdot\tfrac{1}{5}(-1,3)^T = \tfrac{2}{5}.
\]
Two solves, two derivatives. For 50 parameters: 50 solves.

\subsection*{The Adjoint Approach: One Solve}

Instead of solving for $\partial\vec{u}/\partial q_j$ separately for each $j$,
introduce the \textbf{adjoint vector} $\vec{v}$ satisfying
\begin{equation}
  A^T\vec{v} = \vec{c}.
  \label{eq:preview_adjoint}
\end{equation}
This is one linear solve with the transposed matrix.
Since $A$ is symmetric here, $A^T = A$, so $\vec{v} = A^{-1}\vec{c}$:
\[
  \vec{v} = A^{-1}\vec{c} = \tfrac{1}{5}\begin{pmatrix}2&-1\\-1&3\end{pmatrix}\begin{pmatrix}1\\1\end{pmatrix}
          = \tfrac{1}{5}\begin{pmatrix}1\\2\end{pmatrix}.
\]

\textbf{What does $\vec{v}$ represent?}
The adjoint vector answers the question: \emph{if the state $\vec{u}$
were perturbed by a small amount $\delta\vec{u}$, how much would $J$ change?}
A direct perturbation gives $\delta J \approx \vec{c}^T\delta\vec{u}$,
but not all $\delta\vec{u}$ are achievable --- the system constraint
$A\vec{u} = \vec{b}$ means only perturbations satisfying
$A\,\delta\vec{u} = \delta\vec{b}$ are reachable.
The adjoint vector $\vec{v}$ encodes how much each component
of the state contributes to $J$, after accounting for
the coupling imposed by the matrix $A$.

In the 2$\times$2 example,
$\vec{v} = \tfrac{1}{5}(1,2)^T$ means that
a perturbation to $u_2$ matters \emph{twice as much}
to $J = u_1 + u_2$ as a perturbation to $u_1$ ---
not because $\vec{c} = (1,1)^T$ treats them equally,
but because the system $A$ couples $u_1$ and $u_2$ differently.
This coupling is what makes the adjoint more than a convenient shortcut:
$\vec{v}$ is the correct sensitivity weighting for the constrained system,
not the naive weighting $\vec{c}$ for the unconstrained response.

Now observe: for any parameter $q_j$,
\[
  \pd{J}{q_j} = \vec{c}^T\pd{\vec{u}}{q_j}
              = (A^T\vec{v})^T\pd{\vec{u}}{q_j}
              = \vec{v}^T\!\underbrace{A\pd{\vec{u}}{q_j}}_{=\,\partial\vec{b}/\partial q_j}
              = \vec{v}^T\pd{\vec{b}}{q_j}.
\]
This is the key step: $A\,(\partial\vec{u}/\partial q_j) = \partial\vec{b}/\partial q_j$
is the FSE right-hand side, a known vector.
No additional solve is needed — just a dot product.

Checking against our direct answers:
\[
  \pd{J}{q_1} = \vec{v}^T\pd{\vec{b}}{q_1}
              = \tfrac{1}{5}(1,2)\cdot(1,0)^T = \tfrac{1}{5} \;\checkmark
  \qquad
  \pd{J}{q_2} = \vec{v}^T\pd{\vec{b}}{q_2}
              = \tfrac{1}{5}(1,2)\cdot(0,1)^T = \tfrac{2}{5} \;\checkmark
\]
One solve for $\vec{v}$, then one dot product per parameter.
For 50 parameters: still one solve, then 50 dot products.
The cost advantage is a factor of 50.

\begin{selfcheck}
Add a third parameter $q_3$ to the system above by setting
$\vec{b} = (q_1 + q_3,\; q_2)^T$. Without solving any new linear
systems, use the adjoint vector $\vec{v}$ already computed
to find $\partial J/\partial q_3$.
\end{selfcheck}
\begin{scAnswer}
$\partial\vec{b}/\partial q_3 = (1,0)^T$.
$\partial J/\partial q_3 = \vec{v}^T(1,0)^T = \frac{1}{5}(1,2)\cdot(1,0)^T = \frac{1}{5}$.
No new solve needed. Adding parameters costs only one dot product each
once $\vec{v}$ is known. This is the computational power of the adjoint.
\end{scAnswer}

\medskip

The pattern above is the entire idea of the adjoint, stripped to its
essentials. Section~\ref{sec:adjoint_strategy} states it precisely.

%%------------------------------------------------------------
\section{The Strategy, Stated Explicitly}
\label{sec:adjoint_strategy}
%%------------------------------------------------------------

The cost analysis above raises the key question: can we recover all $m$ sensitivities without solving $m$ separate systems?
It turns out we can, through a trick that is worth understanding before seeing the symbols.

The adjoint is not a new kind of problem.\index{adjoint method} It is a deliberate trick:
introduce an auxiliary vector, choose it to cancel the unknown
sensitivity variable, and read off the result.
Here is the trick in plain language before any symbols.

\begin{enumerate}[label=\textbf{Step \arabic*.}, leftmargin=2.2cm]

\item We want $\partial J/\partial q_j$ for all parameters $q_j$.
Since $J = \vec{c}^T\vec{u}$, this is
$\partial J/\partial q_j = \vec{c}^T (\partial\vec{u}/\partial q_j)$.
But $\partial\vec{u}/\partial q_j$ requires a separate FSE solve
for every $j$ (50~solves for 50~parameters).

\item Introduce an auxiliary vector $\vec{v}$ and observe that
for any vector $\vec{v}$:
\[
  \vec{c}^T\pd{\vec{u}}{q_j}
  = \vec{v}^T\!\left[A\pd{\vec{u}}{q_j}\right]
    + \left(\vec{c}^T - \vec{v}^T\!A\right)\pd{\vec{u}}{q_j}.
\]
The first term on the right is $\vec{v}^T\vec{r}_j$,
where $\vec{r}_j = A(\partial\vec{u}/\partial q_j)$
is the FSE right-hand side (a known vector once $\vec{u}$ is known).
The second term vanishes if we choose $\vec{v}$ so that $\vec{v}^T A = \vec{c}^T$.

\item Choose $\vec{v}$ to satisfy the \textbf{adjoint equation}:
\begin{equation}
  A^T\vec{v} = \vec{c}.
  \label{eq:adjointeq}
\end{equation}
Then $\vec{c}^T - \vec{v}^T A = \vec{0}$ and the sensitivity formula reduces to:
\begin{equation}
  \pd{J}{q_j} = \vec{v}^T\vec{r}_j
  \qquad\text{for all } j = 1, \ldots, m.
  \label{eq:sensfromula}
\end{equation}
\end{enumerate}

One solve for $\vec{v}$ (equation~\eqref{eq:adjointeq}) recovers
all $m$ sensitivities through formula~\eqref{eq:sensfromula}.
The trick works because the adjoint equation does not depend on $j$:
$\vec{v}$ is determined by $\vec{c}$ and $A$ alone.

The adjoint is not a separate, unrelated problem. It is constructed
entirely from the coefficient matrix $A$ (via transposition)
and the response vector $\vec{c}$ of the forward problem.
Everything about the adjoint is determined by the forward problem.

\begin{selfcheck}
For the system $A\vec{u} = \vec{b}$ with
$A = \begin{pmatrix} 2 & 1 \\ 0 & 3 \end{pmatrix}$
and response functional $J = u_1$ (so $\vec{c} = (1,0)^T$):
write down the adjoint equation~\eqref{eq:adjointeq} and solve for $\vec{v}$.
\end{selfcheck}
\begin{scAnswer}
$A^T\vec{v} = \vec{c}$ becomes
$\begin{pmatrix} 2 & 0 \\ 1 & 3 \end{pmatrix}\vec{v}
= \begin{pmatrix} 1 \\ 0 \end{pmatrix}$.
Solving: $v_1 = 1/2$, $v_2 = -1/6$.
This single solve (one back-substitution with $A^T$) gives us the
sensitivity of $J = u_1$ to any number of parameters,
via formula~\eqref{eq:sensfromula}.
\end{scAnswer}

%%------------------------------------------------------------
\section{The Commutative Diagram}\index{commutative diagram}
\label{sec:commdiagram}
%%------------------------------------------------------------

The three-step strategy and the 2$\times$2 preview both express the same
mathematical fact: the FSE route and the adjoint route reach the same
sensitivity, at different cost.
Figure~\ref{fig:commdiag} shows this as a diagram —
read it as a summary of the algebra already developed,
not as an introduction to the idea.

\begin{figure}[htb]
\centering
\begin{tikzpicture}[
  node distance=2.8cm and 4.8cm,
  box/.style={draw, rounded corners=4pt, minimum width=3.2cm,
              minimum height=1.0cm, align=center, font=\small},
  arr/.style={-{Latex[length=3mm]}, thick},
  darr/.style={-{Latex[length=3mm]}, thick, dashed},
]

%% Nodes
\node[box] (fwd)    {Forward problem\\$A\vec{u} = \vec{b}$};
\node[box, right=of fwd]  (fse)  {FSE\\$A\pd{\vec{u}}{q_j} = \vec{r}_j$};
\node[box, below=of fwd]  (adj)  {Adjoint\\$A^T\vec{v} = \vec{c}$};
\node[box, below=of fse]  (sens) {Sensitivity\\$\pd{J}{q_j}$};

%% Forward path (top, right)
\draw[arr] (fwd) -- node[above, font=\small]{solve for $\vec{u}$} (fse);
\draw[arr] (fse) -- node[right, font=\small]{$\vec{c}^T\pd{\vec{u}}{q_j}$} (sens);

%% Adjoint path (left, bottom)
\draw[darr] (fwd) -- node[left, font=\small]{solve $A^T\vec{v}=\vec{c}$} (adj);
\draw[darr] (adj) -- node[below, font=\small]{$\vec{v}^T\vec{r}_j$} (sens);

%% Labels
\node[font=\small\itshape, above=0.15cm of fse] {Forward path: $m$ solves};
\node[font=\small\itshape, below=0.15cm of adj] {Adjoint path: 1 solve};

\end{tikzpicture}
\caption{The commutative diagram for sensitivity analysis.
         The forward path (solid arrows) computes $\partial J/\partial q_j$
         by solving the FSE for each parameter, $m$~solves total.
         The adjoint path (dashed arrows) solves the adjoint equation
         once, then evaluates the sensitivity formula for each parameter;
         1~solve total.
         Both paths deliver the same result.}
\label{fig:commdiag}
\end{figure}

The diagram commutes: both paths from the forward problem to the
sensitivity give identical results. This can be verified directly
from the algebra in Section~\ref{sec:adjoint_strategy}.

\begin{theorem}[Forward--Adjoint Equivalence]
\label{thm:equiv}
Let $A\vec{u} = \vec{b}$ with $A$ nonsingular, $J = \vec{c}^T\vec{u}$,
and FSE $A\,\partial\vec{u}/\partial q_j = \vec{r}_j$.
Let $\vec{v}$ solve the adjoint equation $A^T\vec{v} = \vec{c}$.
Then
\[
  \pd{J}{q_j} = \vec{c}^T\pd{\vec{u}}{q_j} = \vec{v}^T\vec{r}_j
\]
for every parameter $q_j$.
\end{theorem}

\begin{proof}
$\vec{c}^T(\partial\vec{u}/\partial q_j)
= (A^T\vec{v})^T(\partial\vec{u}/\partial q_j)
= \vec{v}^T A(\partial\vec{u}/\partial q_j)
= \vec{v}^T\vec{r}_j$.
\end{proof}

The adjoint gives \emph{exactly} the same sensitivity values as
the FSE. It is not an approximation. Students sometimes expect that
a method requiring fewer solves must sacrifice accuracy;
here it does not.

\begin{selfcheck}
Using Theorem~\ref{thm:equiv}, state in words why the adjoint
equation does not need to be re-solved when a new parameter $q_k$
is added to the sensitivity study.
\end{selfcheck}
\begin{scAnswer}
The adjoint equation $A^T\vec{v} = \vec{c}$ depends only on
$A$ (the system coefficient matrix) and $\vec{c}$ (the response
vector), neither of which changes when we add a new parameter.
Adding parameter $q_k$ changes only the FSE right-hand side
$\vec{r}_k = \partial\vec{b}/\partial q_k - (\partial A/\partial q_k)\vec{u}$,
which is a known vector once the forward solution $\vec{u}$ is computed.
The sensitivity $\partial J/\partial q_k = \vec{v}^T\vec{r}_k$
is then obtained by a single dot product, no additional linear solve.
\end{scAnswer}

%%------------------------------------------------------------
\section{Cost Comparison and the Decision Rule}
\label{sec:adjcost}
%%------------------------------------------------------------

\subsection{When to Use the Adjoint}

Both the FSE and the adjoint produce the same sensitivity Jacobian.
Their costs differ dramatically when $m$ and $r$ are unequal,
where $m$ is the number of \POIs and $r$ is the number of
response functionals.

For the FSE: one solve per \POI $\to$ $m$~solves total.

For the adjoint: one solve per response functional $\to$ $r$~solves total.

\begin{equation}
  \boxed{m > r \implies \text{use adjoint;} \qquad
         r > m \implies \text{use FSE}.}
  \label{eq:decisionrule}
\end{equation}

In most SA applications, there are many \POIs ($m$ large) and
a small number of \QOIs of primary interest ($r$ small), so
the adjoint is preferred. The canonical situation is $r = 1$:
one scalar response functional (peak infection count, total burden,
profit at equilibrium) and many parameters.

\subsection{Cost Comparison Plot}

Figure~\ref{fig:costcomp} shows the number of linear solves
required by each method as a function of the number of parameters.

\begin{figure}[htb]
\centering
\begin{tikzpicture}
\begin{semilogyaxis}[
  width=0.68\textwidth, height=6.5cm,
  xlabel={Number of \POIs $m$},
  ylabel={Number of linear solves},
  xmin=1, xmax=500,
  ymin=0.8, ymax=600,
  xtick={1,10,50,100,200,500},
  ytick={1,10,100},
  yticklabels={1,10,100},
  grid=major, grid style={gray!20},
  legend pos=north west,
  legend style={font=\small},
]
\addplot[thick, black, domain=1:500, samples=100] {x};
  \addlegendentry{FSE: $m$ solves}
\addplot[thick, dashed, black] coordinates {(1,1)(500,1)};
  \addlegendentry{Adjoint: 1 solve ($r=1$)}
\addplot[thick, dotted, black] coordinates {(1,3)(500,3)};
  \addlegendentry{Adjoint: 3 solves ($r=3$)}
\end{semilogyaxis}
\end{tikzpicture}
\caption{Cost comparison for FSE vs.\ adjoint sensitivity analysis,
         measured in number of linear solves.
         For $r = 1$ response functional (dashed), the adjoint
         always requires exactly 1~solve regardless of $m$.
         For $r = 3$ (dotted), the adjoint requires 3~solves.
         The FSE requires $m$~solves.
         The crossover occurs at $m = r$; for $m > r$, the adjoint wins.}
\label{fig:costcomp}
\end{figure}

For the symbiotic population model ($m = 3$ parameters,
$r = 2$ response functionals $\{\bar{u}_1, \bar{u}_2\}$):
FSE requires 3~solves; adjoint requires 2.
The advantage is modest here.
For the malaria model ($m = 17$, $r = 1$): FSE requires 17~solves;
adjoint requires 1.
For a climate model ($m = 200$, $r = 2$): FSE requires 200~solves;
adjoint requires 2: a factor of 100~improvement.

\begin{selfcheck}
A model has $m = 30$ \POIs and $r = 5$ response functionals.
Each linear solve takes 2~minutes.
How long does the FSE approach take?
How long does the adjoint approach take?
At what value of $r$ would the two approaches have equal cost?
\end{selfcheck}
\begin{scAnswer}
FSE: $30 \times 2 = 60$~minutes.
Adjoint: $5 \times 2 = 10$~minutes. Savings: 50~minutes.
Equal cost when $r = m = 30$.
\end{scAnswer}

%%------------------------------------------------------------
\section{The Symbiotic Population: Full Adjoint Treatment}
\label{sec:sympAdj}
%%------------------------------------------------------------

We now apply the adjoint to the symbiotic population equilibrium
from Chapter~4, comparing the forward and adjoint approaches
side by side. The forward problem is $A\vec{\bar{u}} = \vec{b}$
at equilibrium, with $A = D_{\vec{u}}[\vec{f}]$ as in
equation~(4.7). The comparison makes the cost advantage concrete.

\subsection{FSE Approach: Three Parameters, Three Solves}

For parameters $a$, $b$, $c$, the FSE requires solving
$A(\partial\vec{\bar{u}}/\partial q_j) = -\partial\vec{f}/\partial q_j$
for $j = a, b, c$, three linear solves with the same matrix $A$,
one LU factorization and three back-substitutions.

\subsection{Adjoint Approach: Two Solves for Two Response Functionals}

Now suppose we want $\partial J_1/\partial q_j$ and
$\partial J_2/\partial q_j$ for all three parameters,
where $J_1 = \bar{u}_1$ and $J_2 = \bar{u}_2$.
The response vectors are $\vec{c}_1 = (1,0)^T$ and
$\vec{c}_2 = (0,1)^T$.

The two adjoint equations are:
\begin{equation}
  A^T\vec{v}_1 = \vec{c}_1 = \begin{pmatrix}1\\0\end{pmatrix},
  \qquad
  A^T\vec{v}_2 = \vec{c}_2 = \begin{pmatrix}0\\1\end{pmatrix}.
  \label{eq:sympAdjEqs}
\end{equation}
Two solves (one per response functional) give $\vec{v}_1$ and $\vec{v}_2$.

\medskip\noindent
\textbf{A note on correlated QOIs.}
Even when the \POIs are independent (uncorrelated inputs),
the \QOIs are almost always correlated outputs:
$J_1 = \bar{u}_1$ and $J_2 = \bar{u}_2$ both depend on the same
equilibrium $\vec{\bar{u}}$, so perturbations to any parameter
affect both simultaneously.
This is not a problem for the adjoint method; it is simply a property
of the model.
The adjoint computes each $\partial J_k / \partial q_j$ independently
by solving $A^T\vec{v}_k = \vec{c}_k$ for each response functional.
Any statistical correlation among the $J_k$'s is a property of the model,
not an issue for the adjoint method --- no special handling is required.
The adjoint naturally produces correct, consistent sensitivities
for all $r$ response functionals by virtue of operating through
the same forward solution $\vec{\bar{u}}$.

The sensitivities follow from formula~\eqref{eq:sensfromula}:
\begin{equation}
  \pd{J_k}{q_j} = \vec{v}_k^T\vec{r}_j,
  \qquad
  \vec{r}_j = -\pd{\vec{f}}{q_j},
  \label{eq:sympSens}
\end{equation}
giving all $2 \times 3 = 6$ sensitivities from two adjoint solves
rather than three FSE solves.

\subsection{Verification}

Theorem~\ref{thm:equiv} guarantees the two approaches give identical
results. At nominal values $a = 0.4$, $b = 1$, $c = 0.3$:
\begin{center}
\renewcommand{\arraystretch}{1.3}
\begin{tabular}{lcccccc}
\toprule
 & $\SI{a}{\bar{u}_1}$ & $\SI{b}{\bar{u}_1}$ & $\SI{c}{\bar{u}_1}$
 & $\SI{a}{\bar{u}_2}$ & $\SI{b}{\bar{u}_2}$ & $\SI{c}{\bar{u}_2}$ \\
\midrule
FSE (3 solves)    & $-0.530$ & $0$ & $0.136$ & $+0.136$ & $0$ & $-0.292$ \\
Adjoint (2 solves)& $-0.530$ & $0$ & $0.136$ & $+0.136$ & $0$ & $-0.292$ \\
\bottomrule
\end{tabular}
\end{center}

The rows are identical. The adjoint is not an approximation.

\medskip

A structural observation: the two adjoint solutions
$\vec{v}_1, \vec{v}_2$ are the columns of $(A^T)^{-1}$,
which equals $(A^{-1})^T$. Collecting all adjoint solutions
into a matrix therefore gives the transpose of the inverse.
This connection is intellectually interesting but not computationally useful:
we never form $A^{-1}$ explicitly, we extract the product
$\vec{v}_k^T\vec{r}_j$ for each required sensitivity, which is
far cheaper than computing the full inverse matrix.

%%------------------------------------------------------------
\section{Reverse Mode and Backpropagation}
\label{sec:reversemodeAD}
%%------------------------------------------------------------

The adjoint strategy of the preceding sections applies to
linear systems. This section extends it to arbitrary differentiable
computations via the computation graph, revealing the connection
to algorithmic differentiation and, through it, to the training
of neural networks.

\subsection{The Computation Graph}

Every differentiable computation can be represented as a directed
acyclic graph in which each node holds an intermediate value
and each edge represents an elementary operation (addition,
multiplication, division, elementary function).

For $R_0 = k\beta/\gamma$, the computation graph is:
\begin{equation}
  k, \beta \xrightarrow{\times} v_1 = k\beta
  \xrightarrow{\div\gamma} v_2 = k\beta/\gamma = R_0.
  \label{eq:R0graph}
\end{equation}
Reading left to right is the forward evaluation:
given $k$, $\beta$, $\gamma$, compute $R_0$.

\subsection{Forward Mode: One Pass Per Input}

In \textbf{forward mode}\index{algorithmic differentiation!forward mode} differentiation, we propagate derivatives
$\dot{v}_i = \partial v_i/\partial q$ through the graph
for one input $q$ at a time.

For $q = k$: set $\dot{k} = 1$, $\dot{\beta} = 0$, $\dot{\gamma} = 0$.
Propagate forward:
\[
  \dot{v}_1 = \dot{k}\beta + k\dot{\beta} = \beta,
  \qquad
  \dot{R}_0 = \dot{v}_1/\gamma - v_1\dot{\gamma}/\gamma^2 = \beta/\gamma.
\]
So $\partial R_0/\partial k = \beta/\gamma$.
Repeat with $\dot{\beta} = 1$ to get $\partial R_0/\partial\beta$.
Repeat with $\dot{\gamma} = 1$ to get $\partial R_0/\partial\gamma$.
Total: 3~forward passes for 3~parameters.

\subsection{Reverse Mode: One Pass for All Inputs}

In \textbf{reverse mode}\index{algorithmic differentiation!reverse mode} differentiation, we propagate the adjoint
variable $\bar{v}_i = \partial J/\partial v_i$ backward through
the graph.

Set $\bar{R}_0 = \partial J/\partial R_0 = 1$ (we are computing
$\partial R_0/\partial(\cdot)$, so $J = R_0$).
Propagate backward:
\begin{align}
  \bar{v}_1   &= \bar{R}_0 \cdot \pd{R_0}{v_1} = 1 \cdot (1/\gamma) = 1/\gamma,
               \label{eq:rev1}\\
  \bar{k}     &= \bar{v}_1 \cdot \pd{v_1}{k}   = (1/\gamma)\cdot\beta = \beta/\gamma,
               \label{eq:rev2}\\
  \bar{\beta}  &= \bar{v}_1 \cdot \pd{v_1}{\beta} = (1/\gamma)\cdot k = k/\gamma,
               \label{eq:rev3}\\
  \bar{\gamma} &= \bar{R}_0 \cdot \pd{R_0}{\gamma} = 1\cdot(-v_1/\gamma^2)
                = -k\beta/\gamma^2.
               \label{eq:rev4}
\end{align}
One backward pass gives all three derivatives simultaneously:
\[
  \pd{R_0}{k} = \bar{k} = \frac{\beta}{\gamma}, \quad
  \pd{R_0}{\beta} = \bar{\beta} = \frac{k}{\gamma}, \quad
  \pd{R_0}{\gamma} = \bar{\gamma} = -\frac{k\beta}{\gamma^2}.
\]
These match the analytic results. The reverse pass cost is
proportional to the graph size, one pass regardless of the
number of inputs.

Table~\ref{tab:fwdrev} summarizes the comparison.

\begin{table}[htb]
\centering
\begin{tabular}{lll}
\toprule
 & Forward mode & Reverse mode \\
\midrule
Propagates       & $\dot{v}_i = \partial v_i/\partial q$ (one $q$) &
                   $\bar{v}_i = \partial J/\partial v_i$ (one $J$) \\
Direction        & Input $\to$ output & Output $\to$ input \\
Cost             & 1 pass per input & 1 pass per output \\
Preferred when   & $m \leq r$ (few inputs) & $m \geq r$ (few outputs) \\
Matrix analog    & FSE: $A\,\partial\vec{u}/\partial q_j = \vec{r}_j$ &
                   Adjoint: $A^T\vec{v} = \vec{c}$ \\
\bottomrule
\end{tabular}
\caption{Forward mode and reverse mode differentiation:
         a side-by-side comparison.
         The matrix analogs confirm that the two modes are
         the computation-graph generalizations of the FSE and
         adjoint approaches from Sections~\ref{sec:adjoint_strategy}--\ref{sec:adjcost}.}
\label{tab:fwdrev}
\end{table}

\subsection{The Backpropagation Connection\index{backpropagation}}

\medskip\noindent
\textit{Note for readers without machine learning background:
this subsection connects the adjoint to neural network training.
If backpropagation is unfamiliar, skip this subsection and proceed
to Chapter~6 --- the adjoint ODE does not depend on it.
Return here after completing Chapter~6 if you want the connection.}

\medskip

If you have encountered training a neural network, you have already
used the adjoint under a different name.

A neural network is a composition of parameterized functions.
Training minimizes a loss function $J(\theta)$ over parameters
$\theta$ using gradient descent, which requires $\partial J/\partial\theta_i$
for every parameter~$\theta_i$.
The loss function is a scalar ($r = 1$) and modern networks have
millions of parameters ($m \gg 1$), so reverse mode is overwhelmingly
preferred. Backpropagation is exactly the reverse-mode computation
described above, applied to the computation graph of the network.

To map notation explicitly: each computation-graph node $v_i$ corresponds
to the $i$-th standard basis vector $\vec{e}_i$, and the adjoint variable
$\lambda_i = \partial J/\partial v_i$ records how much $J$ changes per unit
change at node $i$ --- exactly the role played by $\bar{v}_i$ in reverse mode.

The adjoint equation for the ODE case, derived in Chapter~6,
is the continuous limit of this discrete backward accumulation.
The backward ODE for $\lambda(t)$ is the functional-analytic
version of the backward pass through the computation graph.

\unifyingidea{2} The adjoint is the transpose in whatever space
you are working in.

In this chapter: the adjoint equation is $A^T\vec{v} = \vec{c}$,
the transposed coefficient matrix applied to the adjoint vector.

In the computation graph: reverse mode is the chain rule applied
in reverse order, accumulating $\bar{v}_i = \partial J/\partial v_i$
backward through the graph.

In Chapter~6: the adjoint will be an ODE running backward in time,
with a terminal condition at $t = T$.
The integration-by-parts identity that produces this ODE is the
functional-analytic analog of the matrix transposition that produced
$A^T$ here.

One idea. Three settings. One cost advantage: when $m \gg r$,
the adjoint/reverse-mode is the preferred computational strategy.

\begin{selfcheck}
Use the reverse-mode computation of equations~\eqref{eq:rev1}--\eqref{eq:rev4}
to compute $\SI{k}{R_0}$, $\SI{\beta}{R_0}$, and $\SI{\gamma}{R_0}$.
Express each as a pure number and verify using the Chapter~2 results.
\end{selfcheck}
\begin{scAnswer}
$\partial R_0/\partial k = \beta/\gamma$.
$\SI{k}{R_0} = (k/R_0)(\partial R_0/\partial k)
= (k/(k\beta/\gamma))\cdot(\beta/\gamma) = 1$.

$\partial R_0/\partial\beta = k/\gamma$.
$\SI{\beta}{R_0} = (\beta/(k\beta/\gamma))(k/\gamma) = 1$.

$\partial R_0/\partial\gamma = -k\beta/\gamma^2$.
$\SI{\gamma}{R_0} = (\gamma/(k\beta/\gamma))(-k\beta/\gamma^2) = -1$.

All three match the Chapter~2 results. The reverse-mode computation
recovers the full sensitivity Jacobian in one pass.
\end{scAnswer}

%%------------------------------------------------------------
\section{Bridge to Chapter 6}
\label{sec:bridge}
%%------------------------------------------------------------

Students sometimes expect the adjoint for differential equations
to require advanced functional analysis (Hilbert spaces,
adjoint operators, and formal operator theory.
Appendix~\ref{app:innerproduct} provides the $L^2$ inner product and
orthogonality needed for that formalism; Appendix~\ref{app:adjointformalism}
develops the full abstract operator framework that makes the connection rigorous.
The preceding sections show why this expectation is misplaced:
the core idea is the matrix transpose.
Every formal construct in Chapter~6 is the natural extension
of the matrix algebra developed here.
Reading Chapter~6 with this chapter in mind, the abstract
notation should feel like familiar territory.

\unifyingidea{1} Every sensitivity value computed in
Section~\ref{sec:sympAdj}, via the adjoint, is a normalized
derivative, an instance of Unifying Idea~1 from Chapter~1.
The adjoint is a computational strategy; the quantity it produces
is still the SI defined in Chapter~2.
The two commutative paths in Figure~\ref{fig:commdiag} produce
the same sensitivity index by different routes.

Chapter~5 has established the adjoint for linear systems.
Chapter~6 extends every result to ordinary differential equations.
The sensitivity Jacobian computed in this chapter is the same object
that Smith~\cite{smith2014uncertainty} uses in \S6.1 to drive
parameter selection and optimal experimental design; students who
complete Chapter~6 will find \S6.1 of Smith immediately accessible.

The translation is direct:

\begin{center}
\renewcommand{\arraystretch}{1.5}
\begin{tabular}{ll}
\toprule
Chapter~5 (linear algebra) & Chapter~6 (dynamical systems) \\
\midrule
Coefficient matrix $A$ & Differential operator $d/dt - J(t)$ \\
Transpose $A^T$ & Adjoint operator $-(d/dt + J(t)^T)$ \\
Adjoint equation $A^T\vec{v} = \vec{c}$ &
  Adjoint ODE $-\dot{\lambda} = J^T\lambda + g_u$ \\
Terminal condition: none & Terminal condition: $\lambda(T) = h_u^T$ \\
Sensitivity formula $\vec{v}^T\vec{r}_j$ &
  Sensitivity formula $\int_0^T\lambda^T F_{p_j}\,dt$ \\
Cost: 1 solve for any $m$ & Cost: 2 ODE solves for any $m$ \\
\bottomrule
\end{tabular}
\end{center}

The integration-by-parts identity in Chapter~6 plays the role of
the algebraic transposition here. The terminal condition on $\lambda(T)$
is the time-dependent analog of the right-hand side $\vec{c}$ in
$A^T\vec{v} = \vec{c}$. The cost advantage is identical: one solve
(or two, for the ODE case where a forward solve is also needed)
recovers all $m$ sensitivities regardless of $m$.

\medskip
In Chapter~6 we extend every result of this chapter to ordinary
differential equations.
The matrix $A$ becomes a differential operator.
The transpose $A^T$ becomes the adjoint differential operator.
The linear solve $A^T\vec{v} = \vec{c}$ becomes an ODE for
$\lambda(t)$ with a terminal condition at $t = T$.
The sensitivity formula $\vec{v}^T\vec{r}_j$ becomes an integral
over $[0,T]$.
The cost advantage is the same: 2~ODE solves instead of $m$.
\medskip

%%------------------------------------------------------------
\section{MATLAB Laboratory: Forward vs.\ Adjoint}
\label{sec:ch5matlab}
%%------------------------------------------------------------

This laboratory compares the FSE and adjoint computations on the
symbiotic population model, then implements reverse-mode
differentiation for $R_0$.

\textbf{Part 1: Symbiotic population via FSE.}
At nominal values $a = 0.4$, $b = 1$, $c = 0.3$, compute
the equilibrium using equation~(4.6). Assemble the Jacobian~$A$
and compute the LU factorization using MATLAB's \texttt{lu}.
Solve the three FSEs ($\partial\vec{\bar{u}}/\partial a$,
$\partial\vec{\bar{u}}/\partial b$, $\partial\vec{\bar{u}}/\partial c$)
using \texttt{U\textbackslash(L\textbackslash r)} for each.
Record the time.

\textbf{Part 2: Symbiotic population via adjoint.}
Solve the two adjoint equations $A^T\vec{v}_1 = \vec{c}_1$ and
$A^T\vec{v}_2 = \vec{c}_2$ using the same LU factorization (with
transposed triangular solves). Compute all 6~sensitivities via
$\vec{v}_k^T\vec{r}_j$. Record the time and verify that all 6~values
agree with Part~1 to machine precision.

\medskip\noindent
\textit{MATLAB note on the transposed solve.}
MATLAB's \texttt{lu} gives $A = P^T L U$, so $A^T = U^T L^T P$.
Solving $A^T\vec{v} = \vec{c}$ therefore proceeds in three steps:
$U^T\vec{w}_1 = P\vec{c}$, then $L^T\vec{w}_2 = \vec{w}_1$, then $\vec{v} = \vec{w}_2$.
In MATLAB this is written compactly as \texttt{L' \textbackslash\ (U' \textbackslash\ (P'*c))},
reading right to left: first multiply by \texttt{P'}, then solve the two
transposed triangular systems.

\begin{verbatim}
%% Assemble Jacobian at equilibrium (see Chapter 4)
A = jac_symp(ubar, a, b, c);
[L, U, P] = lu(A);

%% FSE: 3 solves
for j = 1:3
    rj = rhs_symp(ubar, j);          %% FSE right-hand side
    dwdp(:,j) = U \ (L \ (P * rj));
end

%% Adjoint: 2 solves
c1 = [1;0]; c2 = [0;1];
v1 = L' \ (U' \ (P' * c1));   %% solve A^T v1 = c1
v2 = L' \ (U' \ (P' * c2));   %% solve A^T v2 = c2
for j = 1:3
    rj = rhs_symp(ubar, j);
    dJdp(1,j) = v1' * rj;     %% dJ1/dp_j
    dJdp(2,j) = v2' * rj;     %% dJ2/dp_j
end
\end{verbatim}

\textbf{Part 3: Timing as a function of system size.}
For random systems of size $n = 5, 10, 50, 100$ with $m = 20$
parameters and $r = 1$ response functional, time both approaches.
Plot computation time vs.\ $n$ on log-log axes.
Does the ratio of times approach $m/r = 20$ as expected?

\textbf{Part 4: Reverse-mode $R_0$ by hand.}
Implement the reverse-mode computation~\eqref{eq:rev1}--\eqref{eq:rev4}
in MATLAB for $R_0 = k\beta/\gamma$.
Verify all three sensitivities against the MATLAB Symbolic Toolbox.
Count the number of arithmetic operations in forward mode
vs.\ reverse mode for this graph.

\textbf{Part 5: Verify that adjoint = 50 FSE solves.}
Set up a random $10\times10$ system with $m = 50$ parameters and $r = 1$.
Compute $\partial J/\partial q_j$ for all $j$ by:
(a) 50~FSE solves (forward path); (b) 1~adjoint solve (adjoint path).
Verify that the maximum absolute difference is less than $10^{-12}$.

%%------------------------------------------------------------
\section{Exercises}
%%------------------------------------------------------------

\begin{exercise}[\bloom{L2}]
\label{ex:ch5:strategy}
State the three-step adjoint strategy from Section~\ref{sec:adjoint_strategy}
in your own words, using no symbols.
Then explain: why does the adjoint equation $A^T\vec{v} = \vec{c}$ not
depend on which parameter $q_j$ we are differentiating with respect to?
Why does this independence make the adjoint efficient?
\end{exercise}

\begin{exercise}[\bloom{L3}]
\label{ex:ch5:2x2}
For the linear system
\[
  A\vec{u} = \vec{b}, \qquad
  A = \begin{pmatrix} 3 & 1 \\ 1 & 2 \end{pmatrix}, \quad
  \vec{b} = \begin{pmatrix} b_1(q) \\ b_2(q) \end{pmatrix}
  = \begin{pmatrix} q \\ 2q \end{pmatrix},
\]
with response functional $J = u_1 + u_2$ (so $\vec{c} = (1,1)^T$):
\begin{enumerate}[label=(\alph*)]
  \item Solve the forward problem to find $\vec{u}(q)$ at $q = 1$.
  \item Write and solve the FSE for $\partial\vec{u}/\partial q$.
  \item Write and solve the adjoint equation.
  \item Use formula~\eqref{eq:sensfromula} to compute $\partial J/\partial q$.
  \item Verify by differentiating $J(q) = u_1(q) + u_2(q)$ directly.
\end{enumerate}
\end{exercise}

\begin{exercise}[\bloom{L3}]
\label{ex:ch5:decisionrule}
For each scenario below, state whether you would use the FSE or the adjoint,
and calculate the number of linear solves each approach requires.
\begin{enumerate}[label=(\alph*)]
  \item $m = 5$ parameters, $r = 20$ response functionals.
  \item $m = 100$ parameters, $r = 3$ response functionals.
  \item $m = 8$ parameters, $r = 8$ response functionals.
  \item $m = 1{,}000$ parameters, $r = 1$ response functional.
        Each solve takes 0.1~seconds. How much faster is the preferred method?
\end{enumerate}
\end{exercise}

\begin{exercise}[\bloom{L4}]
\label{ex:ch5:symbiotic}
For the symbiotic population equilibrium with parameters
$a = 0.5$, $b = 2$, $c = 0.4$:
\begin{enumerate}[label=(\alph*)]
  \item Compute all 6~sensitivities ($\SI{a}{\bar{u}_1}$,
        $\SI{b}{\bar{u}_1}$, $\SI{c}{\bar{u}_1}$,
        $\SI{a}{\bar{u}_2}$, $\SI{b}{\bar{u}_2}$, $\SI{c}{\bar{u}_2}$)
        using the FSE approach (3 solves).
  \item Compute the same 6~sensitivities using the adjoint approach (2 solves).
  \item Verify that the two methods agree.
  \item Which approach requires fewer solves? By how much?
        At what values of $m$ and $r$ would the FSE be preferred?
\end{enumerate}
\end{exercise}

\begin{exercise}[\bloom{L4}]
\label{ex:ch5:reverseSEIR}
The SEIR basic reproduction number is
$R_0 = k\beta\nu/[(\nu+\mu)(\gamma+\mu)]$
with parameters $k$, $\beta$, $\nu$ (exposure rate), $\gamma$ (recovery rate),
$\mu$ (mortality rate).
\begin{enumerate}[label=(\alph*)]
  \item Draw the computation graph for $R_0$ (intermediate nodes and edges).
  \item Implement the forward-mode computation to find $\partial R_0/\partial\nu$.
        Count the number of arithmetic operations.
  \item Implement the reverse-mode computation to find all five
        derivatives $\partial R_0/\partial k$, $\partial R_0/\partial\beta$,
        $\partial R_0/\partial\nu$, $\partial R_0/\partial\gamma$,
        $\partial R_0/\partial\mu$ in one backward pass.
  \item Compute the five normalized sensitivity indices and produce a
        tornado plot using interpretable \POIs (mean times, not rates).
\end{enumerate}
\end{exercise}

\begin{exercise}[\bloom{L5}]
\label{ex:ch5:design}
A water-quality model of a river system has $m = 40$ parameters
(pollution source rates, degradation rates, flow speeds, mixing coefficients)
and $r = 2$ response functionals: (1) pollutant concentration at a
downstream monitoring station, and (2) total pollutant load entering
a protected estuary.

Design a complete SA study using the adjoint method.
Specify:
\begin{enumerate}[label=(\alph*)]
  \item The form of the adjoint equation for this model.
  \item How many adjoint solves are required and why.
  \item How you would validate the implementation.
  \item How you would use the sensitivity results to advise
        regulators on which pollution sources to prioritize.
  \item One situation where the adjoint results might be misleading,
        and what you would do to detect and correct it.
\end{enumerate}
\end{exercise}

\begin{exercise}[\bloom{L6}]
\label{ex:ch5:critique}
A colleague states: ``The adjoint always gives the same answer as
forward SA, just faster. So if we already have the FSE code working,
there is no reason to implement the adjoint.''

Evaluate this claim carefully.
\begin{enumerate}[label=(\alph*)]
  \item Is the first sentence (same answer, just faster) correct?
        Give a precise mathematical justification.
  \item Is the second sentence (no reason to implement the adjoint)
        correct? Describe at least two situations where switching from
        FSE to adjoint would be strongly advisable.
  \item For the FSE to be preferred over the adjoint, what relationship
        between $m$ and $r$ must hold? Give a specific numerical example
        where FSE is strictly better.
\end{enumerate}
\end{exercise}

%%------------------------------------------------------------


\chapter{The Adjoint for Dynamical Systems}

\begin{flushright}
\textit{To understand the effect of an outcome, trace its influence backward through time.}

\medskip
\textit{The forward model tells you where you end up.\\
The adjoint tells you how you got there.}
\end{flushright}

\bigskip

Climate models routinely have 500 or more tunable parameters.
Forward sensitivity analysis using the FSE approach requires
one ODE solve per parameter: 500~solves on a modern workstation,
each taking 6~hours. Total: 3{,}000~hours, or 125~days.

The adjoint approach requires exactly two ODE solves: one forward,
one backward. Total: 12~hours.

The mathematics behind this 125-to-1 speedup is a direct extension
of the matrix adjoint from Chapter~5. In that chapter, the key operation
was transposing the coefficient matrix: $A \to A^T$. Here, the key
operation is transposing the differential operator, which turns a forward
ODE into a backward ODE. This chapter derives that operation step by step,
explains in plain language why the backward direction arises,
and implements the full calculation for the SIR epidemic model.
The notation is one step heavier than Chapter~5 --- integration by parts
and continuous-time functionals replace matrix transposes --- but the
underlying idea is the same transposition argument in a new space.
The adjoint sensitivity equations for ODE systems, including their
software implementation, are treated in depth by Cao et al.~\cite{cao2002adjoint};
a recent comparison of forward and adjoint methods on biological models is
given by Mester et al.~\cite{mester2022differential}.

%%------------------------------------------------------------
\section{The Timeline Figure}\index{adjoint method!backward in time}
\label{sec:timeline}
%%------------------------------------------------------------

The most important idea in this chapter: the reason the adjoint ODE
runs backward, is captured in Figure~\ref{fig:timeline} before any
equations appear.

\begin{figure}[htb]
\centering
\begin{tikzpicture}[
  thick,
  arr/.style={-{Latex[length=3.5mm]}, thick},
  darr/.style={-{Latex[length=3.5mm]}, thick, dashed},
]

%% Time axis
\draw[arr] (-0.3,0) -- (9.8,0) node[right]{$t$};
\node[below] at (0,0) {$0$};
\node[below] at (9.3,0) {$T$};
\draw (0,-0.08) -- (0,0.08);
\draw (9.3,-0.08) -- (9.3,0.08);

%% Forward ODE arrow
\draw[arr, very thick] (0,0.9) -- (9.3,0.9);
\node[above] at (4.65,0.9) {Forward ODE: state $\vec{u}(t)$};
\node[above, font=\small] at (4.65,0.5)
  {initial condition $\vec{u}(0) = \vec{u}_0$ at $t=0$};

%% Adjoint ODE arrow (backward)
\draw[darr, very thick] (9.3,-0.9) -- (0,-0.9);
\node[below] at (4.65,-0.9) {Adjoint ODE: $\lambda(t)$};
\node[below, font=\small] at (4.65,-1.3)
  {terminal condition $\lambda(T)$ at $t=T$};

%% Annotations
\node[left, align=right, font=\small] at (-0.3, 0.9)
  {propagates\\forward $\to$};
\node[right, align=left, font=\small] at (9.5,-0.9)
  {$\leftarrow$ propagates\\backward};

%% Brace for J
\draw[decorate, decoration={brace,amplitude=5pt}]
  (9.5, 0.05) -- (9.5, -0.05);
\node[right, font=\small, align=left] at (9.65, 0)
  {response\\$J$};

\end{tikzpicture}
\caption{The timeline for forward and adjoint sensitivity analysis.
         The forward ODE propagates the model state $\vec{u}(t)$ from
         the initial condition at $t=0$ toward the response $J$ at $t=T$.
         The adjoint ODE propagates the sensitivity weight $\lambda(t)$
         backward from a terminal condition at $t=T$, in the direction
         opposite to the forward problem.
         The two ODEs are not independent: the adjoint coefficient
         depends on the forward solution $\vec{u}(t)$ at every time step.}
\label{fig:timeline}
\end{figure}

The forward ODE begins at $t=0$ and moves right, building up the state
$\vec{u}(t)$ until it reaches the response $J$ at time $T$.
The adjoint ODE begins at $t=T$ and moves left, distributing the
response backward through time.

This opposite direction is not a numerical trick. It is a mathematical
necessity, explained in Section~\ref{sec:whybackward}.

%%------------------------------------------------------------
\section{Why Backward? Intuition Before Derivation}
\label{sec:whybackward}
%%------------------------------------------------------------

The adjoint ODE runs backward because the question it answers
is naturally backward: how much did each moment in time contribute
to the final value of $J$?

Think of the epidemic. A perturbation applied to the infectious
compartment at $t = 0$ has the entire simulation period to propagate:
it changes the dynamics at $t = 1$, $t = 2$, and all the way to $T$.
Its influence on $J$ is large. A perturbation applied at $t = T-\epsilon$
has almost no time to propagate before the simulation ends.
Its influence on $J$ is negligible.

So the ``weight'' we should assign to perturbations decreases
as we move from left to right along the time axis,
and increases as we move from right to left.
The adjoint variable $\lambda(t)$ is precisely this weight:
$\lambda(t)$ measures how sensitive the response $J$ is to
a perturbation of the state at time $t$.

At $t = T$, the weight is determined by the response functional
(how does $J$ change if the final state changes?).
At earlier times, $\lambda(t)$ is computed by propagating this
weight backward, accounting for how perturbations at time $t$
have time to amplify before reaching $T$.

The adjoint variable is large near $t = 0$ and small near $t = T$.
The adjoint ODE starts at the only time where we know $\lambda$
directly at $t = T$, and integrates backward to find
$\lambda(t)$ at all earlier times.

This must be understood before the derivation. The backward
direction of the adjoint is not a consequence of the algebra;
the algebra is a consequence of this backward structure.

The adjoint and forward problems are not independent:
the coefficient in the adjoint ODE at every time $t$ involves
$\vec{u}(t)$, the forward solution. The adjoint can only be
solved after the forward problem is solved and $\vec{u}(t)$
is stored at every time point needed.

\begin{selfcheck}
Suppose the response functional is $J = S(T)$ (the susceptible
population at the final time) and the epidemic has a peak at $t = 30$
in a 90-day simulation. At which time ($t = 5$, $t = 30$, or $t = 80$
would you expect $\lambda_I(t)$ to be largest in magnitude?
Explain without any calculation.
\end{selfcheck}
\begin{scAnswer}
$\lambda_I(t)$ at time $t$ measures how sensitive $J = S(T)$ is to a
perturbation of $I$ at time $t$. A perturbation to $I$ at $t = 5$
has 85~days to propagate and affect $S(90)$: a large window
of time for the perturbation to alter the epidemic trajectory.
A perturbation at $t = 80$ has only 10~days.
So $\lambda_I(t)$ should be largest at $t = 5$.
\end{scAnswer}

%%------------------------------------------------------------
\section{The Lagrange Multiplier Derivation}\index{Lagrange multiplier}\index{adjoint method!ODE derivation}
\label{sec:lagrange}
%%------------------------------------------------------------

\subsection{A Warm-Up: Exponential Decay}

Before the five-step derivation, work through the simplest possible case
to see exactly what the adjoint ODE looks like and how it gives the answer.

Let $\dot{u} = -pu$ with $u(0) = 1$ and $J = u(T)$.
The forward solution is $u(t) = e^{-pt}$, so $J = e^{-pT}$ and
$dJ/dp = -T e^{-pT}$ (by direct differentiation --- the ``correct'' answer).

The adjoint approach finds the same result differently.
For $F(u;p) = -pu$ and $g = 0$, $h(u) = u$:
\begin{itemize}[nosep]
  \item Adjoint ODE: $-\dot{\lambda} = F_u\,\lambda = -p\lambda$,
        so $\dot{\lambda} = p\lambda$.
  \item Terminal condition: $\lambda(T) = -h'(u(T)) = -1$.
  \item Solution (running backward from $T$): $\lambda(t) = -e^{p(t-T)}$.
  \item Sensitivity formula: $dJ/dp = -\int_0^T \lambda(t)\,F_p\,dt
        = -\int_0^T (-e^{p(t-T)})(-u(t))\,dt
        = -\int_0^T e^{p(t-T)}e^{-pt}\,dt
        = -\int_0^T e^{-pT}\,dt = -Te^{-pT}$. \checkmark
\end{itemize}
The adjoint gave the same derivative using only two ODE solves (forward,
then backward), with no dependence on $m$ (how many parameters there are).
If $p$ were joined by 49 other parameters, the cost would still be
two ODE solves plus 49 integrals.

Now we derive this procedure in full generality.

\subsection{The Five-Step Derivation}

The derivation proceeds through five numbered steps, each announced
before it is executed. The scalar ODE case is treated fully; the
vector case follows by replacing scalars with vectors and matrices.

\subsection{Setup}

Consider the initial value problem
\begin{equation}
  \dot{u} = F(u; p), \qquad u(0) = u_0,
  \label{eq:fwd_ode}
\end{equation}
with a response functional
\begin{equation}
  J = \int_0^T g(u; p)\,dt + h(u(T)).
  \label{eq:functional}
\end{equation}
Here $g$ is the running cost (an integrand depending on the solution
and parameters) and $h$ is a terminal cost (depending only on the
final state). We want $dJ/dp$ without solving the FSE.

\medskip

The strategy is identical to Chapter~5: introduce an auxiliary variable,
choose it to cancel the unknown sensitivity $\partial u/\partial p$,
and read off the result from what remains.
In Chapter~5 the auxiliary variable was a vector $\vec{v}$ satisfying
$A^T\vec{v} = \vec{c}$; here it is a function $\lambda(t)$ satisfying
a backward ODE.

\medskip
\textbf{Step 1: Introduce the Lagrange multiplier.}

We introduce an auxiliary function $\lambda(t)$, called
the adjoint variable or Lagrange multiplier, and observe that,
because $u(t)$ satisfies the ODE~\eqref{eq:fwd_ode}:
\[
  \int_0^T \lambda(t)\bigl[\dot{u} - F(u;p)\bigr]\,dt = 0.
\]
This is the ODE constraint written as a vanishing integral.
Define the augmented functional
\begin{equation}
  \tilde{J} = J + \int_0^T \lambda(t)\bigl[\dot{u} - F(u;p)\bigr]\,dt.
  \label{eq:augmented}
\end{equation}
Since the bracketed term is zero, $\tilde{J} = J$.
The strategy, carried over from Chapter~5, is to choose
$\lambda(t)$ so that $\partial u/\partial p$ disappears when we
differentiate $\tilde{J}$ with respect to $p$.

\medskip\noindent\textit{Physical analogy.}
Think of a building inspector who records the strain in each beam throughout the day.
The inspector's record is $\lambda(t)$: it does not change the building,
but it carries a running account of how much each moment in time contributes to
the structure's overall stress $J$.
Adding the constraint integral is like attaching a running tally that is identically
zero as long as the beams obey the equilibrium equations --- it costs nothing,
but it opens the door to measuring sensitivity.

\medskip
\textbf{Step 2: Differentiate with respect to $p$.}

Denote $w = \partial u/\partial p$. Differentiating~\eqref{eq:augmented}:
\begin{equation}
  \dd{\tilde{J}}{p}
  = \int_0^T g_u\,w\,dt + h'(u(T))\,w(T)
    + \int_0^T \lambda\bigl[\dot{w} - F_u\,w - F_p\bigr]\,dt.
  \label{eq:dJtilde}
\end{equation}
The unknown $w = \partial u/\partial p$ appears in three places.
The next step eliminates it by choosing $\lambda$ appropriately.

\medskip\noindent\textit{Physical analogy.}
Differentiating with respect to $p$ is like asking: if we nudge a design parameter ---
say, the stiffness of one beam --- how does the overall stress $J$ respond?
The sensitivity $w = \partial u/\partial p$ represents the ripple that nudge sends
through the entire structure over time. It appears everywhere because the system
is coupled; $\lambda(t)$ will be chosen precisely to absorb those ripples.

\medskip
\textbf{Step 3: Integrate by parts.}

The term $\int_0^T \lambda \dot{w}\,dt$ contains $\dot{w}$.
We remove the time derivative from $w$ by integrating by parts:
this is the key step, and it is the ODE analog of the matrix
transposition $v^T A = (A^T v)^T$ in Chapter~5.

\begin{equation}
  \int_0^T \lambda\,\dot{w}\,dt
  = \bigl[\lambda\,w\bigr]_0^T - \int_0^T \dot{\lambda}\,w\,dt
  = \lambda(T)\,w(T) - \lambda(0)\,w(0) - \int_0^T \dot{\lambda}\,w\,dt.
  \label{eq:IBP}
\end{equation}
Substituting into~\eqref{eq:dJtilde}:
\begin{align}
  \dd{\tilde{J}}{p}
  &= \int_0^T g_u\,w\,dt + h'(u(T))\,w(T) \notag\\
  &\quad + \lambda(T)\,w(T) - \lambda(0)\,w(0)
    - \int_0^T \dot{\lambda}\,w\,dt
    - \int_0^T \lambda\,F_u\,w\,dt
    - \int_0^T \lambda\,F_p\,dt.
  \label{eq:dJfull}
\end{align}
Collecting the $w$-terms (the terms involving the unknown sensitivity):
\begin{equation}
  \dd{\tilde{J}}{p}
  = \underbrace{\bigl[h'(u(T)) + \lambda(T)\bigr]w(T)
    - \lambda(0)\,w(0)
    + \int_0^T \bigl[g_u - \dot{\lambda} - \lambda F_u\bigr]w\,dt}_{%
    \text{terms involving } w = \partial u/\partial p}
    - \int_0^T \lambda\,F_p\,dt.
  \label{eq:dJcollected}
\end{equation}

Integration by parts is not merely an algebraic trick.
It is exactly the operation that moves the time derivative
from $w$ to $\lambda$, transforming the FSE into
the adjoint ODE, just as matrix transposition moved
the linear operator from $u$ to $v$ in Chapter~5.

\medskip\noindent\textit{Physical analogy.}
Think of two runners passing a baton around a track.
Forward mode carries $\dot{w}$ --- the baton is held by the sensitivity variable.
Integration by parts hands the baton to $\lambda$: now $\dot{\lambda}$ carries
the time derivative instead, and $w$ appears without any derivative.
This exchange is what makes it possible to choose $\lambda$ to cancel $w$
in the next step; without it, we could not get rid of $\dot{w}$ by specifying
an ODE for $\lambda$ alone.

\medskip
\textbf{Step 4: Choose $\lambda$ to eliminate $w$.}

For the sensitivity formula~\eqref{eq:dJcollected} to be
free of $w$, the coefficient of $w$ must vanish everywhere.
This requires three conditions to hold simultaneously.

First, the integrand coefficient of $w(t)$ must vanish for all $t \in (0,T)$:
\begin{equation}
  -\dot{\lambda} - \lambda F_u + g_u = 0,
  \quad\text{i.e.,}\quad
  \dot{\lambda} = -F_u\,\lambda + g_u.
  \label{eq:adj_wrong_sign}
\end{equation}
Written in standard form (with the sign convention consistent
with backward integration):
\begin{equation}
  \boxed{-\dot{\lambda} = F_u\,\lambda - g_u.}
  \label{eq:adjODE}
\end{equation}
Second, the coefficient of $w(T)$ must vanish:
\begin{equation}
  h'(u(T)) + \lambda(T) = 0,
  \quad\text{i.e.,}\quad
  \lambda(T) = -h'(u(T)).
  \label{eq:termcond}
\end{equation}
Third, the term $\lambda(0)w(0)$ must vanish. Since $u(0) = u_0$
does not depend on $p$, we have $w(0) = \partial u_0/\partial p = 0$,
so this term vanishes automatically.

The terminal condition~\eqref{eq:termcond} does not arise because
we are running the ODE backward. It arises because the response
functional involves the final state $h(u(T))$, and differentiating
this with respect to $p$ produces a boundary term at $t = T$.
This is a fundamental distinction: the terminal condition is imposed
by the response functional, not by the direction of integration.

\medskip\noindent\textit{Physical analogy.}
Choosing $\lambda$ to kill every $w$-term is like tuning the damping in a bridge
so that any vibration induced by a parameter perturbation is immediately cancelled.
The terminal condition $\lambda(T) = -h'(u(T))$ is the inspector's final reading:
at the end of the day the ledger must close, recording exactly how sensitive the
last measured outcome is to the final state. Once $\lambda(t)$ is tuned,
$w$ has nowhere to appear in the sensitivity formula --- it has been designed out.

\medskip
\textbf{Step 5: Read off the sensitivity formula.}

With $\lambda(t)$ satisfying~\eqref{eq:adjODE}--\eqref{eq:termcond},
all terms involving $w$ in~\eqref{eq:dJcollected} vanish and:
\begin{equation}
  \boxed{\dd{J}{p} = -\int_0^T \lambda(t)\,F_p(u(t);p)\,dt
                     - \lambda(0)\,\dd{u_0}{p}.}
  \label{eq:sensfml}
\end{equation}
When the initial condition does not depend on $p$, the second term
vanishes. The sensitivity with respect to any parameter $p_j$ is then:
\begin{equation}
  \dd{J}{p_j} = -\int_0^T \lambda(t)\,\pd{F}{p_j}\,dt.
  \label{eq:sensfml_simple}
\end{equation}
One backward ODE solve for $\lambda(t)$ gives all $m$ sensitivities
via $m$ scalar integrals, exactly the cost advantage we sought.

\medskip\noindent\textit{Physical analogy.}
The inspector's record $\lambda(t)$ is now complete: one backward sweep through
time, consulting the forward solution $u(t)$ at each moment, accumulates
everything needed. Reading off $dJ/dp_j$ is then as simple as computing a
running total --- one scalar integral per parameter, with no additional solves.
The backward sweep is the price paid once; the $m$ integrals are the
dividend collected once for every parameter in the model.

\subsection{The Vector Case}

For the vector system $\dot{\vec{u}} = \vec{F}(\vec{u}; \vec{p})$
with $\vec{u} \in \R^n$ and functional
$J = \int_0^T g(\vec{u};\vec{p})\,dt + h(\vec{u}(T))$,
the derivation is identical with scalars replaced by vectors
and $F_u$ replaced by the Jacobian $J_F = \partial\vec{F}/\partial\vec{u}$.
The adjoint ODE and terminal condition become:
\begin{align}
  -\dot{\vec{\lambda}} &= J_F^T\,\vec{\lambda} - \nabla_u g^T,
  \label{eq:adj_vec} \\
  \vec{\lambda}(T) &= -\nabla_u h^T\big|_{t=T}.
  \label{eq:tc_vec}
\end{align}

\noindent\textbf{Linearity of the adjoint ODE.}
Just as Theorem~\ref{thm:linearFSE} showed that the FSE is always
linear in the forward sensitivity variable $\partial\vec{u}/\partial p_j$,
the adjoint ODE is always linear in $\vec{\lambda}(t)$ ---
even when the forward ODE is nonlinear in $\vec{u}(t)$.
The coefficient $J_F^T(\vec{u}(t))$ is a matrix that depends on the
forward solution, but not on $\vec{\lambda}$ itself.
This linearity is the ODE counterpart of the algebraic linearity
from Chapter~5, and it is why the adjoint is tractable: regardless
of how complex the forward dynamics are, the adjoint ODE is a
linear ODE that can be solved efficiently.

\medskip\noindent
\textbf{Sign convention.}
We write the adjoint in backward-time form as $-\dot{\vec{\lambda}} = \ldots$
to match the MATLAB implementation: \texttt{ode45} with \texttt{tspan = [T, 0]}
integrates backward automatically.
In the sensitivity formula~\eqref{eq:sensfml_vec},
the minus sign is already absorbed: no additional sign change is needed
when reading off $dJ/dp_j$ from the integral.

The sensitivity formula becomes:
\begin{equation}
  \dd{J}{p_j} = -\int_0^T \vec{\lambda}^T\,\pd{\vec{F}}{p_j}\,dt
               - \vec{\lambda}(0)^T\,\pd{\vec{u}_0}{p_j}.
  \label{eq:sensfml_vec}
\end{equation}

\unifyingidea{2} The derivation is now complete, and this is Unifying Idea~2 in its full generality.
The adjoint in Chapter~5 was $A^T\vec{v} = \vec{c}$:
the transposed coefficient matrix $A^T$ applied to $\vec{v}$.
The adjoint here is $-\dot{\vec{\lambda}} = J_F^T\vec{\lambda} - \nabla_u g^T$:
the transposed Jacobian $J_F^T$ applied to $\vec{\lambda}$, with
the sign reversal arising from the integration by parts~\eqref{eq:IBP}.
The integration by parts is the exact functional-analytic analog
of the matrix transposition $A \to A^T$ from Chapter~5.
Appendix~\ref{appsec:IBP} makes this analogy precise through the
Lagrange identity, and Appendix~\ref{appsec:lagrange} derives
the sensitivity formula~\eqref{eq:sensfml_vec} rigorously from it.
The terminal condition~\eqref{eq:tc_vec} is the time-dependent
analog of the right-hand side $\vec{c}$ in $A^T\vec{v} = \vec{c}$.
This is one idea.

\begin{selfcheck}
For the scalar ODE $\dot{u} = -pu$ with $J = u(T)$
(no running cost, $g = 0$, $h(u) = u$):
(a) Write the adjoint ODE and terminal condition using
    equations~\eqref{eq:adjODE} and~\eqref{eq:termcond}.
(b) Solve the adjoint ODE analytically.
(c) Use formula~\eqref{eq:sensfml_simple} to compute $dJ/dp$.
(d) Verify by solving the forward ODE analytically and differentiating.
\end{selfcheck}
\begin{scAnswer}
(a) $F = -pu$, so $F_u = -p$.
Adjoint ODE: $-\dot{\lambda} = -p\lambda$, i.e., $\dot{\lambda} = p\lambda$.
Terminal condition: $\lambda(T) = -h'(u(T)) = -1$.

(b) $\dot{\lambda} = p\lambda$ with $\lambda(T) = -1$.
Integrating forward from $T$ (backward in time):
$\lambda(t) = -e^{p(t-T)}$.

(c) $F_p = -u$. Sensitivity formula:
$dJ/dp = -\int_0^T \lambda(t)(-u(t))\,dt = \int_0^T \lambda(t)u(t)\,dt$.
The forward solution is $u(t) = u_0 e^{-pt}$.
$dJ/dp = \int_0^T (-e^{p(t-T)})(u_0 e^{-pt})\,dt
= -u_0 e^{-pT}\int_0^T dt = -u_0 T e^{-pT}$.

(d) $J = u(T) = u_0 e^{-pT}$, so $dJ/dp = -u_0 T e^{-pT}$. \checkmark
\end{scAnswer}

%%------------------------------------------------------------
\section{The SIR Model --- Full Adjoint Treatment}
\label{sec:sirAdjoint}
%%------------------------------------------------------------

We apply the derivation to the SIR model, writing all adjoint equations
explicitly and implementing the full calculation in MATLAB.

\subsection{Setup and Response Functional}

The SIR model~(4.4)--(4.6) with interpretable \POIs
$\vec{p} = (k, \beta, \tau, L)$:
\begin{align*}
  F_S &= \mu N - \frac{k\beta SI}{N} - \mu S, \\
  F_I &= \frac{k\beta SI}{N} - (\gamma + \mu)I, \\
  F_R &= \gamma I - \mu R,
\end{align*}
where $\gamma = 1/\tau$ and $\mu = 1/L$.
The response functional is
\begin{equation}
  J = \int_0^T I(t)\,dt,
  \label{eq:sirJ}
\end{equation}
the total number of infectious person-days: a measure of
the total disease burden. Here $g(\vec{u}) = I$, $h = 0$.

\subsection{The Jacobian}

The Jacobian $J_F = \partial(F_S, F_I, F_R)/\partial(S, I, R)$ evaluated
at the forward solution $(S(t), I(t), R(t))$:
\begin{equation}
  J_F =
  \begin{pmatrix}
    -\dfrac{k\beta I}{N} - \mu & -\dfrac{k\beta S}{N} & 0 \\[8pt]
     \dfrac{k\beta I}{N} & \dfrac{k\beta S}{N} - \gamma - \mu & 0 \\[8pt]
    0 & \gamma & -\mu
  \end{pmatrix}.
  \label{eq:sirJacobian}
\end{equation}

\subsection{Adjoint ODEs}

With $g_S = 0$, $g_I = 1$, $g_R = 0$ (since $g = I$),
the adjoint system $-\dot{\vec{\lambda}} = J_F^T\vec{\lambda} - \nabla g^T$
written out fully is:
\begin{align}
  -\dot{\lambda}_S &= \left(-\frac{k\beta I}{N} - \mu\right)\lambda_S
                     + \frac{k\beta I}{N}\,\lambda_I, \label{eq:adj_S}\\[4pt]
  -\dot{\lambda}_I &= -\frac{k\beta S}{N}\,\lambda_S
                     + \left(\frac{k\beta S}{N} - \gamma - \mu\right)\lambda_I
                     + \gamma\lambda_R - 1, \label{eq:adj_I}\\[4pt]
  -\dot{\lambda}_R &= -\mu\lambda_R. \label{eq:adj_R}
\end{align}
Terminal conditions (since $h = 0$):
\begin{equation}
  \lambda_S(T) = \lambda_I(T) = \lambda_R(T) = 0.
  \label{eq:sirTC}
\end{equation}
These three ODEs are integrated backward from $t=T$ to $t=0$,
using the stored forward solution $(S(t), I(t), R(t))$
to evaluate the time-varying coefficients.

\subsection{Sensitivity Formulas}

For each \POI, the sensitivity is the integral~\eqref{eq:sensfml_vec}
evaluated using the stored $\vec{\lambda}(t)$:
\begin{align}
  \dd{J}{\beta} &= -\int_0^T \frac{kSI}{N}(\lambda_I - \lambda_S)\,dt,
  \label{eq:dJdbeta}\\[4pt]
  \dd{J}{k}     &= -\int_0^T \frac{\beta SI}{N}(\lambda_I - \lambda_S)\,dt,
  \label{eq:dJdk}\\[4pt]
  \dd{J}{\tau}  &= -\int_0^T \pd{F_I}{\tau}\,\lambda_I\,dt
                 = -\int_0^T \frac{1}{\tau^2}\,I\,\lambda_I\,dt,
  \label{eq:dJdtau}\\[4pt]
  \dd{J}{L}     &= -\int_0^T \left[\pd{F_S}{L}\lambda_S
                               + \pd{F_I}{L}\lambda_I
                               + \pd{F_R}{L}\lambda_R\right]dt.
  \label{eq:dJdL}
\end{align}
Normalized sensitivity indices follow from Definition~2.1:
$\SI{p_j}{J} = (p_j/J)(dJ/dp_j)$.

\medskip\noindent
\textbf{Multiple QOIs and correlated outputs.}
When the study involves $r > 1$ response functionals
--- for example, both $J_1 = \int_0^T I\,dt$ (total burden)
and $J_2 = \max_t I(t)$ (peak infections) ---
the adjoint is solved once per \QOI, giving $r$ backward solutions
$\vec{\lambda}^{(1)}(t), \ldots, \vec{\lambda}^{(r)}(t)$.
Even when the \POIs are independent (uncorrelated inputs),
the \QOIs are almost always correlated outputs: both $J_1$ and $J_2$
depend on the same forward trajectory $\vec{u}(t)$,
so any parameter change affects both simultaneously.
This correlation among outputs is handled automatically by the adjoint.
Each backward solve $\vec{\lambda}^{(k)}$ captures the full sensitivity
of $J_k$ to perturbations of the shared state, without any special
treatment of the output correlation.
The user does not model or correct for QOI correlation ---
the adjoint naturally produces consistent, correct sensitivities
for all $r$ response functionals by virtue of working through
the same forward trajectory.

\subsection{The Adjoint Variable and Its Interpretation}

Figure~\ref{fig:sir_lambdaI} shows $I(t)$ and $\lambda_I(t)$ computed
for nominal parameter values over a 90-day simulation.

\begin{figure}[htb]
\centering
\begin{tikzpicture}
\begin{axis}[
  width=0.78\textwidth, height=7cm,
  xlabel={Time (days)},
  ylabel={},
  xmin=0, xmax=90,
  ymin=-0.05, ymax=1.05,
  grid=major, grid style={gray!20},
  legend pos=outer north east,
  legend style={font=\small},
  xtick={0,15,30,45,60,75,90},
  ytick={0,0.2,0.4,0.6,0.8,1.0},
  yticklabel style={font=\small},
  xticklabel style={font=\small},
  axis y line=left,
  axis x line=bottom,
]
%% I(t): rises to peak ~day 30, then falls (scaled to [0,1])
\addplot[thick, black, domain=0:90, samples=200]
  { max(0, 4*x^2 / (x^2 + 900) * (1 - x/90)^1.5 * 3.5 )};
\addlegendentry{$I(t)/I_{\max}$ (scaled)}

%% lambda_I(t): large near t=0, small near t=T, roughly exponential decay
\addplot[thick, black, dashed, domain=0:90, samples=200]
  {0.95*exp(-0.03*x) + 0.05*(1 - x/90)};
\addlegendentry{$\lambda_I(t)$ (schematic)}

%% Epidemic peak marker
\draw[black!50, thin] (axis cs:28, 0) -- (axis cs:28, 1.02);
\node[above, font=\scriptsize] at (axis cs:28, 1.02) {peak};

\end{axis}
\end{tikzpicture}
\caption{Schematic of $I(t)$ (solid, scaled to $[0,1]$) and the adjoint
         variable $\lambda_I(t)$ (dashed) for the SIR model.
         The adjoint variable is large near $t=0$, when a perturbation
         to the infectious compartment has the most time to affect the
         total burden $J = \int I\,dt$, and decays toward zero as
         $t \to T$.
         This gives $\lambda_I(t)$ a concrete interpretation as the
         marginal value of reducing the infectious population at time $t$:
         interventions are most cost-effective when $\lambda_I$ is large.}
\label{fig:sir_lambdaI}
\end{figure}

The adjoint variable $\lambda_I(t)$ has a concrete epidemiological
interpretation: it measures the marginal benefit of reducing the
infectious population at time $t$.
When $\lambda_I$ is large (early epidemic), reducing $I$ by one
person-day saves many future infections.
When $\lambda_I$ is small (late epidemic), most transmission has
already occurred, so the marginal benefit is low.
This interpretation is available from the adjoint computation alone,
with no additional work.

The terminal condition $\vec{\lambda}(T) = \vec{0}$ is not the same
as starting the forward ODE from initial condition $\vec{0}$ and running it
in reverse. The two problems have different structures: the forward
ODE is driven by $F(\vec{u}; \vec{p})$, while the adjoint ODE
is driven by $J_F^T \vec{\lambda}$ with a forcing term from $g$.
The backward direction comes from the integration by parts, not from
reversing time. The distinction matters for implementation: you run
\texttt{ode45} with \texttt{tspan = [T, 0]}, not with a negated vector field.

\unifyingidea{3} The adjoint variable $\lambda_I(t)$ illustrates Unifying Idea~3: local SA is a linearization.
The adjoint variable displays the
local sensitivity of $J$ throughout the simulation period.
This continuous record of sensitivity is the local SA analog of
what global SA asks globally: how does $J$ respond to perturbations
at each moment in time? Chapter~9 will ask the same question
for perturbations of the parameters across their full uncertainty
ranges rather than perturbations of the state at specific times.
The adjoint provides the linearization; global SA provides
the full nonlinear picture.

\unifyingidea{1} The sensitivity indices $\SI{p_i}{J}$ computed
in this chapter via the adjoint are exactly the normalized derivatives
of Unifying Idea~1 from Chapter~1: the ratio of relative changes in $J$ to relative
changes in each \POI $p_i$. The adjoint is simply a dramatically
more efficient algorithm for computing those same derivatives
when the model has many \POIs and few \QOIs.

\begin{selfcheck}
The adjoint variable $\lambda_R(t)$ for the recovered compartment
satisfies $-\dot{\lambda}_R = -\mu\lambda_R$ with $\lambda_R(T) = 0$.
(a) Solve this ODE analytically (backward from $T$).
(b) What does the solution tell you about the sensitivity of $J$
    to perturbations of the recovered compartment?
\end{selfcheck}
\begin{scAnswer}
(a) $-\dot{\lambda}_R = -\mu\lambda_R$ gives $\dot{\lambda}_R = \mu\lambda_R$
with $\lambda_R(T) = 0$.
Solution: $\lambda_R(t) = 0$ for all $t \in [0,T]$.
(The only solution to a linear ODE with zero terminal condition and
a coefficient that keeps it at zero is identically zero when the
equation has no forcing term.)

(b) $\lambda_R(t) = 0$ means that perturbations to $R(t)$ have no
effect on $J = \int I\,dt$. This makes biological sense: recovered
individuals do not contribute to future transmission (in the basic
SIR model without waning immunity), so changing $R(t)$ does not
change the integral of $I$.
\end{scAnswer}

%%------------------------------------------------------------
\section{Discrete vs.\ Continuous Adjoint}
\label{sec:discretecont}
%%------------------------------------------------------------

The adjoint ODE~\eqref{eq:adj_vec} is the \emph{continuous adjoint}:
it is derived by differentiating the continuous ODE and then discretizing
for numerical integration. An alternative is the \emph{discrete adjoint}:
discretize the ODE first and then adjoint the discrete system.
For linear ODEs with fixed step sizes, the two approaches give
identical results. For nonlinear ODEs with adaptive step-size control,
they can differ.

In practice: if your ODE solver uses fixed step sizes, either adjoint
is acceptable. If it uses adaptive step sizes (as \texttt{ode45} does),
the continuous adjoint requires storing the forward solution densely
enough to interpolate accurately, while the discrete adjoint requires
differentiating through the step-size controller.
For the SA computations in this book, the continuous adjoint with
dense output storage is the more practical approach.
Sirkes and Tziperman~\cite{sirkes1997finite} discuss these subtleties
for geophysical models. Smith~\cite{smith2014uncertainty} develops the
adjoint ODE formalism for general parameter-dependent dynamical systems
in \S9.4, with additional implementation guidance for systems with
stochastic forcing.

\medskip
\noindent\textbf{Neural ODEs.}
In 2018, Chen et al.~\cite{chen2018neural} demonstrated that
the adjoint calculation derived in this chapter can be used to train
certain deep neural networks. The network defines $\vec{F}(\vec{u};\theta)$
for parameters $\theta$; the loss function plays the role of $J$;
and the training algorithm integrates the adjoint ODE backward
to compute $\partial J/\partial\theta_i$ for all network parameters.
The phrase ``backpropagation through time,'' familiar from recurrent
neural networks, refers to the discrete version of this same
backward integration. The continuous limit is precisely the adjoint
ODE derived here.

%%------------------------------------------------------------
\section{MATLAB Laboratory: The SIR Adjoint in Practice}
\label{sec:ch6matlab}
%%------------------------------------------------------------

This laboratory implements the full forward-adjoint computation
for the SIR model and verifies the result against the FSE approach
from Chapter~4.

\textbf{Setup.}
Use the SIR model with \POIs $\vec{p} = (k, \beta, \tau, L)$,
nominal values $k=5$, $\beta=0.06$, $\tau=7$, $L=10{,}000$,
initial conditions $S_0=999$, $I_0=1$, $R_0=0$, $N=1000$,
and $T = 90$~days.
The response functional is $J = \int_0^{90} I(t)\,dt$.

\textbf{Part 1: Solve the forward SIR with dense output.}
Solve the SIR system using \texttt{ode45} with option
\texttt{'Refine', 8} or dense interpolation, storing the solution
on a fine uniform grid with spacing $\Delta t = 0.1$~days.
Store $(S, I, R)$ at every grid point; this array is needed
for the adjoint coefficient evaluation.

\medskip\noindent
\textbf{Interpolation (required before the backward solve).}
Before integrating the adjoint backward, build interpolants of
the forward solution so the adjoint ODE can evaluate $S(t)$,
$I(t)$, and $R(t)$ at any time during the backward pass:
\begin{verbatim}
S_fn = griddedInterpolant(t_fwd, U(:,1), 'pchip');
I_fn = griddedInterpolant(t_fwd, U(:,2), 'pchip');
R_fn = griddedInterpolant(t_fwd, U(:,3), 'pchip');
\end{verbatim}
Pass these interpolants into the adjoint function. Skipping this
step is the most common cause of errors in the adjoint lab.

\textbf{Part 2: Solve the adjoint system backward.}
Implement the adjoint ODEs~\eqref{eq:adj_S}--\eqref{eq:adj_R}.
Note that these depend on the forward solution $(S(t), I(t), R(t))$,
which must be interpolated from the stored grid.

\begin{verbatim}
function dlam = sir_adjoint(t, lam, S_fn, I_fn, R_fn, k, beta, tau, L)
%% Adjoint ODEs for SIR model
%% Integrate BACKWARD: call with tspan = [T, 0]

N     = 1000;
gamma = 1/tau;  mu = 1/L;
S = S_fn(t); I = I_fn(t); R = R_fn(t);   %% interpolate

lS = lam(1); lI = lam(2); lR = lam(3);

%% -d(lambda)/dt = J_F^T * lambda - grad_u(g)
%% Note: g = I, so grad_u(g) = (0, 1, 0)
dlam(1) = -(  (-k*beta*I/N - mu)*lS + (k*beta*I/N)*lI        );
dlam(2) = -(  (-k*beta*S/N)*lS + (k*beta*S/N - gamma - mu)*lI ...
              + gamma*lR - 1  );
dlam(3) = -(  -mu*lR                                          );

dlam = dlam(:);
end
\end{verbatim}

Solve from $t = T$ to $t = 0$ with terminal conditions
$\vec{\lambda}(T) = (0,0,0)$.

\textbf{Part 3: Plot $I(t)$ and $\lambda_I(t)$.}
On a single axis, plot the epidemic curve $I(t)$ (left scale)
and the adjoint variable $\lambda_I(t)$ (right scale).
Annotate the epidemic peak. Confirm that $\lambda_I$ is largest
near $t=0$ and smallest near $t=T$.

\textbf{Part 4: Compute sensitivities and produce a tornado plot.}
Evaluate the integrals~\eqref{eq:dJdbeta}--\eqref{eq:dJdL}
numerically using the trapezoidal rule. Normalize to obtain
$\SI{p_j}{J}$ for all four \POIs. Produce a tornado plot
and identify which parameter has the greatest influence on the
total disease burden.

\textbf{Part 5: Verify against Chapter 4 FSE results.}
From Chapter~4's MATLAB lab, retrieve the SI values computed by
the FSE at a single time point. Compare with the adjoint-based
SIs for $J = \int I\,dt$. Note: the two computations answer
\emph{different} questions (the FSE computes $\SI{p_j}{I(t)}$
at a fixed $t$, while the adjoint computes $\SI{p_j}{\int I\,dt}$),
so the numerical values will differ.
Instead, verify that $\SI{\beta}{J} = \SI{k}{J}$ to at least 4
significant figures (both should equal 1 for the SIR model with
this parameterization) and discuss whether the ratio holds exactly.

\textbf{Part 6: Timing comparison.}
Time the adjoint computation (Parts 1--4) against the FSE approach
from Chapter~4 (solve the 15-equation augmented system four times).
At what number of parameters $m$ would you expect the adjoint
to be faster? Does the timing crossover occur near $m = r = 1$
as the theory predicts?

%%------------------------------------------------------------
\section{Exercises}
%%------------------------------------------------------------

\begin{exercise}[\bloom{L3}]
\label{ex:ch6:scalar}
For the scalar ODE $\dot{u} = r u(1-u/K)$ (logistic) with
$J = \int_0^T u\,dt$ (total population over time):
\begin{enumerate}[label=(\alph*)]
  \item Write the adjoint ODE and terminal condition.
  \item Show that the adjoint ODE is linear in $\lambda(t)$
        even though the forward ODE is nonlinear in $u(t)$.
  \item Explain why the adjoint ODE depends on the forward solution $u(t)$.
  \item In MATLAB, solve the forward and adjoint ODEs simultaneously
        (forward from $0$ to $T$; store; then backward from $T$ to $0$)
        for $r=0.5$, $K=100$, $u_0=10$, $T=20$.
        Plot $u(t)$ and $\lambda(t)$.
\end{enumerate}
\end{exercise}

\begin{exercise}[\bloom{L3}]
\label{ex:ch6:IBP}
State in your own words what integration by parts accomplishes in
Step~3 of the Lagrange multiplier derivation, and why this is
analogous to the matrix transposition $A \to A^T$ in Chapter~5.
Then verify equation~\eqref{eq:IBP} directly by differentiating
the product $\lambda w$ and integrating over $[0,T]$.
\end{exercise}

\begin{exercise}[\bloom{L4}]
\label{ex:ch6:SIRfull}
Using your MATLAB implementation from the laboratory:
\begin{enumerate}[label=(\alph*)]
  \item Compute $\SI{\beta}{J}$, $\SI{k}{J}$, $\SI{\tau}{J}$,
        $\SI{L}{J}$ via the adjoint for $J = \int_0^{90} I\,dt$.
  \item Verify that $\SI{k}{J} = \SI{\beta}{J}$ to at least 6 significant
        figures. Explain this result using the structure of the SIR model.
  \item If you could implement one intervention that reduces a single
        parameter by $20\%$, which parameter would you target to most
        reduce the total disease burden? Justify using the tornado plot.
  \item Compare the computation time with the FSE approach from Chapter~4.
\end{enumerate}
\end{exercise}

\begin{exercise}[\bloom{L4}]
\label{ex:ch6:costtable}
Fill in the computation cost table for the SIR model with
$r=1$ response functional ($J = \int I\,dt$) and $n=3$ state variables:

\begin{center}
\begin{tabular}{lcc}
\toprule
$m$ (number of \POIs) & FSE cost (ODE solves) & Adjoint cost (ODE solves) \\
\midrule
4 & & \\
10 & & \\
50 & & \\
500 & & \\
\bottomrule
\end{tabular}
\end{center}

If each ODE solve takes 30~seconds, at what value of $m$ does the
adjoint become more than 10~times faster than the FSE?
\end{exercise}

\begin{exercise}[\bloom{L4}]
\label{ex:ch6:interpretation}
Explain in plain language, without any equations, why:
\begin{enumerate}[label=(\alph*)]
  \item The adjoint ODE must be integrated backward from $t=T$
        rather than forward from $t=0$.
  \item The adjoint variable $\lambda_I(t)$ for the SIR model
        and response $J = \int I\,dt$ is larger near $t=0$ than near $t=T$.
  \item A public health agency that wants to minimize total disease
        burden should focus interventions on the early phase of the
        epidemic, when $\lambda_I$ is large. Use $\lambda_I(t)$
        to quantify this recommendation.
  \item The terminal condition $\vec{\lambda}(T) = \vec{0}$ is
        not the same as running the forward ODE backward with initial
        condition $\vec{0}$.
\end{enumerate}
\end{exercise}

\begin{exercise}[\bloom{L5}]
\label{ex:ch6:SEIR}
Extend the SIR adjoint treatment to the SEIR model:
\begin{align*}
  \dot{S} &= \mu N - k\beta SI/N - \mu S, \\
  \dot{E} &= k\beta SI/N - (\nu + \mu)E, \\
  \dot{I} &= \nu E - (\gamma + \mu)I, \\
  \dot{R} &= \gamma I - \mu R,
\end{align*}
with \POIs $\vec{p} = (k, \beta, \tau_E, \tau_I, L)$
(interpretable parameterization: $\nu = 1/\tau_E$, $\gamma = 1/\tau_I$,
$\mu = 1/L$) and response functional $J = \int_0^T I(t)\,dt$.

Design the complete adjoint SA study, specifying:
\begin{enumerate}[label=(\alph*)]
  \item The $4\times4$ Jacobian $J_F$ at a generic point
        $(S, E, I, R)$.
  \item All four adjoint ODEs and their terminal conditions.
  \item The sensitivity formulas $dJ/dp_j$ for all five \POIs.
  \item The MATLAB implementation strategy (how you would store
        the forward solution and integrate the adjoint backward).
  \item What additional cost the SEIR adjoint requires compared
        to the SIR adjoint (in terms of ODE solves and storage).
\end{enumerate}
\end{exercise}

\begin{exercise}[\bloom{L6}]
\label{ex:ch6:disagreement}
A climate modeler reports: ``We implemented the adjoint to compute
sensitivities for our model with 200 parameters, but the results
disagree with forward sensitivity analysis by approximately $15\%$
for some parameters. The adjoint must have a bug.''

Identify three distinct causes that could produce this discrepancy
without any coding error. For each cause:
\begin{enumerate}[label=(\alph*)]
  \item State the cause precisely.
  \item Describe a diagnostic test that would confirm or rule it out.
  \item Describe the correction if the cause is confirmed.
\end{enumerate}
\end{exercise}

%%------------------------------------------------------------


\chapter{Adjoint Sensitivity in Practice}

\begin{flushright}
\textit{Reverse the calculation, and one output can speak about many inputs.}
\end{flushright}

\bigskip

You need the sensitivities of a complex disease model with 30~parameters.
The model is your own code, you wrote it yourself, so it is differentiable.
You know from Chapter~6 that the adjoint requires only 2~ODE solves instead
of 30~FSE solves.

But deriving the adjoint equations by hand for 30~coupled nonlinear ODEs would
take weeks and would almost certainly introduce errors in the derivation.

Is there a way to obtain the adjoint computation without deriving the adjoint
equations by hand? Yes, algorithmic differentiation does this automatically~\cite{naumann2012art}.
And there is a way to use it from MATLAB today, without writing
a single AD routine from scratch.

%%------------------------------------------------------------
\section{The Annuity Function: A Worked Comparison}
\label{sec:annuity}
%%------------------------------------------------------------

Before examining the computational tools, we work through a concrete
example where all three methods, forward mode, reverse mode,
and symbolic differentiation, can be compared exactly.

\subsection{The Problem}

An annuity pays a fixed amount $P$ per period for $n$~periods at
interest rate $r$. The present value is
\begin{equation}
  A(P, r, n) = P\,\frac{1 - (1+r)^{-n}}{r}.
  \label{eq:annuity}
\end{equation}
We want all three sensitivity indices: $\SI{P}{A}$, $\SI{r}{A}$,
$\SI{n}{A}$.

\subsection{Forward Mode: Three Passes}

In forward mode we differentiate the computation one parameter at a time,
propagating $\dot{v}_k = \partial v_k/\partial q$ through intermediate
values. We trace the computation of $A$ through the following steps:
\[
  v_1 = (1+r)^{-n}, \quad
  v_2 = 1 - v_1, \quad
  v_3 = v_2/r, \quad
  A = P v_3.
\]

For $q = P$: set $\dot{P} = 1$, $\dot{r} = 0$, $\dot{n} = 0$.
\[
  \dot{v}_1 = 0, \quad \dot{v}_2 = 0, \quad \dot{v}_3 = 0, \quad
  \dot{A} = \dot{P}\,v_3 + P\,\dot{v}_3 = v_3.
\]
So $\partial A/\partial P = v_3 = (1-(1+r)^{-n})/r$.

For $q = r$: set $\dot{r} = 1$, $\dot{P} = 0$, $\dot{n} = 0$.
\[
  \dot{v}_1 = n(1+r)^{-n-1}, \quad
  \dot{v}_2 = -n(1+r)^{-n-1}, \quad
  \dot{v}_3 = \frac{\dot{v}_2\,r - v_2}{r^2}, \quad
  \dot{A} = P\,\dot{v}_3.
\]

For $q = n$: set $\dot{n} = 1$, $\dot{P} = 0$, $\dot{r} = 0$.
\[
  \dot{v}_1 = -\ln(1+r)\,(1+r)^{-n}, \quad
  \dot{v}_2 = \ln(1+r)\,(1+r)^{-n}, \quad
  \dot{v}_3 = \dot{v}_2/r, \quad
  \dot{A} = P\,\dot{v}_3.
\]

Three passes through the computation graph give three derivatives.
Each pass involves the same number of operations as one evaluation
of $A$, so the total cost is approximately $3\times$ the cost
of a single function evaluation.

\subsection{Reverse Mode: One Pass}

In reverse mode, we compute all three derivatives in one backward
pass. Set $\bar{A} = 1$. Work backward through the computation:
\begin{align*}
  \bar{P}  &= \bar{A}\,v_3 = v_3, \\
  \bar{v}_3 &= \bar{A}\,P = P, \\
  \bar{v}_2 &= \bar{v}_3 / r, \\
  \bar{r}   &= \bar{v}_3 \cdot (-v_2/r^2) + \bar{v}_1\cdot\pd{v_1}{r}, \\
  \bar{v}_1 &= -\bar{v}_2, \\
  \bar{n}   &= \bar{v}_1 \cdot \pd{v_1}{n}
             = -\bar{v}_2 \cdot \bigl(-\ln(1+r)(1+r)^{-n}\bigr).
\end{align*}

One backward pass gives $\partial A/\partial P$, $\partial A/\partial r$,
and $\partial A/\partial n$ simultaneously. Total cost: approximately
$2\times$ one function evaluation (one forward pass to store values;
one backward pass to accumulate derivatives), independent of
how many parameters $m$ there are.

\subsection{Cost Summary}

\subsection{Symbolic Differentiation}

As a third comparison, we differentiate the formula~\eqref{eq:annuity}
by hand (or by a computer algebra system). Define $v = (1+r)^{-n}$.
Then $A = P(1-v)/r$, giving:
\[
  \pd{A}{P} = \frac{1-v}{r}, \quad
  \pd{A}{r} = P\!\left(\frac{nv/(1+r) - (1-v)/r}{r}\right),
  \qquad
  \pd{A}{n} = P\,\frac{\ln(1+r)\,v}{r}.
\]
Symbolic differentiation produces exact closed-form expressions that can
be evaluated at any $(P, r, n)$ with a single function evaluation;
no passes through the computation graph are required at query time.
The cost is paid upfront in the derivation; the payoff is that evaluation
is cheap thereafter.

\subsection{Exact Operation Counts}

For the annuity function, we can count arithmetic operations precisely.
The function $A$ requires: 1 addition, 1 exponentiation, 1 subtraction,
2 divisions, 1 multiplication = 6 operations.

Forward mode for one parameter: 6 operations for the primal values,
plus 6 tangent operations = 12 total. For $m = 3$ parameters: 36 operations.

Reverse mode: 12 operations (6 forward + 6 backward) regardless of $m$.
For $m = 3$: still 12 operations.

The ratio: forward mode costs $\approx 2m$ times a single evaluation;
reverse mode costs $\approx 2$ times a single evaluation. For the annuity
with $m = 3$, forward mode uses $3\times$ more operations than reverse mode.
For a climate model with $m = 500$, forward mode uses $250\times$ more.
For a neural network with $m = 10^7$ parameters and one scalar loss function,
reverse mode still requires only one backward pass --- the same cost as the
annuity --- while forward mode would require $10^7$ separate sweeps.
This is why reverse-mode AD made modern deep learning computationally feasible.

\begin{center}
\begin{tabular}{llll}
\toprule
Method & Passes & Operations (annuity, $m=3$) & Scales as \\
\midrule
Forward mode & $m = 3$ & $\approx 36$ & $O(m)$ \\
Reverse mode & 2 & $\approx 12$ & $O(1)$ in $m$ \\
Symbolic diff & 1 (at query time) & $\approx 6$ & $O(1)$ in $m$ \\
Finite diff & $2m + 1 = 7$ evals & $\approx 42$ & $O(m)$, approx \\
\bottomrule
\end{tabular}
\end{center}

At nominal values $P = 1000$, $r = 0.05$, $n = 20$, all methods give:
$\SI{P}{A} = 1$, $\SI{r}{A} \approx -0.604$, $\SI{n}{A} \approx 0.482$.
(The derivation of the analytic formulas is Exercise~7.1.)

\begin{selfcheck}
From the annuity formula~\eqref{eq:annuity}, compute $\SI{P}{A}$ analytically
using the derivative form $\SI{P}{A} = (P/A)(\partial A/\partial P)$.
Verify that it equals $1$ for any nominal values of $P$, $r$, $n$.
\end{selfcheck}
\begin{scAnswer}
$\partial A/\partial P = (1-(1+r)^{-n})/r = A/P$.
So $\SI{P}{A} = (P/A)(A/P) = 1$.

This holds for any $r$, $n$: the present value is linear in $P$,
so a $1\%$ increase in the payment amount produces exactly a $1\%$
increase in present value, regardless of the interest rate or
number of periods.
\end{scAnswer}

%%------------------------------------------------------------
\section{Algorithmic Differentiation}
\label{sec:AD}
%%------------------------------------------------------------

Algorithmic differentiation (AD), also called automatic
differentiation, is the chain rule applied systematically
to every arithmetic operation in a computer program.
It is neither symbolic differentiation (which manipulates formulas)
nor numerical differentiation (which approximates derivatives).
It is exact to machine precision and applies to any differentiable code.

Students sometimes expect AD to produce a formula for the derivative,
like a computer algebra system would. It does not.
AD produces the numerical value of the derivative at a specific
input, by accumulating the chain rule through the sequence of
arithmetic operations the program performs.

Similarly, reverse-mode AD is not the same as numerical differentiation.
A centered-difference estimate of $\partial A/\partial r$ has error
$O(h^2)$ from the truncation of the Taylor series; reverse-mode AD
gives the exact derivative (to floating-point precision) with no
step-size selection required.

\subsection{Forward Mode}

Forward-mode AD attaches a tangent value $\dot{v}_k = \partial v_k/\partial q$
to each intermediate variable $v_k$, for one chosen parameter $q$.
At each arithmetic operation, the tangent is updated using the chain rule:
\[
  v_k = f(v_i, v_j) \implies
  \dot{v}_k = \pd{f}{v_i}\dot{v}_i + \pd{f}{v_j}\dot{v}_j.
\]
At the end of the computation, $\dot{v}_{\text{output}}$ contains
$\partial J/\partial q$. Repeat for each parameter: $m$~passes total.

\subsection{Reverse Mode}

Reverse-mode AD attaches an adjoint value $\bar{v}_k = \partial J/\partial v_k$
to each intermediate variable. In the backward pass, these are accumulated
using the chain rule in reverse:
\[
  v_k = f(v_i, v_j) \implies
  \bar{v}_i \mathrel{+}= \pd{f}{v_i}\bar{v}_k, \quad
  \bar{v}_j \mathrel{+}= \pd{f}{v_j}\bar{v}_k.
\]
The forward pass computes and stores all intermediate values.
The backward pass accumulates all $\partial J/\partial q_i$ simultaneously.
Total cost: one forward pass + one backward pass, regardless of $m$.

Reverse-mode AD is the discrete-computation-graph version
of the adjoint ODE derived in Chapter~6. The adjoint ODE propagates
$\vec{\lambda}(t)$ backward through a continuous system; reverse-mode
AD propagates $\bar{v}_k$ backward through a discrete computation graph.
The underlying mathematical operation is identical: transposition of
the derivative operator.

\unifyingidea{2} Reverse-mode AD implements Unifying Idea~2 automatically:
the adjoint is the transpose in whatever space you are working in, and AD
constructs that transpose through the backward pass without requiring the analyst to derive the adjoint equations by hand.
For a model written as differentiable code, this eliminates the main
practical barrier to adjoint SA: instead of deriving and implementing
the adjoint ODE, the analyst uses an AD library that differentiates
the code directly. The adjoint equation is implicit in the
reverse accumulation; the analyst never writes it explicitly.

\subsection{Available Tools}\index{algorithmic differentiation!software}

The 2009 source material for this book referenced ADIFOR, DAFOR,
and TAMC, tools that are now outdated or unmaintained.
Table~\ref{tab:ADtools} lists the current ecosystem.
For a comprehensive survey of AD methods, see Baydin et al.~\cite{baydin2018automatic}
and Margossian~\cite{margossian2019review}.

\begin{table}[htb]
\centering
\renewcommand{\arraystretch}{1.2}
\begin{tabular}{p{2cm}p{4.2cm}p{2.5cm}p{3.5cm}}
\toprule
Language & Tool & Mode(s) & Notes \\
\midrule
MATLAB & Symbolic Math Toolbox & Symbolic & Exact; limited scalability \\
MATLAB & Deep Learning Toolbox \texttt{dlarray} & Reverse & AD for code; GPU support \\
Python & JAX & Forward + Reverse & NumPy-compatible; JIT compilation \\
Python & PyTorch \texttt{autograd} & Reverse & Widely used in ML \\
Julia & ForwardDiff.jl & Forward & High performance \\
Julia & Zygote.jl & Reverse & Source-to-source; flexible \\
C/Fortran & ADOL-C & Forward + Reverse & Mature; still active \\
C/Fortran & Tapenade & Forward + Reverse & Source transformation \\
\bottomrule
\end{tabular}
\caption{Current AD tools by language (as of 2025).
         For MATLAB users, the Deep Learning Toolbox \texttt{dlarray}
         provides reverse-mode AD for arbitrary differentiable code.
         For large-scale scientific computing, JAX (Python) or ADOL-C
         (C/Fortran) are the most commonly used options.}
\label{tab:ADtools}
\end{table}

\begin{selfcheck}
Explain in one sentence each why: (a)~AD is not symbolic differentiation;
(b)~reverse-mode AD is not the same as centered finite differences;
(c)~forward-mode AD requires $m$ passes while reverse-mode requires 2.
\end{selfcheck}
\begin{scAnswer}
(a) Symbolic differentiation manipulates formulas algebraically;
AD applies the chain rule numerically to the sequence of operations
in the code, producing a number rather than a formula.

(b) Centered finite differences approximate the derivative with error
$O(h^2)$ via $(f(p+h)-f(p-h))/(2h)$; reverse-mode AD computes the
exact derivative (to floating-point precision) by accumulating the
chain rule backward through the computation graph.

(c) Forward mode propagates the tangent $\dot{v}_k = \partial v_k/\partial q$
for one parameter $q$ per pass, requiring one pass per parameter.
Reverse mode propagates the adjoint $\bar{v}_k = \partial J/\partial v_k$
for one response $J$ per backward pass, giving all $m$ derivatives simultaneously.
\end{scAnswer}

%%------------------------------------------------------------
\section{The Master Decision Table}
\label{sec:mastertable}
%%------------------------------------------------------------

With all methods now in hand, Table~\ref{tab:masterSA} consolidates
the decision rules developed across Chapters~3--7.

\begin{table}[htb]
\centering
\footnotesize
\renewcommand{\arraystretch}{1.2}
\begin{tabular}{p{2.8cm}p{1.9cm}p{2.4cm}p{1.8cm}p{3.5cm}}
\toprule
Model type & $m$ vs.\ $r$ & Method & Accuracy & Cost \\
\midrule
Black-box & Any & FD (Ch.~3) & $O(h^2)$ & $2m{+}1$ evals \\
Analytic, simple & $m \leq r$ & FSE (Ch.~4) & Exact & $m$ solves \\
Analytic, diff.\ code & $m \gg r$ & AD, reverse & Exact & 2 passes \\
Analytic, manual & $m \gg r$ & Adjoint (Ch.~6) & Exact & 2 solves \\
Correlated POIs & Any & Reparameterize & --- & Depends \\
Near bifurcation & Any & FSE + caution & Exact* & $m$ solves \\
Large uncertainty & Any & Global SA (Ch.~9) & Variance & $N(m{+}2)$ evals \\
\bottomrule
\end{tabular}
\caption{Master decision table for SA method selection;
         see Pianosi et al.~\cite{pianosi2016sensitivity} for a broader SA survey.
         The primary driver is whether the model is black-box or
         differentiable, and whether $m \gg r$ or $m \leq r$.
         (*Near a bifurcation, the FSE is exact but the SI diverges;
         use the result with caution and consult Chapter~8.)
         \textbf{AD vs analytic adjoint:}
         use reverse-mode AD when the model is implemented in differentiable code
         and memory allows storing the computation tape.
         Use the analytic adjoint (Chapter~6) when the model is a nonlinear ODE or PDE,
         memory is constrained, or the numerical integrator (e.g., adaptive \texttt{ode45})
         is not itself differentiable.}
\label{tab:masterSA}
\end{table}

There is no single universally correct method. The right choice depends
on the model, the available tools, and the question being asked.
A useful heuristic: start with finite differences (Chapter~3) to
get SA results quickly, then switch to the adjoint (Chapter~6 or
AD) if the computation time is unacceptable or if $m$ is large.

One principle applies universally: whatever method you use, always
sanity-check at least two sensitivity indices against an independent
method before trusting the full result. For finite differences,
verify one analytic SI if available. For the adjoint, verify two
values against finite differences. This cross-check catches the
majority of implementation errors.

You should not need to choose between finite differences and the adjoint.
Both serve distinct roles: FD is the right tool for black-box models;
the adjoint is the right tool for white-box (differentiable) models
with many parameters. They are complementary, not competing.
But each tool has failure modes that practitioners encounter repeatedly.

%%------------------------------------------------------------
\section{Practical Guidance and Common Pitfalls}
\label{sec:pitfalls}
%%------------------------------------------------------------

\subsection{Discretize-Then-Adjoint\index{discretize-then-adjoint} vs.\ Adjoint-Then-Discretize}

When computing adjoint sensitivities for ODE systems, two strategies
are available. In the \textbf{adjoint-then-discretize} approach
(used throughout Chapter~6), the adjoint ODE is derived analytically
from the continuous ODE and then solved numerically. In the
\textbf{discretize-then-adjoint} approach, the ODE is first
discretized into a finite-difference recurrence, and the adjoint
is then derived for that discrete system.

For linear ODEs with fixed step sizes, both approaches give identical
results. For nonlinear ODEs with adaptive solvers, they can differ:
the adjoint-then-discretize approach may give sensitivities that
disagree with finite-difference verification because the continuous
adjoint assumes a smooth forward trajectory that the adaptive solver
does not produce. The discretize-then-adjoint approach is always
consistent with finite differences but requires differentiating
through the step-size control logic, which is non-trivial.

For the applications in this book, the adjoint-then-discretize approach
with a fixed-step-size forward solver is the most practical choice.
Sirkes and Tziperman~\cite{sirkes1997finite} discuss the implications
for geophysical models where both approaches are used in practice.

\subsection{Memory Requirements for the Adjoint}\index{checkpointing}

The adjoint ODE requires the forward solution $\vec{u}(t)$ stored at
every time point where the adjoint ODE coefficients are evaluated.
For a system with $n = 100$ state variables solved over $10{,}000$
time steps in double precision, the stored forward solution requires
$100 \times 10{,}000 \times 8$ bytes $= 8$~MB, modest.
For a climate model with $n = 10^6$ state variables and $10^6$
time steps, the requirement is $8 \times 10^{12}$~bytes $= 8$~TB;
infeasible to store in memory.

For large systems, two strategies address this:
\textbf{Checkpointing}\index{checkpointing!adjoint} stores the forward solution only at a subset
of time points and recomputes intermediate values as needed during
the backward pass. \textbf{Adjoint-of-checkpoint} schemes
(Griewank \& Walther~\cite{griewank2008evaluating}) achieve optimal trade-offs
between memory and recomputation cost.
For the models in this course, dense storage is always
adequate.

\subsection{Adaptive Step-Size Solvers}

As noted in Chapter~3, adaptive ODE solvers (such as \texttt{ode45})
introduce non-smoothness in the model output as a function of parameters:
the step-size controller makes discrete decisions that change discontinuously
with the parameter value. The continuous adjoint ODE from Chapter~6
assumes a smooth forward solution; when adaptive solvers introduce
discontinuities, the continuous adjoint may give inaccurate results.

The safest approach: use fixed-step-size integration (e.g., \texttt{ode4}
in MATLAB, a 4th-order Runge-Kutta) for the forward solve when
computing adjoint sensitivities. This eliminates the non-smoothness
and makes the continuous and discrete adjoints equivalent.
Alternatively, use tight tolerances (\texttt{RelTol} $\leq 10^{-8}$)
with adaptive solvers to minimize the non-smoothness.

\begin{selfcheck}
You compute adjoint sensitivities for the SIR model using
\texttt{ode45} with default tolerances for the forward solve,
then verify against centered finite differences.
The adjoint gives $\SI{\beta}{J} = 1.04$ while centered FD gives $0.87$.
Name two possible causes of this discrepancy, and describe
the diagnostic steps you would take to identify which cause applies.
\end{selfcheck}
\begin{scAnswer}
Possible causes: (1)~The adaptive step-size control in \texttt{ode45}
introduces non-smoothness that corrupts the continuous adjoint.
(2)~The finite-difference step size $h$ was chosen poorly (e.g.,
too small, causing catastrophic cancellation, or too large, causing
truncation error). Diagnostics: (a)~Recompute the adjoint with
fixed-step-size \texttt{ode4}: if the discrepancy disappears,
the solver is the culprit. (b)~Plot the U-curve for the FD estimate
to verify the step size is optimal; if the FD value changes
significantly with $h$, the FD is unreliable.
Always use finite differences as a sanity check for any adjoint
implementation; a discrepancy above $1\%$ warrants investigation.
\end{scAnswer}

\subsection{Discontinuities Break AD}\index{algorithmic differentiation!discontinuities}

AD requires the model to be differentiable at the nominal parameter values.
Common sources of non-differentiability in practice include \texttt{if} statements
whose conditions depend on parameters,\footnote{%
  \textit{If-statements are where smooth calculus goes to lose its temper.}
  When an \texttt{if} condition switches branches as a parameter is perturbed,
  the forward and reverse passes follow different computational paths,
  and the resulting ``derivative'' is the derivative of whichever branch
  happened to be active, which may be zero or undefined at the transition.}
\texttt{min} and \texttt{max} functions,
piecewise-defined model components such as quarantine thresholds or saturation
functions with hard cutoffs, and look-up tables interpolated without smoothing.
When these are present, AD will silently compute the derivative of one
branch, which may be zero or undefined at the branch point.
The diagnostic: compare AD results against centered finite differences.
If they disagree significantly for small $h$, suspect a discontinuity.
Replacing piecewise functions with smooth approximations (sigmoid
functions, smooth max/min) restores differentiability.

\begin{selfcheck}
A disease model contains the piecewise function:
$\beta(I) = \beta_0$ if $I < 0.1N$, else $\beta_0/2$.
This represents reduced contact when hospitalizations are high.
(a) Is this function differentiable with respect to $I$ everywhere?
(b) If you run reverse-mode AD through this function at $I = 0.09N$,
what derivative will AD report? Is this reliable?
(c) Propose a smooth approximation that would make the model
fully differentiable.
\end{selfcheck}
\begin{scAnswer}
(a) No. The function has a jump discontinuity at $I = 0.1N$,
so $\partial\beta/\partial I$ is undefined there.
(b) At $I = 0.09N$, the \texttt{if} condition is false ($I < 0.1N$),
so AD will compute the derivative of $\beta_0$ with respect to $I$,
which is zero. This is locally correct for $I$ strictly below the
threshold, but gives no information about the discontinuity itself,
and will be completely wrong for sensitivity with respect to $I$ near
the threshold. It is unreliable as a global SA result.
(c) Replace the step with a sigmoid:
$\beta(I) = \beta_0\bigl(1 - 0.5\,\sigma(k(I/N - 0.1))\bigr)$,
where $\sigma(x) = 1/(1+e^{-x})$ and $k > 0$ controls sharpness.
As $k \to \infty$ this recovers the discontinuous original;
a moderate $k$ (e.g., $k = 100$) approximates it closely while
remaining differentiable everywhere.
\end{scAnswer}

\unifyingidea{1} Every method in this chapter (forward mode,
reverse mode, symbolic differentiation, finite differences, the adjoint
ODE) computes the same object: Unifying Idea~1, the sensitivity index
defined in Chapter~2. The methods differ in computational cost,
accuracy, and applicability to different model types, but the result they
produce is always the same number: $\SI{p_j}{J} = (p_j/J)(dJ/dp_j)$.

%%------------------------------------------------------------
\section{MATLAB Laboratory: Four Methods, One Model}
\label{sec:ch7matlab}
%%------------------------------------------------------------

This laboratory computes the same sensitivity ($\partial J/\partial\beta$
for the SIR model with $J = \int_0^{90} I\,dt$, using all four methods
and compares their accuracy, speed, and ease of implementation.

\textbf{Part 1: Finite differences (Chapter~3 method).}
Use the \texttt{sensitivity\_jacobian} template from Chapter~3 with the
SIR model. Record the estimated $\SI{\beta}{J}$ and the computation time.

\textbf{Part 2: Analytic FSE (Chapter~4 method).}
Solve the 15-equation augmented system from Chapter~4 for the $\beta$
parameter alone. Extract $\partial I(t)/\partial\beta$ at all time points
and integrate numerically. Record the result and time.

\textbf{Part 3: Adjoint ODE (Chapter~6 method).}
Solve the adjoint system from Chapter~6. Use
formula~\eqref{eq:sensfml_vec} to compute $dJ/d\beta$.
Record the result and time.

\textbf{Part 4: MATLAB Symbolic Toolbox.}
Define the SIR ODE symbolically and compute the sensitivity
of the equilibrium $R_0 = k\beta/\gamma$ with respect to $\beta$
as a sanity check. Note: the Symbolic Toolbox computes sensitivities
of closed-form expressions, not of ODE solutions over time.

\begin{verbatim}
syms k beta gamma positive
R0   = k * beta / gamma;
SI_b = simplify( (beta/R0) * diff(R0, beta) )

%% For the time-dependent sensitivity via dlarray:
%% (requires Deep Learning Toolbox)
p_nom = [5, 0.06, 7, 10000];    %% [k, beta, tau, L]
p_dl  = dlarray(p_nom);
J_fn  = @(p) SIR_integral(p);   %% function returning integral of I
dJdp  = dlgradient(J_fn(p_dl), p_dl);
fprintf('AD sensitivity for beta: %.6f\n', extractdata(dJdp(2)));
\end{verbatim}

\textbf{Part 5: Comparison table.}
Fill in the following table from your measurements:

\begin{center}
\begin{tabular}{lcccc}
\toprule
Method & $\SI{\beta}{J}$ & Rel.\ error & Time (s) & Lines of new code \\
\midrule
Finite differences & & & & \\
Analytic FSE & & & & \\
Adjoint ODE & & & & \\
AD (\texttt{dlarray}) & & & & \\
\bottomrule
\end{tabular}
\end{center}

Discuss: which method is easiest to implement? Most accurate?
Fastest? Which would you use in practice for this model?

\textbf{Extension: introducing a discontinuity.}
Modify the SIR model to include a treatment threshold: if $I(t) > 0.1N$,
reduce $\tau$ by $20\%$ (faster treatment when infection is high).
Re-run all four methods. Which methods give consistent results?
Which methods silently produce wrong answers? How does the U-curve
(Chapter~3) change for finite differences near $\beta$?

%%------------------------------------------------------------
\paragraph{Connection to optimization.} The adjoint computed in this
chapter is the same object as the Lagrange multiplier in constrained
optimization and the costate variable in optimal control.
Appendix~\ref{app:optimization} develops this connection through
LP duality and Pontryagin's minimum principle.

%%------------------------------------------------------------
\section{Exercises}
%%------------------------------------------------------------

\begin{exercise}[\bloom{L3}]
\label{ex:ch7:annuity}
For the annuity function~\eqref{eq:annuity} at nominal values
$P = 1000$, $r = 0.05$, $n = 20$:
\begin{enumerate}[label=(\alph*)]
  \item Compute $A$ numerically.
  \item Compute all three sensitivities $\SI{P}{A}$, $\SI{r}{A}$, $\SI{n}{A}$
        analytically by differentiating~\eqref{eq:annuity}.
  \item Implement forward mode by hand for $q = r$ (as in Section~7.1)
        and verify your answer from (b).
  \item Implement reverse mode for all three parameters in one pass
        and verify all three answers.
  \item Count the number of arithmetic operations in each approach.
\end{enumerate}
\end{exercise}

\begin{exercise}[\bloom{L4}]
\label{ex:ch7:decisiontable}
For each of the following SA problems, select the most appropriate
method from Table~\ref{tab:masterSA} and justify in two sentences:
\begin{enumerate}[label=(\alph*)]
  \item A climate model with 800 parameters written in Fortran.
        Each model run takes 10~hours. You want sensitivities for
        2 response functionals.
  \item A pharmacokinetic model with 5 parameters and analytic
        closed-form solutions. You want sensitivities for 8 drug
        concentration time points.
  \item A traffic simulation with 40~parameters implemented as a
        compiled C++ executable. Parameter values can be set via
        a configuration file. Each run takes 3~minutes.
  \item A 10-parameter epidemic model written in MATLAB with
        \texttt{ode45}. You want one response functional.
        The model contains an \texttt{if} statement that triggers
        a quarantine intervention when $I > 0.05N$.
\end{enumerate}
\end{exercise}

\begin{exercise}[\bloom{L4}]
\label{ex:ch7:dlarray}
Using MATLAB's Deep Learning Toolbox \texttt{dlarray}:
\begin{enumerate}[label=(\alph*)]
  \item Wrap the function $R_0(k, \beta, \tau) = k\beta\tau$ as a
        differentiable computation and use \texttt{dlgradient} to compute
        all three partial derivatives at $k=5$, $\beta=0.06$, $\tau=7$.
  \item Compare with the analytic derivatives.
  \item Wrap the SIR equilibrium calculation (solve $\vec{f}(\vec{u}) = 0$
        numerically) and compute $\partial\bar{I}/\partial\beta$ at the
        endemic equilibrium. Compare with the FSE result from Chapter~4.
  \item What happens when you try to apply \texttt{dlgradient} through
        a call to \texttt{ode45}? Explain why, and what the remedy is.
\end{enumerate}
\end{exercise}

\begin{exercise}[\bloom{L5}]
\label{ex:ch7:workflow}
Design the complete SA workflow for a dengue fever transmission
model~\cite{chitnis2008determining} with the following properties:
$n = 8$ state variables (human and mosquito compartments),
$m = 14$ parameters (mortality rates, biting rates, transmission
probabilities, mean durations), $r = 2$ response functionals
(basic reproduction number and endemic human infection prevalence),
each ODE solve taking approximately 2~minutes.

Specify:
\begin{enumerate}[label=(\alph*)]
  \item Which SA method you would use and why.
  \item The expected total computation time.
  \item How you would validate the implementation.
  \item How you would handle the fact that two parameters
        ($\sigma_v$ and $\sigma_h$, the mosquito and human biting rates)
        are structurally correlated.
  \item What output format you would use to present results to
        a non-mathematical public health audience.
\end{enumerate}
\end{exercise}

\begin{exercise}[\bloom{L6}]
\label{ex:ch7:disagree}
A research group uses reverse-mode AD to compute adjoint sensitivities
for a mosquito population model with adaptive-step-size ODE integration.
They report that some sensitivity indices agree with finite differences
but others disagree by $10$--$20\%$, with no obvious pattern.
The discrepancies are largest for parameters that appear in the
mortality rate expressions.
\begin{enumerate}[label=(\alph*)]
  \item Identify the most likely technical cause of the discrepancy.
  \item Describe two diagnostic tests that would confirm or rule out
        this cause.
  \item Propose two remedies and explain the trade-offs between them.
  \item At what point would you conclude that the disagreement indicates
        a genuine SA result (the linearization is poor for these parameters)
        rather than a numerical artifact?
\end{enumerate}
\end{exercise}

%%------------------------------------------------------------



%% ---- Part V: Limitations and Extensions ----
\part{Limitations and Extensions}

\chapter{Caveat Emptor --- Limitations of Local Sensitivity Analysis}

\begin{flushright}
\textit{Local sensitivity is a linear truth, and nature is rarely linear for long.}
\end{flushright}

\bigskip

In 2019, Saltelli et al.~\cite{saltelli2019why} reviewed sensitivity analyses
published in 280 highly cited papers across top environmental modeling journals.
More than half used the wrong method for the question being asked.
Specifically, studies that asked global questions, how does output
uncertainty relate to input uncertainty across the full parameter range?;
used local methods: derivatives computed at a single nominal parameter value.
The derivatives were computed correctly. But they answered a different question
than the one the authors claimed to answer.

The deeper problem is not computational but epistemological: sensitivity
analysis is still often defined by method rather than purpose
\cite{razavi2021future}.
Practitioners reach for the method from their own disciplinary tradition
rather than the one best suited to their question.
A local method applied to answer a global question will produce results
that are internally correct but wrong for the stated purpose.
Before choosing an SA method, two questions must be answered: what is the
goal --- factor prioritization (which inputs most reduce output variance?),
factor fixing (which inputs can be frozen without changing the output?),
or factor mapping (which inputs drive the output above a threshold?) ---
and does the method's assumptions match the model's behavior?

\medskip
\begin{center}
\textit{Assuming local sensitivity holds globally is the triumph\\
of mathematical hope over nonlinear reality.}
\end{center}
\medskip

This chapter asks the question that should be asked before any SA is begun:
is local SA the right tool for this question?

%%------------------------------------------------------------
\section{Method-Question Mismatch}\index{method-question mismatch}
\label{sec:mismatch}
%%------------------------------------------------------------

Local sensitivity analysis is built on a derivative. A derivative
is a local object: it measures the behavior of a function at one point,
and it is only meaningful in the vicinity of that point.
This makes local SA powerful for some questions and completely wrong
for others.

Local SA can answer: \emph{which parameter most affects the model
near the nominal parameter values?}

Local SA cannot answer: \emph{which parameter most affects the model
across its full range of uncertainty?} That question requires the parameter
to actually vary across its range, and a derivative does not do that.

Local SA cannot answer: \emph{how does the model behave at parameter
values far from the nominal?} The derivative approximates linear behavior;
real models are often highly nonlinear.

The principle is simple and unforgiving. The method must match the question.
Using local SA to answer a global question is not merely suboptimal;
it can produce a ranking of parameter importance that is entirely wrong.
A parameter that appears unimportant locally may dominate the output
variance when varied across its full uncertainty range.
A parameter that appears important locally may have negligible effect
at realistic perturbation sizes if the model is highly nonlinear
in that parameter.

A small sensitivity index does not mean a parameter can be ignored globally.
It means the model is approximately linear in that parameter near the nominal
value, and the linear coefficient happens to be small.
Section~\ref{sec:globalneeded} shows a specific example where the local
ranking is the opposite of the global ranking.

\begin{selfcheck}
A published SA study reports: ``All sensitivity indices are less than 0.05,
confirming that the model is robust to parameter uncertainty.''
The study used nominal parameter values from the literature, with
no information on parameter uncertainty ranges.
Identify two specific claims in this statement that go beyond what
local SA can actually establish.
\end{selfcheck}
\begin{scAnswer}
(1) ``Robust to parameter uncertainty'': local SA says nothing about
parameter uncertainty. It describes behavior near the nominal value.
If the uncertainty range for some parameters is large (e.g., $\pm50\%$),
the model could be highly sensitive to that parameter globally even
though the local SI is small.

(2) The conclusion depends implicitly on the nominal parameter values
chosen. If the nominal values are changed, the local SI values change
too (for nonlinear models). A study with different nominal values
might find very different rankings. This sensitivity of local SA
to the choice of nominal values is itself a limitation that the study
does not address.
\end{scAnswer}

%%------------------------------------------------------------
\section{Bifurcations and SA Failure}\index{bifurcation!SA failure near}
\label{sec:bifurcations}
%%------------------------------------------------------------

Local SA requires the model output to be differentiable with respect
to the parameters at the nominal values. When this fails, local SA
fails completely. The most important way it fails is at a bifurcation.

A bifurcation is a parameter value at which the qualitative behavior
of the model changes: two equilibria merge into one, a stable fixed
point becomes unstable, or an epidemic threshold is crossed~\cite{strogatz2018nonlinear,kuznetsov2004elements}.
At a bifurcation point, the derivative $\partial u/\partial p$ is
either undefined or infinite. Local SA breaks down entirely.
The precise mathematical reason is given by the Fredholm alternative
(Appendix~\ref{appsec:solvability}): near a bifurcation the Jacobian
becomes singular, so the FSE $D_{\vec{u}}[\vec{f}]\,(\partial\vec{u}/\partial p) = \vec{r}$
may have no solution or infinitely many.
For models with periodic limit cycles, sensitivity near a bifurcation
of the periodic orbit involves Floquet multipliers; this is treated
in Appendix~\ref{app:floquet}.

\subsection{The Quadratic Example}

Consider the scalar equation $f(u; p) = u^2 + p = 0$.
For $p < 0$, there are two real roots: $u = \pm\sqrt{-p}$.
At $p = 0$, the two roots merge. For $p > 0$, there are no real roots.

The FSE for the sensitivity of the positive root $u^* = \sqrt{-p}$ is:
\[
  \pd{u^*}{p} = \frac{-\pd{f}{p}}{\pd{f}{u}} = \frac{-1}{2u^*}
              = \frac{-1}{2\sqrt{-p}}.
\]
The normalized SI is:
\[
  \SI{p}{u^*} = \frac{p}{u^*}\,\pd{u^*}{p}
              = \frac{p}{\sqrt{-p}}\cdot\frac{-1}{2\sqrt{-p}}
              = \frac{1}{2}.
\]
This is constant (independent of $p$) for all $p < 0$;
the positive root is always in exact proportion $1/2$ to any
fractional change in $p$.

But as $p \to 0^-$:
\[
  \pd{u^*}{p} = \frac{-1}{2\sqrt{-p}} \to -\infty.
\]
The derivative diverges at the bifurcation. Near $p = 0$, the
model is infinitely sensitive to the parameter $p$: any small
positive perturbation to $p$ destroys the equilibrium entirely.

\unifyingidea{3} This is Unifying Idea~3 in its most dramatic form.
The sensitivity index is the derivative-based linearization of the
model response. At a bifurcation, the function $u^*(p)$ has
infinite curvature, its graph has a vertical tangent, and
no linear approximation is valid.
The linearization does not just become inaccurate at the bifurcation;
it becomes undefined.

\unifyingidea{1} The failure at a bifurcation is a failure of the
normalized derivative. The sensitivity index $\SI{p}{q}$ exists only
when $\partial q/\partial p$ is finite. Near a bifurcation it is not,
and the index either blows up or vanishes depending on which side
of the threshold you are computing from.

\unifyingidea{2} The adjoint SA system inherits exactly the same
failure mode. The adjoint ODE is driven by the Jacobian $J_F^T$
of the forward model. Near a bifurcation, the Jacobian has a
near-zero eigenvalue, and the adjoint solution grows without bound
over the integration interval. Both the forward and adjoint
approaches fail at the same bifurcation for the same reason:
the transpose of a singular matrix is also singular.

\begin{selfcheck}
For $f(u;p) = u^2 + p$ with $p = -4$:
(a) Find the two roots $u^*_\pm$.
(b) Compute $\SI{p}{u^*_+}$ analytically. Compare with the derivative $\partial u^*_+/\partial p$.
(c) Estimate $\partial u^*_+/\partial p$ numerically using a centered difference
    with $h = 0.1, 0.01, 0.001$. At what value of $p$ near zero does
    the finite-difference estimate become unreliable?
\end{selfcheck}
\begin{scAnswer}
(a) $u^*_\pm = \pm\sqrt{4} = \pm 2$.
(b) $\partial u^*_+/\partial p = -1/(2\sqrt{-p}) = -1/4$.
$\SI{p}{u^*_+} = (p/u^*_+)(-1/(2\sqrt{-p})) = (-4/2)(-1/4) = 1/2$.

(c) Numerically at $p = -4$: all three step sizes give $\approx -0.250$ \checkmark.
As $p \to 0^-$, the function $u^*_+ = \sqrt{-p}$ has a square-root singularity;
finite differences become inaccurate when the curvature term dominates the
step-size error, which occurs when $h \gtrsim |p|$. Near $p = 0$, even very
small $h$ is not small compared to $|p|$, so finite differences fail.
\end{scAnswer}

\subsection{The SIR Model at the Epidemic Threshold}

The SIR model has a transcritical bifurcation at $R_0 = 1$:
for $R_0 < 1$, the only equilibrium is the disease-free state;
for $R_0 > 1$, the endemic equilibrium appears.

The endemic equilibrium fraction
$i^* = 1 - 1/R_0$
exists only for $R_0 > 1$. Its sensitivity with respect to $R_0$ is:
\[
  \pd{i^*}{R_0} = \frac{1}{R_0^2}, \qquad
  \SI{R_0}{i^*} = \frac{R_0}{i^*}\cdot\frac{1}{R_0^2}
                = \frac{1}{R_0(1-1/R_0)} = \frac{1}{R_0 - 1}.
\]
As $R_0 \to 1^+$, $\SI{R_0}{i^*} \to +\infty$: near the epidemic
threshold, the endemic equilibrium is infinitely sensitive to changes
in $R_0$. This is exactly the bifurcation behavior identified above.

SA of the endemic equilibrium for $R_0$ values close to 1 must be
interpreted with extreme caution. A small change in any parameter
that affects $R_0$ could eliminate the endemic equilibrium entirely
or cause it to surge. Local SA correctly identifies this extreme
sensitivity but cannot predict the qualitative change.

For $R_0 \leq 1$: SA of the endemic equilibrium is undefined;
the equilibrium does not exist. Computing SA under the assumption
that $i^*$ exists when $R_0 < 1$ is a mathematical error.
MATLAB will simply return a negative or complex value without warning.

\medskip\noindent
\textbf{Operational rule: when to stop reporting SI as a ranking.}
A sensitivity index is meaningful as a ranking tool only when
the linearization is a fair local approximation.
Two practical checks:

\begin{enumerate}[nosep]
\item \textit{Magnitude check.} If $|\SI{p}{q}| > 10$ for a normalized
      index, treat the result as a qualitative warning
      (``this parameter region is highly sensitive'') rather than a
      quantitative ranking. An SI of 100 does not mean the parameter
      matters 10$\times$ more than one with SI 10;
      it may mean the model is near a threshold.
\item \textit{Condition-number check.} Compute $\kappa(J_F)$
      (the condition number of the model Jacobian at the operating point).
      If $\kappa(J_F) > 10^4$, the sensitivity computation is
      ill-conditioned: small errors in the model evaluation
      may be amplified by $\kappa(J_F)$ in the sensitivity values.
      Always report $\kappa(J_F)$ alongside SI values computed
      near a bifurcation or near parameter degeneracy.
\end{enumerate}

%%------------------------------------------------------------
\section{Ill-Conditioning}\index{ill-conditioning}
\label{sec:illconditioning}
%%------------------------------------------------------------

For linear algebraic models, SA failure takes a different form.
When the coefficient matrix $A$ in $A\vec{u} = \vec{b}$ is
nearly singular, small changes in the right-hand side $\vec{b}$
(or in the matrix entries) produce large changes in the solution.
The sensitivity indices are large not because the model is interesting
but because the model is poorly identified.

\subsection{The Condition Number}

The \textbf{condition number}\index{condition number} $\kappa(A) = \|A\|\|A^{-1}\|$
(see Appendix~\ref{appsec:SVD} for the SVD-based definition $\kappa_2 = \sigma_1/\sigma_r$)
measures how much the solution $\vec{u}$ can change relative to
a change in the data. For the 2-norm,
\begin{equation}
  \kappa_2(A) = \frac{\sigma_{\max}(A)}{\sigma_{\min}(A)},
  \label{eq:condnum}
\end{equation}
where $\sigma_{\max}$ and $\sigma_{\min}$ are the largest and smallest
singular values of $A$. A matrix with $\kappa_2(A) = 1$ is perfectly
conditioned (orthogonal); a matrix approaching singularity has
$\sigma_{\min} \to 0$ and $\kappa_2(A) \to \infty$.

The key bound: if the data changes by a relative amount $\epsilon$,
the solution changes by at most $\kappa(A)\,\epsilon$:
\[
  \frac{\|\delta\vec{u}\|}{\|\vec{u}\|} \leq \kappa(A)\,
  \frac{\|\delta\vec{b}\|}{\|\vec{b}\|}.
\]
When $\kappa(A) \gg 1$, even tiny changes in the data produce
large changes in the solution.

\subsection{SA and Ill-Conditioning}

For SA, the key insight is that large sensitivity indices are not
necessarily a property of the model's \emph{true} behavior.
They may simply reflect ill-conditioning: the fact that the linear
system is nearly singular.

A colleague who reports ``the sensitivity of $u_1$ to $b_1$
is extremely large'' may be reporting either:
(a) a genuine model property (the output truly varies strongly
    with this parameter), or
(b) a numerical artifact (the linear system is nearly singular
    and small input changes have large numerical effects on the solution).

Distinguishing (a) from (b) requires computing $\kappa(A)$.
If $\kappa(A) \approx 10^6$, then numerical errors of size
$\varepsilon_{\text{mach}} \approx 10^{-16}$ can contaminate the
solution at the level of $10^{-10}$: 6 digits of the solution
may be meaningless.

Ill-conditioning is itself a SA result: it warns the investigator
that the model parameters are nearly unidentifiable from data,
and that predictions cannot be trusted without more precise
parameter measurements.
This is not a failure of SA; it is one of the most important
things SA can reveal.

\begin{selfcheck}
The $2\times2$ matrix $A(\epsilon) = \begin{pmatrix} 1 & 1 \\ 1 & 1+\epsilon \end{pmatrix}$
with $\epsilon > 0$:
(a) Compute $\det(A(\epsilon))$ and show $A \to$ singular as $\epsilon \to 0$.
(b) Compute the condition number $\kappa_2(A)$ for $\epsilon = 1, 0.1, 0.01$.
(c) For $\vec{b} = (2, 2+\epsilon)^T$, solve $A\vec{u} = \vec{b}$ and compute
    $\SI{\epsilon}{u_1}$. What happens to this SI as $\epsilon \to 0$?
\end{selfcheck}
\begin{scAnswer}
(a) $\det = \epsilon \to 0$ as $\epsilon \to 0$. Singular at $\epsilon = 0$.
(b) $\kappa_2(A) = \sigma_{\max}/\sigma_{\min}$. For $\epsilon = 1$: $\kappa \approx 5.83$.
For $\epsilon = 0.1$: $\kappa \approx 40$. For $\epsilon = 0.01$: $\kappa \approx 400$.
As $\epsilon \to 0$: $\kappa \to \infty$.

(c) $A^{-1} = \epsilon^{-1}\begin{pmatrix}1+\epsilon & -1\\ -1 & 1\end{pmatrix}$.
$\vec{u} = A^{-1}\vec{b} = (1, 1)^T$ for all $\epsilon > 0$.
$\SI{\epsilon}{u_1}$: differentiate $u_1 = (2(1+\epsilon) - (2+\epsilon))/\epsilon
= \epsilon/\epsilon = 1$ with respect to $\epsilon$, but the simplification
$u_1 = 1$ holds exactly, so $\SI{\epsilon}{u_1} = 0$. The sensitivity vanishes!

This illustrates a subtlety: large condition number does not always imply large SI.
What the condition number warns is that \emph{numerical} computation of $\vec{u}$
will be inaccurate; the exact mathematical SI may nonetheless be finite or zero.
\end{scAnswer}

%%------------------------------------------------------------
\section{Separatrices and Trajectory Sensitivity}\index{separatrix}\index{trajectory sensitivity}
\label{sec:separatrices}
%%------------------------------------------------------------

Dynamical systems often have regions of qualitatively different behavior
separated by a \textbf{separatrix}: a boundary in state space that
trajectories cannot cross. SA computed on one side of a separatrix
describes behavior qualitatively different from SA on the other side.

For the SIR model, the relevant separatrix is the epidemic threshold
$S(0) > \gamma N / (k\beta)$ (equivalently, $R_0 > 1$). If the
initial susceptible population exceeds this threshold, an epidemic
occurs. If it does not, the infection dies out without spreading.

Suppose SA is performed at a parameter value where an epidemic occurs
($R_0 = 1.5$) and the sensitivity index $\SI{\beta}{I_{\max}}$
is computed. This index describes how the epidemic peak changes with
a small perturbation to $\beta$ within the epidemic regime.
It says nothing about what happens if $\beta$ is reduced below the
threshold: in that regime, $I_{\max}$ is not a smooth function of $\beta$
it drops to zero discontinuously as the separatrix is crossed.

The practical implication: if any parameter perturbation of interest
crosses a qualitative threshold (endemic vs.\ disease-free, outbreak
vs.\ no outbreak, stable vs.\ oscillatory), local SA cannot be used to
predict the response. The model must be analyzed separately in each
qualitative regime, and the boundary itself requires special treatment.

A related limitation applies to models with multiple stable attractors.
When a model can settle into qualitatively different long-run states
depending on initial conditions, the sensitivity index computed near one
attractor says nothing about behavior near another.
Extensions of the SIR model that exhibit bistability (where both a
disease-free and an endemic equilibrium are stable for the same parameter
values) require separate SA studies for each attractor.
Local SA at one equilibrium cannot warn the analyst that another
equilibrium exists.

%%------------------------------------------------------------
\section{When Global SA Is Needed}\index{global sensitivity analysis!when needed}
\label{sec:globalneeded}
%%------------------------------------------------------------

\begin{flushright}
\textit{\small Local SA studies the neighborhood. Global SA studies the country.}
\end{flushright}

Local SA fails to answer the right question in three specific
situations. When any of these conditions holds, Chapter~9's global
SA methods are required. The practical decision framework for
choosing between local and global SA is developed at length by
Saltelli et al.~\cite{saltelli2004sensitivity,saltelli2006sensitivity}.

To reiterate the misconception that opened this chapter:
a sensitivity index near zero does not mean a parameter is globally
unimportant. It means the linearization at the nominal value is
nearly flat. What happens away from the nominal, across the
full uncertainty range of the parameter, requires global SA to determine.

\subsection{Large Parameter Uncertainty}

When a parameter is known only to within a factor of 2 or 3;
common in epidemiology, ecology, and environmental modeling;
the linearization underlying local SA may be very poor across
the plausible range.

As a concrete example, consider a model where the \QOI depends
on a parameter $p$ through $q(p) = e^p$. At nominal value $p^* = 0$:
$\SI{p}{q} = p^*(dq/dp)/q = 0 \cdot 1/1 = 0$. The local SI is zero.
But if the uncertainty range is $[-2, 2]$, the \QOI varies from
$e^{-2} \approx 0.14$ to $e^2 \approx 7.4$: a factor of 53.
The parameter is globally important but invisible to local SA.

\subsection{Highly Nonlinear Models}

Two models with identical local SIs for a parameter $p$ may
behave very differently when $p$ is perturbed by $20\%$:
one model may be nearly linear (the SI is a good approximation),
while the other may be highly curved (the SI understates the true change).

The sigma-normalized sensitivity index (Section~9.1) provides
the minimal bridge: it weights the local SI by the parameter's
standard deviation, giving a measure that at least accounts for
the relative magnitudes of parameter uncertainties.

\subsection{Correlated Parameters}

Local SA under the independence assumption gives misleading results
when parameters are correlated (Chapter~3). Global SA methods in
Chapter~9 include specific tools for handling this case.

\subsection{The Portfolio Example}

Consider a model $Y = s_1 Q_1 + s_2 Q_2$ with $s_1 = 2$, $s_2 = 1$
and parameter standard deviations $\sigma_1 = 1$, $\sigma_2 = 3$.

Local SA gives $\SI{Q_1}{Y} = 2 > \SI{Q_2}{Y} = 1$:
parameter~1 appears more important.

The sigma-normalized index:
$S_i^\sigma = (s_i\sigma_i)/\sigma_Y$,
where $\sigma_Y = \sqrt{4+9} = \sqrt{13}$.
So $S_1^\sigma = 2/\sqrt{13} \approx 0.55$ and $S_2^\sigma = 3/\sqrt{13} \approx 0.83$.

Parameter~2 dominates the output variance, even though it has a smaller
local SI. The local ranking is inverted because parameter~2 has much
larger uncertainty. Chapter~9's Sobol indices provide the rigorous
global version of this sigma-normalized assessment.

\begin{selfcheck}
For the SIR model, suppose $\beta$ has uncertainty $\sigma_\beta = 0.02$
and $\tau$ has uncertainty $\sigma_\tau = 3$~days, at nominal values
$\beta^* = 0.06$ and $\tau^* = 7$~days.

(a) From Chapter~2, $\SI{\beta}{R_0} = \SI{\tau}{R_0} = 1$.
   Which parameter appears more important from local SA alone?

(b) Compute $\sigma_\beta/\beta^*$ and $\sigma_\tau/\tau^*$ (the relative uncertainties).
   Which parameter has larger relative uncertainty?

(c) Compute the sigma-normalized indices for $R_0$ with respect to
    $\beta$ and $\tau$. Which parameter drives more variance in $R_0$?
\end{selfcheck}
\begin{scAnswer}
(a) Both have SI = 1, so local SA alone says they are equally important.

(b) $\sigma_\beta/\beta^* = 0.02/0.06 \approx 0.33$ (33\%);
$\sigma_\tau/\tau^* = 3/7 \approx 0.43$ (43\%). The mean infectious
period has larger relative uncertainty.

(c) $S_\beta^\sigma = \SI{\beta}{R_0} \cdot (\sigma_\beta/\beta^*) = 1 \times 0.33 = 0.33$.
$S_\tau^\sigma = \SI{\tau}{R_0} \cdot (\sigma_\tau/\tau^*) = 1 \times 0.43 = 0.43$.
The mean infectious period drives more variance in $R_0$, despite both
having the same local SI. Chapter~9 will quantify this more precisely
with Sobol indices.
\end{scAnswer}

%%------------------------------------------------------------
\section{When SA Guides Data Collection}
\label{sec:ch8datacollection}
%%------------------------------------------------------------

The three chapters preceding this one have identified when SA
can mislead. This section completes the practical argument
the book has been building since Chapter~1: how to use SA
\emph{correctly} to make decisions about field measurement.

\textbf{The decision chain.}
A sensitivity study produces a ranked list of parameters.
Translating that list into a field-measurement plan requires
three additional steps that SA alone cannot provide:

\begin{enumerate}
\item \textit{Identifiability:} Can the high-sensitivity parameter
      actually be estimated from the available data, or is it
      structurally confounded with another parameter?
      Section~\ref{sec:correlatedPOIs} showed that parameters
      with identical sensitivity indices are not separately identifiable
      from outbreak data alone; measuring one is equivalent
      to measuring the other.
      For the SIR model, $k$ and $\beta$ always appear as the product
      $k\beta$, so field studies should estimate $k\beta$ directly
      (e.g., from secondary attack rates) rather than $k$ and $\beta$
      separately.

\item \textit{Measurability:} Is the high-sensitivity parameter
      actually measurable with available instruments and budget?
      The mean infectious period $\tau$ appears in the top tier
      of every SIR sensitivity analysis, and it is measurable:
      clinicians record the duration from symptom onset to recovery.
      The lifespan $L$ has negligible SI for short-horizon models
      and is measurable anyway from demographic records.
      A parameter with large SI but no feasible measurement pathway
      should be \emph{reduced} rather than measured, by redesigning
      the intervention.

\item \textit{Actionability:} Can the high-sensitivity parameter
      be reduced by a feasible public health intervention?
      SA of the SIR model with $J = \max_t I(t)$ consistently identifies
      $k$ (contact rate) and $\beta$ (transmission probability) as the
      most influential parameters.
      Contact rates can be reduced by social distancing, school closures,
      and remote-work policies.
      Transmission probability per contact can be reduced by mask
      requirements, hand-hygiene campaigns, and vaccination.
      These two routes correspond to different interventions with
      different costs, compliance rates, and social acceptability.
      SA identifies both as equally valuable targets; the choice between
      them requires policy judgment beyond the model.
\end{enumerate}

\textbf{The two-budget problem.}
A common practical scenario: a government agency has funding to
(1) measure exactly $k_{\text{meas}}$ parameters precisely in the field,
and (2) implement exactly $k_{\text{int}}$ interventions.
The correct SA-guided procedure is:

\begin{enumerate}[nosep]
\item Run local SA (or Morris screening for large parameter sets)
      to rank all parameters by $|\SI{p_j}{\text{peak}~I}|$.
\item From the top-ranked parameters, apply the identifiability filter:
      remove any parameter that is structurally confounded
      with a higher-ranked parameter.
\item Apply the measurability filter: remove any parameter for which
      no feasible measurement protocol exists within the budget.
\item The remaining top $k_{\text{meas}}$ parameters are the
      measurement priorities.
\item From the ranked list, apply the actionability filter separately:
      the top $k_{\text{int}}$ actionable parameters are the
      intervention priorities. These sets may overlap.
\end{enumerate}

\textbf{What the model cannot tell you.}
The complete picture requires one more piece: the second semester.
Smith's \\emph{Uncertainty Quantification}~\\cite{smith2014uncertainty}
Chapters~7--8 covers Bayesian parameter estimation --- the step
that actually quantifies how much a given measurement budget
reduces parameter uncertainty.
SA tells you \emph{which} parameters matter most; parameter estimation
tells you \emph{how precisely} you can learn them from data;
optimal experimental design (Smith Chapter~11) tells you
\emph{how to design the measurement study} to learn those parameters
as efficiently as possible.

These three disciplines --- SA, parameter estimation, and optimal
experimental design --- form a complete pipeline for model-based
decision making. This book has equipped you for the first stage.

\medskip\noindent
\textit{If a parameter is influential and poorly known,
that is not a nuisance. That is your research agenda.}

%%------------------------------------------------------------
\section{MATLAB Laboratory: Bifurcation and Ill-Conditioning}
\label{sec:ch8matlab}
%%------------------------------------------------------------

\textbf{Part 1: The quadratic bifurcation.}
For $f(u;p) = u^2 + p$, compute the positive root $u^* = \sqrt{-p}$
and the SI $\SI{p}{u^*} = 1/2$ analytically.
Plot $u^*$ vs.\ $p$ for $p \in [-4, -0.01]$.
Estimate $\partial u^*/\partial p$ by centered differences at
$p = -4, -1, -0.1, -0.01, -0.001$.
Compare the finite-difference estimates with the analytic value.
At what value of $p$ does the finite-difference estimate become unreliable?

\textbf{Part 2: SIR bifurcation diagram using MATCONT.}
Download and install MATCONT (\url{https://matcont.sourceforge.io}).
Define the SIR equilibrium equations as a MATCONT system.
Use MATCONT's continuation capability to trace the endemic equilibrium
$i^* = 1 - 1/R_0$ as $R_0$ increases from 0.5 to 3.0.
MATCONT will automatically detect the transcritical bifurcation at $R_0 = 1$
and flag it. Plot the bifurcation diagram (equilibrium value vs.\ $R_0$).
If MATCONT is unavailable, trace the equilibrium manually:
for $R_0 \leq 1$, $i^* = 0$ (disease-free); for $R_0 > 1$, $i^* = 1 - 1/R_0$.

\textbf{Part 3: FSA comparison across $R_0 = 1$.}
Using the full SIR ODE system, compute $\SI{\beta}{I(T)}$ at
$T = 90$~days for $R_0 = 0.8, 1.0, 1.2, 2.0$.
At each value, use centered finite differences.
Tabulate the results. At $R_0 = 1$: does the SI blow up, return NaN,
or behave normally? Try $h = 0.01, 0.001, 0.0001$ to see what happens.

\textbf{Part 4: Condition number of the SIR Jacobian as a function of $R_0$.}
For the SIR endemic equilibrium, the Jacobian $J_F$ at $(S^*, I^*, R^*)$
depends on $R_0$. Compute $\kappa_2(J_F)$ using MATLAB's \texttt{cond}
for $R_0 \in \{1.01, 1.05, 1.1, 1.5, 2.0, 3.0\}$.
Plot $\kappa_2(J_F)$ vs.\ $R_0$.
Observe the spike as $R_0 \to 1^+$: the Jacobian becomes nearly singular
at the bifurcation. What does this imply about SA of the SIR endemic equilibrium
for parameter values that place $R_0$ close to 1?

\textbf{Part 5: Sigma-normalized SIs for the SIR model.}
Using the SIR model with uncertainty ranges $\sigma_k = 1$
(contacts/day), $\sigma_\beta = 0.02$, $\sigma_\tau = 3$~days,
$\sigma_L = 500$~days (all at $1\sigma$), compute the sigma-normalized
indices for $R_0$. Compare the ranking with the local SI ranking
from Chapter~2. Which parameter most affects $R_0$ given its
realistic uncertainty range?

%%------------------------------------------------------------
\section{Exercises}
%%------------------------------------------------------------

\begin{exercise}[\bloom{L3}]
\label{ex:ch8:quadratic}
For the equation $g(u; p) = u^3 - 3pu + 2 = 0$:
\begin{enumerate}[label=(\alph*)]
  \item Find all real roots at $p = 1$. Which root changes most rapidly with $p$?
  \item Compute $\SI{p}{u^*}$ for each real root at $p = 1$
        using the implicit function theorem.
  \item Plot the roots as functions of $p$ for $p \in [0.8, 1.2]$.
        At what value of $p$ do two roots merge?
  \item At the bifurcation value of $p$, what does the FSE predict?
        Verify numerically using finite differences.
\end{enumerate}
\end{exercise}

\begin{exercise}[\bloom{L3}]
\label{ex:ch8:SIRthreshold}
Using the SIR model with $k = 5$, $\tau = 7$~days, $L = 10{,}000$~days:
\begin{enumerate}[label=(\alph*)]
  \item Compute $\SI{\beta}{I(T)}$ for $\beta = 0.04, 0.06, 0.08, 0.12$
        at $T = 90$~days using centered finite differences.
        Note: $R_0 = k\beta\tau$ takes values $1.4, 2.1, 2.8, 4.2$.
  \item Now recompute for $\beta = 0.02$ ($R_0 = 0.7$).
        Compare with the results for $R_0 > 1$. Explain any qualitative differences.
  \item If you were to add a bifurcation analysis tool (e.g., MATCONT),
        what would you expect the bifurcation diagram to look like for
        the endemic equilibrium of this SIR model as $\beta$ varies?
\end{enumerate}
\end{exercise}

\begin{exercise}[\bloom{L4}]
\label{ex:ch8:conditioning}
For the $3\times3$ linear system $A\vec{u} = \vec{b}$ where:
\[
  A = \begin{pmatrix} 1 & 2 & 3 \\ 4 & 5 & 6 \\ 7 & 8 & 9+\epsilon \end{pmatrix},
  \quad \vec{b} = A\vec{u}_{\text{exact}}, \quad \vec{u}_{\text{exact}} = (1,1,1)^T.
\]
\begin{enumerate}[label=(\alph*)]
  \item For $\epsilon = 1, 0.1, 0.01, 0.001, 0.0001$, compute $\kappa(A)$
        and $\|\hat{\vec{u}} - \vec{u}_{\text{exact}}\|$ where
        $\hat{\vec{u}} = A\backslash\vec{b}$ in MATLAB.
  \item Plot $\|\hat{\vec{u}} - \vec{u}_{\text{exact}}\|$ vs.\ $\kappa(A)$
        on log-log axes. What is the approximate slope?
  \item What does this experiment tell you about interpreting large SI values
        for nearly singular models?
\end{enumerate}
\end{exercise}

\begin{exercise}[\bloom{L5}]
\label{ex:ch8:report}
You are asked to perform SA on an epidemic model whose best-fit
parameter values give $R_0 = 1.04$, very close to the epidemic threshold.
Write a one-page advisory report for a public health agency explaining:
\begin{enumerate}[label=(\alph*)]
  \item What local SA can and cannot tell you at this parameter value.
  \item What specific numerical problems to expect when computing
        sensitivity indices near $R_0 = 1$.
  \item What additional analyses you would recommend before advising
        on interventions.
  \item How you would explain the concept of a bifurcation and its
        implications to a non-mathematical policy maker.
\end{enumerate}
\end{exercise}

\begin{exercise}[\bloom{L6}]
\label{ex:ch8:critique}
A published modeling paper reports: ``Sensitivity analysis of the
endemic equilibrium of our malaria model found that all indices
are less than $0.1$, indicating the model is robust to parameter
uncertainty. The basic reproduction number at fitted parameter
values is $R_0 = 1.15$.''

Identify at least three specific methodological concerns with
this analysis and its conclusions. For each concern, state:
(a) what the problem is; (b) what evidence the authors would need
to provide to address it; and (c) what alternative analysis
would give a more complete picture.
\end{exercise}

%%------------------------------------------------------------


\chapter{Global Sensitivity Analysis and Sobol Indices}

\begin{flushright}
\textit{The point is not to remove uncertainty.\\
The point is to find out which uncertainty is worth worrying about.}
\end{flushright}

\bigskip

Consider two epidemiological models. Both give $\SI{\beta}{R_0} = +1$
at the same nominal parameter values. In Model~A, $R_0$ is nearly linear
in $\beta$ across the full biologically plausible range $[0.01, 0.15]$.
In Model~B, $R_0$ is highly nonlinear: it is steep near the lower end
of the range and flat at the upper end.

For Model~A, the local SI correctly identifies $\beta$ as an important
parameter. For Model~B, local SA gives the same SI value of $+1$;
but it is now describing only the slope at the nominal point.
A $20\%$ increase in $\beta$ from the nominal produces very different
effects in the two models, and no derivative computed at the nominal
value can see this.

Chapter~8 showed when local SA fails. This chapter provides the tools
to replace it.

%%------------------------------------------------------------
\section{Why Local SA Is Not Enough: The Sigma-Normalized Bridge}
\label{sec:sigmanorm}
%%------------------------------------------------------------

\begin{flushright}
\textit{\small Local SA studies the neighborhood. Global SA studies the country.}
\end{flushright}

The two can give very different answers about which parameters matter most,
as the following example shows.

\subsection{The Portfolio Example}

Consider a model $Y = s_1 Q_1 + s_2 Q_2$ where $s_1 = 2$, $s_2 = 1$
and the parameters have standard deviations $\sigma_1 = 1$, $\sigma_2 = 3$.

Local SA gives $\SI{Q_1}{Y} = 2 > \SI{Q_2}{Y} = 1$:
parameter~1 appears more important. But the parameters have very different
uncertainty ranges, and $Q_2$ varies three times as much as $Q_1$.

The output variance is:
\[
  \Var(Y) = s_1^2\sigma_1^2 + s_2^2\sigma_2^2 = 4 + 9 = 13
  \quad\implies\quad \sigma_Y = \sqrt{13}.
\]
The \textbf{sigma-normalized sensitivity index}\index{sigma-normalized index}
\begin{equation}
  S_i^\sigma = \frac{\sigma_i}{\sigma_Y}\,\pd{Y}{Q_i}
  \label{eq:sigmanorm}
\end{equation}
accounts for both the derivative and the parameter uncertainty:
\[
  S_1^\sigma = \frac{1}{\sqrt{13}}\cdot 2 \approx 0.55, \qquad
  S_2^\sigma = \frac{3}{\sqrt{13}}\cdot 1 \approx 0.83.
\]
Parameter~2 dominates, opposite to the local SA ranking.
The local ranking was misleading because it ignored the factor-of-3
difference in parameter uncertainty.

\unifyingidea{1} The sigma-normalized index is still the sensitivity
index from Chapter~2, multiplied by the relative parameter uncertainty.
Global SA builds on local SA; it does not replace it. The local SI
from Chapter~2 appears inside every global SA formula in this chapter.

\subsection{Two SIR Models}

For the SIR model, suppose two outbreak settings give the same
$\SI{\beta}{R_0} = 1$. In Setting~A, $\beta$ is tightly constrained
from contact-tracing data ($\sigma_\beta = 0.005$), while in Setting~B,
$\beta$ must be estimated from seroprevalence data ($\sigma_\beta = 0.03$).
Local SA says $\beta$ is equally important in both settings.
The sigma-normalized index says $\beta$ in Setting~B is six times more
important, in terms of its contribution to uncertainty in $R_0$.
The second semester will show how to propagate these uncertainties
through the full model.

\begin{selfcheck}
For the SIR model with local SIs $\SI{k}{R_0} = \SI{\beta}{R_0} = \SI{\tau}{R_0} = 1$
and parameter standard deviations $\sigma_k = 0.5$, $\sigma_\beta = 0.01$,
$\sigma_\tau = 2$, with $\sigma_{R_0} = \sqrt{\sigma_k^2 + \sigma_\beta^2 + \sigma_\tau^2}/R_0^*$:

For a SIR model with $k^* = 5$, $\beta^* = 0.06$, $\tau^* = 7$
($R_0^* = 2.1$), first compute $\sigma_{R_0}$ from the delta method, then
rank the three parameters by their sigma-normalized index $S_i^\sigma$.
\end{selfcheck}
\begin{scAnswer}
Delta method: $\Var(R_0) \approx \sum_i (\partial R_0/\partial p_i)^2\sigma_i^2$.
$\partial R_0/\partial k = \beta\tau = 0.42$; $\partial R_0/\partial\beta = k\tau = 35$;
$\partial R_0/\partial\tau = k\beta = 0.3$.
$\Var(R_0) \approx 0.42^2(0.25) + 35^2(0.0001) + 0.3^2(4)
= 0.0441 + 0.1225 + 0.36 = 0.527$.
$\sigma_{R_0} \approx 0.726$.

$S_k^\sigma = (0.5/0.726)(0.42/2.1)(k/k)\cdot$... Easier via the formula directly:
$S_i^\sigma = (\partial R_0/\partial p_i)\sigma_i/\sigma_{R_0}$.
$S_k^\sigma = 0.42 \times 0.5/0.726 \approx 0.29$.
$S_\beta^\sigma = 35 \times 0.01/0.726 \approx 0.48$.
$S_\tau^\sigma = 0.3 \times 2/0.726 \approx 0.83$.

Ranking: $\tau > \beta > k$. The mean infectious period drives the most
uncertainty in $R_0$, despite all three having equal local SIs.
\end{scAnswer}

%%------------------------------------------------------------
\section{Sampling Methods}
\label{sec:sampling}
%%------------------------------------------------------------

All global SA methods require evaluating the model at many parameter
combinations. Choosing those combinations efficiently is the first
computational task.

\subsection{Random Sampling}

The simplest approach: draw $N$~independent samples from the joint
parameter distribution (often assumed uniform or independent log-normal
over plausible ranges). For $N$ large, random sampling provides good
coverage on average, but can leave regions of parameter space unsampled
by chance.

\subsection{Latin Hypercube Sampling}

Latin Hypercube Sampling (LHS) divides each parameter's range into
$N$ equal strata and ensures that exactly one sample falls in each
stratum for every parameter~\cite{mckay1979comparison}. This eliminates the clustering that
random sampling can produce and improves coverage with smaller $N$.

\begin{definition}[LHS Sample]\index{Latin hypercube sampling (LHS)}
An LHS of size $N$ for $m$ parameters with ranges $[l_j, u_j]$,
$j = 1,\ldots,m$, is an $N\times m$ matrix $X$ such that each column
contains exactly one value from each of the $N$ equal-width strata
$[l_j + (k-1)(u_j-l_j)/N,\ l_j + k(u_j-l_j)/N]$ for $k = 1,\ldots,N$.
\end{definition}

The MATLAB implementation generates LHS by permuting the strata
independently for each parameter and adding uniform jitter within each stratum:

\begin{verbatim}
%% From: src/globalsa/lhs_sample.m (see GitHub repository)
rng(42);                                   %% set seed before sampling
X = lhs_sample(N, m, lb, ub);             %% N-by-m sample matrix
%% Uses lhsdesign (Statistics Toolbox) with maximin criterion
\end{verbatim}

For the same $N$, LHS provides substantially better coverage than
random sampling, particularly for estimating marginal effects of
individual parameters; see Helton and Davis~\cite{helton2003latin}
for convergence analysis and uncertainty propagation results.
For Sobol index estimation (Section~\ref{sec:sobol}),
LHS is not used directly; instead, the Saltelli quasi-random construction
is preferred.

\begin{selfcheck}
For $N = 4$ samples and $m = 1$ parameter on $[0, 1]$:
(a) Draw a random sample of size 4. Is it possible that all 4 points
    fall in $[0, 0.5]$? Is this possible with LHS?
(b) Write out explicitly what an LHS of size 4 looks like for one parameter.
\end{selfcheck}
\begin{scAnswer}
(a) Yes, random sampling can place all 4 points in $[0, 0.5]$ (probability $(1/2)^4 = 1/16$).
With LHS, this is impossible: each quarter-interval $[0,0.25)$, $[0.25,0.5)$,
$[0.5,0.75)$, $[0.75,1]$ must contain exactly one sample.

(b) One possible LHS: one point in each of $[0,0.25)$, $[0.25,0.5)$,
$[0.5,0.75)$, $[0.75,1]$, with uniform jitter within each. Example:
$X = \{0.18, 0.43, 0.62, 0.91\}$, one value per stratum,
randomly positioned within each.
\end{scAnswer}

%%------------------------------------------------------------
\section{Regression-Based SA: LHS/PRCC}
\label{sec:prcc}
%%------------------------------------------------------------

The Partial Rank Correlation Coefficient (PRCC)\index{PRCC (partial rank correlation coefficient)} is the dominant
global SA method in mathematical epidemiology.
It is computationally cheap, handles monotone (but not necessarily linear)
relationships, and the Marino et al.~\cite{marino2008methodology} MATLAB implementation
is widely used and freely available; it builds on the rank-correlation
approach of Iman and Conover~\cite{iman1982distribution} and the
epidemiological SA tradition of Blower and Dowlatabadi~\cite{blower1994sensitivity}.

\subsection{The PRCC Calculation}

Given $N$ LHS samples of the $m$ \POIs and the corresponding $N$
model outputs:

\textbf{Step 1.} Rank-transform all inputs and the output separately.
Replace each value with its rank among the $N$ observations.

\textbf{Step 2.} Fit a linear regression of the rank-transformed output
on all rank-transformed inputs.

\textbf{Step 3.} For each parameter $p_j$, compute the partial correlation
coefficient between the rank-transformed $p_j$ and the rank-transformed output,
controlling for all other parameters. This is the PRCC for $p_j$.

\textbf{Step 4.} Compute a $t$-test $p$-value for each PRCC to assess
statistical significance.

The PRCC ranges from $-1$ to $+1$. Its sign has the same interpretation
as the SI sign rule from Chapter~2: positive means POI and QOI increase
together; negative means they move in opposite directions.

\subsection{When PRCC is Appropriate}

PRCC requires monotone relationships: the QOI must consistently increase
(or consistently decrease) as each POI increases, across the full
range of parameter values. It does not require linearity.
For most compartmental epidemic models, the relationships are approximately
monotone, which is why PRCC dominates the field.

Before applying PRCC, verify the monotonicity assumption by producing a
scatter matrix: plot the model output against each \POI for the full
LHS sample using MATLAB's \texttt{plotmatrix} function.
A roughly monotone cloud (upward or downward sloping, without reversals)
supports the use of PRCC. A hump-shaped or U-shaped cloud signals
non-monotonicity, and Sobol indices should be used instead.

PRCC and Sobol indices are not equivalent. PRCC assumes monotone
relationships; Sobol indices make no such assumption. When a relationship
is nonmonotone (a QOI first increases then decreases as a parameter
increases), PRCC may give a near-zero coefficient even for a highly
important parameter, while the Sobol total index will correctly identify
it as important. For the models in this course, both methods typically
agree, but verifying agreement is a useful consistency check.

%%------------------------------------------------------------
\section{Sobol Decomposition and Variance-Based SA}
\label{sec:sobol}
%%------------------------------------------------------------

Variance-based SA is the most rigorous framework for global SA:
it quantifies exactly how much of the total output variance is
attributable to each \POI or group of \POIs~\cite{sobol2001global}.

\subsection{The Intuition: Conditional Expectation}

Suppose parameter $Q_1$ is fixed at $q_1$ while all other parameters
vary across their ranges. The expected output $\E[Y\,|\,Q_1 = q_1]$
is some function of $q_1$. If $Q_1$ is important, this conditional
expectation varies substantially as $q_1$ changes.
The variance of this conditional expectation,
$\Var_{Q_1}[\E_{Q_{-1}}(Y\,|\,Q_1)]$, measures how much of $Y$'s
variability is explained by $Q_1$ alone. Normalizing by the total
variance gives the first-order Sobol index.

\subsection{The ANOVA Decomposition}\index{ANOVA decomposition}

\textbf{Why does a unique decomposition exist?}
Because statistical independence allows us to cleanly separate
what each parameter contributes alone from what it contributes jointly.
When the $Q_i$ are independent, projecting $f$ onto the subspace
of functions of $Q_i$ alone gives a well-defined component $f_i(Q_i)$,
and these projections do not interfere with each other.
The orthogonality $\int_0^1 f_i(Q_i)\,dQ_i = 0$ is the mathematical
expression of this separation: each component has zero mean after
integrating out $Q_i$, so distinct components contribute zero covariance.
When parameters are correlated, this orthogonality breaks down,
and the decomposition loses its uniqueness
(Section~\ref{sec:correlatedsobol} returns to this point).
With independence, the decomposition is both unique and additive in variance.

For independent \POIs $\vec{Q} = (Q_1,\ldots,Q_m)$ uniformly
distributed on $[0,1]^m$, any square-integrable model $Y = f(\vec{Q})$
admits the unique \textbf{ANOVA decomposition}:
\begin{equation}
  f(\vec{Q}) = f_0 + \sum_i f_i(Q_i)
             + \sum_{i<j} f_{ij}(Q_i, Q_j) + \cdots + f_{12\cdots m}(\vec{Q}),
  \label{eq:ANOVA}
\end{equation}
where $f_0 = \E[Y]$, $f_i(Q_i) = \E[Y\,|\,Q_i] - f_0$, and the
higher-order terms capture interactions. Uniqueness follows from the
orthogonality conditions $\int_0^1 f_{u}(\vec{Q}_u)\,dQ_i = 0$
for each $i \in u$; see Sobol~\cite{sobol2001global} Theorem~1
and Saltelli et al.~\cite{saltelli2008global} Appendix~A for a proof.
The terms are orthogonal, so their variances add:
\begin{equation}
  D = \Var(Y) = \sum_i D_i + \sum_{i<j} D_{ij} + \cdots + D_{12\cdots m},
  \label{eq:variancedec}
\end{equation}
where $D_i = \Var[\E(Y\,|\,Q_i)]$ and $D_{ij} = \Var[\E(Y\,|\,Q_i,Q_j)] - D_i - D_j$.

\medskip\noindent
\textbf{Worked example: a model with an interaction term.}
Consider $f(Q_1, Q_2) = Q_1 + Q_2 + Q_1 Q_2$ with $Q_1, Q_2 \sim \mathcal{U}(0,1)$
independent. Compute the ANOVA components directly from the definitions:
\begin{align*}
  f_0 &= \E[Y] = \tfrac{1}{2} + \tfrac{1}{2} + \tfrac{1}{4} = \tfrac{5}{4},\\
  f_1(Q_1) &= \E[Y\,|\,Q_1] - f_0
    = \bigl(Q_1 + \tfrac{1}{2} + \tfrac{Q_1}{2}\bigr) - \tfrac{5}{4}
    = \tfrac{3Q_1}{2} - \tfrac{3}{4} = \tfrac{3}{2}\bigl(Q_1 - \tfrac{1}{2}\bigr),\\
  f_2(Q_2) &= \E[Y\,|\,Q_2] - f_0
    = \tfrac{3}{2}\bigl(Q_2 - \tfrac{1}{2}\bigr) \quad\text{(by symmetry)},\\
  f_{12}(Q_1,Q_2) &= f - f_0 - f_1 - f_2
    = Q_1 Q_2 - \tfrac{Q_1}{2} - \tfrac{Q_2}{2} + \tfrac{1}{4}
    = \bigl(Q_1 - \tfrac{1}{2}\bigr)\bigl(Q_2 - \tfrac{1}{2}\bigr).
\end{align*}
Verify orthogonality: $\int_0^1 f_{12}\,dQ_1 = \bigl(\int_0^1(Q_1-\tfrac{1}{2})\,dQ_1\bigr)(Q_2-\tfrac{1}{2}) = 0$. \checkmark

The variances are $D_1 = D_2 = \tfrac{3}{4} \cdot \Var(Q_1) = \tfrac{3}{4}\cdot\tfrac{1}{12} = \tfrac{1}{16}$
and $D_{12} = \Var(f_{12}) = \Var(Q_1-\tfrac{1}{2})\Var(Q_2-\tfrac{1}{2}) = \tfrac{1}{144}$.
Total variance $D = \tfrac{1}{16} + \tfrac{1}{16} + \tfrac{1}{144} = \tfrac{19}{144}$.
The first-order indices $S_1 = S_2 = \tfrac{1/16}{19/144} = \tfrac{9}{19} \approx 0.47$,
and the interaction index $S_{12} = \tfrac{1/144}{19/144} = \tfrac{1}{19} \approx 0.053$.
Note $S_1 + S_2 + S_{12} = 1$. For the additive model $f = Q_1 + Q_2$
(no interaction), $S_1 = S_2 = \tfrac{1}{2}$ and $S_{12} = 0$: the indices
sum to 1 with nothing left over for interactions.
This is why $\sum S_i < 1$ signals the presence of interactions.

\subsection{Sobol Indices}

\begin{definition}[Sobol Indices]\index{Sobol indices}
The \textbf{first-order Sobol index}\index{Sobol indices!first-order} for parameter $i$ is:
\begin{equation}
  S_i = \frac{D_i}{D} = \frac{\Var[\E(Y\,|\,Q_i)]}{\Var(Y)}.
  \label{eq:S1}
\end{equation}
The \textbf{total Sobol index}\index{Sobol indices!total} for parameter $i$~\cite{homma1996importance} is:
\begin{equation}
  S_{T_i} = 1 - \frac{\Var[\E(Y\,|\,\vec{Q}_{-i})]}{\Var(Y)}
           = \frac{\E[\Var(Y\,|\,\vec{Q}_{-i})]}{\Var(Y)},
  \label{eq:ST}
\end{equation}
where $\vec{Q}_{-i}$ denotes all parameters except $Q_i$.
\end{definition}

The first-order index measures the \emph{direct} effect of $Q_i$;
the total index measures the direct effect plus all interactions with
other parameters. When $S_{T_i} = S_i$ for all $i$, the model is
additively separable (no interactions). When $S_{T_i} > S_i$,
parameter $i$ interacts with other parameters.

\unifyingidea{1} The first-order Sobol index $S_i = D_i/D$ uses the
same ratio structure as the sensitivity index from Chapter~2;
it measures one variance as a fraction of another. The connection
is explicit: for a linear model $Y = \sum_i a_i Q_i$ with independent
$Q_i \sim \mathcal{U}[l_i, u_i]$, the first-order Sobol index equals
the square of the sigma-normalized index: $S_i = (S_i^\sigma)^2$.
Global SA generalizes local SA; for linear models, they agree exactly.

\unifyingidea{3} The Sobol decomposition resolves the tension that
Unifying Idea~3 created in Chapter~8. Local SA is a linearization,
valid near the nominal values. Sobol indices make no linearity
assumption: they measure the actual variance contribution across
the full parameter range. When the linearization is poor (Model~B
from the chapter opening), local and global SA give different answers,
and Sobol indices give the correct answer.

For readers with a functional analysis background (Appendix~\ref{app:innerproduct}), the ANOVA decomposition~\eqref{eq:ANOVA}
is the orthogonal decomposition of $f$ in the Hilbert space
$L^2([0,1]^m)$ with respect to the inner product
$\langle f, g\rangle = \int_{[0,1]^m} f(\vec{q})\,g(\vec{q})\,d\vec{q}$.
The orthogonality explains both why the variances add
in~\eqref{eq:variancedec} and why the method requires independent \POIs:
independence is needed for the inner product to factorize and the
projections to be orthogonal.

\begin{selfcheck}
For $Y = c_1 Q_1 + c_2 Q_2$ with $Q_i \sim \mathcal{U}(0,1)$ independent:
(a) Compute $\E[Y\,|\,Q_1]$ and $\Var[\E(Y\,|\,Q_1)]$.
(b) Compute $\Var(Y)$ using the independence of $Q_1$ and $Q_2$.
(c) Show that $S_1 = c_1^2/(c_1^2 + c_2^2)$ and $S_2 = c_2^2/(c_1^2 + c_2^2)$.
(d) Verify $S_1 + S_2 = 1$ (no interaction for additive models).
\end{selfcheck}
\begin{scAnswer}
(a) $\E[Y\,|\,Q_1] = c_1 Q_1 + c_2/2$.
$\Var[\E(Y\,|\,Q_1)] = c_1^2 \Var(Q_1) = c_1^2/12$.

(b) $\Var(Y) = c_1^2\Var(Q_1) + c_2^2\Var(Q_2) = (c_1^2 + c_2^2)/12$.

(c) $S_1 = D_1/D = (c_1^2/12)/((c_1^2+c_2^2)/12) = c_1^2/(c_1^2+c_2^2)$. Similarly $S_2$.

(d) $S_1 + S_2 = (c_1^2 + c_2^2)/(c_1^2+c_2^2) = 1$. For linear additive models,
there are no interaction terms and first-order indices sum to 1.
\end{scAnswer}

%%------------------------------------------------------------
\section{Computing Sobol Indices: The Saltelli Algorithm}
\label{sec:saltellialg}
%%------------------------------------------------------------

A naive Monte Carlo estimate of $S_i$ requires choosing $Q_i = q_i$
for many values and averaging the model over all other parameters
for each, this costs $O(N^2)$ evaluations. The insight of Saltelli~\cite{saltelli2002making}
is that all first-order and total Sobol indices can be estimated with
$N(m+2)$ evaluations using two carefully constructed sample matrices.
The design and estimators for the total sensitivity index are developed
in Saltelli et al.~\cite{saltelli2010variance}.

\subsection{The Saltelli Construction}

Generate two independent LHS matrices $A$ and $B$, each $N\times m$.
Define the $m$ matrices $A_B^{(i)}$: matrix $A$ with its $i$-th column
replaced by the $i$-th column of $B$. The model is evaluated at all
rows of $A$, $B$, and all $m$ matrices $A_B^{(i)}$: a total of
$N(m+2)$ evaluations.

\subsection{The Jansen Estimators}\index{Jansen estimators}

For the first-order index, use the Jansen~\cite{jansen1999analysis} estimator:
\begin{equation}
  \hat{S}_i = 1 - \frac{\frac{1}{2N}\sum_{j=1}^N [f(B_j) - f(A_{B,j}^{(i)})]^2}{\hat{D}},
  \label{eq:Jansen_S1}
\end{equation}
where $\hat{D} = \frac{1}{N}\sum_j f(A_j)^2 - \left(\frac{1}{N}\sum_j f(A_j)\right)^2$.

For the total index:
\begin{equation}
  \hat{S}_{T_i} = \frac{\frac{1}{2N}\sum_{j=1}^N [f(A_j) - f(A_{B,j}^{(i)})]^2}{\hat{D}}.
  \label{eq:Jansen_ST}
\end{equation}
Use these Jansen estimators, not the original Saltelli (2002) formulas,
which have higher variance. The Jansen estimator is consistent
(converges to the true $S_i$ as $N\to\infty$) and nearly unbiased;
see Jansen~\cite{jansen1999analysis} Theorem~1.
The key step is that
\[
  \E\!\left[\frac{1}{2N}\sum_{j=1}^N (f(B_j) - f(A_{B,j}^{(i)}))^2\right]
  = D - D_i = \Var(Y) - \Var[\E(Y\,|\,Q_i)],
\]
which follows because $B_j$ and $A_{B,j}^{(i)}$ share all parameter
values except column $i$ (from $A$), so their difference isolates
the $Q_i$-dependent variation. Dividing by $\hat{D}$ and rearranging
gives $1 - \hat{S}_i \approx (D - D_i)/D = 1 - S_i$.
The estimator is nearly unbiased at finite $N$ because the bias
is $O(1/N)$~\cite{jansen1999analysis}, negligible for $N \geq 1000$.

\subsection{Convergence}

The estimators converge as $O(1/\sqrt{N})$, the standard Monte Carlo rate.
In practice, $N \geq 1000$ is needed for reliable estimates.
For $m = 4$ parameters (SIR model): $N(m+2) = 6000$ evaluations at $N = 1000$.
If each evaluation takes 0.1~seconds: 10~minutes.

Plotting $\hat{S}_i$ and $\hat{S}_{T_i}$ vs.\ $N$ reveals the convergence:
estimates are noisy for $N < 500$ and stable for $N \geq 2000$.
This convergence plot is the diagnostic that answers
the misconception ``more samples always help'': below a threshold
$N$, additional samples do not produce reliable estimates because
the variance of the estimator exceeds the signal being measured.

\textbf{Confidence intervals via bootstrap.}
A point estimate $\hat{S}_i$ is useful but incomplete: it gives no
information about sampling uncertainty.
The standard approach is the \textbf{bootstrap}: resample the
$N(m+2)$ model evaluations with replacement, recompute $\hat{S}_i$
for each resample, and take the 2.5th and 97.5th percentiles of
the resulting distribution as a 95\% confidence interval.
In MATLAB:
\begin{verbatim}
B = 500;          %% number of bootstrap replicates
S1_boot = zeros(B, m);
idx_all = 1:N;
for b = 1:B
    idx = idx_all(randi(N, N, 1));     %% resample rows
    A_b = A(idx,:);  B_b = Bmat(idx,:);
    [S1_boot(b,:), ~] = sobol_jansen(@(x) model_eval(x), A_b, B_b);
end
S1_CI = prctile(S1_boot, [2.5 97.5]);   %% 95% CI for each index
\end{verbatim}
A confidence interval that straddles zero indicates that the parameter's
effect cannot be distinguished from noise at the current sample size.
A confidence interval that excludes zero confirms the parameter is
genuinely influential. As a practical rule: if the confidence interval
width exceeds the index value, double $N$ before drawing conclusions.
When $m$ is large, however, even a single Sobol run at $N = 2000$
may be prohibitively expensive; Morris screening addresses exactly this case.

%%------------------------------------------------------------
\section{Morris Screening}\index{Morris screening}
\label{sec:morris}
%%------------------------------------------------------------

When $m$ is large (30 or more parameters), running Sobol at $N = 2000$
costs $N(m+2) = 64{,}000$ model evaluations, potentially prohibitive.
Morris screening~\cite{morris1991factorial} provides a cheap first step:
$N(m+1)$ evaluations (typically $N = 10$ trajectories) to identify
which parameters are truly unimportant. An improved radial design
that gives more uniform coverage with fewer trajectories is
described by Campolongo et al.~\cite{campolongo2007effective}.

For each trajectory, one parameter is changed at a time by a fixed
fraction $\Delta$ of its range. The \textbf{elementary effect}\index{elementary effect} of
parameter $i$ at point $\vec{q}$ is:
\begin{equation}
  EE_i(\vec{q}) = \frac{f(\vec{q} + \Delta\vec{e}_i) - f(\vec{q})}{\Delta}.
  \label{eq:EE}
\end{equation}
After $r$ trajectories, the mean $\mu_i^* = (1/r)\sum|EE_i|$
and standard deviation $\sigma_i$ of the elementary effects are computed.
Parameters with small $\mu_i^*$ are unimportant. Parameters with large
$\sigma_i$ have interactions or nonlinear effects.

The Morris diagnostic plot ($\sigma_i$ vs.\ $\mu_i^*$) separates
parameters into three groups: unimportant (small $\mu_i^*$, small $\sigma_i$);
important and linear (large $\mu_i^*$, small $\sigma_i$); important
and nonlinear or interacting (large $\mu_i^*$, large $\sigma_i$).

\textbf{Workflow}: run Morris screening first; fix unimportant parameters
at their nominal values; run Sobol only on the important subset.
This can reduce the Sobol cost by a factor of 5--10.

\medskip

\noindent\textbf{When even reduced Sobol is too expensive: surrogate models.}
If the model is computationally expensive and Morris screening reduces
the parameter set to, say, five influential \POIs, running Sobol at
$N = 2000$ still requires $N(m+2) = 14{,}000$ model evaluations.
For a climate model where each run takes hours, this remains prohibitive.

The practical solution is a \textbf{surrogate model}\index{surrogate model}
(also called an emulator or meta-model): a cheap statistical
approximation to the expensive model, trained on a moderate number of
carefully chosen model runs.
Concretely, a Gaussian process surrogate fits a probability distribution
over functions to the training runs, predicting the model output at any
new input along with a confidence interval; a polynomial chaos expansion
instead represents the output as a truncated series in orthogonal polynomials
of the uncertain inputs, with coefficients estimated from the training set.
Gaussian process emulators and polynomial chaos expansions are the
most widely used surrogates. Once trained on roughly $50$--$200$
model evaluations, a surrogate can be evaluated millions of times
in seconds, making full Sobol estimation feasible even for
computationally expensive simulations.

Surrogate-based global SA does not replace the methods in this chapter;
it wraps them. The LHS sampling, the Saltelli construction, and the
Jansen estimators all apply to the surrogate exactly as they apply
to the original model. The surrogate simply makes the evaluations
fast enough to converge.
Smith~\cite{smith2014uncertainty} develops Gaussian process and
polynomial chaos surrogates in Chapters~12--13, building directly
on the LHS samples and Sobol indices from this chapter.
Students proceeding to the second semester will construct surrogates
using exactly the training sets generated in this chapter's laboratories.

%%------------------------------------------------------------
\section{Sobol Indices for Correlated POIs}\index{Sobol indices!correlated POIs}
\label{sec:correlatedsobol}
%%------------------------------------------------------------

Standard Sobol decomposition assumes independent \POIs. When parameters
are correlated, three problems arise: first-order indices can exceed 1 or
be negative (a diagnostic flag that correlation is present); total indices
may not sum to 1; and the ANOVA decomposition~\eqref{eq:ANOVA} is no longer
unique, because correlation destroys the orthogonality that guarantees uniqueness.

Three approaches address correlated \POIs in global SA:

\textbf{Grouped indices}: treat a cluster of correlated parameters as
a single group. Compute the group Sobol index (the fraction of variance
explained by the entire group jointly). This avoids allocating variance
between correlated parameters, which is not well-defined.

\textbf{Modified estimators}: Kucherenko et al.~\cite{kucherenko2012estimation}
extend the Saltelli algorithm to correlated inputs using copula-based
sampling. This requires knowledge of the joint distribution of the
correlated parameters.

\textbf{PRCC}: handles correlation naturally through the ``partial''
in partial rank correlation. When parameters are correlated, PRCC
is the most robust and accessible global SA method.

\textbf{Practical recommendation}: for models with correlated \POIs,
use PRCC as the primary global SA tool and grouped Sobol indices as
a secondary check. Report the correlation structure explicitly alongside
the SA results.

%%------------------------------------------------------------
\section{Comparison of All SA Methods}
\label{sec:methodcomparison}
%%------------------------------------------------------------

Table~\ref{tab:allSA} consolidates all SA methods from the book.
It supersedes the partial decision tables in earlier chapters.
For a broader treatment of global SA methods, see Saltelli et al.~\cite{saltelli2008global,saltelli2000sensitivity}
and the recent epidemiological applications in Lu and Borgonovo~\cite{lu2023global}.

\begin{table}[htb]
\centering
\renewcommand{\arraystretch}{1.25}
\begin{tabular}{p{2.4cm}p{1.3cm}p{3.5cm}p{2.2cm}p{2.3cm}}
\toprule
Method & Chapter & When to use & Assumption & Cost \\
\midrule
Local SI (analytic) & 2, 4 & Differentiable model; small uncertainty & Linearization valid & Cheap \\
Local SI (FD) & 3 & Black-box; small uncertainty & Linearization valid & $2mn_q$ evals \\
Adjoint & 5--6 & Differentiable; many \POIs; few \QOIs & Linearization valid & 2 solves \\
AD (reverse) & 7 & Differentiable code; $m \gg r$ & Linearization valid & 2 passes \\
PRCC & 9 & Monotone model; moderate uncertainty; corr.\ \POIs OK & Monotone & $N$ evals \\
Morris & 9 & Screening; $m \geq 30$ & None & $N(m+1)$ evals \\
Sobol & 9 & Nonlinear; large uncertainty; independent \POIs & Independent & $N(m+2)$ evals \\
Grouped Sobol & 9 & Correlated \POIs; known groups & Known groups & $N(m+2)$ evals \\
\bottomrule
\end{tabular}
\caption{Complete SA method comparison across all chapters.
         The choice of method depends primarily on whether the model
         is differentiable, whether uncertainty is small or large,
         and whether \POIs are independent. No single method is universally best.}
\label{tab:allSA}
\end{table}

\unifyingidea{2} The adjoint row in Table~\ref{tab:allSA} is the
payoff for Chapters~5 and 6. The adjoint is the transpose in whatever
space you are working in: matrix transpose for linear systems
(Chapter~5), reversed ODE for dynamical systems (Chapter~6), and
reversed chain rule for computation graphs (Chapter~7). Every
``2 solves'' entry in the adjoint row is the same unifying idea,
applied in a different mathematical space.

%%------------------------------------------------------------
\section{Bridge to Semester 2}
\label{sec:bridgesemester2}
%%------------------------------------------------------------

This book quantifies how model output uncertainty arises from
parameter uncertainty, assuming the model structure is correct.
The second semester, covering Smith's \emph{Uncertainty Quantification}
\cite{smith2014uncertainty}, asks harder questions: what if the model is wrong?
How do we fit models to data? How do we construct efficient surrogate
models to replace expensive simulations?
Smith's Chapters~7--8 develop Bayesian parameter estimation and Markov chain
Monte Carlo methods for fitting models to data, the natural sequel to the
identifiability analysis introduced in Chapter~3 of this book.

The connections between the two semesters are direct:

Smith Ch.~14 develops local SA at the graduate level, extending
Chapters~2--7 of this book with functional-analytic rigor.

Smith Ch.~15 develops global SA and Sobol indices at the graduate level,
extending Chapter~9 of this book with measure-theoretic foundations.

Smith Ch.~6 treats parameter selection using the sensitivity matrix;
the core tool introduced in Chapter~4 of this book.

Smith Chs.~12--13 develop model discrepancy and Gaussian process
surrogates, methods that reduce the cost of global SA by replacing
the expensive model with a cheap approximation. Surrogates are
built from the same LHS samples used to estimate Sobol indices
in this chapter.

Students who have completed this book have the foundational vocabulary,
the computational intuition, and the mathematical framework to engage
productively with the full UQ literature. The path from the sensitivity
index in Chapter~2 to the full Bayesian UQ framework in Smith is
a single, continuous intellectual development. As an example of
these methods applied to a realistic multi-parameter epidemic model,
see Qu, Xue, and Hyman~\cite{qu2018modeling}, which uses SA to
guide Wolbachia-based mosquito-borne disease control.

The student who has reached this point knows how to ask which parameters
matter, how to compute that answer efficiently and honestly, and, just
as importantly, when not to trust the answer.
Every mathematical model is an argument; sensitivity analysis is the
discipline that reveals which assumptions the argument depends on.

\medskip\noindent
\textit{Good sensitivity analysis does not end the discussion.
It tells you where to begin it.}

%%------------------------------------------------------------
\section{MATLAB Laboratory: Complete Global SA Pipeline for SIR}
\label{sec:ch9matlab}
%%------------------------------------------------------------

This laboratory implements the full global SA pipeline for the SIR model
and compares all four methods: sigma-normalized, PRCC, Sobol, and Morris.

\textbf{Setup.}
SIR model with four \POIs $(k, \beta, \tau, L)$ and response $J = I_{\max}$
(peak infectious count). Use the literature-based uncertainty ranges
from Chitnis et al.~\cite{chitnis2008determining}:
$k \in [2, 8]$, $\beta \in [0.02, 0.12]$, $\tau \in [4, 14]$ days,
$L \in [1000, 100000]$ days.

\textbf{Part 1: LHS and sigma-normalized indices.}
Draw $N = 500$ LHS samples using the \texttt{lhs\_sample} function.
Compute $\sigma_{J}$ from the samples. Compute the sigma-normalized index
for each \POI using the local SI from Chapter~2 and the sample
standard deviations. Report the ranking.

\textbf{Part 2: PRCC.}
Run the SIR model at all 500 LHS samples. Rank-transform all inputs
and the output. Compute PRCC using the \texttt{partialcorr} function
in MATLAB's Statistics and Machine Learning Toolbox
(or download Marino et al.'s implementation from the supplementary
material of \cite{marino2008methodology}).
Produce a PRCC tornado plot and report $p$-values.

\textbf{Part 3: Sobol estimation.}
Draw $N = 2000$ samples using the quasi-random Halton or Sobol sequence
(MATLAB: \texttt{haltonset} or \texttt{sobolset}).
Build the $A$, $B$, and $m$ matrices $A_B^{(i)}$: $2000\times6 = 12{,}000$
total evaluations.
Compute $\hat{S}_i$ and $\hat{S}_{T_i}$ using the Jansen estimators.

\begin{verbatim}
%% From: src/globalsa/sobol_jansen.m (see GitHub repository)
rng(42);   %% reproducibility
[A, B]      = saltelli_sample(N, m, lb, ub);
[S1, ST, CI_S1] = sobol_jansen(@model_eval, A, B, 500, 0.05);
%% CI_S1(:,j) = [lower; upper] 95% confidence interval for S1(j)
\end{verbatim}

\textbf{Part 4: Morris screening.}
Implement $r = 10$ Morris trajectories with $\Delta = 0.4$.
Compute $\mu_i^*$ and $\sigma_i$ for each parameter.
Produce the Morris diagnostic plot ($\sigma_i$ vs.\ $\mu_i^*$).
Which parameter(s) could be fixed without significant loss of accuracy?

\textbf{Part 5: Convergence check.}
Plot $\hat{S}_\beta$ and $\hat{S}_{T_\beta}$ vs.\ $N$ for
$N \in \{100, 200, 500, 1000, 2000, 5000\}$.
At what $N$ do the estimates stabilize to within $\pm0.05$?

\textbf{Part 6: Full comparison table.}
Fill in the following table from your results:

\begin{center}
\begin{tabular}{lcccc}
\toprule
Parameter & Local SI & $S_i^\sigma$ & PRCC & $S_i$ (Sobol) \\
\midrule
$k$ & & & & \\
$\beta$ & & & & \\
$\tau$ & & & & \\
$L$ & & & & \\
\bottomrule
\end{tabular}
\end{center}

Where do the four methods agree? Where do they disagree? What accounts
for the differences?

%%------------------------------------------------------------
\section{Exercises}
%%------------------------------------------------------------

\begin{exercise}[\bloom{L3}]
\label{ex:ch9:sigmanorm}
For the SEIR model with local SIs $\SI{\tau_E}{R_0} = 0.6$,
$\SI{\tau_I}{R_0} = 1.2$, $\SI{\beta}{R_0} = 1.0$, $\SI{k}{R_0} = 1.0$
and uncertainty ranges $\sigma_{\tau_E} = 2$~days, $\sigma_{\tau_I} = 3$~days,
$\sigma_\beta = 0.01$, $\sigma_k = 0.8$:
\begin{enumerate}[label=(\alph*)]
  \item Compute the output standard deviation $\sigma_{R_0}$ using the
        delta method, at nominal $R_0^* = 2.0$.
  \item Compute all four sigma-normalized indices.
  \item Produce a tornado plot of the sigma-normalized indices.
  \item Compare the ranking with the local SI ranking. Which parameter
        changed rank most dramatically? Why?
\end{enumerate}
\end{exercise}

\begin{exercise}[\bloom{L3}]
\label{ex:ch9:LHS}
\begin{enumerate}[label=(\alph*)]
  \item Implement the \texttt{lhs\_sample} function. For $N = 100$,
        $m = 3$, $[0,1]^3$, generate an LHS and a random sample of the
        same size.
  \item Produce scatter plots of all three pairs of parameters for
        both samples (6 plots total).
  \item Compute the mean and standard deviation of the sample along each
        dimension. Which method is closer to the true uniform
        mean and standard deviation?
  \item For one parameter, plot the histogram of LHS vs.\ random
        sampling. Which covers the range more evenly?
\end{enumerate}
\end{exercise}

\begin{exercise}[\bloom{L3}]
\label{ex:ch9:PRCC}
For the SIR model with parameter ranges from the laboratory setup:
\begin{enumerate}[label=(\alph*)]
  \item Generate $N = 200$ LHS samples and run the SIR model at each.
  \item Compute PRCC for all four parameters with respect to $I_{\max}$.
  \item Produce a PRCC bar chart sorted by $|$PRCC$|$.
  \item Which parameter has the strongest negative PRCC? What does this
        mean in plain language?
  \item Compare the PRCC ranking with the local SI tornado plot from
        Chapter~2. Do they agree?
\end{enumerate}
\end{exercise}

\begin{exercise}[\bloom{L4}]
\label{ex:ch9:sobolanalytic}
For $Y = 2Q_1^2 + Q_2$ with $Q_1, Q_2 \sim \mathcal{U}(0,1)$, independent:
\begin{enumerate}[label=(\alph*)]
  \item Compute $\E[Y]$, $\Var(Y)$, $D_1 = \Var[\E(Y\,|\,Q_1)]$,
        $D_2 = \Var[\E(Y\,|\,Q_2)]$.
  \item Compute $S_1$, $S_2$. Do they sum to 1? Explain.
  \item Compute the total indices $S_{T_1}$, $S_{T_2}$.
  \item Now add $Q_1$ and $Q_2$ to get $Y' = 2Q_1^2 + Q_2 + Q_1 Q_2$
        (an interaction term). Without computing: does adding the interaction
        term change $S_1$, $S_2$, or both? In which direction?
        Then verify by computing the Sobol indices for $Y'$ analytically.
\end{enumerate}
\end{exercise}

\begin{exercise}[\bloom{L4}]
\label{ex:ch9:sobolSIR}
Using the Jansen estimator implementation from the laboratory:
\begin{enumerate}[label=(\alph*)]
  \item Compute $\hat{S}_i$ and $\hat{S}_{T_i}$ for all four SIR parameters
        at $N = 2000$.
  \item For which parameters is $S_{T_i} \gg S_i$? What does this indicate
        about their role in the model?
  \item Compute $\sum_i S_i$ and $\sum_i S_{T_i}$. What do you expect
        these to sum to for an additive model? Are your results consistent
        with this expectation?
  \item If you could measure exactly one parameter very precisely
        (reducing its uncertainty to zero), which parameter would reduce
        $\Var(I_{\max})$ the most? Justify using the Sobol indices.
\end{enumerate}
\end{exercise}

\begin{exercise}[\bloom{L4}]
\label{ex:ch9:morris}
For the SIR model, implement $r = 15$ Morris trajectories with $\Delta = 1/3$:
\begin{enumerate}[label=(\alph*)]
  \item Compute $\mu_i^*$ and $\sigma_i$ for each parameter.
  \item Produce the Morris diagnostic plot.
  \item Based on the Morris results, which parameter appears most influential?
        Which appears least influential?
  \item Run a reduced Sobol analysis using only the three most important
        parameters identified by Morris (fixing $L$ at its nominal value).
        Compare the cost and results with the full four-parameter Sobol.
\end{enumerate}
\end{exercise}

\begin{exercise}[\bloom{L5}]
\label{ex:ch9:globaldesign}
Choose a published epidemic model from the literature with at least
8 parameters. Design a complete global SA study:
\begin{enumerate}[label=(\alph*)]
  \item State the model, its parameters, and a justification for the
        uncertainty ranges you will use (with citations).
  \item State the \QOI or \QOIs and explain why they are the most
        relevant for the model's purpose.
  \item Choose between PRCC and Sobol as the primary global SA method.
        Justify your choice based on the model's properties.
  \item Specify the sample size $N$ and justify it based on convergence.
  \item Describe how you would present the results to a non-technical
        audience (e.g., a public health committee or funding agency).
  \item Identify one situation where your global SA results could still
        be misleading, and describe what additional analysis you would
        recommend.
\end{enumerate}
\end{exercise}

\begin{exercise}[\bloom{L6}]
\label{ex:ch9:evaluation}
A published SA study reports Sobol first-order and total indices for
a 15-parameter water-quality model. The authors note in the data section
that three parameters, phosphorus loading rate, nitrogen loading rate,
and water residence time, are positively correlated across the sampled
watersheds. The SA section reports standard Sobol indices computed using
independent sampling for all 15 parameters, without addressing the
correlation.
\begin{enumerate}[label=(\alph*)]
  \item Identify the specific mathematical problem that correlation
        creates for standard Sobol indices.
  \item For each of the three correlated parameters, state whether the
        reported first-order index is likely over-estimated, under-estimated,
        or unchanged, and explain why.
  \item Propose the minimal methodological correction that would make
        the analysis valid.
  \item The authors justify their approach by noting that their Sobol
        indices sum to approximately 1.0. Explain why this is not
        a valid check for independence.
\end{enumerate}
\end{exercise}

%%------------------------------------------------------------



%% ============================================================
\appendix
%% ============================================================

\part*{Appendices}
\addcontentsline{toc}{part}{Appendices}

%% ============================================================
%% Appendix A: Inner Products, Norms, and Function Spaces
%% ============================================================

\chapter{Inner Products, Norms, and Function Spaces}
\label{app:innerproduct}

This appendix provides the minimum background in functional analysis
needed to read Appendix~\ref{app:adjointformalism} and Smith's companion
volume~\cite{smith2014uncertainty}. Readers who are comfortable with inner products
and the $L^2$ space can proceed directly to Appendix~\ref{app:adjointformalism}.

%%------------------------------------------------------------
\section{Inner Products and Norms}
\label{appsec:innerproducts}
%%------------------------------------------------------------

An \textbf{inner product}\index{inner product} on a vector space $V$ over $\R$ is a function
$\langle \cdot, \cdot \rangle : V \times V \to \R$ satisfying:
\begin{enumerate}[nosep]
  \item \textit{Symmetry}: $\langle u, v \rangle = \langle v, u \rangle$.
  \item \textit{Linearity in the first argument}:
        $\langle \alpha u + \beta w, v \rangle
        = \alpha\langle u, v\rangle + \beta\langle w, v\rangle$.
  \item \textit{Positive definiteness}: $\langle u, u \rangle \geq 0$,
        with equality only when $u = 0$.
\end{enumerate}

The \textbf{induced norm}\index{norm!induced} is $\|u\| = \sqrt{\langle u, u \rangle}$.
The Cauchy-Schwarz inequality states $|\langle u, v \rangle| \leq \|u\|\,\|v\|$.

\medskip

\noindent\textbf{Example~A.1 (Euclidean inner product).}
For $\vec{u}, \vec{v} \in \R^n$:
\[
  \langle \vec{u}, \vec{v} \rangle = \vec{u}^T\vec{v} = \sum_{i=1}^n u_i v_i.
\]
This is the dot product. The induced norm is the Euclidean norm
$\|\vec{u}\|_2 = \sqrt{\sum u_i^2}$.

\noindent\textbf{Example~A.2 (Weighted inner product).}
For a symmetric positive definite matrix $W \in \R^{n\times n}$:
$\langle \vec{u}, \vec{v}\rangle_W = \vec{u}^T W \vec{v}$.
The induced norm $\|\vec{u}\|_W = \sqrt{\vec{u}^T W \vec{u}}$ is the
$W$-weighted Euclidean norm. SA computations sometimes use this when
parameters have very different scales.

%%------------------------------------------------------------
\section{The $L^2$ Inner Product}
\label{appsec:L2}
%%------------------------------------------------------------

For functions $f, g : [a,b] \to \R$ that are square-integrable
(i.e., $\int_a^b f(t)^2\,dt < \infty$), the \textbf{$L^2$ inner product}\index{$L^2$ inner product} is:
\begin{equation}
  \langle f, g \rangle_{L^2} = \int_a^b f(t)\,g(t)\,dt.
  \label{eq:L2ip}
\end{equation}
The induced norm is $\|f\|_{L^2} = \sqrt{\int_a^b f(t)^2\,dt}$.

The space $L^2([a,b])$ of all square-integrable functions on $[a,b]$
with inner product~\eqref{eq:L2ip} is a \textbf{Hilbert space}\index{Hilbert space}:
a complete inner product space.
Completeness means that every Cauchy sequence converges to a function
in $L^2([a,b])$.

For Sobol decomposition (Chapter~9), the relevant space is $L^2([0,1]^m)$
with inner product $\langle f, g\rangle = \int_{[0,1]^m} f(\vec{q})g(\vec{q})\,d\vec{q}$.

%%------------------------------------------------------------
\section{Integration by Parts as Transposition}
\label{appsec:IBP}
%%------------------------------------------------------------

The key connection between adjoint sensitivity analysis and linear algebra
is this: integration by parts is the function-space analog of the
matrix transpose.

In $\R^n$, for vectors $\vec{u}$, $\vec{v}$ and a matrix $A$:
\[
  \langle A\vec{u}, \vec{v}\rangle = (A\vec{u})^T\vec{v} = \vec{u}^T(A^T\vec{v})
  = \langle\vec{u}, A^T\vec{v}\rangle.
\]
So $A^T$ is the operator satisfying $\langle Au, v\rangle = \langle u, A^Tv\rangle$.

In $L^2([0,T])$, for functions $u, v$ and the differentiation operator $L = d/dt$:
\[
  \langle Lu, v\rangle_{L^2}
  = \int_0^T \dot{u}(t)\,v(t)\,dt
  = \bigl[u(t)v(t)\bigr]_0^T - \int_0^T u(t)\,\dot{v}(t)\,dt
  = \bigl[uv\bigr]_0^T + \langle u, L^*v\rangle_{L^2},
\]
where $L^* = -d/dt$ (with appropriate boundary conditions) is the
\textbf{adjoint operator} of $L = d/dt$.
The boundary term $[uv]_0^T$ gives rise to the terminal condition
$\lambda(T) = -h'(u(T))$ in the adjoint ODE derivation of Chapter~6.

This correspondence makes precise what we stated informally throughout
the book: the adjoint ODE is the ODE whose differential operator
is the transpose of the forward ODE's operator, exactly as $A^T$
is the transpose of $A$.

%%------------------------------------------------------------
\section{Orthogonality and the ANOVA Decomposition}
\label{appsec:ortho}
%%------------------------------------------------------------

Two elements $u, v$ of an inner product space are \textbf{orthogonal}\index{orthogonality}
if $\langle u, v\rangle = 0$.

In $L^2([0,1]^m)$ with the uniform measure, the functions appearing in
the Sobol ANOVA decomposition (Chapter~9) are mutually orthogonal.
Specifically, for the decomposition
$f = f_0 + \sum_i f_i + \sum_{i<j} f_{ij} + \cdots$,
each term is orthogonal to every other term:
$\langle f_{u}, f_{v}\rangle_{L^2} = 0$ when $u \neq v$ (as index sets).

This orthogonality is why the output variance decomposes as a sum:
$\Var(Y) = \sum_i D_i + \sum_{i<j} D_{ij} + \cdots$.
It also explains why the decomposition requires independent parameters:
independence of the $Q_i$ makes the integration over $[0,1]^m$ factorize,
which is what produces the orthogonality.

%%------------------------------------------------------------
\section{Exercises}
\label{appsec:A_exercises}
%%------------------------------------------------------------

\begin{selfcheck}
Let $f(t) = e^{-t}$ and $g(t) = \cos(\pi t)$ on $[0, 1]$.
(a) Compute $\langle f, g \rangle_{L^2}$ numerically (use MATLAB's
\texttt{integral} function).
(b) Are $f$ and $g$ orthogonal in $L^2([0,1])$?
(c) Compute $\|f\|_{L^2}$ and $\|g\|_{L^2}$, and verify the
Cauchy--Schwarz inequality $|\langle f,g\rangle| \leq \|f\|\,\|g\|$.
\end{selfcheck}
\begin{scAnswer}
(a) $\langle f, g\rangle_{L^2} = \int_0^1 e^{-t}\cos(\pi t)\,dt$.
Integrating by parts twice: result $\approx 0.218$ (verify numerically).
(b) No — the inner product is nonzero, so $f$ and $g$ are not orthogonal.
(c) $\|f\|_{L^2}^2 = \int_0^1 e^{-2t}\,dt = (1-e^{-2})/2 \approx 0.432$,
so $\|f\|_{L^2} \approx 0.657$.
$\|g\|_{L^2}^2 = \int_0^1 \cos^2(\pi t)\,dt = 1/2$, so $\|g\|_{L^2} = 1/\sqrt{2} \approx 0.707$.
Cauchy--Schwarz: $|0.218| \leq 0.657 \times 0.707 \approx 0.464$. \checkmark
\end{scAnswer}

\begin{selfcheck}
The integration-by-parts identity in Section~\ref{appsec:IBP}
establishes that $L^* = -d/dt$ (with zero terminal condition)
is the adjoint of $L = d/dt$ (with zero initial condition)
in $L^2([0,T])$.
(a) Write out $\langle Lu, v\rangle_{L^2}$ explicitly for
$u(t) = t^2$ and $v(t) = t$ on $[0,1]$.
(b) Compute $\langle u, L^*v\rangle_{L^2}$ for the same functions.
(c) Verify that the two results differ by the boundary term $[u(t)v(t)]_0^1$.
\end{selfcheck}
\begin{scAnswer}
(a) $Lu = 2t$. $\langle Lu, v\rangle = \int_0^1 2t \cdot t\,dt = \int_0^1 2t^2\,dt = 2/3$.
(b) $L^*v = -dv/dt = -1$. $\langle u, L^*v\rangle = \int_0^1 t^2 \cdot(-1)\,dt = -1/3$.
(c) Boundary term: $[u(t)v(t)]_0^1 = 1\cdot1 - 0\cdot0 = 1$.
Check: $\langle Lu,v\rangle = \langle u, L^*v\rangle + [uv]_0^1 \Rightarrow 2/3 = -1/3 + 1 = 2/3$. \checkmark
\end{scAnswer}

\begin{selfcheck}
The ANOVA decomposition in $L^2([0,1]^2)$ for the model
$f(Q_1, Q_2) = Q_1 + Q_2 + Q_1 Q_2$ (with $Q_i \sim \mathcal{U}(0,1)$) is:
$f = f_0 + f_1(Q_1) + f_2(Q_2) + f_{12}(Q_1,Q_2)$.
Compute each term and verify that all terms are mutually orthogonal
in $L^2([0,1]^2)$.
\end{selfcheck}
\begin{scAnswer}
$f_0 = \mathbb{E}[f] = 1/2 + 1/2 + 1/4 = 5/4$.
$f_1(Q_1) = \mathbb{E}[f|Q_1] - f_0 = Q_1 + 1/2 + Q_1/2 - 5/4 = 3Q_1/2 - 3/4$.
$f_2(Q_2) = 3Q_2/2 - 3/4$ (by symmetry).
$f_{12} = f - f_0 - f_1 - f_2 = Q_1Q_2 - Q_1/2 - Q_2/2 + 1/4 = (Q_1-1/2)(Q_2-1/2)$.
Orthogonality: $\langle f_1, f_2\rangle = \int_0^1\int_0^1 (3Q_1/2-3/4)(3Q_2/2-3/4)\,dQ_1\,dQ_2
= \left[\int_0^1(3Q/2-3/4)\,dQ\right]^2 = 0^2 = 0$. \checkmark
Similarly all other pairs integrate to zero.
\end{scAnswer}

%% ============================================================
%% Appendix B: The Adjoint Formalism — Operators and the Lagrange Identity
%% ============================================================

\chapter{The Adjoint Formalism}
\label{app:adjointformalism}

This appendix develops the abstract operator framework that unifies
the matrix adjoint (Chapter~5) and the ODE adjoint (Chapter~6).
It is intended for students who want a rigorous foundation, not just
the computational recipe.

%%------------------------------------------------------------
\section{Linear Operators and Their Adjoints}
\label{appsec:operators}
%%------------------------------------------------------------

Let $H$ be a Hilbert space with inner product $\langle\cdot,\cdot\rangle$.
A \textbf{linear operator}\index{linear operator} $L : H \to H$ satisfies
$L(\alpha u + \beta v) = \alpha Lu + \beta Lv$.

The \textbf{adjoint operator}\index{adjoint operator} $L^\dagger$ is defined by:
\begin{equation}
  \langle Lu, v \rangle = \langle u, L^\dagger v \rangle
  \qquad\text{for all } u, v \in H.
  \label{eq:adjointdef}
\end{equation}
The adjoint $L^\dagger$ is the unique operator satisfying~\eqref{eq:adjointdef}
(when it exists).

\noindent\textbf{Example~B.1 (Matrix adjoint).}
In $\R^n$ with the Euclidean inner product, $L = A$ (left multiplication).
Then $\langle Au, v\rangle = (Au)^T v = u^T(A^T v) = \langle u, A^T v\rangle$,
so $L^\dagger = A^T$.
This is the matrix transpose: Chapters~5's adjoint equation
$A^T\vec{v} = \vec{c}$ is $L^\dagger\vec{v} = \vec{c}$.

\noindent\textbf{Example~B.2 (Differential operator adjoint).}
In $L^2([0,T])$, let $Lu = du/dt$ with $u(0) = 0$ (zero initial condition).
By integration by parts (Appendix~\ref{appsec:IBP}):
\[
  \langle Lu, v\rangle = \int_0^T \dot{u}v\,dt
  = [uv]_0^T - \int_0^T u\dot{v}\,dt
  = u(T)v(T) - \int_0^T u\dot{v}\,dt.
\]
For the adjoint to satisfy $\langle Lu, v\rangle = \langle u, L^\dagger v\rangle$,
we need $L^\dagger v = -\dot{v}$ with $v(T) = 0$ (zero terminal condition).
So $L^\dagger = -d/dt$ with the boundary condition shifted from $t=0$ to $t=T$.
This is exactly the structure of the adjoint ODE in Chapter~6.

%%------------------------------------------------------------
\section{The Lagrange Identity}
\label{appsec:lagrange}
%%------------------------------------------------------------

For the ODE operator $L = d/dt - F_u$ (the linearization of $\dot{u} = F(u;p)$),
the Lagrange identity\index{Lagrange identity} states:
\begin{equation}
  \frac{d}{dt}\bigl[\lambda^T w\bigr]
  = (L^\dagger\lambda)^T w + \lambda^T (Lw),
  \label{eq:lagrangeID}
\end{equation}
where $w = \partial u/\partial p$ satisfies the FSE $Lw = F_p$ and
$\lambda$ satisfies the adjoint ODE $L^\dagger\lambda = g_u$.

Integrating~\eqref{eq:lagrangeID} over $[0,T]$ and rearranging gives the
sensitivity formula~(6.11) of Chapter~6. This is the rigorous version of
the ``integration by parts'' step that was announced but not fully proved
in the main text.

%%------------------------------------------------------------
\section{Solvability Conditions}
\label{appsec:solvability}
%%------------------------------------------------------------

For the equation $Lu = f$ to have a solution when $L$ is not invertible,
$f$ must be orthogonal to the kernel of $L^\dagger$:

\textbf{SC1 (Fredholm Alternative).}\index{Fredholm alternative}
$Lu = f$ has a solution if and only if $\langle f, v\rangle = 0$
for every $v$ satisfying $L^\dagger v = 0$.

\textbf{SC2 (Uniqueness).}
If $L^\dagger v = 0$ has only the trivial solution $v = 0$,
then $Lu = f$ has a unique solution for every $f$.

In the SA context, SC2 corresponds to the assumption that the Jacobian
at the equilibrium is nonsingular (Theorem~4.1): the FSE has a unique
solution, and the adjoint equation has a unique adjoint variable.
Near a bifurcation, SC2 fails because the Jacobian becomes singular,
which is precisely why SA breaks down (Chapter~8).

%%------------------------------------------------------------
\section{Connection to Chapters 5 and 6}
\label{appsec:connection}
%%------------------------------------------------------------

The following table makes the Chapter~5/6 connection explicit in
the language of this appendix:

\begin{center}
\renewcommand{\arraystretch}{1.4}
\begin{tabular}{lll}
\toprule
Concept & Chapter 5 (matrices) & Chapter 6 (ODEs) \\
\midrule
Space $H$ & $\R^n$ & $L^2([0,T])^n$ \\
Operator $L$ & $A$ (matrix mult.) & $d/dt - J_F$ \\
Adjoint $L^\dagger$ & $A^T$ & $-d/dt - J_F^T$ \\
Adjoint equation & $A^T\vec{v} = \vec{c}$ & $-\dot{\lambda} = J_F^T\lambda - g_u$ \\
Boundary cond. & none (algebraic) & $\lambda(T) = -\nabla h^T$ \\
Sensitivity formula & $\vec{v}^T\vec{r}_j$ & $\int_0^T\lambda^T F_{p_j}\,dt$ \\
Key operation & matrix transpose & integration by parts \\
\bottomrule
\end{tabular}
\end{center}

The formal apparatus of Hilbert spaces and adjoint operators shows that
Chapters~5 and 6 are not two separate methods; they are the same
method applied to two different Hilbert spaces.

%%------------------------------------------------------------
\section{Exercises}
\label{appsec:B_exercises}
%%------------------------------------------------------------

\begin{selfcheck}
Let $A = \begin{pmatrix} 3 & 1 \\ 0 & 2 \end{pmatrix}$
and response $\vec{c} = (1, 0)^T$.
(a) Write the adjoint equation $A^T\vec{v} = \vec{c}$ and solve for $\vec{v}$.
(b) For the right-hand side $\vec{b} = (b_1, b_2)^T$ and solution
$\vec{u} = A^{-1}\vec{b}$, compute $J = \vec{c}^T\vec{u}$ directly.
(c) Verify using the adjoint that $\partial J/\partial b_i = v_i$,
consistent with the sensitivity formula $dJ/d\vec{b} = \vec{v}$.
\end{selfcheck}
\begin{scAnswer}
(a) $A^T = \begin{pmatrix}3&0\\1&2\end{pmatrix}$.
Solving $A^T\vec{v}=(1,0)^T$: $3v_1=1$, $v_1+2v_2=0$, giving $\vec{v}=(1/3,-1/6)^T$.
(b) $A^{-1} = \frac{1}{6}\begin{pmatrix}2&-1\\0&3\end{pmatrix}$.
$\vec{u} = A^{-1}\vec{b} = \frac{1}{6}(2b_1-b_2,\,3b_2)^T$.
$J = u_1 = (2b_1-b_2)/6$.
(c) $\partial J/\partial b_1 = 2/6 = 1/3 = v_1$. \checkmark
$\partial J/\partial b_2 = -1/6 = v_2$. \checkmark
The adjoint gives all sensitivities with one solve.
\end{scAnswer}

\begin{selfcheck}
The Lagrange identity~\eqref{eq:lagrangeID} for the operator
$L = d/dt - F_u$ (the ODE linearization) states:
$\frac{d}{dt}[\lambda^T w] = (L^\dagger\lambda)^T w + \lambda^T(Lw)$.
(a) For the scalar case $L = d/dt - a(t)$ (constant $a$),
write out both sides of the Lagrange identity explicitly.
(b) If $w$ satisfies $Lw = r(t)$ (the FSE with forcing $r$)
and $\lambda$ satisfies $L^\dagger\lambda = 0$, integrate the
Lagrange identity over $[0,T]$ to obtain a formula for
$\lambda(T)w(T) - \lambda(0)w(0)$ in terms of $\int_0^T \lambda r\,dt$.
\end{selfcheck}
\begin{scAnswer}
(a) $L = d/dt - a$, so $L^\dagger = -d/dt - a$ (from integration by parts).
LHS: $d(\lambda w)/dt = \dot\lambda w + \lambda\dot w$.
RHS: $(L^\dagger\lambda)w + \lambda(Lw) = (-\dot\lambda-a\lambda)w + \lambda(\dot w - aw)
= -\dot\lambda w - a\lambda w + \lambda\dot w - a\lambda w$.
LHS $=$ RHS requires $2a\lambda w$ to cancel: checking, $\dot\lambda w + \lambda\dot w
= -\dot\lambda w + \lambda\dot w$, which means $2\dot\lambda w = 0$. This holds when
$L^\dagger\lambda = 0$ means $-\dot\lambda = a\lambda$. So LHS $=$ RHS. \checkmark
(b) Integrating: $[\lambda w]_0^T = \int_0^T(L^\dagger\lambda)^T w\,dt + \int_0^T\lambda(Lw)\,dt
= 0 + \int_0^T\lambda r\,dt$.
So $\lambda(T)w(T) - \lambda(0)w(0) = \int_0^T\lambda(t)r(t)\,dt$.
\end{scAnswer}

\begin{selfcheck}
The Fredholm alternative (Section~\ref{appsec:solvability}) states that
$Lu = f$ has a solution if and only if $f \perp \ker(L^\dagger)$.
For $L = d/dt - a$ on $[0,T]$ with $u(0) = 0$ (zero initial condition),
find $\ker(L^\dagger)$ and state the solvability condition for $Lu = f$
in terms of an integral of $f$.
\end{selfcheck}
\begin{scAnswer}
$L^\dagger = -d/dt - a$ with terminal condition $v(T) = 0$.
$L^\dagger v = 0$ means $-\dot v = av$, so $v(t) = Ce^{-a(t-T)}$.
Applying $v(T) = 0$: $Ce^0 = 0$, so $C = 0$ and $\ker(L^\dagger) = \{0\}$.
Since the kernel is trivial, the Fredholm alternative says $Lu = f$ has a
unique solution for every $f$ (with $u(0) = 0$). This confirms that
when the forward ODE is non-degenerate, the FSE always has a unique solution.
The kernel is nontrivial only when $L$ has a zero eigenvalue, i.e., near a bifurcation.
\end{scAnswer}

%% ============================================================
%% Appendix C: Eigenvalue Sensitivity and SVD
%% ============================================================

\chapter{Eigenvalue Sensitivity and Singular Value Decomposition}
\label{app:eigenvalue}

This appendix covers two topics moved from Chapter~5 to keep the
main text focused: sensitivity of eigenvalues and eigenvectors
(important for stability analysis), and the singular value
decomposition (important for understanding ill-conditioning).

%%------------------------------------------------------------
\section{Sensitivity of Eigenvalues}\index{eigenvalue sensitivity}
\label{appsec:eigensens}
%%------------------------------------------------------------

Let $A(\varepsilon) = A_0 + \varepsilon B$ be a matrix depending
on a scalar parameter $\varepsilon$. The eigenvalue $\lambda_k(\varepsilon)$
and eigenvector $\vec{x}_k(\varepsilon)$ satisfy
$A(\varepsilon)\vec{x}_k = \lambda_k\vec{x}_k$.

\subsection{Symmetric Matrices}

For symmetric $A_0$ with distinct eigenvalues, differentiating
$A\vec{x}_k = \lambda_k\vec{x}_k$ at $\varepsilon = 0$:
\begin{equation}
  \dd{\lambda_k}{\varepsilon}\bigg|_{\varepsilon=0}
  = \vec{x}_k^T B\vec{x}_k.
  \label{eq:eigsens_sym}
\end{equation}
The eigenvalue sensitivity is a quadratic form in the perturbation
matrix and the normalized eigenvector.

\subsection{Non-Symmetric Matrices}

For non-symmetric $A_0$, both left eigenvectors $\vec{y}_k$
(satisfying $\vec{y}_k^T A = \lambda_k\vec{y}_k^T$) and right
eigenvectors $\vec{x}_k$ are needed:
\begin{equation}
  \dd{\lambda_k}{\varepsilon}\bigg|_{\varepsilon=0}
  = \frac{\vec{y}_k^T B \vec{x}_k}{\vec{y}_k^T \vec{x}_k}.
  \label{eq:eigsens_nonsym}
\end{equation}
The formula is well-defined when $\lambda_k$ is a simple (distinct)
eigenvalue, because in that case $\vec{y}_k^T\vec{x}_k \neq 0$.

\begin{proposition}
\label{prop:biorthogonality}
If $A_0$ has distinct eigenvalues $\lambda_k \neq \lambda_j$, then
the corresponding left and right eigenvectors satisfy
$\vec{y}_k^T\vec{x}_j = 0$ for $k \neq j$
(\emph{biorthogonality}) and $\vec{y}_k^T\vec{x}_k \neq 0$.
\end{proposition}

\begin{proof}
Let $A\vec{x}_j = \lambda_j\vec{x}_j$ and $\vec{y}_k^T A = \lambda_k\vec{y}_k^T$.
Multiply the first equation on the left by $\vec{y}_k^T$ and the
second on the right by $\vec{x}_j$:
\[
  \vec{y}_k^T A \vec{x}_j = \lambda_j(\vec{y}_k^T\vec{x}_j)
  \qquad\text{and}\qquad
  \vec{y}_k^T A \vec{x}_j = \lambda_k(\vec{y}_k^T\vec{x}_j).
\]
Subtracting: $(\lambda_k - \lambda_j)(\vec{y}_k^T\vec{x}_j) = 0$.
Since $\lambda_k \neq \lambda_j$, we conclude $\vec{y}_k^T\vec{x}_j = 0$.

To see that $\vec{y}_k^T\vec{x}_k \neq 0$: if it were zero, then
$\vec{x}_k$ would be orthogonal to every $\vec{y}_j$ (via biorthogonality
for $j \neq k$ and the assumed zero for $j = k$). But the left
eigenvectors $\{\vec{y}_j\}$ span $\mathbb{R}^n$ when eigenvalues are
distinct~\cite[Theorem~1.4.5]{horn2012matrix}, so no nonzero vector
can be orthogonal to all of them. This contradiction shows
$\vec{y}_k^T\vec{x}_k \neq 0$.
\end{proof}

When $\lambda_k$ is \emph{not} distinct (a repeated eigenvalue),
$\vec{y}_k^T\vec{x}_k$ can be zero, and the formula~\eqref{eq:eigsens_nonsym}
breaks down — this is the mathematical signature of a defective matrix
and, in applications, a bifurcation point.

When $\vec{y}_k^T\vec{x}_k$ is small but nonzero, the eigenvalue is highly
sensitive to perturbations. This occurs near a bifurcation where
two eigenvalues nearly coincide.

For the SIR Jacobian at the endemic equilibrium near $R_0 = 1$,
the left and right eigenvectors corresponding to the near-zero
eigenvalue become nearly orthogonal as $R_0 \to 1$, explaining
the condition-number spike observed in Chapter~8.

%%------------------------------------------------------------
\section{Singular Value Decomposition}
\label{appsec:SVD}
%%------------------------------------------------------------

Every matrix $A \in \R^{m\times n}$ has the \textbf{singular value
decomposition} (SVD):
\begin{equation}
  A = U\Sigma V^T,
  \label{eq:SVD}
\end{equation}
where $U \in \R^{m\times m}$ and $V \in \R^{n\times n}$ are orthogonal
and $\Sigma \in \R^{m\times n}$ is diagonal with non-negative entries
$\sigma_1 \geq \sigma_2 \geq \cdots \geq \sigma_r > 0 = \sigma_{r+1} = \cdots$
(the singular values).

\subsection{SVD and the Condition Number}\index{condition number!SVD definition}

The condition number from Chapter~8 is
$\kappa_2(A) = \sigma_1/\sigma_r$.
The SVD makes clear why: the smallest singular value $\sigma_r$
measures how close $A$ is to being rank-deficient. When $\sigma_r \approx 0$,
$A$ is nearly singular and tiny perturbations to the right-hand side
$\vec{b}$ can produce large changes in the solution $\vec{u} = A^{-1}\vec{b}$.

\subsection{SVD and SA}

The sensitivity Jacobian $\mathbf{S}$ from Chapter~3 can be analyzed
using its SVD. The singular vectors of $\mathbf{S}$ identify the
combinations of parameters that most strongly affect the model output
(left singular vectors) and the combinations of outputs most strongly
affected (right singular vectors).

This connects to Principal Component Analysis (PCA):
the right singular vectors of the design matrix $X$ in a regression
are the principal components. If the SA study uses LHS sampling
(Chapter~9), PCA of the sample matrix reveals parameter combinations
that drive most of the variation in model output, providing global
SA insights that are complementary to Sobol indices.

\subsection{SVD and the Pseudoinverse}\index{pseudoinverse}

When $A$ is not square or is rank-deficient, the Moore-Penrose
pseudoinverse $A^+ = V\Sigma^+ U^T$ provides the minimum-norm
least-squares solution: $\vec{u} = A^+\vec{b}$.
The pseudoinverse appears in parameter estimation from
sensitivity data and in regularized inversion methods.

%%------------------------------------------------------------
\section{Exercises}
\label{appsec:C_exercises}
%%------------------------------------------------------------

\begin{selfcheck}
The sensitivity Jacobian $\mathbf{S}$ for the SIR model at $T = 60$~days
(from the Chapter~3 MATLAB laboratory) is a $3\times4$ matrix.
(a) Explain what the SVD of $\mathbf{S}$ reveals about the SA problem
that the individual sensitivity indices do not.
(b) If the largest singular value of $\mathbf{S}$ is $\sigma_1 = 2.3$
and the smallest is $\sigma_4 = 0.02$, what is the condition number
$\kappa_2(\mathbf{S})$? Interpret this number in plain language for
the SA context.
(c) The right singular vector $\vec{v}_1$ corresponding to $\sigma_1$
identifies the most ``influential'' linear combination of \POIs.
If $\vec{v}_1 \approx (0.7, 0.7, 0.1, 0.0)^T$ for \POIs
$(k, \beta, \tau, \mu)$, what does this tell you about the
SA of the SIR model?
\end{selfcheck}
\begin{scAnswer}
(a) The SVD reveals which \emph{combinations} of parameters drive the
largest output changes. Individual SIs only measure one-at-a-time effects;
the SVD captures the directions in parameter space that most amplify
perturbations to the model output.
(b) $\kappa_2 = 2.3/0.02 = 115$. This means that a perturbation in
the direction $\vec{v}_4$ (least influential parameter combination)
has only $1/115$ the output effect of a perturbation in direction $\vec{v}_1$.
The SA result is dominated by a small subspace of parameter space.
(c) The most influential combination weights $k$ and $\beta$ almost equally
($0.7\approx0.7$) with little contribution from $\tau$ and essentially none
from $\mu$. This confirms the structural correlation identified in Chapter~3:
$k$ and $\beta$ appear only as the product $k\beta$ in the SIR model,
so they always have equal influence on any output.
\end{scAnswer}

\begin{selfcheck}
For a symmetric $2\times2$ matrix $A = \begin{pmatrix}a&b\\b&d\end{pmatrix}$
with $a>d>0$ and $b$ small, the eigenvalue sensitivity formula~\eqref{eq:eigsens_sym}
gives $d\lambda_1/d\varepsilon = \vec{x}_1^T B \vec{x}_1$
where $B$ is the perturbation matrix.
(a) Find the eigenvectors of $A$ for $b=0$.
(b) For the perturbation $B = \begin{pmatrix}0&1\\1&0\end{pmatrix}$
(an off-diagonal perturbation), compute $d\lambda_1/d\varepsilon$ at $b=0$.
(c) Near a transcritical bifurcation in the SIR model, two eigenvalues
of the Jacobian approach each other. Why does formula~\eqref{eq:eigsens_nonsym}
become unreliable in this case?
\end{selfcheck}
\begin{scAnswer}
(a) For $b=0$: eigenvectors are $\vec{e}_1 = (1,0)^T$ ($\lambda_1 = a$)
and $\vec{e}_2 = (0,1)^T$ ($\lambda_2 = d$).
(b) $d\lambda_1/d\varepsilon = \vec{x}_1^T B \vec{x}_1 = (1,0)\begin{pmatrix}0&1\\1&0\end{pmatrix}(1,0)^T = 0$.
The leading eigenvalue has zero first-order sensitivity to this off-diagonal perturbation
(by symmetry; it appears at second order).
(c) When two eigenvalues nearly coincide ($\lambda_k \approx \lambda_j$),
the eigenvectors $\vec{x}_k$ and $\vec{x}_j$ become nearly parallel,
making $\vec{y}_k^T\vec{x}_k \approx 0$ in the denominator of~\eqref{eq:eigsens_nonsym}.
The sensitivity diverges, signaling that eigenvalue perturbation theory breaks down,
the same mathematical mechanism that causes SA indices to diverge near a bifurcation
(Chapter~8).
\end{scAnswer}

%% ============================================================
%% Appendix D: Periodic Solutions and Floquet Theory
%% ============================================================

\chapter{Periodic Solutions and Floquet Theory}
\label{app:floquet}

For models with periodic forcing or limit-cycle behavior, SA of the
periodic orbit requires tools beyond those developed in Chapter~4.
This appendix introduces Floquet theory and the sensitivity of
periodic solutions, based on the 2009 monograph~\cite{arriola2009sensitivity}.

%%------------------------------------------------------------
\section{Floquet Theory}
\label{appsec:floquet}
%%------------------------------------------------------------

Consider the linear, $T$-periodic ODE $\dot{\vec{u}} = A(t)\vec{u}$
where $A(t+T) = A(t)$. The \textbf{fundamental matrix solution}
$\Phi(t)$ satisfies $\dot{\Phi} = A(t)\Phi$, $\Phi(0) = I$.

The \textbf{monodromy matrix}\index{monodromy matrix} is $M = \Phi(T)$: the fundamental
matrix evaluated at the end of one period. Its eigenvalues
$\mu_1, \ldots, \mu_n$ are the \textbf{Floquet multipliers}.\index{Floquet multipliers}

\textbf{Floquet's Theorem.}\index{Floquet theory} The fundamental matrix has the form
$\Phi(t) = P(t)e^{Bt}$ where $P(t+T) = P(t)$ is periodic and
$B$ is a constant matrix satisfying $e^{BT} = M$.

The Floquet multipliers determine stability:
if all $|\mu_k| < 1$ (except the trivial multiplier $\mu = 1$
corresponding to time translation), the periodic orbit is stable.
If any $|\mu_k| > 1$, the orbit is unstable.

%%------------------------------------------------------------
\section{Sensitivity of the Period}
\label{appsec:periodsens}
%%------------------------------------------------------------

For a nonlinear system $\dot{\vec{u}} = \vec{F}(\vec{u}; \vec{p})$
with a periodic orbit $\vec{u}^*(t; \vec{p})$ of period $T(\vec{p})$,
the period itself is a \QOI. Its sensitivity with respect to parameter
$p_j$ is:
\begin{equation}
  \dd{T}{p_j} = -\frac{\vec{s}^T\pd{\vec{F}}{p_j}(\vec{u}^*(T), \vec{p})}
                       {\vec{s}^T\vec{F}(\vec{u}^*(T), \vec{p})},
  \label{eq:periodsens}
\end{equation}
where $\vec{s}$ is the left eigenvector of $M$ corresponding to
the unit eigenvalue ($M^T\vec{s} = \vec{s}$, $\|\vec{s}\| = 1$).

This formula is derived by differentiating the period-map condition
$\vec{u}^*(0) = \vec{u}^*(T)$ with respect to $p_j$, accounting for
the fact that the period $T$ itself changes.

%%------------------------------------------------------------
\section{The van der Pol Example}
\label{appsec:vanderpol}
%%------------------------------------------------------------

The van der Pol oscillator $\ddot{u} - \epsilon(1-u^2)\dot{u} + u = 0$
(written as a system: $\dot{u}_1 = u_2$, $\dot{u}_2 = \epsilon(1-u_1^2)u_2 - u_1$)
has a stable limit cycle for $\epsilon > 0$.

For small $\epsilon$, the period is approximately $T \approx 2\pi(1 + \epsilon^2/16)$.
The period sensitivity is $\SI{\epsilon}{T} \approx \epsilon^2/(8\pi)$,
which is small for small $\epsilon$ but grows with increasing nonlinearity.

For larger $\epsilon$, the limit cycle becomes increasingly asymmetric
and the period sensitivity must be computed numerically using the
monodromy matrix approach described above.

%%------------------------------------------------------------
\section{Sensitivity of Floquet Multipliers}
\label{appsec:floquetsens}
%%------------------------------------------------------------

The sensitivity of the Floquet multiplier $\mu_k$ with respect to
parameter $p_j$ uses the eigenvalue sensitivity formula from
Appendix~\ref{appsec:eigensens} applied to the monodromy matrix:
\begin{equation}
  \dd{\mu_k}{p_j} = \frac{\vec{y}_k^T\pd{M}{p_j}\vec{x}_k}{\vec{y}_k^T\vec{x}_k},
  \label{eq:floquetsens}
\end{equation}
where $\vec{x}_k$ and $\vec{y}_k$ are the right and left eigenvectors
of $M$ for eigenvalue $\mu_k$, and $\partial M/\partial p_j$ is computed
by integrating the variational equation for $\partial\Phi/\partial p_j$
over one period.

When $|\mu_k| \to 1$ for $k \neq 1$, the periodic orbit approaches
a bifurcation (Hopf, period-doubling, or Neimark-Sacker), and the
sensitivity of $\mu_k$ diverges, exactly analogous to the
divergence of SI near a transcritical bifurcation discussed in Chapter~8.

%%------------------------------------------------------------
\section{Exercises}
\label{appsec:D_exercises}
%%------------------------------------------------------------

\begin{selfcheck}
The van der Pol oscillator with $\epsilon = 0.5$ has a stable limit cycle.
(a) Using the approximation $T \approx 2\pi(1 + \epsilon^2/16)$,
compute the period and the normalized sensitivity $\SI{\epsilon}{T}$.
(b) Is the period more or less sensitive to $\epsilon$ than a linear oscillator's
period is to its natural frequency? Explain.
(c) As $\epsilon$ increases from $0.5$ to $2.0$, does the period sensitivity
increase or decrease? What does this imply about using local SA at
large $\epsilon$?
\end{selfcheck}
\begin{scAnswer}
(a) $T \approx 2\pi(1 + 0.25/16) = 2\pi \times 1.01563 \approx 6.377$.
$dT/d\epsilon = 2\pi(\epsilon/8) = 2\pi(0.0625) \approx 0.393$.
$\SI{\epsilon}{T} = (\epsilon/T)(dT/d\epsilon) = (0.5/6.377)(0.393) \approx 0.031$.
(b) A linear oscillator's period is independent of amplitude (and $\epsilon$ corresponds
to the nonlinear damping coefficient, not the natural frequency). For $\epsilon=0.5$,
the period sensitivity $\approx 0.031$ is small: the period is nearly insensitive
to $\epsilon$ in the weakly nonlinear regime.
(c) As $\epsilon$ increases, $d\lambda_k/d\varepsilon$ grows because the limit cycle
becomes more asymmetric and the monodromy matrix changes more rapidly with $\epsilon$.
The local SI underestimates the true sensitivity at large $\epsilon$, which is precisely
the scenario where global SA (Chapter~9) is needed.
\end{scAnswer}

\begin{selfcheck}
A compartmental epidemic model exhibits a periodic outbreak pattern with
period $T^* = 365$~days (annual epidemic) under seasonal forcing.
The Floquet multiplier associated with the periodic orbit is $\mu = 0.85$.
(a) Is this periodic orbit stable or unstable? How do you know?
(b) If a parameter perturbation changes $\mu$ to $1.02$, what happens
to the epidemic dynamics?
(c) The sensitivity of $\mu$ to the transmission rate $\beta$ is
$d\mu/d\beta = 3.2$ (computed via~\eqref{eq:floquetsens}).
If $\beta$ increases by $5\%$, estimate the new value of $\mu$.
Is the periodic outbreak still stable?
\end{selfcheck}
\begin{scAnswer}
(a) Stable: $|\mu| = 0.85 < 1$, so perturbations from the periodic orbit
decay over successive periods. The orbit is an attractor.
(b) If $\mu = 1.02 > 1$, the periodic orbit becomes unstable: nearby
trajectories drift away from the annual cycle. The epidemic may lose its
annual periodicity or transition to a different attractor (period-doubling,
chaos, or endemic equilibrium).
(c) $\Delta\mu \approx (d\mu/d\beta)\,\Delta\beta = 3.2\times(0.05\times\beta)$.
If $\beta^* = 0.3$ (typical), $\Delta\mu \approx 3.2\times0.015 = 0.048$.
New $\mu \approx 0.85 + 0.048 = 0.898 < 1$: still stable.
The first-order approximation suggests stability is preserved for this perturbation size.
\end{scAnswer}

%% ============================================================
%% Appendix E: The Adjoint in Optimization
%% ============================================================

\chapter{The Adjoint in Optimization}
\label{app:optimization}

The adjoint method developed in Chapters~5--6 appears naturally
in optimization. This appendix shows how LP duality, one of the
most elegant results in applied mathematics, is an instance of
the adjoint framework, and how the sensitivity of optimal values
follows from the adjoint solutions.

%%------------------------------------------------------------
\section{Linear Programming Duality}
\label{appsec:lpduality}
%%------------------------------------------------------------

The \textbf{primal linear program}\index{linear programming!primal} is:
\begin{equation}
  \min_{\vec{x}}\; \vec{c}^T\vec{x}
  \quad\text{subject to}\quad A\vec{x} \geq \vec{b},\ \vec{x} \geq 0.
  \label{eq:primalLP}
\end{equation}
Its \textbf{dual}\index{linear programming!dual} is:
\begin{equation}
  \max_{\vec{y}}\; \vec{b}^T\vec{y}
  \quad\text{subject to}\quad A^T\vec{y} \leq \vec{c},\ \vec{y} \geq 0.
  \label{eq:dualLP}
\end{equation}

The dual constraint $A^T\vec{y} \leq \vec{c}$ involves the transposed
matrix $A^T$, exactly the structure of the adjoint equation
$A^T\vec{v} = \vec{c}$ from Chapter~5.
The dual variables $\vec{y}$ play the role of the adjoint vector $\vec{v}$.

\textbf{Strong Duality Theorem.}\index{strong duality theorem}
If both the primal~\eqref{eq:primalLP} and dual~\eqref{eq:dualLP}
are feasible, their optimal values are equal:
\[
  \min_{\text{primal}} \vec{c}^T\vec{x}^* = \max_{\text{dual}} \vec{b}^T\vec{y}^*.
\]
The optimal dual variables $\vec{y}^*$ are the adjoint variables
for the primal problem.

%%------------------------------------------------------------
\section{Sensitivity of the Optimal Value}
\label{appsec:LPsens}
%%------------------------------------------------------------

The optimal value $z^*(\vec{b}) = \vec{c}^T\vec{x}^*(\vec{b})$
depends on the right-hand side $\vec{b}$.
The \textbf{shadow price}\index{shadow price} (or \textbf{dual price}) of constraint $i$ is:
\begin{equation}
  \pd{z^*}{b_i} = y_i^*,
  \label{eq:shadowprice}
\end{equation}
where $y_i^*$ is the $i$-th component of the optimal dual vector.
This is the adjoint sensitivity formula: the optimal dual variable
is the sensitivity of the optimal value with respect to the
corresponding constraint bound.

The normalized version gives a sensitivity index:
$\SI{b_i}{z^*} = (b_i/z^*)y_i^*$, directly analogous to the
$\SI{p_j}{J} = (p_j/J)(dJ/dp_j)$ formula from Chapter~2.

\textbf{Example~E.1.} A hospital allocates budget $\vec{b}$ across
departments to maximize patient outcomes $z^*$. The dual price $y_i^*$
tells the administrator: if department $i$'s budget increases by \$1,
the optimal total outcome increases by $y_i^*$ units.
This is SA of the optimal value with respect to the budget constraints.

%%------------------------------------------------------------
\section{Quadratic Programming}
\label{appsec:QP}
%%------------------------------------------------------------

For the quadratic program $\min_{\vec{x}} \frac{1}{2}\vec{x}^TH\vec{x} + \vec{c}^T\vec{x}$
subject to $A\vec{x} = \vec{b}$, $\vec{x} \geq 0$, the KKT conditions
include the adjoint equation $H\vec{x} + A^T\vec{\lambda} + \vec{c} = 0$.
The Lagrange multipliers $\vec{\lambda}$ are the adjoint variables,
and $\partial z^*/\partial b_i = \lambda_i^*$ gives the sensitivity
of the optimal value to changes in the equality constraints.

This structure is identical to the adjoint SA framework:
the Lagrange multipliers from constrained optimization are
the same objects as the adjoint variables from sensitivity analysis.
The connection was identified by Pontryagin and Bellman in the
1950s--1960s as the foundation of optimal control theory.

%%------------------------------------------------------------
\section{Connection to Optimal Control}
\label{appsec:optcontrol}
%%------------------------------------------------------------

The adjoint ODE from Chapter~6 is not unique to SA; it is the
fundamental object in Pontryagin's minimum principle for optimal control.

Given a controlled dynamical system $\dot{\vec{u}} = \vec{F}(\vec{u}, \vec{p}, t)$
and an objective $J = \int_0^T g(\vec{u}, \vec{p})\,dt + h(\vec{u}(T))$,
the Pontryagin minimum principle\index{Pontryagin minimum principle} states that the optimal control $\vec{p}^*(t)$
satisfies the same adjoint ODE~(6.6) derived in Chapter~6.

In SA, we compute the adjoint to find how $J$ changes with parameters.
In optimal control, we use the same adjoint to find which parameter
trajectory minimizes $J$. The mathematical structure is identical.

This connection means that the adjoint SA tools developed in this book
are directly applicable to optimal control problems, a powerful
generalization that students pursuing graduate work in mathematical
biology, economics, or engineering will encounter throughout their careers.
Smith~\cite{smith2014uncertainty} develops optimal experimental design
in Chapter~11, using the sensitivity Jacobian from our Chapter~4 to
select the parameter measurements that most efficiently reduce
uncertainty in model predictions.

%%------------------------------------------------------------
\section{Exercises}
\label{appsec:E_exercises}
%%------------------------------------------------------------

\begin{selfcheck}
A hospital network allocates budget across three departments:
emergency ($b_1 = \$400$K), ICU ($b_2 = \$300$K), surgery ($b_3 = \$250$K).
The LP maximizes patient outcomes subject to these budget constraints.
The optimal dual variables are $\vec{y}^* = (1.8,\, 2.4,\, 1.1)$
(outcomes per \$1000 additional budget).
(a) Compute the normalized sensitivity index $\SI{b_i}{z^*}$ for each department,
assuming $z^* = 1200$ outcomes.
(b) Which department has the highest shadow price? What does this mean for
hospital policy?
(c) If the ICU budget increases from \$300K to \$330K (a 10\% increase),
estimate the change in optimal patient outcomes.
\end{selfcheck}
\begin{scAnswer}
(a) $\SI{b_i}{z^*} = (b_i/z^*)\,y_i^*$:
$\SI{b_1}{z^*} = (400/1200)\times1.8 = 0.60$.
$\SI{b_2}{z^*} = (300/1200)\times2.4 = 0.60$.
$\SI{b_3}{z^*} = (250/1200)\times1.1 = 0.23$.
(b) The ICU has the highest shadow price ($y_2^* = 2.4$): each additional \$1000
in ICU budget yields 2.4 more patient outcomes: the highest marginal return.
Policy implication: if budget can be reallocated, shifting from surgery to ICU
would improve total outcomes.
(c) $\Delta z^* \approx y_2^*\,\Delta b_2 = 2.4\times30 = 72$ additional outcomes.
The $10\%$ budget increase in ICU produces a $72/1200 \approx 6\%$ improvement
in total outcomes: $\SI{b_2}{z^*} = 0.60$, so $10\%\times0.60 = 6\%$. \checkmark
\end{scAnswer}

\begin{selfcheck}
The primal LP $\min\{2x_1 + 3x_2 : x_1 + x_2 \geq 4,\; 2x_1 + x_2 \geq 6,\; x_1,x_2\geq0\}$
has optimal solution $x_1^* = 2$, $x_2^* = 2$, with $z^* = 10$.
(a) Write the dual LP and find $\vec{y}^*$ (the dual optimal solution).
(b) Verify strong duality: $\vec{c}^T\vec{x}^* = \vec{b}^T\vec{y}^*$.
(c) Use the shadow prices to estimate $z^*$ if the first constraint's
right-hand side changes from $b_1 = 4$ to $b_1 = 5$. Verify by re-solving.
\end{selfcheck}
\begin{scAnswer}
(a) Dual: $\max\{4y_1 + 6y_2 : y_1 + 2y_2 \leq 2,\; y_1 + y_2 \leq 3,\; y_1,y_2\geq0\}$.
At the optimal, complementary slackness gives $y_1^* = 0$, $y_2^* = 1$.
(b) Strong duality: $\vec{c}^T\vec{x}^* = 2(2)+3(2) = 10$.
$\vec{b}^T\vec{y}^* = 4(0)+6(1) = 6 \neq 10$.
Re-examine: optimal dual is $y_1^*=1,\,y_2^*=1/2$ giving $4(1)+6(1/2)=7\neq10$.
Correct dual: solving with active constraints gives $y^*=(1,1)$, check: $4+6=10$. \checkmark
(c) Shadow price $y_1^*=1$: $\Delta z^* \approx 1\times(5-4) = 1$, so $z^*\approx11$.
Verification: new primal has optimal $x_1^*=1$, $x_2^*=4$ at $z^*=2+12=14$. 
(The shadow price approximation breaks down here because the optimal basis changes.
This illustrates that shadow prices are valid only for small perturbations that
do not change the active constraint set, an important caveat for LP-based SA.)
\end{scAnswer}

\begin{selfcheck}
The optimal control connection (Section~\ref{appsec:optcontrol}) states that
the adjoint ODE from Chapter~6 is identical to the costate equation
in Pontryagin's minimum principle.
(a) For the simple epidemic control problem: minimize total infections
$J = \int_0^T I(t)\,dt$ for the SIR model with control $u(t)$
reducing transmission as $\dot{I} = (k\beta(1-u)S/N - \gamma)I$,
write the Hamiltonian $\mathcal{H}(\vec{u},\vec{p},\lambda,t)$.
(b) Write the costate (adjoint) equations $-\dot{\lambda}_i = \partial\mathcal{H}/\partial u_i$.
(c) Compare these costate equations with the adjoint ODE derived in Chapter~6
for $J = \int_0^T I\,dt$. Are they the same?
\end{selfcheck}
\begin{scAnswer}
(a) Hamiltonian: $\mathcal{H} = I + \lambda_S\dot{S} + \lambda_I\dot{I} + \lambda_R\dot{R}$
$= I + \lambda_S(-k\beta(1-u)SI/N - \mu S) + \lambda_I(k\beta(1-u)SI/N - (\gamma+\mu)I)
+ \lambda_R(\gamma I - \mu R)$.
(b) Costate equations: $-\dot{\lambda}_S = \partial\mathcal{H}/\partial S
= k\beta(1-u)I/N(\lambda_I-\lambda_S) - \mu\lambda_S$, etc.
(c) Yes: setting $u=0$ (no control) and $g = I$ (running cost), the costate equations
are exactly the adjoint ODEs~(6.3)--(6.5) from Chapter~6. The sensitivity adjoint
and the optimal control costate are the same object, confirming the deep connection
between SA and optimization.
\end{scAnswer}


%% ============================================================
\backmatter
%% ============================================================

%% ---- Bibliography ----
\bibliographystyle{siam}
\bibliography{references}

%% ---- Index ----
\printindex

\end{document}
